\input{preamble}

% OK, start here.
%
\begin{document}

\title{Cohomologie des schémas}


\maketitle

\phantomsection
\label{section-phantom}

\tableofcontents

\section{Introduction}
\label{section-introduction}

\noindent
Dans ce chapitre, nous démontrons d'abord plusieurs résultats sur la cohomologie des
faisceaux quasi-cohérents. Une référence fondamentale est \cite{EGA}.
Cela fait, nous développerons la cohomologie des
faisceaux cohérents dans le cadre noethérien. Voir \cite{FAC}.







\section{Cohomologie de {\v C}ech des faisceaux quasi-cohérents}
\label{section-cech-quasi-coherent}

\noindent
Soit $X$ un schéma.
Soit $U \subset X$ un ouvert affine.
Rappelons qu'un {\it recouvrement standard} de $U$ est un recouvrement
de la forme $\mathcal{U} : U = \bigcup_{i = 1}^n D(f_i)$,
où $f_1, \ldots, f_n \in \Gamma(U, \mathcal{O}_X)$ engendrent
l'idéal unité, voir
Schémas, définition \ref{schemes-definition-standard-covering}.

\begin{lemma}
\label{lemma-cech-cohomology-quasi-coherent-trivial}
Soit $X$ un schéma.
Soit $\mathcal{F}$ un $\mathcal{O}_X$-module quasi-cohérent.
Soit $\mathcal{U} : U = \bigcup_{i = 1}^n D(f_i)$ un recouvrement
standard d'un ouvert affine de $X$.
Alors $\check{H}^p(\mathcal{U}, \mathcal{F}) = 0$ pour
tout $p > 0$.
\end{lemma}

\begin{proof}
Écrivons $U = \Spec(A)$ pour un anneau $A$.
Autrement dit, $f_1, \ldots, f_n$ sont des éléments de $A$
qui engendrent l'idéal unité de $A$.
Écrivons $\mathcal{F}|_U = \widetilde{M}$ pour un $A$-module $M$.
Clairement, le complexe de {\v C}ech
$\check{\mathcal{C}}^\bullet(\mathcal{U}, \mathcal{F})$
s'identifie au complexe
$$
\prod\nolimits_{i_0} M_{f_{i_0}} \to
\prod\nolimits_{i_0i_1} M_{f_{i_0}f_{i_1}} \to
\prod\nolimits_{i_0i_1i_2} M_{f_{i_0}f_{i_1}f_{i_2}} \to
\ldots
$$
Il s'agit de montrer que le complexe augmenté
\begin{equation}
\label{equation-extended}
0 \to
M \to
\prod\nolimits_{i_0} M_{f_{i_0}} \to
\prod\nolimits_{i_0i_1} M_{f_{i_0}f_{i_1}} \to
\prod\nolimits_{i_0i_1i_2} M_{f_{i_0}f_{i_1}f_{i_2}} \to
\ldots
\end{equation}
(dont nous avons étudié la troncature dans
Algèbre, lemme \ref{algebra-lemma-cover-module}) est exact.
Il suffit de montrer que (\ref{equation-extended})
est exact après localisation en un idéal premier $\mathfrak p$, voir
Algèbre, lemme \ref{algebra-lemma-characterize-zero-local}.
En fait, nous montrerons que le complexe augmenté localisé
en $\mathfrak p$ est homotope à zéro.

\medskip\noindent
Il existe un indice $i$ tel que $f_i \not \in \mathfrak p$.
Choisissons et fixons un tel indice $i_{\text{fix}}$. Remarquons que
$M_{f_{i_{\text{fix}}}, \mathfrak p} = M_{\mathfrak p}$. De même,
dans une localisation en un produit $f_{i_0} \ldots f_{i_p}$ et en $\mathfrak p$,
nous pouvons omettre tout $f_{i_j}$ tel que $i_j = i_{\text{fix}}$.
Définissons une homotopie
$$
h :
\prod\nolimits_{i_0 \ldots i_{p + 1}}
M_{f_{i_0} \ldots f_{i_{p + 1}}, \mathfrak p}
\longrightarrow
\prod\nolimits_{i_0 \ldots i_p}
M_{f_{i_0} \ldots f_{i_p}, \mathfrak p}
$$
par la formule
$$
h(s)_{i_0 \ldots i_p} = s_{i_{\text{fix}} i_0 \ldots i_p}
$$
(Ceci est ``dual'' de l'homotopie dans la démonstration de
Cohomologie, lemme \ref{cohomology-lemma-homology-complex}.)
Autrement dit, $h : \prod_{i_0} M_{f_{i_0}, \mathfrak p} \to M_\mathfrak p$
est la projection sur le facteur
$M_{f_{i_{\text{fix}}}, \mathfrak p} = M_{\mathfrak p}$ et, en général,
l'application $h$ est la projection sur les facteurs
$M_{f_{i_{\text{fix}}} f_{i_1} \ldots f_{i_{p + 1}}, \mathfrak p}
= M_{f_{i_1} \ldots f_{i_{p + 1}}, \mathfrak p}$. Calculons
\begin{align*}
(dh + hd)(s)_{i_0 \ldots i_p}
& =
\sum\nolimits_{j = 0}^p
(-1)^j
h(s)_{i_0 \ldots \hat i_j \ldots i_p}
+
d(s)_{i_{\text{fix}} i_0 \ldots i_p}\\
& =
\sum\nolimits_{j = 0}^p
(-1)^j
s_{i_{\text{fix}} i_0 \ldots \hat i_j \ldots i_p}
+
s_{i_0 \ldots i_p}
+
\sum\nolimits_{j = 0}^p
(-1)^{j + 1}
s_{i_{\text{fix}} i_0 \ldots \hat i_j \ldots i_p} \\
& =
s_{i_0 \ldots i_p}
\end{align*}
Cela montre que l'application identique est homotope à zéro, comme voulu.
\end{proof}

\noindent
Le lemme suivant dit notamment que pour tout schéma affine
$X$ et tout faisceau quasi-cohérent $\mathcal{F}$ sur $X$, on a
$$
H^p(X, \mathcal{F}) = 0
$$
pour tout $p > 0$.

\begin{lemma}
\label{lemma-quasi-coherent-affine-cohomology-zero}
\begin{slogan}
Théorème d'annulation de Serre : la cohomologie supérieure s'annule sur les schémas
affines pour les modules quasi-cohérents.
\end{slogan}
Soit $X$ un schéma.
Soit $\mathcal{F}$ un $\mathcal{O}_X$-module
quasi-cohérent. Pour tout ouvert affine $U \subset X$ nous avons
$H^p(U, \mathcal{F}) = 0$ pour tout $p > 0$.
\end{lemma}

\begin{proof}
Nous allons appliquer Cohomologie,
lemme \ref{cohomology-lemma-cech-vanish-basis}.
Comme base $\mathcal{B}$ pour la topologie de $X$, nous utiliserons les
ouverts affines de $X$.
Comme ensemble $\text{Cov}$ de recouvrements ouverts, nous allons utiliser les recouvrements ouverts
standards des ouverts affines de $X$.
Nous vérifions que les conditions (1), (2) et (3) de
Cohomologie, lemme \ref{cohomology-lemma-cech-vanish-basis}
sont satisfaites. Notez que l'intersection d'ouverts standards d'un ouvert affine est encore un
ouvert standard. Par conséquent, la propriété (1) est
vérifiée. Ces recouvrements forment un système cofinal de recouvrements ouverts de tout élément
de $\mathcal{B}$, voir
Schémas, lemme \ref{schemes-lemma-standard-open}.
Donc (2) est
vérifiée. Enfin, la condition (3) du lemme découle du
lemme \ref{lemma-cech-cohomology-quasi-coherent-trivial}.
\end{proof}

\noindent
Voici une version relative de l'annulation de la cohomologie des faisceaux quasi-cohérents sur
les schémas affines.

\begin{lemma}
\label{lemma-relative-affine-vanishing}
Soit $f : X \to S$ un morphisme de schémas.
Soit $\mathcal{F}$ un $\mathcal{O}_X$-module quasi-cohérent.
Si $f$ est affine, alors $R^if_*\mathcal{F} = 0$ pour tout $i > 0$.
\end{lemma}

\begin{proof}
D'après Cohomologie,
lemme \ref{cohomology-lemma-describe-higher-direct-images}
le faisceau
$R^if_*\mathcal{F}$ est le faisceau associé au préfaisceau
$V \mapsto H^i(f^{-1}(V), \mathcal{F}|_{f^{-1}(V)})$. Par
hypothèse, chaque fois que $V$ est affine, $f^{-1}(V)$ est
affine ; voir Morphismes, définition \ref{morphisms-definition-affine}.
Par le lemme \ref{lemma-quasi-coherent-affine-cohomology-zero} nous concluons que
$H^i(f^{-1}(V), \mathcal{F}|_{f^{-1}(V)}) = 0$
chaque fois que $V$ est affine. Puisque $S$ possède une base constituée d'ouverts affines, le
faisceau est nul.
\end{proof}

\begin{lemma}
\label{lemma-relative-affine-cohomology}
Soit $f : X \to S$ un morphisme affine de schémas.
Soit $\mathcal{F}$ un $\mathcal{O}_X$-module quasi-cohérent.
Alors $H^i(X, \mathcal{F}) = H^i(S, f_*\mathcal{F})$ pour tout $i \geq 0$.
\end{lemma}

\begin{proof}
Cela résulte du lemme \ref{lemma-relative-affine-vanishing}
et de la suite spectrale de Leray. Voir
Cohomologie, lemme \ref{cohomology-lemma-apply-Leray}.
\end{proof}

\noindent
Les deux lemmes suivants précisent quand la cohomologie de {\v C}ech peut
être utilisée pour calculer la cohomologie de modules quasi-cohérents.

\begin{lemma}
\label{lemma-affine-diagonal}
Soit $X$ un schéma. Les éléments suivants sont équivalents
\begin{enumerate}
\item $X$ a une diagonale affine $\Delta : X \to X \times X$,
\item pour tous ouverts affines $U, V \subset X$, l'intersection
$U \cap V$ est affine, et
\item il existe un recouvrement ouvert $\mathcal{U} : X = \bigcup_{i \in I} U_i$
tel que $U_{i_0 \ldots i_p}$ est ouvert affine pour tous les $p \ge 0$ et tous les
$i_0, \ldots, i_p \in I$.
\end{enumerate}
En particulier, ces conditions sont satisfaites si $X$ est séparé.
\end{lemma}

\begin{proof}
Supposons que $X$ ait une diagonale affine. Soient $U, V \subset X$ des ouverts affines.
Alors $U \cap V = \Delta^{-1}(U \times V)$ est affine. Ainsi (2) est
vraie. Il est immédiat que (2) implique (3). Inversement, s'il existe un
recouvrement de $X$ comme dans (3), alors $X \times X = \bigcup U_i \times U_{i'}$
est un recouvrement ouvert affine, et on voit que
$\Delta^{-1}(U_i \times U_{i'}) = U_i \cap U_{i'}$
est affine. Alors $\Delta$ est un morphisme affine par
Morphismes, lemme \ref{morphisms-lemma-characterize-affine}.
L'assertion finale découle de Schémas, Lemme
\ref{schemes-lemma-characterize-separated}.
\end{proof}

\begin{lemma}
\label{lemma-cech-cohomology-quasi-coherent}
Soit $X$ un schéma.
Soit $\mathcal{U} : X = \bigcup_{i \in I} U_i$ un recouvrement ouvert tel que
$U_{i_0 \ldots i_p}$ soit ouvert affine pour tout $p \ge 0$ et tout
$i_0, \ldots, i_p \in I$.
Dans ce cas, pour tout faisceau quasi-cohérent $\mathcal{F}$, on a
$$
\check{H}^p(\mathcal{U}, \mathcal{F}) = H^p(X, \mathcal{F})
$$
comme $\Gamma(X, \mathcal{O}_X)$-modules, pour tout $p$.
\end{lemma}

\begin{proof}
D'après le
lemme \ref{lemma-quasi-coherent-affine-cohomology-zero},
c'est un cas particulier de
Cohomologie, Lemme
\ref{cohomology-lemma-cech-spectral-sequence-application}.
\end{proof}







\section{Annulation de la cohomologie}
\label{section-vanishing}

\noindent
Nous avons vu que, sur un schéma affine, les groupes de cohomologie supérieurs
de tout faisceau quasi-cohérent s'annulent
(lemme \ref{lemma-quasi-coherent-affine-cohomology-zero}).
Il s'avère que cela caractérise aussi
les schémas affines. Nous donnons deux versions.

\begin{lemma}
\label{lemma-quasi-compact-h1-zero-covering}
\begin{reference}
\cite{Serre-criterion}, \cite[II, Théorème 5.2.1 (d') et IV (1.7.17)]{EGA}
\end{reference}
\begin{slogan}
Le critère d'affinité de Serre.
\end{slogan}
Soit $X$ un schéma.
Supposons que
\begin{enumerate}
\item $X$ soit quasi-compact,
\item pour tout faisceau quasi-cohérent d'idéaux
$\mathcal{I} \subset \mathcal{O}_X$ nous avons $H^1(X, \mathcal{I}) = 0$.
\end{enumerate}
Alors $X$ est affine.
\end{lemma}

\begin{proof}
Soit $x \in X$ un point fermé. Soit $U \subset X$ un voisinage ouvert
affine de $x$. Écrivons $U = \Spec(A)$ et soit
$\mathfrak m \subset A$ l'idéal maximal correspondant à $x$.
Posons $Z = X \setminus U$ et $Z' = Z \cup \{x\}$.
Par Schémas, lemme \ref{schemes-lemma-reduced-closed-subscheme},
il existe des faisceaux quasi-cohérents d'idéaux
$\mathcal{I}$, resp.\ $\mathcal{I}'$, définissant
les sous-schémas fermés réduits $Z$, resp.\ $Z'$.
Considérons la suite exacte courte
$$
0 \to \mathcal{I}' \to \mathcal{I} \to \mathcal{I}/\mathcal{I}' \to 0.
$$
Puisque $x$ est un point fermé de $X$ et $x \not \in Z$ nous voyons que
$\mathcal{I}/\mathcal{I}'$ est supporté en $x$. En fait, la restriction
de $\mathcal{I}/\mathcal{I'}$ à $U$ correspond au $A$-module
$A/\mathfrak m$. On voit donc que $\Gamma(X, \mathcal{I}/\mathcal{I'})
= A/\mathfrak m$. Puisque par hypothèse $H^1(X, \mathcal{I}') = 0$
nous voyons qu'il existe une section globale $f \in \Gamma(X, \mathcal{I})$
qui s'envoie sur l'élément $1 \in A/\mathfrak m$ en tant que section de
$\mathcal{I}/\mathcal{I'}$. Manifestement, nous avons
$x \in X_f \subset U$. Cela implique que $X_f = D(f_A)$ où
$f_A$ est l'image de $f$ dans $A = \Gamma(U, \mathcal{O}_X)$.
En particulier, $X_f$ est affine.

\medskip\noindent
Considérons la réunion $W = \bigcup X_f$ portant sur les $f \in \Gamma(X, \mathcal{O}_X)$
tels que $X_f$ soit affine. Il est clair que $W$ est ouvert dans $X$.
D'après les arguments ci-dessus, chaque point fermé de
$X$ est contenu dans $W$. Le sous-ensemble fermé $X \setminus W$ de $X$
est également quasi-compact
(voir Topologie, lemme \ref{topology-lemma-closed-in-quasi-compact}).
Il a donc un point fermé s'il n'est pas vide (voir
Topologie, lemme \ref{topology-lemma-quasi-compact-closed-point}).
Cela contredirait le fait que tous les points fermés appartiennent à
$W$. On conclut donc $X = W$.

\medskip\noindent
Choisissons un nombre fini de $f_1, \ldots, f_n \in \Gamma(X, \mathcal{O}_X)$
tels que $X = X_{f_1} \cup \ldots \cup X_{f_n}$ et tels que chaque
$X_{f_i}$ soit affine. C'est possible comme nous l'avons vu ci-dessus.
Par Propriétés, lemme \ref{properties-lemma-characterize-affine}
pour terminer la preuve il suffit
de montrer que $f_1, \ldots, f_n$ engendrent l'idéal unité dans
$\Gamma(X, \mathcal{O}_X)$. Considérons la suite exacte courte
$$
\xymatrix{
0 \ar[r] &
\mathcal{F} \ar[r] &
\mathcal{O}_X^{\oplus n} \ar[rr]^{f_1, \ldots, f_n} & &
\mathcal{O}_X \ar[r] &
0
}
$$
La flèche définie par $f_1, \ldots, f_n$ est surjective puisque
l'ouverture $X_{f_i}$ recouvrent $X$. Notons $\mathcal{F}$ le noyau de
cette application surjective.
Remarquons que $\mathcal{F}$ possède une filtration
$$
0 = \mathcal{F}_0 \subset \mathcal{F}_1 \subset
\ldots \subset \mathcal{F}_n = \mathcal{F}
$$
dont chaque sous-quotient $\mathcal{F}_i/\mathcal{F}_{i - 1}$ est
isomorphe à un faisceau quasi-cohérent d'idéaux. Plus
précisément, on peut prendre $\mathcal{F}_i$ égal à l'intersection de
$\mathcal{F}$ avec les $i$ premiers facteurs directs de
$\mathcal{O}_X^{\oplus n}$.
L'hypothèse
du lemme implique que $H^1(X, \mathcal{F}_i/\mathcal{F}_{i - 1}) = 0$
pour tout $i$. Cela implique que
$H^1(X, \mathcal{F}_2) = 0$, car ce groupe se place entre
$H^1(X, \mathcal{F}_1)$ et $H^1(X, \mathcal{F}_2/\mathcal{F}_1)$.
En continuant ainsi on en déduit que $H^1(X, \mathcal{F}) = 0$.
Nous concluons donc que l'application
$$
\xymatrix{
\bigoplus\nolimits_{i = 1, \ldots, n} \Gamma(X, \mathcal{O}_X)
\ar[rr]^{f_1, \ldots, f_n} & &
\Gamma(X, \mathcal{O}_X)
}
$$
est surjective comme souhaité.
\end{proof}

\noindent
Notons que, si $X$ est un schéma noethérien, tout faisceau quasi-cohérent
d'idéaux est automatiquement un faisceau cohérent d'idéaux, donc un faisceau quasi-cohérent
d'idéaux de type fini. Par conséquent, le
lemme précédent et le lemme suivant s'appliquent tous deux dans ce cas.

\begin{lemma}
\label{lemma-quasi-separated-h1-zero-covering}
\begin{reference}
\cite{Serre-criterion}, \cite[II, Théorème 5.2.1]{EGA}
\end{reference}
\begin{slogan}
Le critère d'affinité de Serre.
\end{slogan}
Soit $X$ un schéma. Supposons que
\begin{enumerate}
\item $X$ soit quasi-compact,
\item $X$ soit quasi-séparé, et
\item $H^1(X, \mathcal{I}) = 0$ pour tout faisceau
quasi-cohérent d'idéaux $\mathcal{I}$ de type fini.
\end{enumerate}
Alors $X$ est affine.
\end{lemma}

\begin{proof}
Par
Propriétés, lemme \ref{properties-lemma-quasi-coherent-colimit-finite-type},
tout faisceau quasi-cohérent d'idéaux est une limite inductive filtrante
de faisceaux quasi-cohérents d'idéaux de type fini.
Par Cohomologie, lemme \ref{cohomology-lemma-quasi-separated-cohomology-colimit}
prenant la cohomologie sur $X$ commute aux limites inductives
filtrantes. On voit donc que $H^1(X, \mathcal{I}) = 0$
pour tout faisceau quasi-cohérent d'idéaux sur $X$. Autrement
dit, le lemme \ref{lemma-quasi-compact-h1-zero-covering} s'applique.
\end{proof}

\noindent
Les arguments précédents fournissent une condition suffisante pour qu'un faisceau
inversible soit ample. Nous avertissons toutefois le lecteur que cette
condition n'est pas nécessaire.

\begin{lemma}
\label{lemma-quasi-compact-h1-zero-invertible}
Soit $X$ un schéma. Soit $\mathcal{L}$ un $\mathcal{O}_X$-module
inversible. Supposons que
\begin{enumerate}
\item $X$ soit quasi-compact,
\item pour tout faisceau quasi-cohérent d'idéaux
$\mathcal{I} \subset \mathcal{O}_X$
il existe un $n \geq 1$ tel que
$H^1(X, \mathcal{I} \otimes_{\mathcal{O}_X} \mathcal{L}^{\otimes n}) = 0$.
\end{enumerate}
Alors $\mathcal{L}$ est ample.
\end{lemma}

\begin{proof}
Cela se démontre exactement de la même manière que le
lemme \ref{lemma-quasi-compact-h1-zero-covering}.
Soit $x \in X$ un point fermé. Soit $U \subset X$ un voisinage ouvert
affine de $x$ tel que $\mathcal{L}|_U \cong \mathcal{O}_U$.
Écrivons $U = \Spec(A)$ et soit
$\mathfrak m \subset A$ l'idéal maximal correspondant à $x$.
Posons $Z = X \setminus U$ et $Z' = Z \cup \{x\}$.
Par Schémas, lemme \ref{schemes-lemma-reduced-closed-subscheme},
il existe des faisceaux quasi-cohérents d'idéaux
$\mathcal{I}$, resp.\ $\mathcal{I}'$, définissant
les sous-schémas fermés réduits $Z$, resp.\ $Z'$.
Considérons la suite exacte courte
$$
0 \to \mathcal{I}' \to \mathcal{I} \to \mathcal{I}/\mathcal{I}' \to 0.
$$
Pour chaque $n \geq 1$ nous obtenons une suite exacte courte
$$
0 \to \mathcal{I}' \otimes_{\mathcal{O}_X} \mathcal{L}^{\otimes n}
\to \mathcal{I} \otimes_{\mathcal{O}_X} \mathcal{L}^{\otimes n} \to
\mathcal{I}/\mathcal{I}' \otimes_{\mathcal{O}_X} \mathcal{L}^{\otimes n} \to 0.
$$
D'après notre hypothèse, nous pouvons choisir $n$ tel que
$H^1(X, \mathcal{I}' \otimes_{\mathcal{O}_X} \mathcal{L}^{\otimes n}) = 0$.
Puisque $x$ est un point fermé de $X$ et $x \not \in Z$, nous voyons que
$\mathcal{I}/\mathcal{I}'$ est supporté en $x$. En fait, la restriction
de $\mathcal{I}/\mathcal{I'}$ à $U$ correspond au $A$-module
$A/\mathfrak m$. Puisque $\mathcal{L}$ est trivial sur $U$,
nous voyons que la restriction de
$\mathcal{I}/\mathcal{I}' \otimes_{\mathcal{O}_X} \mathcal{L}^{\otimes n}$
à $U$ correspond également au $A$-module $A/\mathfrak m$.
Ainsi, nous voyons que
$\Gamma(X, \mathcal{I}/\mathcal{I'} \otimes_{\mathcal{O}_X}
\mathcal{L}^{\otimes n}) = A/\mathfrak m$.
Par notre choix de $n$, nous voyons qu'il existe une section globale
$s \in \Gamma(X, \mathcal{I} \otimes_{\mathcal{O}_X} \mathcal{L}^{\otimes n})$
qui s'envoie sur l'élément $1 \in A/\mathfrak m$. Manifestement, nous avons
$x \in X_s \subset U$ parce que $s$ s'annule aux points de $Z$.
Cela implique que $X_s = D(f)$ où
$f \in A$ est l'image de $s$ dans $A \cong \Gamma(U, \mathcal{L}^{\otimes n})$.
En particulier, $X_s$ est affine.

\medskip\noindent
Considérons la réunion $W = \bigcup X_s$ portant sur les
$s \in \Gamma(X, \mathcal{L}^{\otimes n})$ pour $n \geq 1$
tels que $X_s$ soit affine. Il est clair que $W$ est ouvert dans $X$.
D'après les arguments ci-dessus, chaque point fermé de
$X$ est contenu dans $W$. Le sous-ensemble fermé $X \setminus W$ de $X$
est également quasi-compact
(voir Topologie, lemme \ref{topology-lemma-closed-in-quasi-compact}).
Il possède donc un point fermé s'il n'est pas vide (voir
Topologie, lemme \ref{topology-lemma-quasi-compact-closed-point}).
Cela contredirait le fait que tous les points fermés appartiennent à
$W$. Nous concluons donc $X = W$. Cela signifie que $\mathcal{L}$
est ample par Propriétés, définition \ref{properties-definition-ample}.
\end{proof}

\noindent
Il existe une variante du lemme \ref{lemma-quasi-compact-h1-zero-invertible}
pour les faisceaux d'idéaux de type fini ; nous la formulerons et la démontrerons ici
si nous en avons besoin.

\begin{lemma}
\label{lemma-criterion-affine-morphism}
Soit $f : X \to Y$ un morphisme quasi-compact, avec $X$ et $Y$ quasi-séparés.
Si $R^1f_*\mathcal{I} = 0$ pour tout faisceau quasi-cohérent d'idéaux
$\mathcal{I}$ sur $X$, alors $f$ est affine.
\end{lemma}

\begin{proof}
Soit $V \subset Y$ un sous-schéma ouvert affine. Il suffit de montrer que
$U = f^{-1}(V)$ est affine. Le morphisme d'inclusion $V \to Y$ est quasi-compact
d'après Schémas, lemme \ref{schemes-lemma-quasi-compact-permanence}.
Par conséquent, le changement de base $U \to X$ est quasi-compact, voir
Schémas, lemme \ref{schemes-lemma-quasi-compact-preserved-base-change}.
Ainsi, tout faisceau quasi-cohérent d'idéaux $\mathcal{I}$ sur $U$
s'étend à un faisceau quasi-cohérent d'idéaux sur $X$, voir
Propriétés, lemme \ref{properties-lemma-extend-trivial}.
Puisque la formation de $R^1f_*$ est locale sur $Y$
(Cohomologie, section \ref{cohomology-section-locality}),
nous concluons que $R^1(U \to V)_*\mathcal{I} = 0$ par l'hypothèse du
lemme. Ainsi, par la suite spectrale de Leray
(Cohomologie, lemme \ref{cohomology-lemma-Leray}),
nous concluons que $H^1(U, \mathcal{I}) = H^1(V, (U \to V)_*\mathcal{I})$.
Puisque $(U \to V)_*\mathcal{I}$ est quasi-cohérent d'après
Schémas, lemme \ref{schemes-lemma-push-forward-quasi-coherent}, nous avons
$H^1(V, (U \to V)_*\mathcal{I}) = 0$ d'après le
lemme \ref{lemma-quasi-coherent-affine-cohomology-zero}.
Ainsi, on conclut que $U$ est affine d'après le
lemme \ref{lemma-quasi-compact-h1-zero-covering}.
\end{proof}















\section{Quasi-cohérence des images directes supérieures}
\label{section-quasi-coherence}

\noindent
Nous avons vu que les groupes de cohomologie supérieure d'un module quasi-cohérent sur un schéma
affine sont nuls. Pour les schémas quasi-compacts et (quasi-)séparés $X$, il
en résulte l'annulation des groupes de cohomologie des
faisceaux quasi-cohérents au-delà d'un certain degré. Toutefois, $X$ n'est pas nécessairement
de dimension cohomologique finie, car celle-ci est définie par l'annulation de la cohomologie
de {\it tous} les $\mathcal{O}_X$-modules.

\begin{lemma}[Principe de récurrence]
\label{lemma-induction-principle}
\begin{reference}
\cite[Proposition 3.3.1]{BvdB}
\end{reference}
Soit $X$ un schéma quasi-compact et quasi-séparé. Soit $P$ une propriété
des ouverts quasi-compacts de $X$. Supposons que
\begin{enumerate}
\item $P$ est vérifiée pour tout ouvert affine de $X$,
\item si $U$ est un ouvert quasi-compact, $V$ un ouvert affine, et si
$P$ est vérifiée pour $U$, $V$ et $U \cap V$, alors
$P$ est vérifiée pour $U \cup V$.
\end{enumerate}
Alors $P$ est vérifiée pour tout ouvert quasi-compact de $X$
et en particulier pour $X$.
\end{lemma}

\begin{proof}
Montrons d'abord par récurrence que $P$ est vérifiée pour les ouverts quasi-compacts {\it séparés}
$W \subset X$. En effet, un tel ouvert s'écrit
$W = U_1 \cup \ldots \cup U_n$ et l'on raisonne par récurrence sur $n$ en utilisant
la propriété (2) avec $U = U_1 \cup \ldots \cup U_{n - 1}$ et $V = U_n$.
Cela est possible car
$U \cap V = (U_1 \cap U_n) \cup \ldots \cup (U_{n - 1} \cap U_n)$
est aussi une réunion de $n - 1$ sous-schémas ouverts affines, d'après
Schémas, lemme \ref{schemes-lemma-characterize-separated},
appliqué aux ouverts affines $U_i$ et $U_n$ de $W$.
Cela étant, pour tout ouvert quasi-compact $W \subset X$, on
peut raisonner par récurrence sur le nombre d'ouverts affines nécessaires pour recouvrir $W$,
en employant le même procédé et en observant que l'ouvert quasi-compact
$U_i \cap U_n$ est séparé, comme sous-schéma ouvert du schéma affine $U_n$.
\end{proof}

\begin{lemma}
\label{lemma-vanishing-nr-affines}
\begin{slogan}
Pour les schémas à diagonale affine, la cohomologie des
modules quasi-cohérents s'annule en degrés supérieurs au nombre d'ouverts affines
nécessaires pour former un recouvrement.
\end{slogan}
Soit $X$ un schéma quasi-compact à diagonale affine (par exemple
si $X$ est séparé).
Soit $t = t(X)$ le nombre minimal d'ouverts affines nécessaires pour
recouvrir $X$. Alors $H^n(X, \mathcal{F}) = 0$ pour tout $n \geq t$ et tout
faisceau quasi-cohérent $\mathcal{F}$.
\end{lemma}

\begin{proof}
Première démonstration.
Raisonnons par récurrence sur $t$.
Si $t = 1$, le résultat découle du
lemme \ref{lemma-quasi-coherent-affine-cohomology-zero}.
Si $t > 1$, écrivons $X = U \cup V$, où $V$ est un ouvert affine et
$U = U_1 \cup \ldots \cup U_{t - 1}$ une réunion de $t - 1$ ouverts affines.
Notons que, dans ce cas,
$U \cap V =  (U_1 \cap V) \cup \ldots (U_{t - 1} \cap V)$
est aussi une réunion de $t - 1$ sous-schémas ouverts affines.
En effet, la diagonale étant affine, l'intersection de deux
ouverts affines est affine, voir lemme \ref{lemma-affine-diagonal}.
Appliquons la suite exacte longue de Mayer-Vietoris
$$
0 \to
H^0(X, \mathcal{F}) \to
H^0(U, \mathcal{F}) \oplus H^0(V, \mathcal{F}) \to
H^0(U \cap V, \mathcal{F}) \to
H^1(X, \mathcal{F}) \to \ldots
$$
voir Cohomologie, lemme \ref{cohomology-lemma-mayer-vietoris}.
Par récurrence, les groupes $H^i(U, \mathcal{F})$,
$H^i(V, \mathcal{F})$, $H^i(U \cap V, \mathcal{F})$ sont nuls pour
$i \geq t - 1$. Il en résulte immédiatement que $H^i(X, \mathcal{F})$
est nul pour $i \geq t$.

\medskip\noindent
Deuxième démonstration.
Soit $\mathcal{U} : X = \bigcup_{i = 1}^t U_i$ un recouvrement ouvert affine
fini. Puisque $X$ est à diagonale affine, les intersections multiples
$U_{i_0 \ldots i_p}$ sont toutes affines, voir
lemme \ref{lemma-affine-diagonal}.
D'après le lemme \ref{lemma-cech-cohomology-quasi-coherent}, les groupes de cohomologie de {\v C}ech
$\check{H}^p(\mathcal{U}, \mathcal{F})$
coïncident avec les groupes de cohomologie. D'après
Cohomologie, lemme \ref{cohomology-lemma-alternating-usual},
les groupes de cohomologie de {\v C}ech se calculent au moyen du complexe alterné de
{\v C}ech $\check{\mathcal{C}}_{alt}^\bullet(\mathcal{U}, \mathcal{F})$.
Le recouvrement ayant $t$ éléments, on voit immédiatement
que $\check{\mathcal{C}}_{alt}^p(\mathcal{U}, \mathcal{F}) = 0$
pour tout $p \geq t$. Le résultat s'ensuit.
\end{proof}

\begin{lemma}
\label{lemma-affine-diagonal-universal-delta-functor}
Soit $X$ un schéma quasi-compact à diagonale affine
(par exemple si $X$ est séparé). Alors
\begin{enumerate}
\item pour tout $\mathcal{O}_X$-module quasi-cohérent $\mathcal{F}$,
il existe une injection $\mathcal{F} \to \mathcal{F}'$
de $\mathcal{O}_X$-modules
quasi-cohérents telle que $H^p(X, \mathcal{F}') = 0$ pour tout $p \geq 1$, et
\item $\{H^n(X, -)\}_{n \geq 0}$
est un $\delta$-foncteur universel de $\QCoh(\mathcal{O}_X)$ dans
$\textit{Ab}$.
\end{enumerate}
\end{lemma}

\begin{proof}
Soit $X = \bigcup U_i$ un recouvrement ouvert affine fini.
Posons $U = \coprod U_i$ et notons $j : U \to X$
le morphisme induisant les immersions ouvertes données $U_i \to X$.
Puisque $U$ est un schéma affine et que $X$ est à diagonale affine,
le morphisme $j$ est affine, voir
Morphismes, lemme \ref{morphisms-lemma-affine-permanence}.
Pour tout $\mathcal{O}_X$-module $\mathcal{F}$, il
existe un morphisme canonique $\mathcal{F} \to j_*j^*\mathcal{F}$.
Ce morphisme est injectif, comme on le vérifie sur les fibres :
si $x \in U_i$, on a une factorisation
$$
\mathcal{F}_x \to (j_*j^*\mathcal{F})_x
\to (j^*\mathcal{F})_{x'} = \mathcal{F}_x
$$
où $x' \in U$ est le point $x$ considéré comme point de $U_i \subset U$.
Si maintenant $\mathcal{F}$ est quasi-cohérent, alors $j^*\mathcal{F}$
est quasi-cohérent sur le schéma affine $U$, donc sa cohomologie
supérieure est nulle d'après le
lemme \ref{lemma-quasi-coherent-affine-cohomology-zero}. On
a alors $H^p(X, j_*j^*\mathcal{F}) = 0$ pour
$p > 0$ d'après le lemme \ref{lemma-relative-affine-cohomology},
puisque $j$ est affine. Cela prouve (1).
Enfin, on voit que le morphisme
$H^p(X, \mathcal{F}) \to H^p(X, j_*j^*\mathcal{F})$
est nul, et l'assertion (2) découle de
Homologie, lemme \ref{homology-lemma-efface-implies-universal}.
\end{proof}

\begin{lemma}
\label{lemma-vanishing-nr-affines-quasi-separated}
Soit $X$ un schéma quasi-compact et quasi-séparé.
Soit $X = U_1 \cup \ldots \cup U_n$ un recouvrement ouvert
dont chaque $U_i$ est quasi-compact et séparé (par exemple affine).
Posons
$$
d = \max\nolimits_{I \subset \{1, \ldots, n\}}
\left(|I| + t(\bigcap\nolimits_{i \in I} U_i) - 1\right)
$$
où $t(U)$ est le nombre minimal d'ouverts affines nécessaires pour recouvrir
le schéma $U$. Alors $H^p(X, \mathcal{F}) = 0$ pour tout $p \geq d$ et tout
faisceau quasi-cohérent $\mathcal{F}$.
\end{lemma}

\begin{proof}
Notons que, puisque $X$ est quasi-séparé et les $U_i$ quasi-compacts, les nombres
$t(\bigcap_{i \in I} U_i)$ sont finis. Raisonnons par récurrence sur $n$.
Si $n = 1$, le résultat découle du lemme \ref{lemma-vanishing-nr-affines}.
Si $n > 1$, écrivons $X = U \cup V$, avec $U = U_1 \cup \ldots \cup U_{n - 1}$
et $V = U_n$. Appliquons la suite exacte longue de Mayer-Vietoris
$$
0 \to
H^0(X, \mathcal{F}) \to
H^0(U, \mathcal{F}) \oplus H^0(V, \mathcal{F}) \to
H^0(U \cap V, \mathcal{F}) \to
H^1(X, \mathcal{F}) \to \ldots
$$
voir Cohomologie, lemme \ref{cohomology-lemma-mayer-vietoris}.
Pour achever la démonstration lorsque $q \geq d$, montrons que
$H^q(V, \mathcal{F})$, $H^q(U, \mathcal{F})$ et
$H^{q - 1}(U \cap V, \mathcal{F})$ sont
nuls. Le cas $n = 1$ donne $H^q(V, \mathcal{F}) = 0$ pour
$q \geq t(V) = t(U_n)$. Comme $t(V) \leq d$, on obtient l'annulation voulue. Par
l'hypothèse de récurrence, on a $H^q(U, \mathcal{F}) = 0$ pour
$$
q \geq \max\nolimits_{I \subset \{1, \ldots, n - 1\}}
\left(|I| + t(\bigcap\nolimits_{i \in I} U_i) - 1\right)
$$
L'entier de droite étant inférieur ou égal à $d$, on
obtient l'annulation voulue. Enfin, appliquons l'hypothèse de récurrence à l'ouvert
$U \cap V = (U_1 \cap U_n) \cup \ldots \cup (U_{n - 1} \cap U_n)$ pour
obtenir l'annulation $H^q(U \cap V, \mathcal{F}) = 0$ pour
$$
q \geq \max\nolimits_{I \subset \{1, \ldots, n - 1\}}
\left(|I| + t(U_n \cap \bigcap\nolimits_{i \in I} U_i) - 1\right)
$$
L'entier de droite étant strictement inférieur à $d$,
le lemme s'ensuit.
\end{proof}

\begin{lemma}
\label{lemma-quasi-coherence-higher-direct-images}
Soit $f : X \to S$ un morphisme de schémas.
Supposons que $f$ soit quasi-séparé et quasi-compact.
\begin{enumerate}
\item Pour tout $\mathcal{O}_X$-module quasi-cohérent $\mathcal{F}$,
les images directes supérieures $R^pf_*\mathcal{F}$ sont quasi-cohérentes sur $S$.
\item Si $S$ est quasi-compact, il existe un entier $n = n(X, S, f)$
tel que $R^pf_*\mathcal{F} = 0$ pour tout $p \geq n$ et tout
faisceau quasi-cohérent $\mathcal{F}$ sur $X$.
\item En fait, si $S$ est quasi-compact, on peut trouver $n = n(X, S, f)$
tel que, pour tout
morphisme de schémas $S' \to S$, on ait $R^p(f')_*\mathcal{F}' = 0$
pour $p \geq n$ et tout faisceau quasi-cohérent $\mathcal{F}'$
sur $X'$. Ici $f' : X' = S' \times_S X \to S'$ est le changement de base de $f$.
\end{enumerate}
\end{lemma}

\begin{proof}
Prouvons d'abord (1). Notons que, sous les hypothèses du lemme, le faisceau
$R^0f_*\mathcal{F} = f_*\mathcal{F}$ est quasi-cohérent d'après
Schémas, lemme \ref{schemes-lemma-push-forward-quasi-coherent}.
D'après
Cohomologie, lemme \ref{cohomology-lemma-localize-higher-direct-images},
la formation des images directes supérieures commute à la
restriction aux sous-schémas ouverts. La quasi-cohérence étant locale sur $S$,
on se ramène au cas traité au paragraphe suivant.

\medskip\noindent
Démonstration de (1) lorsque $S$ est affine. Utilisons le principe de récurrence.
Puisque $f$ est quasi-compact et quasi-séparé, $X$
est quasi-compact et quasi-séparé. Pour un ouvert quasi-compact $U \subset X$
et $a = f|_U$, notons $P(U)$ la propriété suivante :
$R^pa_*\mathcal{F}$ est quasi-cohérent sur $S$ pour tout module quasi-cohérent
$\mathcal{F}$ sur $U$ et tout $p \geq 0$. Comme $P(X)$ est l'assertion (1), il suffit
de vérifier les conditions (1) et (2) du lemme \ref{lemma-induction-principle}.
Si $U$ est affine, alors $P(U)$ est vérifiée car $R^pa_*\mathcal{F} = 0$ pour
$p \geq 1$ (d'après le lemme \ref{lemma-relative-affine-vanishing} et
Morphismes, lemme \ref{morphisms-lemma-morphism-affines-affine}),
et nous avons déjà établi le résultat pour $p = 0$ au premier
paragraphe. Soient ensuite $U \subset X$ un ouvert quasi-compact et $V \subset X$
un ouvert affine; supposons $P(U)$, $P(V)$ et $P(U \cap V)$ vérifiées.
Posons $a = f|_U$, $b = f|_V$, $c = f|_{U \cap V}$ et $g = f|_{U \cup V}$.
Pour tout $\mathcal{O}_{U \cup V}$-module quasi-cohérent $\mathcal{F}$,
on dispose alors de la suite de Mayer-Vietoris relative
$$
0 \to
g_*\mathcal{F} \to
a_*(\mathcal{F}|_U) \oplus b_*(\mathcal{F}|_V) \to
c_*(\mathcal{F}|_{U \cap V}) \to
R^1g_*\mathcal{F} \to \ldots
$$
voir Cohomologie, lemme \ref{cohomology-lemma-relative-mayer-vietoris}.
D'après $P(U)$, $P(V)$ et $P(U \cap V)$, les faisceaux
$R^pa_*(\mathcal{F}|_U)$, $R^pb_*(\mathcal{F}|_V)$ et
$R^pc_*(\mathcal{F}|_{U \cap V})$ sont tous quasi-cohérents.
Les résultats sur les faisceaux quasi-cohérents de
Schémas, section \ref{schemes-section-quasi-coherent},
impliquent alors que chacun des faisceaux
$R^pg_*\mathcal{F}$ est quasi-cohérent : il est le terme médian d'une suite exacte
courte dont le terme de gauche est le conoyau d'un morphisme
de faisceaux quasi-cohérents et le terme de droite le noyau d'un morphisme de faisceaux quasi-cohérents.
D'où $P(U \cup V)$, ce qui achève la démonstration de (1).

\medskip\noindent
Prouvons ensuite (3), et a fortiori (2). Choisissons un recouvrement ouvert affine
fini $S = \bigcup_{j = 1, \ldots m} S_j$. Pour chaque $j$,
choisissons un recouvrement ouvert affine fini
$f^{-1}(S_j) = \bigcup_{i = 1, \ldots t_j} U_{ji} $.
Soit
$$
d_j = \max\nolimits_{I \subset \{1, \ldots, t_j\}}
\left(|I| + t(\bigcap\nolimits_{i \in I} U_{ji})\right)
$$
l'entier obtenu dans le
lemme \ref{lemma-vanishing-nr-affines-quasi-separated}.
Nous affirmons que $n(X, S, f) = \max d_j$ convient.

\medskip\noindent
En effet, soient $S' \to S$ un morphisme de schémas et
$\mathcal{F}'$ un faisceau quasi-cohérent sur $X' = S' \times_S X$.
Montrons que $R^pf'_*\mathcal{F}' = 0$ pour $p \geq n(X, S, f)$.
Cette question étant locale sur $S'$, on peut supposer que $S'$ est affine
et s'envoie dans $S_j$ pour un certain $j$. Alors $X' = S' \times_{S_j} f^{-1}(S_j)$
est recouvert par les ouverts affines $S' \times_{S_j} U_{ji}$, $i = 1, \ldots t_j$,
et les intersections
$$
\bigcap\nolimits_{i \in I} S' \times_{S_j} U_{ji} =
S' \times_{S_j} \bigcap\nolimits_{i \in I} U_{ji}
$$
sont recouvertes par le même nombre d'ouverts affines qu'avant le changement de base. En appliquant
le
lemme \ref{lemma-vanishing-nr-affines-quasi-separated},
on obtient $H^p(X', \mathcal{F}') = 0$. La première partie de la démonstration
montre déjà que chaque $R^qf'_*\mathcal{F}'$ est quasi-cohérent,
donc a ses groupes de cohomologie supérieure nuls sur le schéma affine $S'$;
ainsi $H^0(S', R^pf'_*\mathcal{F}') = H^p(X', \mathcal{F}') = 0$
d'après Cohomologie, lemme \ref{cohomology-lemma-apply-Leray}.
Puisque $R^pf'_*\mathcal{F}'$ est quasi-cohérent,
on en déduit que $R^pf'_*\mathcal{F}' = 0$.
\end{proof}

\begin{lemma}
\label{lemma-quasi-coherence-higher-direct-images-application}
Soit $f : X \to S$ un morphisme de schémas.
Supposons que $f$ soit quasi-séparé et quasi-compact.
Supposons $S$ affine.
Pour tout $\mathcal{O}_X$-module quasi-cohérent $\mathcal{F}$,
on a
$$
H^q(X, \mathcal{F}) = H^0(S, R^qf_*\mathcal{F})
$$
pour tout $q \in \mathbf{Z}$.
\end{lemma}

\begin{proof}
Considérons la suite spectrale de Leray $E_2^{p, q} = H^p(S, R^qf_*\mathcal{F})$
convergeant vers $H^{p + q}(X, \mathcal{F})$, voir
Cohomologie, lemme \ref{cohomology-lemma-Leray}.
D'après le lemme \ref{lemma-quasi-coherence-higher-direct-images},
les faisceaux $R^qf_*\mathcal{F}$ sont quasi-cohérents.
D'après le lemme \ref{lemma-quasi-coherent-affine-cohomology-zero},
on a $E_2^{p, q} = 0$ lorsque $p > 0$.
La suite spectrale dégénère donc en $E_2$, ce qui
conclut. Voir aussi
Cohomologie, lemme \ref{cohomology-lemma-apply-Leray} (2),
pour le principe général.
\end{proof}








\section{Cohomologie et changement de base, I}
\label{section-cohomology-and-base-change}

\noindent
Soit $f : X \to S$ un morphisme de schémas.
Soit $\mathcal{F}$ un faisceau quasi-cohérent sur $X$.
Soit de plus $g : S' \to S$ un morphisme quelconque de schémas. Notons
$X' = X_{S'} = S' \times_S X$ le changement de base de $X$ et
$f' : X' \to S'$ le changement de base de $f$.
Notons aussi $g' : X' \to X$ la projection,
et posons $\mathcal{F}' = (g')^*\mathcal{F}$.
La situation est représentée par le diagramme suivant :
\begin{equation}
\label{equation-base-change-diagram}
\vcenter{
\xymatrix{
\mathcal{F}' = (g')^*\mathcal{F} &
X' \ar[r]_{g'} \ar[d]_{f'} &
X \ar[d]^f &
\mathcal{F} \\
Rf'_*\mathcal{F}' &
S' \ar[r]^g &
S &
Rf_*\mathcal{F}
}
}
\end{equation}
Voici le cas le plus simple de la propriété de changement de base envisagée.

\begin{lemma}
\label{lemma-affine-base-change}
Soit $f : X \to S$ un morphisme de schémas.
Soit $\mathcal{F}$ un $\mathcal{O}_X$-module quasi-cohérent.
Supposons $f$ affine. Dans
ce cas, $f_*\mathcal{F} \cong Rf_*\mathcal{F}$ est
un faisceau quasi-cohérent et, pour tout diagramme de changement de base
(\ref{equation-base-change-diagram})
on a
$$
g^*f_*\mathcal{F} = f'_*(g')^*\mathcal{F}.
$$
\end{lemma}

\begin{proof}
L'annulation des images directes supérieures est donnée par le
lemme \ref{lemma-relative-affine-vanishing}.
L'énoncé est local sur $S$ et $S'$. On peut donc
supposer $X = \Spec(A)$, $S = \Spec(R)$,
$S' = \Spec(R')$ et $\mathcal{F} = \widetilde{M}$,
pour un certain $A$-module $M$.
Utilisons Schémas, lemme \ref{schemes-lemma-widetilde-pullback},
pour décrire les images inverses et directes de $\mathcal{F}$. On
a $X' = \Spec(R' \otimes_R A)$ et
$\mathcal{F}'$ est le faisceau quasi-cohérent associé
à $(R' \otimes_R A) \otimes_A M$.
Le lemme se ramène donc à
l'égalité
$$
(R' \otimes_R A) \otimes_A M = R' \otimes_R M
$$
de $R'$-modules.
\end{proof}

\noindent
Dans de nombreuses situations, il suffit de connaître le cas particulier suivant
de la cohomologie et du changement de base. Il découle immédiatement
des résultats plus forts de la
section \ref{section-cohomology-and-base-change-derived},
mais son importance justifie une démonstration distincte.

\begin{lemma}[Changement de base plat]
\label{lemma-flat-base-change-cohomology}
Considérons un diagramme cartésien de schémas
$$
\xymatrix{
X' \ar[d]_{f'} \ar[r]_{g'} & X \ar[d]^f \\
S' \ar[r]^g & S
}
$$
Soit $\mathcal{F}$ un $\mathcal{O}_X$-module
quasi-cohérent, d'image inverse $\mathcal{F}' = (g')^*\mathcal{F}$.
Supposons que $g$ soit plat et que $f$ soit quasi-compact et quasi-séparé.
Pour tout $i \geq 0$,
\begin{enumerate}
\item le morphisme de changement de base
de Cohomologie, lemme \ref{cohomology-lemma-base-change-map-flat-case},
est un isomorphisme
$$
g^*R^if_*\mathcal{F} \longrightarrow R^if'_*\mathcal{F}',
$$
\item si $S = \Spec(A)$ et $S' = \Spec(B)$, alors
$H^i(X, \mathcal{F}) \otimes_A B = H^i(X', \mathcal{F}')$.
\end{enumerate}
\end{lemma}

\begin{proof}
On peut appliquer Cohomologie, lemme \ref{cohomology-lemma-base-change-map-flat-case}, dans
(1) car $g'$ est plat d'après
Morphismes, lemme \ref{morphisms-lemma-base-change-flat}.
Cela étant, l'assertion (1) découle de (2). En effet,
l'assertion (1) est locale sur $S'$, et l'on peut donc supposer $S$
et $S'$ affines. Autrement dit, on a $S = \Spec(A)$
et $S' = \Spec(B)$ comme dans
(2). Puisque $R^if_*\mathcal{F}$ est quasi-cohérent
(lemme \ref{lemma-quasi-coherence-higher-direct-images}),
c'est le $\mathcal{O}_S$-module quasi-cohérent associé au
$A$-module $H^0(S, R^if_*\mathcal{F}) = H^i(X, \mathcal{F})$
(égalité donnée par
le lemme \ref{lemma-quasi-coherence-higher-direct-images-application}). De
même, $R^if'_*\mathcal{F}'$ est le
$\mathcal{O}_{S'}$-module quasi-cohérent associé au $B$-module
$H^i(X', \mathcal{F}')$. L'image inverse par $g$ correspondant au
foncteur $- \otimes_A B$ sur les
modules (Schémas, lemme \ref{schemes-lemma-widetilde-pullback}),
il suffit de prouver (2).

\medskip\noindent
Soit $A \to B$ un homomorphisme plat d'anneaux.
Soit $X$ un schéma quasi-compact et quasi-séparé sur $A$.
Soit $\mathcal{F}$ un $\mathcal{O}_X$-module quasi-cohérent.
Posons $X_B = X \times_{\Spec(A)} \Spec(B)$ et notons
$\mathcal{F}_B$ l'image inverse de $\mathcal{F}$.
Il s'agit de montrer que le morphisme
$$
H^i(X, \mathcal{F}) \otimes_A B \longrightarrow H^i(X_B, \mathcal{F}_B)
$$
(donné par la référence citée dans l'énoncé du lemme)
est un isomorphisme.

\medskip\noindent
Lorsque $X$ est séparé, choisissons un recouvrement ouvert affine
$\mathcal{U} : X = U_1 \cup \ldots \cup U_t$ et rappelons que
$$
\check{H}^p(\mathcal{U}, \mathcal{F}) = H^p(X, \mathcal{F}),
$$
voir
lemme \ref{lemma-cech-cohomology-quasi-coherent}. Pour
le recouvrement $\mathcal{U}_B : X_B = (U_1)_B \cup \ldots \cup (U_t)_B$ obtenu
par changement de base, chaque $(U_i)_B$ est encore
affine et $X_B$ est séparé. On obtient donc
$$
\check{H}^p(\mathcal{U}_B, \mathcal{F}_B) = H^p(X_B, \mathcal{F}_B).
$$
On a la relation suivante entre les complexes de {\v C}ech
$$
\check{\mathcal{C}}^\bullet(\mathcal{U}_B, \mathcal{F}_B) =
\check{\mathcal{C}}^\bullet(\mathcal{U}, \mathcal{F}) \otimes_A B
$$
comme il résulte du
lemme \ref{lemma-affine-base-change}.
Puisque $A \to B$ est plat, cette égalité subsiste en passant à la cohomologie.

\medskip\noindent
Lorsque $X$ est quasi-séparé, choisissons un recouvrement ouvert affine
$\mathcal{U} : X = U_1 \cup \ldots \cup U_t$. Utilisons la suite spectrale reliant la cohomologie de
{\v C}ech à la cohomologie, donnée
dans Cohomologie, lemme \ref{cohomology-lemma-cech-spectral-sequence}.
Le lecteur qui souhaite éviter cette suite spectrale peut utiliser
Mayer-Vietoris et raisonner par récurrence sur $t$, comme dans la démonstration du
lemme \ref{lemma-quasi-coherence-higher-direct-images}.
La suite spectrale a pour terme $E_2$
$E_2^{p, q} = \check{H}^p(\mathcal{U}, \underline{H}^q(\mathcal{F}))$
et converge vers $H^{p + q}(X, \mathcal{F})$.
De même, on dispose d'une suite spectrale de terme $E_2$
$E_2^{p, q} = \check{H}^p(\mathcal{U}_B, \underline{H}^q(\mathcal{F}_B))$
qui converge vers $H^{p + q}(X_B, \mathcal{F}_B)$.
Puisque les intersections $U_{i_0 \ldots i_p}$ sont quasi-compactes
et séparées, le résultat du deuxième paragraphe de la démonstration donne
$\check{H}^p(\mathcal{U}_B, \underline{H}^q(\mathcal{F}_B)) =
\check{H}^p(\mathcal{U}, \underline{H}^q(\mathcal{F})) \otimes_A B$.
La platitude de $A \to B$ permet de conclure que
$H^i(X, \mathcal{F}) \otimes_A B \to H^i(X_B, \mathcal{F}_B)$
est un isomorphisme pour tout $i$, ce qui conclut.
\end{proof}

\begin{lemma}[Changement de base fini localement libre]
\label{lemma-finite-locally-free-base-change-cohomology}
Considérons un diagramme cartésien de schémas
$$
\xymatrix{
Y \ar[d]_{g} \ar[r]_h & X \ar[d]^f \\
\Spec(B) \ar[r] & \Spec(A)
}
$$
Soit $\mathcal{F}$ un $\mathcal{O}_X$-module
quasi-cohérent, d'image inverse $\mathcal{G} = h^*\mathcal{F}$.
Si $B$ est un $A$-module localement libre de type fini, alors
$H^i(X, \mathcal{F}) \otimes_A B = H^i(Y, \mathcal{G})$.
\end{lemma}

\noindent
{\bf Attention} : N'utilisez ce lemme que si vous comprenez la différence entre
celui-ci et le lemme \ref{lemma-flat-base-change-cohomology}.

\begin{proof}
Lorsque $X$ est séparé, choisissons un recouvrement ouvert affine
$\mathcal{U} : X = \bigcup_{i \in I} U_i$ et rappelons que
$$
\check{H}^p(\mathcal{U}, \mathcal{F}) = H^p(X, \mathcal{F}),
$$
voir
lemme \ref{lemma-cech-cohomology-quasi-coherent}.
Soit $\mathcal{V} : Y = \bigcup_{i \in I} g^{-1}(U_i)$
le recouvrement ouvert affine correspondant de $Y$.
Les ouverts $V_i = g^{-1}(U_i) = U_i \times_{\Spec(A)} \Spec(B)$
sont affines et $Y$ est séparé. On obtient donc
$$
\check{H}^p(\mathcal{V}, \mathcal{G}) = H^p(Y, \mathcal{G}).
$$
Nous affirmons que le morphisme de complexes de {\v C}ech
$$
\check{\mathcal{C}}^\bullet(\mathcal{U}, \mathcal{F}) \otimes_A B
\longrightarrow
\check{\mathcal{C}}^\bullet(\mathcal{V}, \mathcal{G})
$$
est un isomorphisme. En effet, puisque $B$ est de présentation finie comme $A$-module,
le produit tensoriel avec $B$ sur $A$ commute aux produits, voir
Algèbre, proposition \ref{algebra-proposition-fp-tensor}.
Il suffit donc de montrer que les morphismes
$\Gamma(U_{i_0 \ldots i_p}, \mathcal{F}) \otimes_A B \to
\Gamma(V_{i_0 \ldots i_p}, \mathcal{G})$
sont des isomorphismes, ce qui résulte du
lemme \ref{lemma-affine-base-change}.
Puisque $A \to B$ est plat, cela reste vrai en passant à la cohomologie.

\medskip\noindent
Dans le cas général, on raisonne exactement de même en utilisant un recouvrement
ouvert affine $\mathcal{U} : X = \bigcup_{i \in I} U_i$ et le
recouvrement correspondant $\mathcal{V} : Y = \bigcup_{i \in I} V_i$,
avec $V_i = g^{-1}(U_i)$ comme ci-dessus. Utilisons la suite spectrale reliant la cohomologie de
{\v C}ech à la cohomologie, donnée
dans Cohomologie, lemme \ref{cohomology-lemma-cech-spectral-sequence}.
La suite spectrale a pour terme $E_2$
$E_2^{p, q} = \check{H}^p(\mathcal{U}, \underline{H}^q(\mathcal{F}))$
et converge vers $H^{p + q}(X, \mathcal{F})$.
De même, on dispose d'une suite spectrale de terme $E_2$
$E_2^{p, q} = \check{H}^p(\mathcal{V}, \underline{H}^q(\mathcal{G}))$
qui converge vers $H^{p + q}(Y, \mathcal{G})$.
Les intersections $U_{i_0 \ldots i_p}$ étant séparées, le résultat
du paragraphe précédent donne des isomorphismes
$\Gamma(U_{i_0 \ldots i_p}, \underline{H}^q(\mathcal{F})) \otimes_A B
\to \Gamma(V_{i_0 \ldots i_p}, \underline{H}^q(\mathcal{G}))$.
Le foncteur $- \otimes_A B$ commutant aux produits et étant exact, on en déduit
que
$\check{H}^p(\mathcal{U}, \underline{H}^q(\mathcal{F})) \otimes_A B
\to \check{H}^p(\mathcal{V}, \underline{H}^q(\mathcal{G}))$
est un isomorphisme. La platitude de $A \to B$ permet de conclure que
$H^i(X, \mathcal{F}) \otimes_A B \to H^i(Y, \mathcal{G})$
est un isomorphisme pour tout $i$, ce qui conclut.
\end{proof}









\section{Limites inductives et images directes supérieures}
\label{section-colimits}

\noindent
On trouvera des résultats généraux de cette nature dans
Cohomologie, section \ref{cohomology-section-limits},
Faisceaux, lemme \ref{sheaves-lemma-directed-colimits-sections}, et
Modules, lemme \ref{modules-lemma-finite-presentation-quasi-compact-colimit}.

\begin{lemma}
\label{lemma-colimit-cohomology}
Soit $f : X \to S$ un morphisme quasi-compact et quasi-séparé de schémas.
Soit $\mathcal{F} = \colim \mathcal{F}_i$ une limite inductive filtrante
de faisceaux abéliens sur $X$. Alors, pour tout $p \geq 0$, on a
$$
R^pf_*\mathcal{F} = \colim R^pf_*\mathcal{F}_i.
$$
\end{lemma}

\begin{proof}
Nous appliquons Cohomologie, Lemme
\ref{cohomology-lemma-higher-direct-image-colimit}.
Comme les ouverts affines forment une base de la topologie de $S$,
il suffit de montrer que, pour tout ouvert affine $U \subset S$, on a
$H^p(f^{-1}U, \mathcal{F}) = \colim H^p(f^{-1}U, \mathcal{F}_i)$.
Puisque $f^{-1}U$ est quasi-compact et quasi-séparé, cela résulte
de
Cohomologie, lemme \ref{cohomology-lemma-quasi-separated-cohomology-colimit} (la
base des ouverts affines de $f^{-1}U$
vérifie les hypothèses de ce lemme).
\end{proof}










\section{Cohomologie et changement de base, II}
\label{section-cohomology-and-base-change-derived}

\noindent
Soit $f : X \to S$ un morphisme de schémas et soit $\mathcal{F}$
un $\mathcal{O}_X$-module quasi-cohérent. Si $f$ est quasi-compact
et quasi-séparé, nous souhaitons représenter $Rf_*\mathcal{F}$
par un complexe de faisceaux quasi-cohérents sur $S$. Cela résulte
du fait que les faisceaux $R^if_*\mathcal{F}$ sont quasi-cohérents
lorsque $S$ est quasi-compact et à diagonale affine, en utilisant
l'équivalence de $D_\QCoh(S)$ et de
$D(\QCoh(\mathcal{O}_S))$ ; voir
Catégories dérivées des schémas, Proposition
\ref{perfect-proposition-quasi-compact-affine-diagonal}.

\medskip\noindent
Dans cette section, nous suivrons une autre approche, qui produit un complexe
explicite doté d'une bonne propriété de changement de base. La construction
est particulièrement simple si $f$ et $S$ sont séparés ou, plus généralement,
à diagonale affine. Comme il s'agit de loin
du cas le plus fréquemment utilisé, nous le traitons séparément.

\begin{lemma}
\label{lemma-separated-case-relative-cech}
Soit $f : X \to S$ un morphisme de schémas.
Soit $\mathcal{F}$ un $\mathcal{O}_X$-module quasi-cohérent.
Supposons que $X$ soit quasi-compact et que $X$ et $S$ aient une diagonale affine
(par exemple, si $X$ et $S$ sont séparés).
On peut alors calculer $Rf_*\mathcal{F}$ comme suit :
\begin{enumerate}
\item Choisir un recouvrement ouvert affine fini
$\mathcal{U} : X = \bigcup_{i = 1, \ldots, n} U_i$.
\item Pour $i_0, \ldots, i_p \in \{1, \ldots, n\}$, noter
$f_{i_0 \ldots i_p} : U_{i_0 \ldots i_p} \to S$ la restriction de $f$
à l'intersection $U_{i_0 \ldots i_p} = U_{i_0} \cap \ldots \cap U_{i_p}$.
\item Noter $\mathcal{F}_{i_0 \ldots i_p}$ la restriction
de $\mathcal{F}$ à $U_{i_0 \ldots i_p}$.
\item Poser
$$
\check{\mathcal{C}}^p(\mathcal{U}, f, \mathcal{F}) =
\bigoplus\nolimits_{i_0 \ldots i_p}
f_{i_0 \ldots i_p *} \mathcal{F}_{i_0 \ldots i_p}
$$
et définir les différentielles
$d : \check{\mathcal{C}}^p(\mathcal{U}, f, \mathcal{F})
\to \check{\mathcal{C}}^{p + 1}(\mathcal{U}, f, \mathcal{F})$
comme dans Cohomologie, Équation (\ref{cohomology-equation-d-cech}).
\end{enumerate}
Le complexe $\check{\mathcal{C}}^\bullet(\mathcal{U}, f, \mathcal{F})$
est alors un complexe de faisceaux quasi-cohérents sur $S$, muni d'un
isomorphisme
$$
\check{\mathcal{C}}^\bullet(\mathcal{U}, f, \mathcal{F})
\longrightarrow
Rf_*\mathcal{F}
$$
dans $D^{+}(S)$. Cet isomorphisme est fonctoriel en le faisceau
quasi-cohérent $\mathcal{F}$.
\end{lemma}

\begin{proof}
Considérons la résolution
$\mathcal{F} \to {\mathfrak C}^\bullet(\mathcal{U}, \mathcal{F})$
de Cohomologie, lemme \ref{cohomology-lemma-covering-resolution}.
On a une égalité de complexes
$\check{\mathcal{C}}^\bullet(\mathcal{U}, f, \mathcal{F}) =
f_*{\mathfrak C}^\bullet(\mathcal{U}, \mathcal{F})$
de $\mathcal{O}_S$-modules
quasi-cohérents. Les morphismes $j_{i_0 \ldots i_p} : U_{i_0 \ldots i_p} \to X$
et $f_{i_0 \ldots i_p} : U_{i_0 \ldots i_p} \to S$
sont affines d'après Morphismes, lemme \ref{morphisms-lemma-affine-permanence},
et le lemme \ref{lemma-affine-diagonal}. Par
conséquent, $R^qj_{i_0 \ldots i_p *}\mathcal{F}_{i_0 \ldots i_p}$
ainsi que $R^qf_{i_0 \ldots i_p *}\mathcal{F}_{i_0 \ldots i_p}$
sont nuls pour $q > 0$ (lemme \ref{lemma-relative-affine-vanishing}). En
utilisant $f \circ j_{i_0 \ldots i_p} = f_{i_0 \ldots i_p}$ et
la suite spectrale de
Cohomologie, lemme \ref{cohomology-lemma-relative-Leray},
on en déduit que
$R^qf_*(j_{i_0 \ldots i_p *}\mathcal{F}_{i_0 \ldots i_p}) = 0$
pour $q > 0$.
Comme les termes du complexe
${\mathfrak C}^\bullet(\mathcal{U}, \mathcal{F})$ sont des sommes
directes finies des faisceaux $j_{i_0 \ldots i_p *}\mathcal{F}_{i_0 \ldots i_p}$,
le lemme d'acyclicité de Leray
(Catégories dérivées, lemme \ref{derived-lemma-leray-acyclicity})
donne
$$
Rf_* \mathcal{F} = f_*{\mathfrak C}^\bullet(\mathcal{U}, \mathcal{F}) =
\check{\mathcal{C}}^\bullet(\mathcal{U}, f, \mathcal{F})
$$
comme souhaité.
\end{proof}

\noindent
Considérons maintenant ce qui se passe après un changement de base.

\begin{lemma}
\label{lemma-base-change-complex}
Avec les notations du diagramme (\ref{equation-base-change-diagram}), supposons
que $f : X \to S$ et $\mathcal{F}$ vérifient les hypothèses du
lemme \ref{lemma-separated-case-relative-cech}. Choisissons un recouvrement
ouvert affine fini $\mathcal{U} : X = \bigcup U_i$ de $X$.
Il existe un isomorphisme canonique
$$
g^*\check{\mathcal{C}}^\bullet(\mathcal{U}, f, \mathcal{F})
\longrightarrow
Rf'_*\mathcal{F}'
$$
dans $D^{+}(S')$. De plus, si $S' \to S$ est affine, on a en fait
$$
g^*\check{\mathcal{C}}^\bullet(\mathcal{U}, f, \mathcal{F})
=
\check{\mathcal{C}}^\bullet(\mathcal{U}', f', \mathcal{F}')
$$
avec $\mathcal{U}' : X' = \bigcup U_i'$, où
$U_i' = (g')^{-1}(U_i) = U_{i, S'}$ est également affine.
\end{lemma}

\begin{proof}
En fait, on peut définir $U_i' = (g')^{-1}(U_i) = U_{i, S'}$,
que $S'$ soit ou non affine sur $S$.
Soit $\mathcal{U}' : X' = \bigcup U_i'$ le recouvrement induit de $X'$.
Nous affirmons alors que
$$
g^*\check{\mathcal{C}}^\bullet(\mathcal{U}, f, \mathcal{F})
=
\check{\mathcal{C}}^\bullet(\mathcal{U}', f', \mathcal{F}')
$$
où $\check{\mathcal{C}}^\bullet(\mathcal{U}', f', \mathcal{F}')$
est défini exactement comme dans le
lemme \ref{lemma-separated-case-relative-cech}.
Cela résulte du cas des morphismes affines
(lemme \ref{lemma-affine-base-change}) en raisonnant localement sur $S'$.
De plus, exactement comme dans la démonstration du
lemme \ref{lemma-separated-case-relative-cech},
on voit qu'il existe un isomorphisme
$$
\check{\mathcal{C}}^\bullet(\mathcal{U}', f', \mathcal{F}')
\longrightarrow
Rf'_*\mathcal{F}'
$$
dans $D^{+}(S')$, car les morphismes $U_i' \to X'$ et $U_i' \to S'$
restent affines (ce sont des changements de base de morphismes affines). Nous
omettons les détails.
\end{proof}

\noindent
Le lemme précédent affirme que le complexe
$$
\mathcal{K}^\bullet = \check{\mathcal{C}}^\bullet(\mathcal{U}, f, \mathcal{F})
$$
est un complexe borné inférieurement de faisceaux quasi-cohérents sur $S$ qui calcule
{\it universellement} les images directes supérieures de $f : X \to S$.
Il s'agit d'une propriété de ce complexe particulier,
qui n'est pas préservée lorsqu'on remplace
$\check{\mathcal{C}}^\bullet(\mathcal{U}, f, \mathcal{F})$ par
un complexe quasi-isomorphe en général ! Autrement dit, cet énoncé
n'a pas de sens dans la catégorie dérivée. La
raison en est que l'image inverse $g^*\mathcal{K}^\bullet$ n'est
{\it pas} égale, en général, à l'image inverse dérivée $Lg^*\mathcal{K}^\bullet$
de $\mathcal{K}^\bullet$ !

\medskip\noindent
Voici un cas plus général où cet énoncé peut être démontré.
Remarquons que l'hypothèse de séparation de $S$ est inoffensive dans la
plupart des applications, car le résultat sert habituellement à établir
une propriété locale de l'image dérivée
totale. La démonstration est nettement plus délicate et
utilise les hyperrecouvrements ; c'est un bon exemple de leur emploi.

\begin{lemma}
\label{lemma-hypercoverings}
Soit $f : X \to S$ un morphisme de schémas.
Soit $\mathcal{F}$ un faisceau quasi-cohérent sur $X$.
Supposons que $f$ soit quasi-compact et quasi-séparé, et
que $S$ soit quasi-compact et
séparé. Il existe un complexe $\mathcal{K}^\bullet$,
borné inférieurement, de $\mathcal{O}_S$-modules quasi-cohérents, possédant la
propriété suivante : pour tout morphisme
$g : S' \to S$, le complexe $g^*\mathcal{K}^\bullet$
représente $Rf'_*\mathcal{F}'$, avec les notations du
diagramme (\ref{equation-base-change-diagram}).
\end{lemma}

\begin{proof}
(Si $f$ est lui aussi séparé, voir le
lemme \ref{lemma-base-change-complex}.)
Les hypothèses impliquent notamment que $X$
est quasi-compact et quasi-séparé en tant que schéma.
Soit $\mathcal{B}$ l'ensemble des ouverts affines de $X$. D'après
Hyperrecouvrements,
lemme \ref{hypercovering-lemma-quasi-separated-quasi-compact-hypercovering},
il existe un hyperrecouvrement $K = (I, \{U_i\})$ tel que chaque
$I_n$ soit fini et chaque $U_i$ un ouvert affine de $X$. D'après
Hyperrecouvrements, lemme \ref{hypercovering-lemma-cech-spectral-sequence},
il existe une suite spectrale dont la page $E_2$ est
$$
E_2^{p, q} = \check{H}^p(K, \underline{H}^q(\mathcal{F}))
$$
et qui converge vers $H^{p + q}(X, \mathcal{F})$. Remarquons que
$\check{H}^p(K, \underline{H}^q(\mathcal{F}))$ est le $p$-ième groupe
de cohomologie du complexe
$$
\prod\nolimits_{i \in I_0} H^q(U_i, \mathcal{F})
\to
\prod\nolimits_{i \in I_1} H^q(U_i, \mathcal{F})
\to
\prod\nolimits_{i \in I_2} H^q(U_i, \mathcal{F})
\to \ldots
$$
Chaque $U_i$ étant affine, ce groupe est nul sauf si $q = 0$,
auquel cas on obtient
$$
\prod\nolimits_{i \in I_0} \mathcal{F}(U_i)
\to
\prod\nolimits_{i \in I_1} \mathcal{F}(U_i)
\to
\prod\nolimits_{i \in I_2} \mathcal{F}(U_i)
\to \ldots
$$
On en conclut que $R\Gamma(X, \mathcal{F})$ est calculé par
ce complexe.

\medskip\noindent
Pour tout $n$ et tout $i \in I_n$, notons $f_i : U_i \to S$ la restriction de
$f$ à $U_i$. Puisque $S$ est séparé et $U_i$ affine, ce morphisme est
affine. Considérons le complexe de faisceaux quasi-cohérents
$$
\mathcal{K}^\bullet = (
\prod\nolimits_{i \in I_0} f_{i, *}\mathcal{F}|_{U_i}
\to
\prod\nolimits_{i \in I_1} f_{i, *}\mathcal{F}|_{U_i}
\to
\prod\nolimits_{i \in I_2} f_{i, *}\mathcal{F}|_{U_i}
\to \ldots )
$$
sur $S$. Comme dans
Hyperrecouvrements, lemme \ref{hypercovering-lemma-cech-spectral-sequence},
on obtient un morphisme $\mathcal{K}^\bullet \to Rf_*\mathcal{F}$ dans
$D(\mathcal{O}_S)$ en choisissant une résolution injective de $\mathcal{F}$
(nous omettons les détails). Considérons un schéma affine quelconque $V$ et un morphisme
$g : V \to S$. Le changement de base $X_V$ possède alors l'hyperrecouvrement
$K_V = (I, \{U_{i, V}\})$ obtenu par changement de base. De plus,
$g^*f_{i, *}\mathcal{F} = f_{i, V, *}(g')^*\mathcal{F}|_{U_{i, V}}$.
Les arguments précédents montrent donc que $\Gamma(V, g^*\mathcal{K}^\bullet)$
calcule $R\Gamma(X_V, (g')^*\mathcal{F})$.
Cela achève la démonstration du lemme, puisqu'il suffit d'établir l'égalité des complexes
localement pour la topologie de Zariski sur $S'$.
\end{proof}

\noindent
Le lemme suivant est une variante du changement de base plat.

\begin{lemma}
\label{lemma-base-change-tensor-with-flat}
Considérons un diagramme cartésien de schémas
$$
\xymatrix{
X' \ar[d]_{f'} \ar[r]_{g'} & X \ar[d]^f \\
S' \ar[r]^g & S
}
$$
Soit $\mathcal{F}$ un $\mathcal{O}_X$-module quasi-cohérent.
Soit $\mathcal{G}$ un $\mathcal{O}_{S'}$-module quasi-cohérent
et plat sur $S$. Supposons que $f$ soit quasi-compact et quasi-séparé.
Pour tout $i \geq 0$, il existe une identification
$$
\mathcal{G} \otimes_{\mathcal{O}_{S'}} g^*R^if_*\mathcal{F} =
R^if'_*\left((f')^*\mathcal{G}
\otimes_{\mathcal{O}_{X'}} (g')^*\mathcal{F}\right)
$$
\end{lemma}

\begin{proof}
Construisons un morphisme du membre de gauche vers celui de droite. On dispose d'abord du morphisme de changement de base
$Lg^*Rf_*\mathcal{F} \to Rf'_*L(g')^*\mathcal{F}$. On dispose aussi du morphisme
d'adjonction $\mathcal{G} \to Rf'_*L(f')^*\mathcal{G}$. À l'aide du cup-produit relatif,
on obtient
\begin{align*}
\mathcal{G}
\otimes_{\mathcal{O}_{S'}}^\mathbf{L}
Lg^*Rf_*\mathcal{F}
& \to
Rf'_*L(f')^*\mathcal{G}
\otimes_{\mathcal{O}_{S'}}^\mathbf{L}
Rf'_*L(g')^*\mathcal{F} \\
& \to
Rf'_*\left(L(f')^*\mathcal{G}
\otimes_{\mathcal{O}_{X'}}^\mathbf{L}
L(g')^*\mathcal{F}\right) \\
& \to
Rf'_*\left((f')^*\mathcal{G}
\otimes_{\mathcal{O}_{X'}} (g')^*\mathcal{F}\right)
\end{align*}
où la flèche centrale est donnée par le cup-produit relatif ; voir
Cohomologie, remarque \ref{cohomology-remark-cup-product}.
La source du composé est
$$
\mathcal{G} \otimes_{\mathcal{O}_{S'}}^\mathbf{L} Lg^*Rf_*\mathcal{F} =
\mathcal{G} \otimes_{g^{-1}\mathcal{O}_S}^\mathbf{L} g^{-1}Rf_*\mathcal{F}
$$
d'après Cohomologie, lemme \ref{cohomology-lemma-variant-derived-pullback}.
Comme $\mathcal{G}$ est plat en tant que faisceau de $g^{-1}\mathcal{O}_S$-modules
et que $g^{-1}$ est un foncteur exact, il s'agit d'un complexe dont le
$i$-ième faisceau de cohomologie est
$\mathcal{G} \otimes_{g^{-1}\mathcal{O}_S} g^{-1}R^if_*\mathcal{F} =
\mathcal{G} \otimes_{\mathcal{O}_{S'}} g^*R^if_*\mathcal{F}$.
On obtient ainsi, globalement, les morphismes du membre de gauche vers celui de droite
dans l'égalité du lemme. Pour montrer que ce morphisme est un isomorphisme, on peut
raisonner localement sur $S'$. Nous pouvons donc supposer, et supposons, que $S$ et $S'$
sont des schémas affines.

\medskip\noindent
Démonstration dans le cas où $S$ et $S'$ sont affines. Écrivons $S = \Spec(A)$ et $S' = \Spec(B)$,
et supposons que $\mathcal{G}$ corresponde au $B$-module $N$, supposé plat sur
$A$. Puisque $S$ est affine, $X$ est quasi-compact et quasi-séparé.
Nous achevons la démonstration à l'aide d'un hyperrecouvrement ; si $X$ est séparé
ou à diagonale affine, on peut utiliser un recouvrement de {\v C}ech.
Soit $\mathcal{B}$ l'ensemble des ouverts affines de $X$. D'après
Hyperrecouvrements,
lemme \ref{hypercovering-lemma-quasi-separated-quasi-compact-hypercovering},
il existe un hyperrecouvrement $K = (I, \{U_i\})$ de $X$ tel que chaque
$I_n$ soit fini et chaque $U_i$ un ouvert affine de $X$. D'après
Hyperrecouvrements, lemme \ref{hypercovering-lemma-cech-spectral-sequence},
il existe une suite spectrale dont la page $E_2$ est
$$
E_2^{p, q} = \check{H}^p(K, \underline{H}^q(\mathcal{F}))
$$
et qui converge vers $H^{p + q}(X, \mathcal{F})$. Comme chaque $U_i$ est affine et
que $\mathcal{F}$ est quasi-cohérent, la valeur de
$\underline{H}^q(\mathcal{F})$ sur $U_i$ est nulle pour $q > 0$. La suite
spectrale dégénère donc, et on en conclut que les modules de
cohomologie $H^q(X, \mathcal{F})$ sont calculés par
$$
\prod\nolimits_{i \in I_0} \mathcal{F}(U_i)
\to
\prod\nolimits_{i \in I_1} \mathcal{F}(U_i)
\to
\prod\nolimits_{i \in I_2} \mathcal{F}(U_i)
\to \ldots
$$
Remarquons ensuite que le changement de base de notre hyperrecouvrement à
$S'$ est un hyperrecouvrement de $X' = S' \times_S X$.
Les schémas $S' \times_S U_i$ sont eux aussi affines,
et l'on a
$$
\left((f')^*\mathcal{G} \otimes_{\mathcal{O}_{S'}} (g')^*\mathcal{F}\right)
(S' \times_S U_i) =
N \otimes_A \mathcal{F}(U_i)
$$
On en conclut que les modules de cohomologie
$H^q(X', (f')^*\mathcal{G} \otimes_{\mathcal{O}_{S'}} (g')^*\mathcal{F})$
sont calculés par
$$
N \otimes_A \left(
\prod\nolimits_{i \in I_0} \mathcal{F}(U_i)
\to
\prod\nolimits_{i \in I_1} \mathcal{F}(U_i)
\to
\prod\nolimits_{i \in I_2} \mathcal{F}(U_i)
\to \ldots
\right)
$$
Puisque $N$ est plat sur $A$, on en déduit que
$$
H^q(X', (f')^*\mathcal{G} \otimes_{\mathcal{O}_{S'}} (g')^*\mathcal{F})
=
N \otimes_A H^q(X, \mathcal{F})
$$
Il s'agit précisément de la traduction algébrique de l'énoncé
à démontrer, ce qui achève la démonstration.
\end{proof}







\section{Cohomologie de l'espace projectif}
\label{section-cohomology-projective-space}

\noindent
Dans cette section, nous calculons la cohomologie des faisceaux structuraux
tordus sur $\mathbf{P}^n_S$, où $S$ est un
schéma. Rappelons que $\mathbf{P}^n_S$ a été défini comme le produit fibré
$
\mathbf{P}^n_S = S \times_{\Spec(\mathbf{Z})} \mathbf{P}^n_{\mathbf{Z}}
$
dans Constructions, définition \ref{constructions-definition-projective-space}.
On a montré l'égalité
$$
\mathbf{P}^n_S = \underline{\text{Proj}}_S(\mathcal{O}_S[T_0, \ldots, T_n])
$$
dans Constructions, lemme \ref{constructions-lemma-projective-space-bundle}.
En particulier, l'espace projectif est un cas particulier de fibré projectif.
Si $S = \Spec(R)$ est affine, on a
$$
\mathbf{P}^n_S = \mathbf{P}^n_R = \text{Proj}(R[T_0, \ldots, T_n]).
$$
Toutes ces identifications sont compatibles entre elles ainsi qu'avec les
constructions des faisceaux structuraux tordus $\mathcal{O}_{\mathbf{P}^n_S}(d)$.

\medskip\noindent
Avant d'énoncer le résultat, introduisons quelques notations.
Soit $R$ un anneau.
Rappelons que $R[T_0, \ldots, T_n]$ est une
$R$-algèbre graduée dans laquelle chaque $T_i$ est homogène de degré $1$.
Notons $(R[T_0, \ldots, T_n])_d$ la composante homogène de degré $d$.
C'est un $R$-module libre de type fini de rang $\binom{n + d}{d}$
si $d \geq 0$, et nul
sinon. Il possède une base formée des monômes $T_0^{e_0} \ldots T_n^{e_n}$
tels que $\sum e_i = d$. Nous utiliserons aussi la notation suivante :
$R[\frac{1}{T_0}, \ldots, \frac{1}{T_n}]$ désigne l'anneau $\mathbf{Z}$-gradué
dans lequel $\frac{1}{T_i}$ est de degré $-1$. En particulier, le module
$\mathbf{Z}$-gradué sur $R[\frac{1}{T_0}, \ldots, \frac{1}{T_n}]$
$$
\frac{1}{T_0 \ldots T_n} R[\frac{1}{T_0}, \ldots, \frac{1}{T_n}]
$$
qui figure dans l'énoncé ci-dessous est nul en degrés
$\geq -n$, est libre de générateur $\frac{1}{T_0 \ldots T_n}$
en degré $-n - 1$, et est libre de rang $(-1)^n\binom{n + d}{d}$ pour
$d \leq -n - 1$.

\begin{lemma}
\label{lemma-cohomology-projective-space-over-ring}
\begin{reference}
\cite[III Proposition 2.1.12]{EGA}
\end{reference}
Soit $R$ un anneau.
Soit $n \geq 0$ un entier.
On a
$$
H^q(\mathbf{P}^n, \mathcal{O}_{\mathbf{P}^n_R}(d)) =
\left\{
\begin{matrix}
(R[T_0, \ldots, T_n])_d & \text{si} & q = 0,\ d \geq 0 \\
\left(\frac{1}{T_0 \ldots T_n} R[\frac{1}{T_0}, \ldots, \frac{1}{T_n}]\right)_d
& \text{si} & q = n,\ d < 0 \\
0 & \text{si} & q, n, d\text{ hors des cas ci-dessus}
\end{matrix}
\right.
$$
comme $R$-modules.
\end{lemma}

\begin{proof}
Utilisons le recouvrement ouvert affine standard
$$
\mathcal{U} : \mathbf{P}^n_R = \bigcup\nolimits_{i = 0}^n D_{+}(T_i)
$$
pour calculer la cohomologie au moyen du complexe de {\v C}ech.
Le lemme \ref{lemma-cech-cohomology-quasi-coherent}
le permet, puisque toute intersection d'un nombre fini d'ouverts affines $D_{+}(T_i)$ est encore un
ouvert affine standard (voir
Constructions, section \ref{constructions-section-proj}).
On peut même utiliser le complexe de {\v C}ech alterné ou ordonné, d'après
Cohomologie, lemmes \ref{cohomology-lemma-ordered-alternating} et
\ref{cohomology-lemma-alternating-usual}.

\medskip\noindent
Nous munissons $\{0, \ldots, n\}$ de son ordre
usuel. Le complexe considéré a donc pour termes
$$
\check{\mathcal{C}}_{ord}^p(\mathcal{U}, \mathcal{O}_{\mathbf{P}_R}(d))
=
\bigoplus\nolimits_{i_0 < \ldots < i_p}
(R[T_0, \ldots, T_n, \frac{1}{T_{i_0} \ldots T_{i_p}}])_d
$$
De plus, les différentielles sont données par la formule usuelle
$$
d(s)_{i_0 \ldots i_{p + 1}} =
\sum\nolimits_{j = 0}^{p + 1} (-1)^j s_{i_0 \ldots \hat i_j \ldots i_{p + 1}}
$$
voir Cohomologie, section \ref{cohomology-section-alternating-cech}.
Notons que chaque terme de ce complexe possède une
$\mathbf{Z}^{n + 1}$-graduation naturelle : on déclare que le monôme
$T_0^{e_0} \ldots T_n^{e_n}$ est homogène de poids
$(e_0, \ldots, e_n) \in \mathbf{Z}^{n + 1}$. Il est clair que la
différentielle ci-dessus respecte cette graduation. On a donc
$$
\check{\mathcal{C}}_{ord}^\bullet(\mathcal{U}, \mathcal{O}_{\mathbf{P}_R}(d))
=
\bigoplus\nolimits_{\vec{e} \in \mathbf{Z}^{n + 1}}
\check{\mathcal{C}}^\bullet(\vec{e})
$$
où tous les facteurs directs de droite n'interviennent pas nécessairement (voir
ci-dessous). Pour calculer les modules de
cohomologie du complexe, il suffit donc de calculer la cohomologie de ses
composantes homogènes, puis d'en prendre la somme directe.

\medskip\noindent
Fixons $\vec{e} = (e_0, \ldots, e_n) \in \mathbf{Z}^{n + 1}$. Pour que ce
poids intervienne dans le complexe ci-dessus, il faut supposer
$e_0 + \ldots + e_n = d$ (sinon, il intervient bien entendu pour un autre tordu
du faisceau structural). Sous cette hypothèse, posons
$$
NEG(\vec{e}) = \{i \in \{0, \ldots, n\} \mid e_i < 0\}.
$$
Avec cette notation, le facteur direct de poids $\vec{e}$
$\check{\mathcal{C}}^\bullet(\vec{e})$ du complexe de {\v C}ech ci-dessus a
pour termes
$$
\check{\mathcal{C}}^p(\vec{e})
=
\bigoplus\nolimits_{i_0 < \ldots < i_p,
\ NEG(\vec{e}) \subset \{i_0, \ldots, i_p\}}
R \cdot T_0^{e_0} \ldots T_n^{e_n}
$$
Autrement dit, les termes correspondant aux $i_0 < \ldots < i_p$ tels
que $NEG(\vec{e})$ ne soit pas contenu dans $\{i_0 \ldots i_p\}$ sont nuls.
La différentielle du complexe $\check{\mathcal{C}}^\bullet(\vec{e})$
est encore donnée par la même formule que ci-dessus.

\medskip\noindent
Supposons que $NEG(\vec{e}) = \{0, \ldots, n\}$, c'est-à-dire que tous les
exposants $e_i$ soient négatifs.
Dans ce cas, le complexe $\check{\mathcal{C}}^\bullet(\vec{e})$ ne
possède qu'un terme, à savoir $\check{\mathcal{C}}^n(\vec{e}) =
R \cdot \frac{1}{T_0^{-e_0} \ldots T_n^{-e_n}}$. On a donc, dans ce
cas,
$$
H^q(\check{\mathcal{C}}^\bullet(\vec{e})) =
\left\{
\begin{matrix}
R \cdot \frac{1}{T_0^{-e_0} \ldots T_n^{-e_n}} & \text{si} & q = n \\
0 & \text{si} & \text{sinon}
\end{matrix}
\right.
$$
La somme directe de tous ces termes donne manifestement
$$
\left(\frac{1}{T_0 \ldots T_n} R[\frac{1}{T_0}, \ldots, \frac{1}{T_n}]\right)_d
$$
en degré $n$, comme l'indique le lemme. De plus, ces termes ne contribuent
pas à la cohomologie en d'autres degrés (conformément, là encore, à l'énoncé du
lemme).

\medskip\noindent
Supposons $NEG(\vec{e}) = \emptyset$. Dans ce cas, le complexe
$\check{\mathcal{C}}^\bullet(\vec{e})$ possède un facteur direct $R$ correspondant
à chaque $i_0 < \ldots < i_p$.
Comparons le complexe $\check{\mathcal{C}}^\bullet(\vec{e})$
à un autre complexe. Considérons le recouvrement ouvert affine
$$
\mathcal{V} : \Spec(R) = \bigcup\nolimits_{i \in \{0, \ldots, n\}} V_i
$$
où $V_i = \Spec(R)$ pour tout $i$. Considérons le complexe de
{\v C}ech alterné
$$
\check{\mathcal{C}}_{ord}^\bullet(\mathcal{V}, \mathcal{O}_{\Spec(R)})
$$
Le même raisonnement que ci-dessus montre qu'il calcule la cohomologie du
faisceau structural sur $\Spec(R)$. Ainsi,
$H^p(
\check{\mathcal{C}}_{ord}^\bullet(\mathcal{V}, \mathcal{O}_{\Spec(R)})
) = R$ si $p = 0$, et vaut $0$ lorsque $p > 0$.
Pour ces faits, voir le
lemme \ref{lemma-cech-cohomology-quasi-coherent-trivial} et sa
démonstration. Notons que
$\check{\mathcal{C}}_{ord}^\bullet(\mathcal{V}, \mathcal{O}_{\Spec(R)})$
possède lui aussi un facteur direct $R$ pour chaque $i_0 < \ldots < i_p$, et a exactement la même
différentielle que $\check{\mathcal{C}}^\bullet(\vec{e})$. Autrement dit, ces
complexes sont isomorphes et ont donc la même cohomologie. On
en déduit que
$$
H^q(\check{\mathcal{C}}^\bullet(\vec{e})) =
\left\{
\begin{matrix}
R \cdot T_0^{e_0} \ldots T_n^{e_n} & \text{si} & q = 0 \\
0 & \text{si} & \text{sinon}
\end{matrix}
\right.
$$
lorsque $NEG(\vec{e}) = \emptyset$.
La somme directe de tous ces termes donne manifestement
$$
(R[T_0, \ldots, T_n])_d
$$
en degré $0$, comme l'indique le lemme. De plus, ces termes ne contribuent
pas à la cohomologie en d'autres degrés (conformément, là encore, à l'énoncé du
lemme).

\medskip\noindent
Pour achever la démonstration du lemme, il reste à montrer que les complexes
$\check{\mathcal{C}}^\bullet(\vec{e})$ sont acycliques lorsque
$NEG(\vec{e})$ n'est ni vide ni égal à $\{0, \ldots, n\}$.
Choisissons un indice $i_{\text{fix}} \not \in NEG(\vec{e})$ (il en existe un).
Considérons le morphisme
$$
h :
\check{\mathcal{C}}^{p + 1}(\vec{e})
\to
\check{\mathcal{C}}^p(\vec{e})
$$
défini, pour $i_0 < \ldots < i_p$, par
$$
h(s)_{i_0 \ldots i_p} =
\left\{
\begin{matrix}
0 & \text{si} & p \not \in \{0, \ldots, n - 1\} \\
0 & \text{si} & i_{\text{fix}} \in \{i_0, \ldots, i_p\} \\
s_{i_{\text{fix}} i_0 \ldots i_p} & \text{si} & i_{\text{fix}} < i_0 \\
(-1)^a s_{i_0 \ldots i_{a - 1} i_{\text{fix}} i_a \ldots i_p} &
\text{si} & i_{a - 1} < i_{\text{fix}} < i_a \\
(-1)^{p + 1} s_{i_0 \ldots i_p i_{\text{fix}}} &
\text{si} & i_p < i_{\text{fix}}
\end{matrix}
\right.
$$
On comparera avec la démonstration du
lemme \ref{lemma-cech-cohomology-quasi-coherent-trivial}.
Cette définition a un sens car
$$
NEG(\vec{e}) \subset \{i_0, \ldots, i_p\}
\Leftrightarrow
NEG(\vec{e}) \subset \{i_{\text{fix}}, i_0, \ldots, i_p\}
$$
Le même calcul combinatoire\footnote{
Par exemple, supposons que $i_0 < \ldots < i_p$ soit tel que
$i_{\text{fix}} \not \in \{i_0, \ldots, i_p\}$ et que
$i_{a - 1} < i_{\text{fix}} < i_a$ pour un certain $1 \leq a \leq p$. Alors
\begin{align*}
&
(dh + hd)(s)_{i_0 \ldots i_p} \\
& =
\sum\nolimits_{j = 0}^p
(-1)^j
h(s)_{i_0 \ldots \hat i_j \ldots i_p}
+
(-1)^a d(s)_{i_0 \ldots i_{a - 1} i_{\text{fix}} i_a \ldots i_p}\\
& =
\sum\nolimits_{j = 0}^{a - 1}
(-1)^{j + a - 1}
s_{i_0 \ldots \hat i_j \ldots i_{a - 1} i_{\text{fix}} i_a \ldots i_p} +
\sum\nolimits_{j = a}^p
(-1)^{j + a}
s_{i_0 \ldots i_{a - 1} i_{\text{fix}} i_a \ldots \hat i_j \ldots i_p}
+ \\
&
\sum\nolimits_{j = 0}^{a - 1}
(-1)^{a + j}
s_{i_0 \ldots \hat i_j \ldots i_{a - 1} i_{\text{fix}} i_a \ldots i_p} +
(-1)^{2a} s_{i_0 \ldots i_p} +
\sum\nolimits_{j = a}^p
(-1)^{a + j + 1}
s_{i_0 \ldots i_{a - 1} i_{\text{fix}} i_a \ldots \hat i_j \ldots i_p}
\\
& =
s_{i_0 \ldots i_p}
\end{align*}
comme voulu. Les autres cas sont analogues.}
que dans la démonstration du lemme \ref{lemma-cech-cohomology-quasi-coherent-trivial}
montre que
$$
(hd + dh)(s)_{i_0 \ldots i_p}
=
s_{i_0 \ldots i_p}
$$
On voit donc que le morphisme identité du complexe
$\check{\mathcal{C}}^\bullet(\vec{e})$ est homotope à zéro,
ce qui entraîne son acyclicité.
\end{proof}

\noindent
Dans le lemme suivant, nous utiliserons l'accouplement de
$R$-modules libres
$$
R[T_0, \ldots, T_n]
\times
\frac{1}{T_0 \ldots T_n} R[\frac{1}{T_0}, \ldots, \frac{1}{T_n}]
\longrightarrow
R
$$
défini par la règle
$$
(f, g)
\longmapsto
\text{coefficient de }
\frac{1}{T_0 \ldots T_n}
\text{ dans }fg.
$$
Autrement dit, l'accouplement de l'élément de base $T_0^{e_0} \ldots T_n^{e_n}$ avec
l'élément de base $T_0^{d_0} \ldots T_n^{d_n}$ vaut $1$ si et seulement
si $e_i + d_i = -1$ pour tout $i$, et vaut zéro dans tous les autres
cas. Cet accouplement donne une identification
$$
\left(\frac{1}{T_0 \ldots T_n} R[\frac{1}{T_0}, \ldots, \frac{1}{T_n}]\right)_d
=
\Hom_R((R[T_0, \ldots, T_n])_{-n - 1 - d}, R)
$$
On peut donc reformuler le résultat du
lemme \ref{lemma-cohomology-projective-space-over-ring} sous la forme
\begin{equation}
\label{equation-identify}
H^q(\mathbf{P}^n, \mathcal{O}_{\mathbf{P}^n_R}(d)) =
\left\{
\begin{matrix}
(R[T_0, \ldots, T_n])_d & \text{si} & q = 0 \\
0 & \text{si} & q \not = 0, n \\
\Hom_R((R[T_0, \ldots, T_n])_{-n - 1 - d}, R)
& \text{si} & q = n
\end{matrix}
\right.
\end{equation}

\begin{lemma}
\label{lemma-identify-functorially}
Les identifications de l'Équation (\ref{equation-identify}) sont compatibles
au changement de base par les homomorphismes\ d'anneaux $R \to R'$.
De plus, pour tout $f \in R[T_0, \ldots, T_n]$ homogène
de degré $m$, le morphisme de multiplication par $f$
$$
\mathcal{O}_{\mathbf{P}^n_R}(d)
\longrightarrow
\mathcal{O}_{\mathbf{P}^n_R}(d + m)
$$
induit, via les identifications de
l'Équation (\ref{equation-identify}), le morphisme sur la cohomologie donné par la multiplication par
$f$ sur $H^0$ et par la transposée de la multiplication par $f$
$$
(R[T_0, \ldots, T_n])_{-n - 1 - (d + m)}
\longrightarrow
(R[T_0, \ldots, T_n])_{-n - 1 - d}
$$
sur $H^n$.
\end{lemma}

\begin{proof}
Soit $R \to R'$ un homomorphisme d'anneaux.
Soit $\mathcal{U}$ le recouvrement ouvert affine standard de $\mathbf{P}^n_R$,
et soit $\mathcal{U}'$ le recouvrement ouvert affine standard de
$\mathbf{P}^n_{R'}$. Notons que $\mathcal{U}'$ est l'image inverse du recouvrement
$\mathcal{U}$ par le morphisme canonique
$\mathbf{P}^n_{R'} \to \mathbf{P}^n_R$. Il existe donc
un morphisme de complexes de {\v C}ech
$$
\gamma :
\check{\mathcal{C}}_{ord}^\bullet(\mathcal{U},
\mathcal{O}_{\mathbf{P}_R}(d))
\longrightarrow
\check{\mathcal{C}}_{ord}^\bullet(\mathcal{U}',
\mathcal{O}_{\mathbf{P}_{R'}}(d))
$$
compatible avec le morphisme induit sur la cohomologie, d'après
Cohomologie, lemme \ref{cohomology-lemma-functoriality-cech}.
Les calculs de la démonstration du
lemme \ref{lemma-cohomology-projective-space-over-ring}
montrent que ce morphisme de complexes de {\v C}ech est compatible avec les identifications des groupes
de cohomologie en question. (En effet, les éléments de base du complexe
de {\v C}ech sur $R$ sont simplement envoyés sur les éléments de base correspondants
du complexe de {\v C}ech sur $R'$.) D'où la première assertion du lemme.

\medskip\noindent
Fixons maintenant l'anneau $R$ et considérons deux polynômes homogènes
$f, g \in R[T_0, \ldots, T_n]$ de même degré $m$.
La cohomologie étant un foncteur additif, il est clair que le morphisme
induit par la multiplication par $f + g$ est la somme des morphismes
induits par la multiplication par $f$ et
par $g$. De plus, la cohomologie étant un foncteur,
un résultat analogue vaut pour la multiplication par un produit $fg$, où
$f, g$ sont homogènes (mais pas nécessairement de même degré). Pour
vérifier la deuxième assertion du lemme, il suffit donc de
la prouver lorsque $f = x \in R$ ou lorsque $f = T_i$.
Dans le cas de la multiplication par un élément $x \in R$, le
résultat découle du fait que tous les groupes ou complexes de cohomologie considérés
possèdent une structure de $R$-module ou de complexe de $R$-modules.
Enfin, considérons la multiplication par $T_i$
comme morphisme $\mathcal{O}_{\mathbf{P}^n_R}$-linéaire
$$
\mathcal{O}_{\mathbf{P}^n_R}(d)
\longrightarrow
\mathcal{O}_{\mathbf{P}^n_R}(d + 1)
$$
L'assertion concernant $H^0$ est claire. Pour celle concernant $H^n$,
considérons la multiplication par $T_i$ comme morphisme de complexes de {\v C}ech
$$
\check{\mathcal{C}}_{ord}^\bullet(\mathcal{U},
\mathcal{O}_{\mathbf{P}_R}(d))
\longrightarrow
\check{\mathcal{C}}_{ord}^\bullet(\mathcal{U},
\mathcal{O}_{\mathbf{P}_R}(d + 1))
$$
Utilisons les notations introduites dans la démonstration du
lemme \ref{lemma-cohomology-projective-space-over-ring}.
Considérons l'effet de la multiplication par $T_i$
sur les décompositions
$$
\check{\mathcal{C}}_{ord}^\bullet(\mathcal{U}, \mathcal{O}_{\mathbf{P}_R}(d))
=
\bigoplus\nolimits_{\vec{e} \in \mathbf{Z}^{n + 1}, \ \sum e_i = d}
\check{\mathcal{C}}^\bullet(\vec{e})
$$
et
$$
\check{\mathcal{C}}_{ord}^\bullet(\mathcal{U},
\mathcal{O}_{\mathbf{P}_R}(d + 1))
=
\bigoplus\nolimits_{\vec{e} \in \mathbf{Z}^{n + 1}, \ \sum e_i = d + 1}
\check{\mathcal{C}}^\bullet(\vec{e})
$$
Il est clair qu'elle envoie le sous-complexe
$\check{\mathcal{C}}^\bullet(\vec{e})$ dans le sous-complexe
$\check{\mathcal{C}}^\bullet(\vec{e} + \vec{b}_i)$, où
$\vec{b}_i = (0, \ldots, 0, 1, 0, \ldots, 0))$ est le $i$-ième vecteur de
base. Autrement dit, elle envoie le facteur direct de $H^n$ correspondant à
$\vec{e}$, avec $e_i < 0$ et $\sum e_i = d$,
dans le facteur direct de $H^n$ correspondant à
$\vec{e} + \vec{b}_i$ (qui est nul si $e_i + b_i \geq 0$).
On voit aisément que cela correspond exactement à l'action de la
transposée de la multiplication par $T_i$, considérée comme morphisme
$$
(R[T_0, \ldots, T_n])_{-n - 1 - (d + 1)}
\longrightarrow
(R[T_0, \ldots, T_n])_{-n - 1 - d}
$$
Cela prouve le lemme.
\end{proof}

\noindent
Avant d'énoncer la version relative, introduisons quelques notations.
Rappelons que $\mathcal{O}_S[T_0, \ldots, T_n]$ est un
$\mathcal{O}_S$-module gradué dans lequel chaque $T_i$ est homogène de degré $1$.
Notons $(\mathcal{O}_S[T_0, \ldots, T_n])_d$ la composante homogène de degré $d$.
C'est un faisceau localement libre de type fini de rang $\binom{n + d}{d}$ sur $S$.

\begin{lemma}
\label{lemma-cohomology-projective-space-over-base}
Soit $S$ un schéma.
Soit $n \geq 0$ un entier. Considérons
le morphisme structural
$$
f : \mathbf{P}^n_S \longrightarrow S.
$$
On a
$$
R^qf_*(\mathcal{O}_{\mathbf{P}^n_S}(d)) =
\left\{
\begin{matrix}
(\mathcal{O}_S[T_0, \ldots, T_n])_d & \text{si} & q = 0 \\
0 & \text{si} & q \not = 0, n \\
\SheafHom_{\mathcal{O}_S}(
(\mathcal{O}_S[T_0, \ldots, T_n])_{- n - 1 - d}, \mathcal{O}_S)
& \text{si} & q = n
\end{matrix}
\right.
$$
\end{lemma}

\begin{proof}
Démonstration omise. Indication : cela résulte de la compatibilité des identifications
(\ref{equation-identify}) au changement de base affine, donnée par
le lemme \ref{lemma-identify-functorially}.
\end{proof}

\noindent
Énonçons maintenant la version pour les fibrés projectifs associés aux faisceaux localement libres de
type fini. Soit $S$ un schéma. Soit $\mathcal{E}$
un $\mathcal{O}_S$-module localement libre de type fini de rang constant $n + 1$, voir
Modules, section \ref{modules-section-locally-free}.
On considère alors $\text{Sym}(\mathcal{E})$ comme un
$\mathcal{O}_S$-module gradué dont $\mathcal{E}$ est la composante homogène de degré $1$.
Et $\text{Sym}^d(\mathcal{E})$ est la composante homogène de degré $d$.
C'est un faisceau localement libre de type fini de rang $\binom{n + d}{d}$ sur $S$.
Rappelons notre convention :
$$
\pi :
\mathbf{P}(\mathcal{E})
=
\underline{\text{Proj}}_S(\text{Sym}(\mathcal{E}))
\longrightarrow
S
$$
et l'existence de morphismes naturels
$\text{Sym}^d(\mathcal{E}) \to \pi_*\mathcal{O}_{\mathbf{P}(\mathcal{E})}(d)$.

\begin{lemma}
\label{lemma-cohomology-projective-bundle}
Soit $S$ un schéma. Soit $n \geq 1$.
Soit $\mathcal{E}$
un $\mathcal{O}_S$-module localement libre de type fini de rang constant $n + 1$.
Considérons le morphisme structural
$$
\pi : \mathbf{P}(\mathcal{E}) \longrightarrow S.
$$
On a
$$
R^q\pi_*(\mathcal{O}_{\mathbf{P}(\mathcal{E})}(d)) =
\left\{
\begin{matrix}
\text{Sym}^d(\mathcal{E}) & \text{si} & q = 0 \\
0 & \text{si} & q \not = 0, n \\
\SheafHom_{\mathcal{O}_S}(
\text{Sym}^{- n - 1 - d}(\mathcal{E})
\otimes_{\mathcal{O}_S}
\wedge^{n + 1}\mathcal{E},
\mathcal{O}_S)
& \text{si} & q = n
\end{matrix}
\right.
$$
Ces identifications sont compatibles au changement de base et
aux isomorphismes de faisceaux localement libres.
\end{lemma}

\begin{proof}
Considérons le morphisme canonique
$$
\pi^*\mathcal{E} \longrightarrow \mathcal{O}_{\mathbf{P}(\mathcal{E})}(1)
$$
et prenons son tordu de degré opposé à $1$ pour obtenir
$$
\pi^*(\mathcal{E})(-1) \longrightarrow \mathcal{O}_{\mathbf{P}(\mathcal{E})}
$$
C'est un morphisme surjectif d'un faisceau localement libre de rang $n + 1$ sur le
faisceau structural. Le complexe de Koszul correspondant est
donc exact (Compléments d'algèbre, Lemme
\ref{more-algebra-lemma-homotopy-koszul-abstract}).
Autrement dit, on dispose d'un complexe exact
$$
0 \to
\pi^*(\wedge^{n + 1}\mathcal{E})(-n - 1) \to
\ldots \to
\pi^*(\wedge^i\mathcal{E})(-i) \to
\ldots \to
\pi^*\mathcal{E}(-1) \to
\mathcal{O}_{\mathbf{P}(\mathcal{E})} \to 0
$$
Nous plaçons le terme $\pi^*(\wedge^i\mathcal{E})(-i)$ en
degré $-i$.
Calculons les images directes supérieures de ce complexe
acyclique à l'aide de la première suite spectrale de
Catégories dérivées, lemme \ref{derived-lemma-two-ss-complex-functor}.
On dispose ainsi d'une suite spectrale de termes
$$
E_1^{p, q} = R^q\pi_*\left(\pi^*(\wedge^{-p}\mathcal{E})(p)\right)
$$
convergeant vers zéro ! La formule de projection
(Cohomologie, lemme \ref{cohomology-lemma-projection-formula})
donne
$$
E_1^{p, q} = \wedge^{-p} \mathcal{E} \otimes_{\mathcal{O}_S}
R^q\pi_*\left(\mathcal{O}_{\mathbf{P}(\mathcal{E})}(p)\right).
$$
Notons que, localement sur $S$, le faisceau $\mathcal{E}$ est trivial,
c'est-à-dire isomorphe à $\mathcal{O}_S^{\oplus n + 1}$; localement sur
$S$, le morphisme $\mathbf{P}(\mathcal{E}) \to S$ s'identifie donc
à $\mathbf{P}^n_S \to S$. Ainsi,
localement sur $S$, on peut utiliser les résultats des Lemmes
\ref{lemma-cohomology-projective-space-over-ring},
\ref{lemma-identify-functorially} ou
\ref{lemma-cohomology-projective-space-over-base}.
Il en résulte que $E_1^{p, q} = 0$, sauf si $(p, q)$ vaut $(0, 0)$
ou $(-n - 1, n)$. Les termes non nuls sont
\begin{align*}
E_1^{0, 0} & = \pi_*\mathcal{O}_{\mathbf{P}(\mathcal{E})} = \mathcal{O}_S \\
E_1^{-n - 1, n} & =
R^n\pi_*\left(\pi^*(\wedge^{n + 1}\mathcal{E})(-n - 1)\right) =
\wedge^{n + 1}\mathcal{E} \otimes_{\mathcal{O}_S}
R^n\pi_*\left(\mathcal{O}_{\mathbf{P}(\mathcal{E})}(-n - 1)\right)
\end{align*}
Il ne peut donc y avoir qu'une seule différentielle
non nulle dans la suite spectrale, à savoir le morphisme
$d_{n + 1}^{-n - 1, n} : E_{n + 1}^{-n - 1, n} \to E_{n + 1}^{0, 0}$,
qui est nécessairement un isomorphisme (puisque la suite spectrale converge vers
le faisceau $0$). Ainsi $E_1^{p, q} = E_{n + 1}^{p, q}$,
et l'on obtient un isomorphisme canonique
$$
\wedge^{n + 1}\mathcal{E} \otimes_{\mathcal{O}_S}
R^n\pi_*\left(\mathcal{O}_{\mathbf{P}(\mathcal{E})}(-n - 1)\right) =
R^n\pi_*\left(\pi^*(\wedge^{n + 1}\mathcal{E})(-n - 1)\right)
\xrightarrow{d_{n + 1}^{-n - 1, n}}
\mathcal{O}_S
$$
Puisque $\wedge^{n + 1}\mathcal{E}$ est un faisceau
inversible, il en résulte que
$R^n\pi_*\mathcal{O}_{\mathbf{P}(\mathcal{E})}(-n - 1)$ est lui
aussi inversible et canoniquement isomorphe à l'inverse de
$\wedge^{n + 1}\mathcal{E}$. Nous avons donc prouvé le cas
$d = - n - 1$ du lemme.

\medskip\noindent
En travaillant localement sur $S$, on déduit immédiatement des calculs de
cohomologie des lemmes \ref{lemma-cohomology-projective-space-over-ring},
\ref{lemma-identify-functorially} ou
\ref{lemma-cohomology-projective-space-over-base} les assertions d'annulation
du lemme. De plus, le résultat concernant $R^0\pi_*$ est lui
aussi clair, puisqu'il existe des morphismes canoniques
$\text{Sym}^d(\mathcal{E}) \to \pi_* \mathcal{O}_{\mathbf{P}(\mathcal{E})}(d)$
pour tout $d$. Il reste à vérifier la description de
$R^n\pi_*\mathcal{O}_{\mathbf{P}(\mathcal{E})}(d)$
lorsque $d < -n - 1$. Pour cela, considérons le morphisme
$$
\pi^*(\text{Sym}^{-d - n - 1}(\mathcal{E}))
\otimes_{\mathcal{O}_{\mathbf{P}(\mathcal{E})}}
\mathcal{O}_{\mathbf{P}(\mathcal{E})}(d)
\longrightarrow
\mathcal{O}_{\mathbf{P}(\mathcal{E})}(-n - 1)
$$
En appliquant $R^n\pi_*$ et la formule de projection (voir ci-dessus), on obtient un morphisme
$$
\text{Sym}^{-d - n - 1}(\mathcal{E})
\otimes_{\mathcal{O}_S}
R^n\pi_*(\mathcal{O}_{\mathbf{P}(\mathcal{E})}(d))
\longrightarrow
R^n\pi_*\mathcal{O}_{\mathbf{P}(\mathcal{E})}(-n - 1) =
(\wedge^{n + 1}\mathcal{E})^{\otimes -1}
$$
(la dernière égalité a été établie
ci-dessus). Les calculs locaux des Lemmes
\ref{lemma-cohomology-projective-space-over-ring},
\ref{lemma-identify-functorially} ou
\ref{lemma-cohomology-projective-space-over-base}
montrent à nouveau que ce morphisme induit un accouplement parfait entre
$R^n\pi_*(\mathcal{O}_{\mathbf{P}(\mathcal{E})}(d))$ et
$\text{Sym}^{-d - n - 1}(\mathcal{E}) \otimes \wedge^{n + 1}(\mathcal{E})$,
comme voulu.
\end{proof}

















\section{Faisceaux cohérents sur les schémas localement noethériens}
\label{section-coherent-sheaves}

\noindent
Nous avons défini la notion de module cohérent sur un espace annelé quelconque dans
Modules, section \ref{modules-section-coherent}.
Bien que l'on puisse considérer des faisceaux cohérents sur des schémas
non noethériens, nous supposerons toujours le schéma de base localement noethérien lorsque
nous étudierons des faisceaux cohérents. Voici une caractérisation de ces faisceaux
sur les schémas localement noethériens.

\begin{lemma}
\label{lemma-coherent-Noetherian}
Soit $X$ un schéma localement noethérien.
Soit $\mathcal{F}$ un $\mathcal{O}_X$-module.
Les conditions suivantes sont équivalentes :
\begin{enumerate}
\item $\mathcal{F}$ est cohérent,
\item $\mathcal{F}$ est un $\mathcal{O}_X$-module quasi-cohérent de type fini,
\item $\mathcal{F}$ est un $\mathcal{O}_X$-module de présentation finie,
\item pour tout ouvert affine $\Spec(A) = U \subset X$, on a
$\mathcal{F}|_U = \widetilde M$, où $M$ est un $A$-module de type fini, et
\item il existe un recouvrement ouvert affine $X = \bigcup U_i$,
$U_i = \Spec(A_i)$, tel que, pour chaque indice,
$\mathcal{F}|_{U_i} = \widetilde M_i$, où $M_i$ est un $A_i$-module de type fini.
\end{enumerate}
En particulier, $\mathcal{O}_X$ est cohérent, tout
$\mathcal{O}_X$-module inversible est cohérent et, plus
généralement, tout $\mathcal{O}_X$-module localement libre de type fini est cohérent.
\end{lemma}

\begin{proof}
Les implications (1) $\Rightarrow$ (2) et (1) $\Rightarrow$ (3) sont
vraies en général, voir
Modules, lemme \ref{modules-lemma-coherent-finite-presentation}.
Si $\mathcal{F}$ est de présentation finie, alors $\mathcal{F}$ est
quasi-cohérent, voir
Modules, lemme \ref{modules-lemma-finite-presentation-quasi-coherent}. On
a donc aussi (3) $\Rightarrow$ (2).

\medskip\noindent
Supposons que $\mathcal{F}$ soit un $\mathcal{O}_X$-module quasi-cohérent
de type
fini. D'après Propriétés, lemme \ref{properties-lemma-finite-type-module},
sur tout ouvert affine
$\Spec(A) = U \subset X$, on a $\mathcal{F}|_U = \widetilde M$,
où $M$ est un $A$-module de type fini. Puisque $A$ est noethérien,
$M$ admet une présentation finie
$$
A^{\oplus m} \to A^{\oplus n} \to M \to 0.
$$
Ainsi $\mathcal{F}$ est de présentation finie d'après
Propriétés, lemme \ref{properties-lemma-finite-presentation-module}.
Autrement dit, (2) $\Rightarrow$ (3).

\medskip\noindent
D'après Modules, lemme \ref{modules-lemma-coherent-structure-sheaf}, il
suffit de montrer que $\mathcal{O}_X$ est cohérent pour prouver que (3) implique
(1). Il faut donc montrer que, pour tout ouvert $U \subset X$ et
toute famille finie de sections $f_i \in \mathcal{O}_X(U)$,
$i = 1, \ldots, n$, le noyau du morphisme
$\bigoplus_{i = 1, \ldots, n} \mathcal{O}_U \to \mathcal{O}_U$
est de type fini. Cette propriété étant locale, il suffit
de la vérifier au voisinage de tout $x \in U$.
On peut donc supposer $U = \Spec(A)$ affine. Dans ce cas,
$f_1, \ldots, f_n \in A$ sont des éléments de $A$. Puisque $A$ est
noethérien, voir
Propriétés, lemme \ref{properties-lemma-locally-Noetherian},
le noyau $K$ du morphisme $\bigoplus_{i = 1, \ldots, n} A \to A$
est un $A$-module de type fini. Voir par
exemple Algèbre, lemme \ref{algebra-lemma-Noetherian-basic}.
Le foncteur\ $\widetilde{ }$\ étant exact, voir
Schémas, lemme \ref{schemes-lemma-spec-sheaves},
on obtient une suite exacte
$$
\widetilde K \to
\bigoplus\nolimits_{i = 1, \ldots, n} \mathcal{O}_U \to
\mathcal{O}_U
$$
et, en appliquant à
nouveau Propriétés, lemme \ref{properties-lemma-finite-type-module},
on voit que $\widetilde K$ est de type fini. On en
déduit l'équivalence de (1), (2) et (3).

\medskip\noindent
Il résulte de
Propriétés, lemme \ref{properties-lemma-finite-type-module},
que (2) implique (4). Il est immédiat que (4) implique (5).
La discussion de
Schémas, section \ref{schemes-section-quasi-coherent},
montre que (5) implique la
quasi-cohérence de $\mathcal{F}$, et il est clair que (5)
implique que $\mathcal{F}$ soit de type fini. Ainsi (5) implique
(2), ce qui conclut.
\end{proof}

\begin{lemma}
\label{lemma-coherent-abelian-Noetherian}
Soit $X$ un schéma localement noethérien.
La catégorie des $\mathcal{O}_X$-modules cohérents est
abélienne. Plus précisément, le noyau et le conoyau d'un morphisme de
$\mathcal{O}_X$-modules cohérents sont cohérents. Toute extension
de faisceaux cohérents est cohérente.
\end{lemma}

\begin{proof}
C'est une reformulation de
Modules, lemme \ref{modules-lemma-coherent-abelian},
dans un cas particulier.
\end{proof}

\noindent
Le lemme suivant n'est pas toujours vrai pour la catégorie des
$\mathcal{O}_X$-modules cohérents sur un espace annelé quelconque $X$.

\begin{lemma}
\label{lemma-coherent-Noetherian-quasi-coherent-sub-quotient}
Soit $X$ un schéma localement noethérien.
Soit $\mathcal{F}$ un $\mathcal{O}_X$-module
cohérent. Tout sous-module quasi-cohérent de $\mathcal{F}$ est cohérent.
Tout module quotient quasi-cohérent de $\mathcal{F}$ est cohérent.
\end{lemma}

\begin{proof}
On peut supposer $X$ affine, disons $X = \Spec(A)$.
Propriétés, lemme \ref{properties-lemma-locally-Noetherian},
implique que $A$ est noethérien. Le lemme \ref{lemma-coherent-Noetherian}
ramène alors l'énoncé à l'algèbre. Le fait
algébrique correspondant est qu'un quotient ou un sous-module d'un $A$-module de
type fini est un $A$-module de type fini, voir par
exemple Algèbre, lemme \ref{algebra-lemma-Noetherian-basic}.
\end{proof}

\begin{lemma}
\label{lemma-tensor-hom-coherent}
Soit $X$ un schéma localement noethérien.
Soient $\mathcal{F}$, $\mathcal{G}$ des $\mathcal{O}_X$-modules cohérents.
Les $\mathcal{O}_X$-modules $\mathcal{F} \otimes_{\mathcal{O}_X} \mathcal{G}$
et $\SheafHom_{\mathcal{O}_X}(\mathcal{F}, \mathcal{G})$ sont
cohérents.
\end{lemma}

\begin{proof}
On montre dans
Modules, lemme \ref{modules-lemma-internal-hom-locally-kernel-direct-sum}, que
$\SheafHom_{\mathcal{O}_X}(\mathcal{F}, \mathcal{G})$ est cohérent. Le
résultat pour les produits tensoriels est donné
par Modules, lemme \ref{modules-lemma-tensor-product-permanence}.
\end{proof}

\begin{lemma}
\label{lemma-local-isomorphism}
Soit $X$ un schéma localement noethérien.
Soient $\mathcal{F}$, $\mathcal{G}$ des $\mathcal{O}_X$-modules cohérents.
Soit $\varphi : \mathcal{G} \to \mathcal{F}$ un homomorphisme
de $\mathcal{O}_X$-modules. Soit $x \in X$.
\begin{enumerate}
\item Si $\mathcal{F}_x = 0$, il existe un voisinage ouvert
$U \subset X$ de $x$ tel que $\mathcal{F}|_U = 0$.
\item Si $\varphi_x : \mathcal{G}_x \to \mathcal{F}_x$ est injectif,
il existe un voisinage ouvert $U \subset X$ de $x$ tel que
$\varphi|_U$ soit injectif.
\item Si $\varphi_x : \mathcal{G}_x \to \mathcal{F}_x$ est surjectif,
il existe un voisinage ouvert $U \subset X$ de $x$ tel que
$\varphi|_U$ soit surjectif.
\item Si $\varphi_x : \mathcal{G}_x \to \mathcal{F}_x$ est bijectif,
il existe un voisinage ouvert $U \subset X$ de $x$ tel que
$\varphi|_U$ soit un isomorphisme.
\end{enumerate}
\end{lemma}

\begin{proof}
Voir Modules, Lemmes
\ref{modules-lemma-finite-type-surjective-on-stalk},
\ref{modules-lemma-finite-type-stalk-zero} et
\ref{modules-lemma-finite-type-to-coherent-injective-on-stalk}.
\end{proof}

\begin{lemma}
\label{lemma-map-stalks-local-map}
Soit $X$ un schéma localement noethérien.
Soient $\mathcal{F}$, $\mathcal{G}$ des $\mathcal{O}_X$-modules cohérents.
Soit $x \in X$. Supposons
que $\psi : \mathcal{G}_x \to \mathcal{F}_x$ soit un morphisme de
$\mathcal{O}_{X, x}$-modules.
Il existe alors un voisinage ouvert $U \subset X$ de $x$ et un morphisme
$\varphi : \mathcal{G}|_U \to \mathcal{F}|_U$ tels que
$\varphi_x = \psi$.
\end{lemma}

\begin{proof}
Compte tenu du lemme \ref{lemma-coherent-Noetherian},
c'est une reformulation de
Modules, lemme \ref{modules-lemma-stalk-internal-hom}.
\end{proof}

\begin{lemma}
\label{lemma-coherent-support-closed}
Soit $X$ un schéma localement noethérien. Soit $\mathcal{F}$ un
$\mathcal{O}_X$-module cohérent. Alors $\text{Supp}(\mathcal{F})$ est fermé, et
$\mathcal{F}$ provient d'un faisceau cohérent sur le support schématique
de $\mathcal{F}$, voir
Morphismes, définition \ref{morphisms-definition-scheme-theoretic-support}.
\end{lemma}

\begin{proof}
Soit $i : Z \to X$ le support schématique de $\mathcal{F}$, et
soit $\mathcal{G}$ le faisceau quasi-cohérent de type fini sur $Z$
tel que $i_*\mathcal{G} \cong \mathcal{F}$.
Puisque $Z = \text{Supp}(\mathcal{F})$, le support est fermé.
Le schéma $Z$ est localement noethérien d'après
Morphismes, lemmes \ref{morphisms-lemma-immersion-locally-finite-type}
et \ref{morphisms-lemma-finite-type-noetherian}.
Enfin, $\mathcal{G}$ est un $\mathcal{O}_Z$-module cohérent d'après
le lemme \ref{lemma-coherent-Noetherian}.
\end{proof}

\begin{lemma}
\label{lemma-i-star-equivalence}
Soit $i : Z \to X$ une immersion fermée de schémas localement noethériens.
Soit $\mathcal{I} \subset \mathcal{O}_X$ le faisceau quasi-cohérent d'idéaux
définissant $Z$. Le foncteur $i_*$ induit une équivalence entre
la catégorie des $\mathcal{O}_X$-modules cohérents annulés par $\mathcal{I}$
et la catégorie des $\mathcal{O}_Z$-modules cohérents.
\end{lemma}

\begin{proof}
Le foncteur est pleinement fidèle d'après
Morphismes, lemme \ref{morphisms-lemma-i-star-equivalence}.
Soit $\mathcal{F}$ un $\mathcal{O}_X$-module
cohérent annulé par $\mathcal{I}$. D'après
Morphismes, lemme \ref{morphisms-lemma-i-star-equivalence},
on peut écrire $\mathcal{F} = i_*\mathcal{G}$,
où $\mathcal{G}$ est un faisceau quasi-cohérent sur $Z$. D'après
Modules, lemme \ref{modules-lemma-i-star-reflects-finite-type},
le faisceau $\mathcal{G}$ est de type fini.
Ainsi $\mathcal{G}$ est cohérent d'après le
lemme \ref{lemma-coherent-Noetherian}.
Le foncteur est donc aussi essentiellement surjectif, comme voulu.
\end{proof}

\begin{lemma}
\label{lemma-finite-pushforward-coherent}
Soit $f : X \to Y$ un morphisme de schémas.
Soit $\mathcal{F}$ un $\mathcal{O}_X$-module quasi-cohérent.
Supposons $f$ fini et $Y$ localement noethérien.
Alors $R^pf_*\mathcal{F} = 0$ pour $p > 0$, et
$f_*\mathcal{F}$ est cohérent si $\mathcal{F}$ est cohérent.
\end{lemma}

\begin{proof}
Les images directes supérieures sont nulles d'après le
lemme \ref{lemma-relative-affine-vanishing}, et
parce qu'un morphisme fini est affine (par
définition). Notons que les hypothèses impliquent que $X$ est lui aussi localement
noethérien (voir Morphismes, lemme \ref{morphisms-lemma-finite-type-noetherian}),
de sorte que l'énoncé a un sens.
Soit $\Spec(A) = V \subset Y$ un ouvert affine. D'après
Morphismes, définition \ref{morphisms-definition-integral},
on a $f^{-1}(V) = \Spec(B)$, avec $A \to B$ fini. Le
lemme \ref{lemma-coherent-Noetherian}
ramène l'énoncé au fait algébrique suivant :
si $M$ est un $B$-module de type fini, alors $M$ est aussi de
type fini comme $A$-module, voir
Algèbre, lemme \ref{algebra-lemma-finite-module-over-finite-extension}.
\end{proof}

\noindent
Dans la situation du lemme, les images directes supérieures sont elles
aussi cohérentes puisqu'elles sont
nulles. Nous montrerons qu'il en est toujours ainsi pour un morphisme propre
entre schémas localement noethériens
(proposition \ref{proposition-proper-pushforward-coherent}).

\begin{lemma}
\label{lemma-coherent-support-dimension-0}
Soit $X$ un schéma localement noethérien. Soit $\mathcal{F}$
un faisceau cohérent tel que $\dim(\text{Supp}(\mathcal{F})) \leq 0$.
Alors $\mathcal{F}$ est engendré par ses sections globales, et
$H^i(X, \mathcal{F}) = 0$ pour $i > 0$.
\end{lemma}

\begin{proof}
D'après le lemme \ref{lemma-coherent-support-closed}, on a
$\mathcal{F} = i_*\mathcal{G}$, où $i : Z \to X$ est l'inclusion
du support schématique de $\mathcal{F}$ et où $\mathcal{G}$
est un $\mathcal{O}_Z$-module cohérent. Puisque la dimension de $Z$ est
$0$, le schéma $Z$ est une réunion disjointe de schémas affines (Propriétés, Lemme
\ref{properties-lemma-locally-Noetherian-dimension-0}).
Ainsi $\mathcal{G}$ est engendré par ses sections globales et les groupes
de cohomologie supérieure de $\mathcal{G}$ sont nuls
(lemme \ref{lemma-quasi-coherent-affine-cohomology-zero}). Par
conséquent, $\mathcal{F} = i_*\mathcal{G}$ est engendré par ses
sections globales. Puisque les cohomologies de $\mathcal{F}$ et de $\mathcal{G}$ coïncident (le
lemme \ref{lemma-relative-affine-cohomology} s'applique, car une
immersion fermée est
affine), les groupes de cohomologie supérieure de $\mathcal{F}$ sont nuls.
\end{proof}

\begin{lemma}
\label{lemma-pushforward-coherent-on-open}
Soit $X$ un schéma. Soit $j : U \to X$ l'inclusion d'un ouvert.
Soit $T \subset X$ un sous-ensemble fermé contenu dans $U$.
Si $\mathcal{F}$ est un $\mathcal{O}_U$-module cohérent
tel que $\text{Supp}(\mathcal{F}) \subset T$, alors
$j_*\mathcal{F}$ est un $\mathcal{O}_X$-module cohérent.
\end{lemma}

\begin{proof}
Considérons le recouvrement ouvert $X = U \cup (X \setminus T)$.
Alors $j_*\mathcal{F}|_U = \mathcal{F}$ est cohérent, et
$j_*\mathcal{F}|_{X \setminus T} = 0$ est également cohérent.
Ainsi $j_*\mathcal{F}$ est cohérent.
\end{proof}














\section{Faisceaux cohérents sur les schémas noethériens}
\label{section-coherent-quasi-compact}

\noindent
Dans cette section, nous donnons quelques propriétés des faisceaux cohérents sur les
schémas noethériens.

\begin{lemma}
\label{lemma-acc-coherent}
Soit $X$ un schéma noethérien.
Soit $\mathcal{F}$ un $\mathcal{O}_X$-module
cohérent. Les sous-modules quasi-cohérents
de $\mathcal{F}$ vérifient la condition de chaîne ascendante. Autrement dit, pour toute suite
$$
\mathcal{F}_1 \subset \mathcal{F}_2 \subset \ldots \subset \mathcal{F}
$$
de sous-modules quasi-cohérents, on a
$\mathcal{F}_n = \mathcal{F}_{n + 1} = \ldots $ pour un certain $n \geq 0$.
\end{lemma}

\begin{proof}
Choisissons un recouvrement ouvert affine fini.
Sur chaque ouvert du recouvrement, la suite stationne d'après
Algèbre, lemme \ref{algebra-lemma-Noetherian-basic}.
Le lemme s'ensuit.
\end{proof}

\begin{lemma}
\label{lemma-power-ideal-kills-sheaf}
Soit $X$ un schéma noethérien.
Soit $\mathcal{F}$ un faisceau cohérent sur $X$.
Soit $\mathcal{I} \subset \mathcal{O}_X$ un faisceau
quasi-cohérent d'idéaux correspondant à un sous-schéma fermé $Z \subset X$.
Il existe un entier $n \geq 0$ tel que $\mathcal{I}^n\mathcal{F} = 0$
si et seulement si $\text{Supp}(\mathcal{F}) \subset Z$ ensemblistement.
\end{lemma}

\begin{proof}
Cela résulte immédiatement
d'Algèbre, lemme \ref{algebra-lemma-Noetherian-power-ideal-kills-module},
puisque $X$ admet un recouvrement fini par des spectres d'anneaux noethériens.
\end{proof}

\begin{lemma}[Artin-Rees]
\label{lemma-Artin-Rees}
Soit $X$ un schéma noethérien.
Soit $\mathcal{F}$ un faisceau cohérent sur $X$.
Soit $\mathcal{G} \subset \mathcal{F}$ un sous-faisceau quasi-cohérent.
Soit $\mathcal{I} \subset \mathcal{O}_X$ un faisceau quasi-cohérent d'idéaux.
Il
existe alors un entier $c \geq 0$ tel que, pour tout $n \geq c$, on
ait
$$
\mathcal{I}^{n - c}(\mathcal{I}^c\mathcal{F} \cap \mathcal{G})
=
\mathcal{I}^n\mathcal{F} \cap \mathcal{G}
$$
\end{lemma}

\begin{proof}
Cela résulte immédiatement
d'Algèbre, lemme \ref{algebra-lemma-Artin-Rees},
puisque $X$ admet un recouvrement fini par des spectres d'anneaux noethériens.
\end{proof}

\begin{lemma}
\label{lemma-directed-colimit-coherent}
Soit $X$ un schéma noethérien. Tout $\mathcal{O}_X$-module
quasi-cohérent est la limite inductive filtrante de ses sous-modules cohérents.
\end{lemma}

\begin{proof}
C'est une reformulation de Propriétés, Lemme
\ref{properties-lemma-quasi-coherent-colimit-finite-type},
compte tenu du fait qu'un
$\mathcal{O}_X$-module quasi-cohérent de type fini est cohérent d'après
le lemme \ref{lemma-coherent-Noetherian}.
\end{proof}

\begin{lemma}
\label{lemma-homs-over-open}
Soit $X$ un schéma noethérien.
Soit $\mathcal{F}$ un $\mathcal{O}_X$-module quasi-cohérent.
Soit $\mathcal{G}$ un $\mathcal{O}_X$-module cohérent.
Soit $\mathcal{I} \subset \mathcal{O}_X$ un faisceau quasi-cohérent d'idéaux.
Notons $Z \subset X$ le sous-schéma fermé correspondant et
posons $U = X \setminus Z$.
Il existe un isomorphisme canonique
$$
\colim_n \Hom_{\mathcal{O}_X}(\mathcal{I}^n\mathcal{G}, \mathcal{F})
\longrightarrow
\Hom_{\mathcal{O}_U}(\mathcal{G}|_U, \mathcal{F}|_U).
$$
En particulier, on a un isomorphisme
$$
\colim_n \Hom_{\mathcal{O}_X}(
\mathcal{I}^n, \mathcal{F})
\longrightarrow
\Gamma(U, \mathcal{F}).
$$
\end{lemma}

\begin{proof}
Montrons d'abord que le deuxième morphisme est un isomorphisme. Il est injectif d'après
Propriétés, lemme \ref{properties-lemma-sections-over-quasi-compact-open}.
Puisque $\mathcal{F}$ est la réunion de ses sous-modules cohérents, voir
Propriétés, lemme \ref{properties-lemma-quasi-coherent-colimit-finite-type}
(et lemme \ref{lemma-coherent-Noetherian}),
on peut supposer $\mathcal{F}$ cohérent pour prouver la surjectivité.
Notons $\mathcal{F}_n$ le sous-faisceau quasi-cohérent de $\mathcal{F}$
formé des sections annulées par $\mathcal{I}^n$,
voir Propriétés, lemme \ref{properties-lemma-sections-over-quasi-compact-open}.
Puisque $\mathcal{F}_1 \subset \mathcal{F}_2 \subset \ldots$, on a
$\mathcal{F}_n = \mathcal{F}_{n + 1} = \ldots $ pour un certain $n \geq 0$
d'après le lemme \ref{lemma-acc-coherent}. Posons $\mathcal{H} = \mathcal{F}_n$
pour cet $n$. D'après Artin-Rees (lemme \ref{lemma-Artin-Rees}),
il existe un entier $c \geq 0$ tel que
$\mathcal{I}^m\mathcal{F} \cap \mathcal{H}
\subset \mathcal{I}^{m - c}\mathcal{H}$. En prenant $m = n + c$, on obtient
$\mathcal{I}^m\mathcal{F} \cap \mathcal{H} \subset \mathcal{I}^n\mathcal{H}
= 0$. Si l'on pose $\mathcal{F}' = \mathcal{I}^m\mathcal{F}$, on
voit donc que $\mathcal{F}' \cap \mathcal{F}_n = 0$ et
$\mathcal{F}'|_U = \mathcal{F}|_U$. En particulier, le sous-faisceau
$(\mathcal{F}')_N$ des sections annulées par $\mathcal{I}^N$ est nul
pour tout $N \geq 0$. D'après
Propriétés, lemme \ref{properties-lemma-sections-over-quasi-compact-open},
on en
déduit que la flèche horizontale supérieure du diagramme commutatif
suivant est une bijection :
$$
\xymatrix{
\colim_n \Hom_{\mathcal{O}_X}(
\mathcal{I}^n, \mathcal{F}')
\ar[r] \ar[d] &
\Gamma(U, \mathcal{F}') \ar[d] \\
\colim_n \Hom_{\mathcal{O}_X}(
\mathcal{I}^n, \mathcal{F})
\ar[r] &
\Gamma(U, \mathcal{F})
}
$$
La flèche verticale de droite étant elle aussi bijective, on en déduit que
la flèche horizontale inférieure est surjective, comme voulu.

\medskip\noindent
Montrons ensuite que la première flèche du lemme est une bijection. D'après
le lemme \ref{lemma-coherent-Noetherian}, le faisceau $\mathcal{G}$
est de présentation finie, donc le faisceau
$\mathcal{H} = \SheafHom_{\mathcal{O}_X}(\mathcal{G}, \mathcal{F})$
est quasi-cohérent, voir
Schémas, section \ref{schemes-section-quasi-coherent}.
Par définition, on a
$$
\mathcal{H}(U)
=
\Hom_{\mathcal{O}_U}(\mathcal{G}|_U, \mathcal{F}|_U)
$$
Choisissons un élément $\psi$ dans le but de la première flèche du
lemme, c'est-à-dire $\psi \in \mathcal{H}(U)$. Le résultat que nous venons de prouver s'applique
à $\mathcal{H}$; il existe donc un entier $n \geq 0$ et un morphisme
$\varphi : \mathcal{I}^n \to \mathcal{H}$ qui redonne
$\psi$ par restriction à $U$. D'après
Modules, lemme \ref{modules-lemma-internal-hom},
$\varphi$ correspond à un morphisme
$$
\varphi :
\mathcal{I}^n \otimes_{\mathcal{O}_X} \mathcal{G}
\longrightarrow
\mathcal{F}.
$$
C'est presque le morphisme cherché, à ceci près que sa source est
le produit tensoriel de $\mathcal{I}^n$ et de $\mathcal{G}$
et non leur produit. Nous allons montrer que, quitte à augmenter $n$,
cette différence disparaît. Considérons la suite exacte courte
$$
0 \to \mathcal{K} \to
\mathcal{I}^n \otimes_{\mathcal{O}_X} \mathcal{G} \to
\mathcal{I}^n\mathcal{G} \to 0
$$
où $\mathcal{K}$ est le noyau. Notons que
$\mathcal{I}^n\mathcal{K} = 0$ (démonstration omise). En appliquant
à nouveau Artin-Rees, on obtient
$$
\mathcal{K}
\cap
\mathcal{I}^m(\mathcal{I}^n \otimes_{\mathcal{O}_X} \mathcal{G})
=
0
$$
pour un certain $m$ suffisamment grand. Autrement dit, le morphisme
$$
\mathcal{I}^m(\mathcal{I}^n \otimes_{\mathcal{O}_X} \mathcal{G})
\longrightarrow
\mathcal{I}^{n + m}\mathcal{G}
$$
est un isomorphisme. Soit $\varphi'$ la restriction de
$\varphi$ à ce sous-module, considérée comme morphisme
$\mathcal{I}^{m + n}\mathcal{G} \to \mathcal{F}$.
Alors $\varphi'$ définit un élément
de la source de la première flèche du lemme, dont l'image par cette
flèche est $\psi$. Nous avons donc prouvé la surjectivité de la
flèche. Nous omettons la démonstration de son injectivité.
\end{proof}

\begin{lemma}
\label{lemma-extend-coherent}
Soit $X$ un schéma localement noethérien. Soient $\mathcal{F}$, $\mathcal{G}$
des $\mathcal{O}_X$-modules cohérents. Soit $U \subset X$ un ouvert, et
soit $\varphi : \mathcal{F}|_U \to \mathcal{G}|_U$ un morphisme de
$\mathcal{O}_U$-modules. Il existe alors un sous-module
cohérent $\mathcal{F}' \subset \mathcal{F}$ coïncidant avec
$\mathcal{F}$ sur $U$, tel que $\varphi$ se prolonge en
$\varphi' : \mathcal{F}' \to \mathcal{G}$.
\end{lemma}

\begin{proof}
Soit $\mathcal{I} \subset \mathcal{O}_X$ le faisceau cohérent d'idéaux
définissant la structure réduite induite sur $X \setminus U$.
Si $X$ est noethérien, le lemme \ref{lemma-homs-over-open} montre que
l'on peut prendre $\mathcal{F}' = \mathcal{I}^n\mathcal{F}$
pour un certain $n$. Le cas général s'en déduit au moyen du lemme de Zorn.

\medskip\noindent
Considérons l'ensemble des triplets $(U', \mathcal{F}', \varphi')$ tels que
$U \subset U' \subset X$ soit ouvert, que $\mathcal{F}' \subset \mathcal{F}|_{U'}$
soit un sous-faisceau cohérent coïncidant avec $\mathcal{F}$ sur $U$, et que
$\varphi' : \mathcal{F}' \to \mathcal{G}|_{U'}$ se restreigne à
$\varphi$ sur $U$. Posons
$(U'', \mathcal{F}'', \varphi'') \geq (U', \mathcal{F}', \varphi')$
si et seulement si $U'' \supset U'$, $\mathcal{F}''|_{U'} = \mathcal{F}'$
et $\varphi''|_{U'} = \varphi'$.
Il est clair que, pour une famille totalement ordonnée
de triplets $(U_i, \mathcal{F}_i, \varphi_i)$, on
peut recoller les $\mathcal{F}_i$ en un sous-faisceau $\mathcal{F}'$ de $\mathcal{F}$
sur $U' = \bigcup U_i$, et prolonger $\varphi$ en un morphisme
$\varphi' : \mathcal{F}' \to \mathcal{G}|_{U'}$.
Tout sous-ensemble totalement ordonné de triplets admet donc un
majorant. Enfin, soit $(U', \mathcal{F}', \varphi')$
un triplet tel que $U' \not = X$. On peut alors choisir un
ouvert affine $W \subset X$ qui n'est pas contenu dans $U'$.
D'après le résultat du premier paragraphe, on peut prolonger
le sous-faisceau $\mathcal{F}'|_{W \cap U'}$ et la restriction
$\varphi'|_{W \cap U'}$ en un sous-faisceau $\mathcal{F}'' \subset \mathcal{F}|_W$
et en un morphisme $\varphi'' : \mathcal{F}'' \to \mathcal{G}|_W$.
La coïncidence de $(\mathcal{F}', \varphi')$ et de
$(\mathcal{F}'', \varphi'')$ sur $W \cap U'$ signifie précisément que
l'on peut prolonger le triplet en $(U' \cup W, \mathcal{F}''', \varphi''')$.
Ainsi, tout triplet maximal $(U', \mathcal{F}', \varphi')$
(dont l'existence résulte du lemme de Zorn) vérifie $U' = X$, ce
qui achève la démonstration.
\end{proof}









\section{Profondeur}
\label{section-depth}

\noindent
Dans cette section, nous étudions brièvement la profondeur et la propriété
$(S_k)$ pour les modules cohérents sur les schémas localement
noethériens. Rappelons que cette notion a déjà été abordée pour les
schémas localement noethériens dans Propriétés, section \ref{properties-section-Rk}.

\begin{definition}
\label{definition-depth}
Soit $X$ un schéma localement noethérien.
Soit $\mathcal{F}$ un $\mathcal{O}_X$-module cohérent.
Soit $k \geq 0$ un entier.
\begin{enumerate}
\item On dit que $\mathcal{F}$ est de {\it profondeur $k$ en un point}
$x$ de $X$ si $\text{depth}_{\mathcal{O}_{X, x}}(\mathcal{F}_x) = k$.
\item On dit que $X$ est de {\it profondeur $k$ en un point} $x$ de $X$ si
$\text{depth}(\mathcal{O}_{X, x}) = k$.
\item On dit que $\mathcal{F}$ vérifie la propriété {\it $(S_k)$} si
$$
\text{depth}_{\mathcal{O}_{X, x}}(\mathcal{F}_x)
\geq \min(k, \dim(\text{Supp}(\mathcal{F}_x)))
$$
pour tout $x \in X$.
\item On dit que $X$ vérifie la propriété {\it $(S_k)$} si $\mathcal{O}_X$ vérifie la
propriété $(S_k)$.
\end{enumerate}
\end{definition}

\noindent
Tout faisceau cohérent vérifie la condition $(S_0)$. La
condition $(S_1)$ équivaut à l'absence de points associés
immergés, voir Diviseurs, lemme \ref{divisors-lemma-S1-no-embedded}.

\begin{lemma}
\label{lemma-hom-into-depth}
Soit $X$ un schéma localement noethérien. Soient $\mathcal{F}$, $\mathcal{G}$
des $\mathcal{O}_X$-modules cohérents et $x \in X$.
\begin{enumerate}
\item Si $\mathcal{G}_x$ est de profondeur $\geq 1$, alors
$\SheafHom_{\mathcal{O}_X}(\mathcal{F}, \mathcal{G})_x$
est de profondeur $\geq 1$.
\item Si $\mathcal{G}_x$ est de profondeur $\geq 2$, alors
$\Hom_{\mathcal{O}_X}(\mathcal{F}, \mathcal{G})_x$ est de profondeur $\geq 2$.
\end{enumerate}
\end{lemma}

\begin{proof}
Remarquons que $\SheafHom_{\mathcal{O}_X}(\mathcal{F}, \mathcal{G})$ est
un $\mathcal{O}_X$-module cohérent d'après le lemme \ref{lemma-tensor-hom-coherent}.
Les modules cohérents sont de présentation finie
(lemme \ref{lemma-coherent-Noetherian}); le passage aux fibres commute donc à
la formation de $\SheafHom$ et de $\Hom$, voir
Modules, lemme \ref{modules-lemma-stalk-internal-hom}.
On se ramène ainsi au cas des modules de type fini sur
les anneaux locaux, traité dans Compléments d'algèbre, lemme \ref{more-algebra-lemma-hom-into-depth}.
\end{proof}

\begin{lemma}
\label{lemma-hom-into-S2}
Soit $X$ un schéma localement noethérien. Soient $\mathcal{F}$, $\mathcal{G}$
des $\mathcal{O}_X$-modules cohérents.
\begin{enumerate}
\item Si $\mathcal{G}$ vérifie la propriété $(S_1)$, alors
$\SheafHom_{\mathcal{O}_X}(\mathcal{F}, \mathcal{G})$ vérifie la propriété $(S_1)$.
\item Si $\mathcal{G}$ vérifie la propriété $(S_2)$, alors
$\SheafHom_{\mathcal{O}_X}(\mathcal{F}, \mathcal{G})$ vérifie la propriété $(S_2)$.
\end{enumerate}
\end{lemma}

\begin{proof}
Cela résulte immédiatement du lemme \ref{lemma-hom-into-depth}
et des définitions.
\end{proof}

\noindent
Nous avons vu dans Propriétés, lemme \ref{properties-lemma-scheme-CM-iff-all-Sk},
qu'un schéma localement noethérien
est de Cohen-Macaulay si et seulement s'il vérifie $(S_k)$ pour tout $k$.
Il est donc naturel d'introduire la définition suivante, qui équivaut
à demander que toutes les fibres soient des modules de Cohen-Macaulay.

\begin{definition}
\label{definition-Cohen-Macaulay}
Soit $X$ un schéma localement noethérien.
Soit $\mathcal{F}$ un $\mathcal{O}_X$-module cohérent.
On dit que $\mathcal{F}$ est {\it de Cohen-Macaulay} si et seulement s'il
vérifie $(S_k)$ pour tout $k \geq 0$.
\end{definition}

\begin{lemma}
\label{lemma-Cohen-Macaulay-over-regular}
Soit $X$ un schéma régulier. Soit $\mathcal{F}$ un
$\mathcal{O}_X$-module cohérent. Les conditions suivantes sont équivalentes :
\begin{enumerate}
\item $\mathcal{F}$ est de Cohen-Macaulay et $\text{Supp}(\mathcal{F}) = X$,
\item $\mathcal{F}$ est localement libre de type fini de rang $> 0$.
\end{enumerate}
\end{lemma}

\begin{proof}
Soit $x \in X$. Si (2) est vérifiée, alors $\mathcal{F}_x$ est un
$\mathcal{O}_{X, x}$-module libre de rang $> 0$. Ainsi
$\text{depth}(\mathcal{F}_x) = \dim(\mathcal{O}_{X, x})$,
car un anneau local régulier est de Cohen-Macaulay
(Algèbre, lemme \ref{algebra-lemma-regular-ring-CM}).
Réciproquement, si (1) est vérifiée, alors $\mathcal{F}_x$ est un
module de Cohen-Macaulay maximal sur $\mathcal{O}_{X, x}$
(Algèbre, définition \ref{algebra-definition-maximal-CM}).
Ainsi $\mathcal{F}_x$ est libre d'après
Algèbre, lemme \ref{algebra-lemma-regular-mcm-free}.
\end{proof}









\section{Dévissage des faisceaux cohérents}
\label{section-devissage}

\noindent
Soit $X$ un schéma noethérien. Considérons un sous-schéma fermé intègre
$i : Z \to X$. Il est souvent commode de considérer les faisceaux cohérents de
la forme $i_*\mathcal{G}$, où $\mathcal{G}$ est un faisceau cohérent sur
$Z$. Nous nous intéressons en particulier au cas où $\mathcal{G}$
est sans torsion et de rang $1$. Par exemple, $\mathcal{G}$ peut être
un faisceau d'idéaux non nul sur $Z$, ou, plus particulièrement,
$\mathcal{G} = \mathcal{O}_Z$.

\medskip\noindent
Tout au long de cette section, nous utiliserons le fait qu'un faisceau
cohérent est la même chose qu'un faisceau quasi-cohérent de type fini, et
qu'un sous-quotient quasi-cohérent d'un faisceau cohérent est cohérent, voir
section \ref{section-coherent-sheaves}.
Le support d'un faisceau cohérent est fermé, voir
Modules, lemme \ref{modules-lemma-support-finite-type-closed}.

\begin{lemma}
\label{lemma-prepare-filter-support}
Soit $X$ un schéma noethérien.
Soit $\mathcal{F}$ un faisceau cohérent sur $X$.
Supposons $\text{Supp}(\mathcal{F}) = Z \cup Z'$, avec $Z$, $Z'$ fermés.
Il existe alors une suite exacte courte de faisceaux cohérents
$$
0 \to \mathcal{G}' \to \mathcal{F} \to \mathcal{G} \to 0
$$
avec $\text{Supp}(\mathcal{G}') \subset Z'$ et
$\text{Supp}(\mathcal{G}) \subset Z$.
\end{lemma}

\begin{proof}
Soit $\mathcal{I} \subset \mathcal{O}_X$ le faisceau d'idéaux définissant la
structure réduite induite de sous-schéma fermé sur $Z$, voir
Schémas, lemme \ref{schemes-lemma-reduced-closed-subscheme}.
Considérons les sous-faisceaux
$\mathcal{G}'_n = \mathcal{I}^n\mathcal{F}$ et les
quotients $\mathcal{G}_n = \mathcal{F}/\mathcal{I}^n\mathcal{F}$.
Pour tout $n$, on a une suite exacte courte
$$
0 \to \mathcal{G}'_n \to \mathcal{F} \to \mathcal{G}_n \to 0
$$
Pour tout point $x$ de $Z' \setminus Z$, on a
$\mathcal{I}_x = \mathcal{O}_{X, x}$
et donc $\mathcal{G}_{n, x} = 0$. Ainsi
$\text{Supp}(\mathcal{G}_n) \subset Z$. Notons que $X \setminus Z'$
est un schéma noethérien. D'après le lemme \ref{lemma-power-ideal-kills-sheaf},
il existe donc un entier $n$ tel que
$\mathcal{G}'_n|_{X \setminus Z'} =
\mathcal{I}^n\mathcal{F}|_{X \setminus Z'} = 0$.
Pour un tel $n$, on a $\text{Supp}(\mathcal{G}'_n) \subset Z'$. Il
suffit donc de poser
$\mathcal{G}' = \mathcal{G}'_n$ et $\mathcal{G} = \mathcal{G}_n$
pour conclure.
\end{proof}

\begin{lemma}
\label{lemma-prepare-filter-irreducible}
Soit $X$ un schéma noethérien.
Soit $i : Z \to X$ un sous-schéma fermé intègre.
Soit $\xi \in Z$ son point générique.
Soit $\mathcal{F}$ un faisceau cohérent sur $X$.
Supposons que $\mathcal{F}_\xi$ soit annulé par
$\mathfrak m_\xi$. Il existe alors un entier
$r \geq 0$, un faisceau cohérent d'idéaux $\mathcal{I} \subset \mathcal{O}_Z$
et un morphisme injectif de faisceaux cohérents
$$
i_*\left(\mathcal{I}^{\oplus r}\right) \to \mathcal{F}
$$
qui est un isomorphisme au voisinage de $\xi$.
\end{lemma}

\begin{proof}
Soit $\mathcal{J} \subset \mathcal{O}_X$ le faisceau d'idéaux de $Z$.
Soit $\mathcal{F}' \subset \mathcal{F}$ le sous-faisceau des
sections locales de $\mathcal{F}$ annulées par
$\mathcal{J}$. C'est un faisceau quasi-cohérent d'après
Propriétés, lemme \ref{properties-lemma-sections-annihilated-by-ideal}. De
plus, $\mathcal{F}'_\xi = \mathcal{F}_\xi$, car
$\mathcal{J}_\xi = \mathfrak m_\xi$, et d'après l'assertion (3)
de Propriétés, lemme \ref{properties-lemma-sections-annihilated-by-ideal}.
Le lemme \ref{lemma-local-isomorphism} montre que
$\mathcal{F}' \to \mathcal{F}$
induit un isomorphisme au voisinage de $\xi$.
On peut donc remplacer $\mathcal{F}$ par $\mathcal{F}'$ et supposer
que $\mathcal{F}$ est annulé par $\mathcal{J}$.

\medskip\noindent
Supposons $\mathcal{J}\mathcal{F} = 0$. D'après le
lemme \ref{lemma-i-star-equivalence}, on peut écrire
$\mathcal{F} = i_*\mathcal{G}$,
où $\mathcal{G}$ est un faisceau cohérent sur $Z$. Supposons que l'on puisse trouver un morphisme
$\mathcal{I}^{\oplus r} \to \mathcal{G}$ qui soit un
isomorphisme au voisinage du point générique $\xi$ de $Z$.
En appliquant $i_*$ (qui est exact à gauche), on obtient alors le résultat du
lemme. Nous sommes donc ramenés au cas $X = Z$.

\medskip\noindent
Supposons que $Z = X$ soit un schéma noethérien intègre de point générique $\xi$.
Dans ce cas, $\mathcal{O}_{X, \xi} = \kappa(\xi)$ est, pour $X$, le corps
des fonctions.
Puisque $\mathcal{F}_\xi$ est un $\mathcal{O}_\xi$-module de type fini,
l'entier $r = \dim_{\kappa(\xi)} \mathcal{F}_\xi$ est fini.
Les faisceaux $\mathcal{O}_X^{\oplus r}$ et $\mathcal{F}$
ont donc des fibres isomorphes en $\xi$.
D'après le lemme \ref{lemma-map-stalks-local-map}, il existe un ouvert non
vide $U \subset X$ et un morphisme
$\psi : \mathcal{O}_X^{\oplus r}|_U \to \mathcal{F}|_U$
qui est un isomorphisme
en $\xi$, et donc au voisinage de $\xi$, d'après
le lemme \ref{lemma-local-isomorphism}.
D'après Schémas, lemme \ref{schemes-lemma-reduced-closed-subscheme},
il existe un faisceau quasi-cohérent d'idéaux
$\mathcal{I} \subset \mathcal{O}_X$
dont le sous-schéma fermé associé $Y \subset X$ est le complémentaire
de $U$.
D'après le lemme \ref{lemma-homs-over-open}, il existe un entier $n \geq 0$ et un morphisme
$\mathcal{I}^n(\mathcal{O}_X^{\oplus r}) \to \mathcal{F}$
qui redonne $\psi$ sur $U$. Puisque
$\mathcal{I}^n(\mathcal{O}_X^{\oplus r}) = (\mathcal{I}^n)^{\oplus r}$,
on obtient un morphisme comme dans le lemme. Il est injectif, car $X$ est
intègre et il est injectif au point générique de $X$
(démonstration facile omise).
\end{proof}

\begin{lemma}
\label{lemma-coherent-filter}
Soit $X$ un schéma noethérien.
Soit $\mathcal{F}$ un faisceau cohérent sur $X$.
Il existe une filtration
$$
0 = \mathcal{F}_0 \subset \mathcal{F}_1 \subset
\ldots \subset \mathcal{F}_m = \mathcal{F}
$$
par des sous-faisceaux cohérents telle que, pour chaque $j = 1, \ldots, m$,
il existe un sous-schéma fermé intègre $Z_j \subset X$ et
un faisceau cohérent d'idéaux non nul $\mathcal{I}_j \subset \mathcal{O}_{Z_j}$
tels que
$$
\mathcal{F}_j/\mathcal{F}_{j - 1}
\cong (Z_j \to X)_* \mathcal{I}_j
$$
\end{lemma}

\begin{proof}
Considérons l'ensemble
$$
\mathcal{T} =
\left\{
\begin{matrix}
Z \subset X
\text{ fermé tel qu'il existe un faisceau cohérent }
\mathcal{F} \\
\text{ avec }
\text{Supp}(\mathcal{F}) = Z
\text{ pour lequel le lemme est faux}
\end{matrix}
\right\}
$$
Nous voulons montrer que $\mathcal{T}$ est vide. Sinon,
puisque $X$ est noethérien, on peut choisir un élément minimal
$Z \in \mathcal{T}$. Il existe alors un faisceau
cohérent $\mathcal{F}$ sur $X$, de support $Z$, pour lequel le lemme
est faux. Il est clair que $Z \not = \emptyset$, car le
seul faisceau de support vide est le faisceau nul, pour lequel le
lemme est vrai (avec $m = 0$).

\medskip\noindent
Si $Z$ n'est pas irréductible, on peut écrire $Z = Z_1 \cup Z_2$,
où $Z_1, Z_2$ sont fermés et strictement contenus dans $Z$.
On peut alors appliquer le lemme \ref{lemma-prepare-filter-support}
pour obtenir une suite exacte courte de faisceaux cohérents
$$
0 \to
\mathcal{G}_1 \to
\mathcal{F} \to
\mathcal{G}_2 \to 0
$$
avec $\text{Supp}(\mathcal{G}_i) \subset Z_i$. Par minimalité de
$Z$, chacun des $\mathcal{G}_i$ admet une filtration comme dans l'énoncé
du lemme. En considérant la filtration induite sur $\mathcal{F}$,
on aboutit à une contradiction. On en déduit
que $Z$ est irréductible.

\medskip\noindent
Supposons $Z$ irréductible. Soit $\mathcal{J}$ le faisceau d'idéaux définissant
la structure réduite induite de sous-schéma fermé sur $Z$,
voir Schémas, lemme \ref{schemes-lemma-reduced-closed-subscheme}.
D'après le lemme \ref{lemma-power-ideal-kills-sheaf}, il existe un
entier $n \geq 0$ tel que $\mathcal{J}^n\mathcal{F} = 0$. On obtient donc
une filtration
$$
0 = \mathcal{J}^n\mathcal{F} \subset \mathcal{J}^{n - 1}\mathcal{F}
\subset \ldots \subset \mathcal{J}\mathcal{F} \subset \mathcal{F}
$$
dont chaque sous-quotient successif est annulé par $\mathcal{J}$.
Si chacun de ces sous-quotients admet une filtration comme dans l'énoncé du lemme, il
en va de même pour $\mathcal{F}$. Autrement dit, on peut
supposer que $\mathcal{J}$ annule $\mathcal{F}$.

\medskip\noindent
Lorsque $Z$ est irréductible et $\mathcal{J}\mathcal{F} = 0$,
on peut appliquer le lemme \ref{lemma-prepare-filter-irreducible}.
On obtient ainsi une suite exacte courte
$$
0 \to
i_*(\mathcal{I}^{\oplus r}) \to
\mathcal{F} \to
\mathcal{Q} \to 0
$$
où $\mathcal{Q}$ est le quotient.
Puisque $\mathcal{Q}$ est nul au voisinage de $\xi$,
d'après le lemme qui vient d'être cité, le support de $\mathcal{Q}$
est strictement contenu dans $Z$. Ainsi $\mathcal{Q}$
admet une filtration du type voulu par minimalité de $Z$. Mais
il en va alors de même pour $\mathcal{F}$, d'où la contradiction finale.
\end{proof}

\begin{lemma}
\label{lemma-property-initial}
Soit $X$ un schéma noethérien.
Soit $\mathcal{P}$ une propriété des faisceaux cohérents sur $X$. Supposons que
\begin{enumerate}
\item pour toute suite exacte courte de faisceaux cohérents
$$
0 \to \mathcal{F}_1 \to \mathcal{F} \to \mathcal{F}_2 \to 0
$$
si les $\mathcal{F}_i$, $i = 1, 2$, vérifient la propriété $\mathcal{P}$,
il en va de même pour $\mathcal{F}$;
\item pour tout sous-schéma fermé intègre $Z \subset X$
et tout faisceau quasi-cohérent d'idéaux
$\mathcal{I} \subset \mathcal{O}_Z$, on ait
$\mathcal{P}$ pour $i_*\mathcal{I}$.
\end{enumerate}
Alors la propriété $\mathcal{P}$ est vérifiée par tout faisceau cohérent
sur $X$.
\end{lemma}

\begin{proof}
Notons d'abord que, si $\mathcal{F}$ est un faisceau cohérent muni d'une filtration
$$
0 = \mathcal{F}_0 \subset \mathcal{F}_1 \subset
\ldots \subset \mathcal{F}_m = \mathcal{F}
$$
par des sous-faisceaux cohérents telle que chaque $\mathcal{F}_i/\mathcal{F}_{i - 1}$
vérifie la propriété $\mathcal{P}$, il en va de même pour $\mathcal{F}$.
Cela résulte de la condition (1) sur $\mathcal{P}$.
D'autre part, le lemme \ref{lemma-coherent-filter}
permet de filtrer tout $\mathcal{F}$
avec des sous-quotients successifs comme dans
(2). Le lemme s'ensuit.
\end{proof}

\begin{lemma}
\label{lemma-property-irreducible}
Soit $X$ un schéma noethérien. Soit $Z_0 \subset X$ un sous-ensemble fermé irréductible
de point générique $\xi$. Soit $\mathcal{P}$ une propriété des faisceaux
cohérents sur $X$ dont le support est contenu dans $Z_0$, telle que
\begin{enumerate}
\item pour toute suite exacte courte de faisceaux cohérents, si
deux d'entre eux vérifient la propriété $\mathcal{P}$, il en va de même pour
le troisième;
\item pour tout sous-schéma fermé intègre $Z \subset Z_0 \subset X$,
$Z \not = Z_0$, et tout faisceau quasi-cohérent d'idéaux
$\mathcal{I} \subset \mathcal{O}_Z$, on ait
$\mathcal{P}$ pour $(Z \to X)_*\mathcal{I}$;
\item il existe un faisceau cohérent $\mathcal{G}$ sur $X$ tel que
\begin{enumerate}
\item $\text{Supp}(\mathcal{G}) = Z_0$,
\item $\mathcal{G}_\xi$ soit annulé par $\mathfrak m_\xi$,
\item $\dim_{\kappa(\xi)} \mathcal{G}_\xi = 1$, et
\item la propriété $\mathcal{P}$ soit vérifiée par $\mathcal{G}$.
\end{enumerate}
\end{enumerate}
Alors la propriété $\mathcal{P}$ est vérifiée par tout faisceau cohérent
$\mathcal{F}$ sur $X$ dont le support est contenu dans $Z_0$.
\end{lemma}

\begin{proof}
Notons d'abord que, si $\mathcal{F}$ est un faisceau cohérent de support
contenu dans $Z_0$, muni d'une filtration
$$
0 = \mathcal{F}_0 \subset \mathcal{F}_1 \subset
\ldots \subset \mathcal{F}_m = \mathcal{F}
$$
par des sous-faisceaux cohérents telle que chaque $\mathcal{F}_i/\mathcal{F}_{i - 1}$
vérifie la propriété $\mathcal{P}$, il en va de même pour $\mathcal{F}$. De même, si $\mathcal{F}$
vérifie la propriété $\mathcal{P}$ et que tous les
$\mathcal{F}_i/\mathcal{F}_{i - 1}$ sauf un vérifient $\mathcal{P}$,
le dernier la vérifie aussi. Cela résulte de l'hypothèse (1).

\medskip\noindent
On en déduit d'abord que tout faisceau cohérent dont le support
est strictement contenu dans $Z_0$ vérifie la propriété $\mathcal{P}$. En effet, un
tel faisceau admet une filtration (voir lemme \ref{lemma-coherent-filter})
dont les sous-quotients vérifient $\mathcal{P}$ d'après (2).

\medskip\noindent
Soit $\mathcal{G}$ comme dans (3). D'après le lemme \ref{lemma-prepare-filter-irreducible},
il existe un faisceau d'idéaux $\mathcal{I}$ sur $Z_0$, un
entier $r \geq 1$ et une suite exacte courte
$$
0 \to
\left((Z_0 \to X)_*\mathcal{I}\right)^{\oplus r} \to
\mathcal{G} \to
\mathcal{Q} \to 0
$$
où le support de $\mathcal{Q}$ est strictement contenu dans $Z_0$.
D'après (3)(c), on a $r = 1$. Puisque $\mathcal{Q}$ vérifie lui aussi la propriété $\mathcal{P}$,
on en déduit que $(Z_0 \to X)_*\mathcal{I}$ vérifie la propriété
$\mathcal{P}$.

\medskip\noindent
Supposons ensuite que $\mathcal{I}' \not = 0$ soit un autre
faisceau quasi-cohérent d'idéaux sur $Z_0$. On peut considérer l'intersection
$\mathcal{I}'' = \mathcal{I}' \cap \mathcal{I}$, ce qui
donne deux suites exactes courtes
$$
0 \to
(Z_0 \to X)_*\mathcal{I}'' \to
(Z_0 \to X)_*\mathcal{I} \to
\mathcal{Q} \to 0
$$
et
$$
0 \to
(Z_0 \to X)_*\mathcal{I}'' \to
(Z_0 \to X)_*\mathcal{I}' \to
\mathcal{Q}' \to 0.
$$
Notons que les supports des faisceaux cohérents $\mathcal{Q}$ et
$\mathcal{Q}'$ sont strictement contenus dans $Z_0$.
Ainsi $\mathcal{Q}$ et $\mathcal{Q}'$ vérifient la propriété $\mathcal{P}$
(voir ci-dessus). La condition (1) entraîne alors
que $(Z_0 \to X)_*\mathcal{I}''$ et $(Z_0 \to X)_*\mathcal{I}'$
vérifient eux aussi $\mathcal{P}$.

\medskip\noindent
Pour achever la démonstration, il suffit de noter que tout faisceau cohérent
$\mathcal{F}$ sur $X$ dont le support est contenu dans $Z_0$ admet une filtration (voir
à nouveau lemme \ref{lemma-coherent-filter}) dont les sous-quotients
vérifient tous la propriété $\mathcal{P}$ d'après ce qui précède.
\end{proof}

\begin{lemma}
\label{lemma-property}
Soit $X$ un schéma noethérien.
Soit $\mathcal{P}$ une propriété des faisceaux cohérents sur $X$ telle que
\begin{enumerate}
\item pour toute suite exacte courte de faisceaux cohérents, si
deux d'entre eux vérifient la propriété $\mathcal{P}$, il en va de même pour
le troisième;
\item pour tout sous-schéma fermé intègre $Z \subset X$
de point générique $\xi$, il
existe un faisceau cohérent $\mathcal{G}$ tel que
\begin{enumerate}
\item $\text{Supp}(\mathcal{G}) = Z$,
\item $\mathcal{G}_\xi$ soit annulé par $\mathfrak m_\xi$,
\item $\dim_{\kappa(\xi)} \mathcal{G}_\xi = 1$, et
\item la propriété $\mathcal{P}$ soit vérifiée par $\mathcal{G}$.
\end{enumerate}
\end{enumerate}
Alors la propriété $\mathcal{P}$ est vérifiée par tout faisceau cohérent
sur $X$.
\end{lemma}

\begin{proof}
D'après le lemme \ref{lemma-property-initial}, il suffit de montrer que,
pour tout sous-schéma fermé intègre $Z \subset X$ et tout faisceau
quasi-cohérent d'idéaux $\mathcal{I} \subset \mathcal{O}_Z$, la propriété $\mathcal{P}$
est vérifiée par $(Z \to X)_*\mathcal{I}$. Sinon, puisque $X$ est noethérien,
il existe un sous-schéma fermé intègre minimal $Z_0 \subset X$ tel que
$\mathcal{P}$ ne soit pas vérifiée par $(Z_0 \to X)_*\mathcal{I}_0$ pour un
certain faisceau quasi-cohérent d'idéaux $\mathcal{I}_0 \subset \mathcal{O}_{Z_0}$, mais
que $\mathcal{P}$ soit vérifiée par $(Z \to X)_*\mathcal{I}$ pour tout sous-schéma
fermé intègre $Z \subset Z_0$, $Z \not = Z_0$, et tout
faisceau quasi-cohérent d'idéaux $\mathcal{I} \subset \mathcal{O}_Z$. Puisque l'existence
de $\mathcal{G}$ pour $Z_0$ est assurée par (2), le
lemme \ref{lemma-property-irreducible} exclut cette éventualité.
\end{proof}

\begin{lemma}
\label{lemma-property-irreducible-higher-rank-cohomological}
Soit $X$ un schéma noethérien. Soit $Z_0 \subset X$ un sous-ensemble fermé
irréductible de point générique $\xi$. Soit $\mathcal{P}$ une propriété des
faisceaux cohérents sur $X$ telle que
\begin{enumerate}
\item pour toute suite exacte courte de faisceaux cohérents
$$
0 \to \mathcal{F}_1 \to \mathcal{F} \to \mathcal{F}_2 \to 0
$$
si les $\mathcal{F}_i$, $i = 1, 2$, vérifient la propriété $\mathcal{P}$,
il en va de même pour $\mathcal{F}$;
\item si $\mathcal{P}$ est vérifiée par $\mathcal{F}^{\oplus r}$ pour un
certain $r \geq 1$, elle est vérifiée par $\mathcal{F}$;
\item pour tout sous-schéma fermé intègre $Z \subset Z_0 \subset X$,
$Z \not = Z_0$, et tout faisceau quasi-cohérent d'idéaux
$\mathcal{I} \subset \mathcal{O}_Z$, on ait
$\mathcal{P}$ pour $(Z \to X)_*\mathcal{I}$;
\item il existe un faisceau cohérent $\mathcal{G}$ tel que
\begin{enumerate}
\item $\text{Supp}(\mathcal{G}) = Z_0$,
\item $\mathcal{G}_\xi$ soit annulé par $\mathfrak m_\xi$, et
\item pour tout faisceau quasi-cohérent d'idéaux
$\mathcal{J} \subset \mathcal{O}_X$ tel que
$\mathcal{J}_\xi = \mathcal{O}_{X, \xi}$, il existe un sous-faisceau
quasi-cohérent $\mathcal{G}' \subset \mathcal{J}\mathcal{G}$ avec
$\mathcal{G}'_\xi = \mathcal{G}_\xi$ et tel que
$\mathcal{P}$ soit vérifiée par $\mathcal{G}'$.
\end{enumerate}
\end{enumerate}
Alors la propriété $\mathcal{P}$ est vérifiée par tout faisceau cohérent
$\mathcal{F}$ sur $X$ dont le support est contenu dans $Z_0$.
\end{lemma}

\begin{proof}
Notons que, si $\mathcal{F}$ est un faisceau cohérent muni d'une filtration
$$
0 = \mathcal{F}_0 \subset \mathcal{F}_1 \subset
\ldots \subset \mathcal{F}_m = \mathcal{F}
$$
par des sous-faisceaux cohérents telle que chaque $\mathcal{F}_i/\mathcal{F}_{i - 1}$
vérifie la propriété $\mathcal{P}$, il en va de même pour $\mathcal{F}$.
Cela résulte de l'hypothèse (1).

\medskip\noindent
On en déduit d'abord que tout faisceau cohérent dont le support
est strictement contenu dans $Z_0$ vérifie la propriété $\mathcal{P}$. En effet, un
tel faisceau admet une filtration (voir lemme \ref{lemma-coherent-filter})
dont les sous-quotients vérifient $\mathcal{P}$ d'après (3).

\medskip\noindent
Notons $i : Z_0 \to X$ l'immersion fermée. Considérons
un faisceau cohérent $\mathcal{G}$ comme dans (4). D'après
le lemme \ref{lemma-prepare-filter-irreducible},
il existe un faisceau d'idéaux $\mathcal{I}$ sur $Z_0$ et
une suite exacte courte
$$
0 \to
i_*\mathcal{I}^{\oplus r} \to
\mathcal{G} \to
\mathcal{Q} \to 0
$$
où le support de $\mathcal{Q}$ est strictement contenu dans $Z_0$.
En particulier, $r > 0$ et $\mathcal{I}$ est non
nul, car le support de $\mathcal{G}$ est égal à $Z_0$.
Soit $\mathcal{I}' \subset \mathcal{I}$ un faisceau quasi-cohérent d'idéaux
non nul sur $Z_0$ contenu dans $\mathcal{I}$.
On obtient alors aussi une suite exacte courte
$$
0 \to
i_*(\mathcal{I}')^{\oplus r} \to
\mathcal{G} \to
\mathcal{Q}' \to 0
$$
où $\mathcal{Q}'$ a son support strictement contenu dans $Z_0$.
Soit $\mathcal{J} \subset \mathcal{O}_X$ un faisceau quasi-cohérent
d'idéaux définissant le support de $\mathcal{Q}'$ (par exemple
l'idéal correspondant à la structure réduite induite de sous-schéma
fermé sur le support de $\mathcal{Q}'$). Alors
$\mathcal{J}_\xi = \mathcal{O}_{X, \xi}$. D'après le
lemme \ref{lemma-power-ideal-kills-sheaf},
on a $\mathcal{J}^n\mathcal{Q}' = 0$ pour un certain $n$.
Ainsi $\mathcal{J}^n\mathcal{G} \subset i_*(\mathcal{I}')^{\oplus r}$.
L'hypothèse (4)(c) du lemme assure l'existence d'un
sous-faisceau quasi-cohérent $\mathcal{G}' \subset \mathcal{J}^n\mathcal{G}$
tel que $\mathcal{G}'_\xi = \mathcal{G}_\xi$
et qui vérifie la propriété $\mathcal{P}$.
On obtient donc une suite exacte courte
$$
0 \to \mathcal{G}' \to
i_*(\mathcal{I}')^{\oplus r} \to
\mathcal{Q}'' \to 0
$$
où $\mathcal{Q}''$ a son support strictement contenu dans $Z_0$.
D'après nos remarques initiales et la condition (1) du lemme, on
en déduit que $i_*(\mathcal{I}')^{\oplus r}$ vérifie
$\mathcal{P}$. Ainsi $i_*\mathcal{I}'$ vérifie
$\mathcal{P}$ d'après (2). Enfin, pour un faisceau
quasi-cohérent d'idéaux quelconque $\mathcal{I}'' \subset \mathcal{O}_{Z_0}$, on peut poser
$\mathcal{I}' = \mathcal{I}'' \cap \mathcal{I}$, ce qui
donne une suite exacte courte
$$
0 \to
i_*(\mathcal{I}') \to
i_*(\mathcal{I}'') \to
\mathcal{Q}''' \to 0
$$
où $\mathcal{Q}'''$ a son support strictement contenu dans $Z_0$.
On en déduit que la propriété $\mathcal{P}$ est également vérifiée par
$i_*\mathcal{I}''$.

\medskip\noindent
Pour achever la démonstration, il suffit de noter que tout faisceau cohérent
$\mathcal{F}$ sur $X$ dont le support est contenu dans $Z_0$ admet une filtration (voir
à nouveau lemme \ref{lemma-coherent-filter}) dont les sous-quotients
vérifient tous la propriété $\mathcal{P}$ d'après ce qui précède.
\end{proof}

\begin{lemma}
\label{lemma-property-higher-rank-cohomological}
Soit $X$ un schéma noethérien.
Soit $\mathcal{P}$ une propriété des faisceaux cohérents sur $X$ telle que
\begin{enumerate}
\item pour toute suite exacte courte de faisceaux cohérents
$$
0 \to \mathcal{F}_1 \to \mathcal{F} \to \mathcal{F}_2 \to 0
$$
si les $\mathcal{F}_i$, $i = 1, 2$, vérifient la propriété $\mathcal{P}$,
il en va de même pour $\mathcal{F}$;
\item si $\mathcal{P}$ est vérifiée par $\mathcal{F}^{\oplus r}$ pour un
certain $r \geq 1$, elle est vérifiée par $\mathcal{F}$;
\item pour tout sous-schéma fermé intègre $Z \subset X$
de point générique $\xi$, il
existe un faisceau cohérent $\mathcal{G}$ tel que
\begin{enumerate}
\item $\text{Supp}(\mathcal{G}) = Z$,
\item $\mathcal{G}_\xi$ soit annulé par $\mathfrak m_\xi$, et
\item pour tout faisceau quasi-cohérent d'idéaux
$\mathcal{J} \subset \mathcal{O}_X$ tel que
$\mathcal{J}_\xi = \mathcal{O}_{X, \xi}$, il existe un sous-faisceau
quasi-cohérent $\mathcal{G}' \subset \mathcal{J}\mathcal{G}$ avec
$\mathcal{G}'_\xi = \mathcal{G}_\xi$ et tel que
$\mathcal{P}$ soit vérifiée par $\mathcal{G}'$.
\end{enumerate}
\end{enumerate}
Alors la propriété $\mathcal{P}$ est vérifiée par tout faisceau cohérent
sur $X$.
\end{lemma}

\begin{proof}
Cela résulte du lemme \ref{lemma-property-irreducible-higher-rank-cohomological}
exactement comme le lemme \ref{lemma-property} résulte du
lemme \ref{lemma-property-irreducible}.
\end{proof}








\section{Morphismes finis et schémas affines}
\label{section-finite-affine}

\noindent
Dans cette section, nous utilisons les résultats des sections précédentes
pour montrer que l'image d'un schéma affine noethérien par un morphisme fini
est affine. Nous verrons plus loin que ce résultat vaut plus
généralement (voir Limites, lemme \ref{limits-lemma-affine} et
proposition \ref{limits-proposition-affine}).

\begin{lemma}
\label{lemma-finite-morphism-Noetherian}
Soit $f : Y \to X$ un morphisme de schémas.
Supposons $f$ fini et surjectif, et $X$ localement noethérien.
Soit $Z \subset X$ un sous-schéma fermé intègre de
point générique $\xi$. Il
existe alors un faisceau cohérent $\mathcal{F}$ sur $Y$
tel que le support de $f_*\mathcal{F}$ soit égal à $Z$
et que $(f_*\mathcal{F})_\xi$ soit annulé par $\mathfrak m_\xi$.
\end{lemma}

\begin{proof}
Notons que $Y$ est localement noethérien d'après
Morphismes, lemme \ref{morphisms-lemma-finite-type-noetherian}.
Puisque $f$ est surjectif, la fibre $Y_\xi$ est non vide.
Choisissons $\xi' \in Y$ d'image $\xi$. Posons $Z' = \overline{\{\xi'\}}$.
On peut considérer $Z' \subset Y$ comme un sous-schéma fermé réduit,
voir Schémas, lemme \ref{schemes-lemma-reduced-closed-subscheme}.
Ainsi le faisceau $\mathcal{F} = (Z' \to Y)_*\mathcal{O}_{Z'}$
est un faisceau cohérent sur $Y$ (voir
lemme \ref{lemma-finite-pushforward-coherent}).
Considérons le diagramme commutatif
$$
\xymatrix{
Z' \ar[r]_{i'} \ar[d]_{f'} &
Y \ar[d]^f \\
Z \ar[r]^i &
X
}
$$
On a $f_*\mathcal{F} = i_*f'_*\mathcal{O}_{Z'}$.
La fibre de $f_*\mathcal{F}$ en $\xi$ est donc celle
de $f'_*\mathcal{O}_{Z'}$ en $\xi$. Puisque $Z'$ est
intègre de point générique $\xi'$, on voit que
$\xi'$ est le seul point de $Z'$ au-dessus de $\xi$, voir
Algèbre, lemmes \ref{algebra-lemma-finite-is-integral} et
\ref{algebra-lemma-integral-no-inclusion}.
La fibre de $f'_*\mathcal{O}_{Z'}$ en $\xi$
est donc égale à $\mathcal{O}_{Z', \xi'} = \kappa(\xi')$. En particulier,
la fibre de $f_*\mathcal{F}$ en $\xi$ est non nulle.
Cela, joint au fait que $f_*\mathcal{F}$ est
de la forme $i_*f'_*(\text{quelque chose})$, entraîne le lemme.
\end{proof}

\begin{lemma}
\label{lemma-affine-morphism-projection-ideal}
Soit $f : Y \to X$ un morphisme de schémas.
Soit $\mathcal{F}$ un faisceau quasi-cohérent sur $Y$.
Soit $\mathcal{I}$ un faisceau quasi-cohérent d'idéaux sur $X$.
Si le morphisme $f$ est affine, alors
$\mathcal{I}f_*\mathcal{F} = f_*(f^{-1}\mathcal{I}\mathcal{F})$.
\end{lemma}

\begin{proof}
Précisons cette notation. Puisque $f^{-1}$ est un
foncteur exact, $f^{-1}\mathcal{I}$ est un
faisceau d'idéaux de $f^{-1}\mathcal{O}_X$. Par le morphisme
$f^\sharp : f^{-1}\mathcal{O}_X \to \mathcal{O}_Y$, il opère sur
$\mathcal{F}$. Alors $f^{-1}\mathcal{I}\mathcal{F}$ est le sous-faisceau
engendré par les sommes de sections locales de la forme $as$, où $a$
est une section locale de $f^{-1}\mathcal{I}$ et $s$ une section locale
de $\mathcal{F}$. C'est un sous-$\mathcal{O}_Y$-module quasi-cohérent
de $\mathcal{F}$, car c'est aussi l'image d'un morphisme naturel
$f^*\mathcal{I} \otimes_{\mathcal{O}_Y} \mathcal{F} \to \mathcal{F}$.

\medskip\noindent
Cela étant, la démonstration est immédiate. La question est locale, et
on peut donc supposer $X$ affine. Puisque $f$ est affine,
$Y$ est lui aussi affine. On peut donc écrire
$Y = \Spec(B)$, $X = \Spec(A)$, $\mathcal{F} = \widetilde{M}$
et $\mathcal{I} = \widetilde{I}$. Dans ce cas, l'assertion du
lemme se ramène à l'égalité
$$
I(M_A) = ((IB)M)_A
$$
où $M_A$ désigne le $A$-module sous-jacent au $B$-module $M$.
\end{proof}

\begin{lemma}
\label{lemma-image-affine-finite-morphism-affine-Noetherian}
Soit $f : Y \to X$ un morphisme de schémas. Supposons
que
\begin{enumerate}
\item $f$ soit fini,
\item $f$ soit surjectif,
\item $Y$ soit affine, et
\item $X$ soit noethérien.
\end{enumerate}
Alors $X$ est affine.
\end{lemma}

\begin{proof}
Nous allons montrer que, sous les hypothèses du lemme, pour tout
$\mathcal{O}_X$-module cohérent $\mathcal{F}$, on a $H^1(X, \mathcal{F}) = 0$.
Cela entraînera en particulier que $H^1(X, \mathcal{I}) = 0$
pour tout faisceau quasi-cohérent d'idéaux de $\mathcal{O}_X$. On en
déduira que $X$ est affine, soit par le
lemme \ref{lemma-quasi-compact-h1-zero-covering}, soit
par le lemme \ref{lemma-quasi-separated-h1-zero-covering}.

\medskip\noindent
Soit $\mathcal{P}$ la propriété des faisceaux cohérents
$\mathcal{F}$ sur $X$ définie par
$$
\mathcal{P}(\mathcal{F}) \Leftrightarrow H^1(X, \mathcal{F}) = 0.
$$
Nous allons appliquer le lemme \ref{lemma-property-higher-rank-cohomological}.
Il faut donc vérifier les conditions (1), (2) et (3) de ce lemme pour $\mathcal{P}$.
La propriété (1) résulte de la suite exacte longue de cohomologie associée à
une suite exacte courte de faisceaux. La propriété (2) résulte de ce que
$H^1(X, -)$ est un foncteur additif. Pour vérifier (3), soit $Z \subset X$ un
sous-schéma fermé intègre de point générique $\xi$.
Soit $\mathcal{F}$ un faisceau cohérent sur $Y$ tel que
le support de $f_*\mathcal{F}$ soit égal à $Z$
et que $(f_*\mathcal{F})_\xi$ soit annulé par $\mathfrak m_\xi$,
voir lemme \ref{lemma-finite-morphism-Noetherian}. Nous affirmons
que $\mathcal{G} = f_*\mathcal{F}$ convient. Il suffit de vérifier la
condition (3)(c) du lemme \ref{lemma-property-higher-rank-cohomological}.
Supposons donc que $\mathcal{J} \subset \mathcal{O}_X$ soit un
faisceau quasi-cohérent d'idéaux tel que
$\mathcal{J}_\xi = \mathcal{O}_{X, \xi}$.
Un morphisme fini est affine; d'après le
lemme \ref{lemma-affine-morphism-projection-ideal}, on a donc
$\mathcal{J}\mathcal{G} = f_*(f^{-1}\mathcal{J}\mathcal{F})$.
De plus, comme on l'a remarqué dans la démonstration du
lemme \ref{lemma-affine-morphism-projection-ideal}, le faisceau
$f^{-1}\mathcal{J}\mathcal{F}$ est un $\mathcal{O}_Y$-module quasi-cohérent.
Puisque $Y$ est affine, on a $H^1(Y, f^{-1}\mathcal{J}\mathcal{F}) = 0$,
voir lemme \ref{lemma-quasi-coherent-affine-cohomology-zero}.
Puisque $f$ est fini, donc affine, on a
$$
H^1(X, \mathcal{J}\mathcal{G}) =
H^1(X, f_*(f^{-1}\mathcal{J}\mathcal{F})) =
H^1(Y, f^{-1}\mathcal{J}\mathcal{F}) = 0
$$
d'après le lemme \ref{lemma-relative-affine-cohomology}.
Le sous-faisceau quasi-cohérent $\mathcal{G}' = \mathcal{J}\mathcal{G}$
vérifie donc $\mathcal{P}$. Cela établit la propriété (3)(c) du
lemme \ref{lemma-property-higher-rank-cohomological}, comme voulu.
\end{proof}












\section{Faisceaux cohérents sur Proj, I}
\label{section-coherent-proj}

\noindent
Dans cette section, nous étudions les faisceaux cohérents sur $\text{Proj}(A)$,
où $A$ est un anneau gradué noethérien engendré par $A_1$ sur $A_0$.
Dans la section suivante, nous examinerons le cas où $A$ n'est pas engendré par ses
éléments de degré $1$. Commençons par un résultat d'ensemble pour
l'espace projectif sur un anneau noethérien.

\begin{lemma}
\label{lemma-coherent-projective}
Soit $R$ un anneau noethérien.
Soit $n \geq 0$ un entier. Pour
tout faisceau cohérent $\mathcal{F}$ sur $\mathbf{P}^n_R$,
on a les propriétés suivantes :
\begin{enumerate}
\item Il existe $r \geq 0$,
$d_1, \ldots, d_r \in \mathbf{Z}$ et une surjection
$$
\bigoplus\nolimits_{j = 1, \ldots, r}
\mathcal{O}_{\mathbf{P}^n_R}(d_j)
\longrightarrow
\mathcal{F}.
$$
\item On a $H^i(\mathbf{P}^n_R, \mathcal{F}) = 0$ sauf si
$0 \leq i \leq n$.
\item Pour tout $i$, le groupe de cohomologie $H^i(\mathbf{P}^n_R, \mathcal{F})$
est un $R$-module de type fini.
\item Si $i > 0$, alors
$H^i(\mathbf{P}^n_R, \mathcal{F}(d)) = 0$ pour tout $d$ assez grand.
\item Pour tout $k \in \mathbf{Z}$, le $R[T_0, \ldots, T_n]$-module gradué
$$
\bigoplus\nolimits_{d \geq k} H^0(\mathbf{P}^n_R, \mathcal{F}(d))
$$
est un $R[T_0, \ldots, T_n]$-module de type fini.
\end{enumerate}
\end{lemma}

\begin{proof}
Nous utiliserons le fait que $\mathcal{O}_{\mathbf{P}^n_R}(1)$ est un faisceau inversible
ample sur
le schéma $\mathbf{P}^n_R$. Cela résulte directement de la définition,
puisque $\mathbf{P}^n_R$ est recouvert par les ouverts affines standard $D_{+}(T_i)$.
Ainsi, d'après
Propriétés, proposition \ref{properties-proposition-characterize-ample},
tout $\mathcal{O}_{\mathbf{P}^n_R}$-module
quasi-cohérent de type fini est quotient d'une somme directe finie de puissances tensorielles de
$\mathcal{O}_{\mathbf{P}^n_R}(1)$. D'autre part, les faisceaux cohérents
coïncident avec les faisceaux quasi-cohérents de type fini sur l'espace
projectif sur $R$, d'après le lemme \ref{lemma-coherent-Noetherian}. D'où (1).

\medskip\noindent
L'espace projectif de dimension $n$, $\mathbf{P}^n_R$, est recouvert par $n + 1$ ouverts
affines, à savoir les ouverts standard $D_{+}(T_i)$, $i = 0, \ldots, n$, voir Constructions,
lemme \ref{constructions-lemma-standard-covering-projective-space}.
Ainsi, pour tout faisceau
quasi-cohérent $\mathcal{F}$ sur $\mathbf{P}^n_R$,
on a $H^i(\mathbf{P}^n_R, \mathcal{F}) = 0$ pour $i \geq n + 1$,
voir lemme \ref{lemma-vanishing-nr-affines}. D'où (2).

\medskip\noindent
Démontrons (3) et (4) simultanément pour tous les faisceaux cohérents
sur $\mathbf{P}^n_R$, par récurrence descendante sur $i$. Le résultat est
vrai pour $i \geq n + 1$ d'après (2). Supposons-le connu pour
$i + 1$ et montrons-le pour $i$. (Si $i = 0$, l'assertion
(4) est sans objet.) Soit $\mathcal{F}$ un faisceau cohérent sur
$\mathbf{P}^n_R$. Choisissons une surjection comme dans (1) et notons
$\mathcal{G}$ son noyau, de sorte que l'on ait une suite exacte courte
$$
0 \to \mathcal{G} \to
\bigoplus\nolimits_{j = 1, \ldots, r}
\mathcal{O}_{\mathbf{P}^n_R}(d_j)
\to
\mathcal{F} \to 0
$$
D'après le lemme \ref{lemma-coherent-abelian-Noetherian},
$\mathcal{G}$ est cohérent. La suite exacte
longue de cohomologie donne une suite exacte
$$
H^i(\mathbf{P}^n_R, \bigoplus\nolimits_{j = 1, \ldots, r}
\mathcal{O}_{\mathbf{P}^n_R}(d_j))
\to
H^i(\mathbf{P}^n_R, \mathcal{F})
\to
H^{i + 1}(\mathbf{P}^n_R, \mathcal{G}).
$$
Par hypothèse de récurrence, le $R$-module de droite est de type fini et, d'après
le lemme \ref{lemma-cohomology-projective-space-over-ring}, le
$R$-module de gauche est de type fini. Puisque $R$ est noethérien, il en résulte aussitôt
que $H^i(\mathbf{P}^n_R, \mathcal{F})$ est un $R$-module
de type fini. Cela établit l'étape de récurrence pour l'assertion (3).
Puisque $\mathcal{O}_{\mathbf{P}^n_R}(d)$ est
inversible, tordre un faisceau sur $\mathbf{P}^n_R$ définit un foncteur exact (car
ce foncteur est donné par le produit tensoriel avec un faisceau inversible, voir
Constructions, définition \ref{constructions-definition-twist}).
Autrement dit, pour tout $d \in \mathbf{Z}$, la suite
$$
0 \to \mathcal{G}(d) \to
\bigoplus\nolimits_{j = 1, \ldots, r}
\mathcal{O}_{\mathbf{P}^n_R}(d_j + d)
\to
\mathcal{F}(d) \to 0
$$
est exacte courte. La suite de cohomologie qui en résulte est
$$
H^i(\mathbf{P}^n_R, \bigoplus\nolimits_{j = 1, \ldots, r}
\mathcal{O}_{\mathbf{P}^n_R}(d_j + d))
\to
H^i(\mathbf{P}^n_R, \mathcal{F}(d))
\to
H^{i + 1}(\mathbf{P}^n_R, \mathcal{G}(d)).
$$
Par hypothèse de récurrence, le module de droite est nul
pour $d \gg 0$ et, d'après le calcul du
lemme \ref{lemma-cohomology-projective-space-over-ring},
le module de gauche est nul dès que $d \geq -\min\{d_j\}$
et $i \geq 1$. D'où l'étape de récurrence pour l'assertion (4).
Cela achève la démonstration de (3) et (4).

\medskip\noindent
Pour démontrer (5), notons que, pour tout $d$
assez grand, l'application
$$
H^0(\mathbf{P}^n_R, \bigoplus\nolimits_{j = 1, \ldots, r}
\mathcal{O}_{\mathbf{P}^n_R}(d_j + d))
\to
H^0(\mathbf{P}^n_R, \mathcal{F}(d))
$$
est surjective, d'après l'annulation de $H^1(\mathbf{P}^n_R, \mathcal{G}(d))$
que nous venons de démontrer. Autrement dit, le module
$$
M_k
=
\bigoplus\nolimits_{d \geq k} H^0(\mathbf{P}^n_R, \mathcal{F}(d))
$$
est, pour $k$ assez grand, un quotient du module correspondant
$$
N_k
=
\bigoplus\nolimits_{d \geq k} H^0(\mathbf{P}^n_R,
\bigoplus\nolimits_{j = 1, \ldots, r}
\mathcal{O}_{\mathbf{P}^n_R}(d_j + d)
)
$$
Lorsque $k$ est assez petit (par exemple\ $k < -d_j$ pour tout $j$), on a
$$
N_k = \bigoplus\nolimits_{j = 1, \ldots, r}
R[T_0, \ldots, T_n](d_j)
$$
d'après les calculs de la section \ref{section-cohomology-projective-space}.
Il est donc en particulier de type fini.
Soit $k \in \mathbf{Z}$ quelconque.
Choisissons $k_{-} \ll k \ll k_{+}$.
Considérons le diagramme
$$
\xymatrix{
N_{k_{-}} & N_{k_{+}} \ar[d] \ar[l] \\
M_k & M_{k_{+}} \ar[l]
}
$$
où la flèche verticale est l'application surjective ci-dessus
et où les flèches horizontales sont les
inclusions évidentes. D'après ce qui précède, $N_{k_{-}}$ est
un $R[T_0, \ldots, T_n]$-module de type fini. Ainsi $N_{k_{+}}$ est
un $R[T_0, \ldots, T_n]$-module de type
fini, car c'est un sous-module d'un module de type fini et l'anneau
$R[T_0, \ldots, T_n]$ est noethérien. La flèche
verticale étant surjective, $M_{k_{+}}$ est
un $R[T_0, \ldots, T_n]$-module de type fini. Le quotient
$M_k/M_{k_{+}}$ est de type fini comme $R$-module, car c'est une
somme directe finie des $R$-modules de type fini
$H^0(\mathbf{P}^n_R, \mathcal{F}(d))$ pour $k \leq d < k_{+}$.
On utilise ici l'assertion (3) pour $i = 0$. Ainsi
$M_k/M_{k_{+}}$ est a fortiori un $R[T_0, \ldots, T_n]$-module
de type fini. Autrement dit, $M_k$ est extension de deux
$R[T_0, \ldots, T_n]$-modules de type fini, ce qui conclut.
\end{proof}

\begin{lemma}
\label{lemma-coherent-on-proj}
Soit $A$ un anneau gradué tel que $A_0$ soit noethérien et que
$A$ soit engendré par un nombre fini d'éléments de $A_1$ sur $A_0$.
Posons $X = \text{Proj}(A)$. Alors $X$ est un schéma noethérien.
Soit $\mathcal{F}$ un $\mathcal{O}_X$-module cohérent.
\begin{enumerate}
\item Il existe $r \geq 0$,
$d_1, \ldots, d_r \in \mathbf{Z}$ et une surjection
$$
\bigoplus\nolimits_{j = 1, \ldots, r} \mathcal{O}_X(d_j)
\longrightarrow \mathcal{F}.
$$
\item Pour tout $p$, le groupe de cohomologie $H^p(X, \mathcal{F})$ est un
$A_0$-module de type fini.
\item Si $p > 0$, alors $H^p(X, \mathcal{F}(d)) = 0$ pour tout $d$ assez grand.
\item Pour tout $k \in \mathbf{Z}$, le $A$-module gradué
$$
\bigoplus\nolimits_{d \geq k} H^0(X, \mathcal{F}(d))
$$
est un $A$-module de type fini.
\end{enumerate}
\end{lemma}

\begin{proof}
Par hypothèse, il existe une surjection de $A_0$-algèbres graduées
$$
A_0[T_0, \ldots, T_n] \longrightarrow A
$$
où $\deg(T_j) = 1$ pour $j = 0, \ldots, n$. D'après Constructions, Lemme
\ref{constructions-lemma-surjective-graded-rings-generated-degree-1-map-proj},
elle définit une immersion fermée $i : X \to \mathbf{P}^n_{A_0}$
telle que $i^*\mathcal{O}_{\mathbf{P}^n_{A_0}}(1) = \mathcal{O}_X(1)$.
En particulier, $X$ est noethérien, étant un sous-schéma fermé du schéma
noethérien $\mathbf{P}^n_{A_0}$. Nous affirmons que les résultats du lemme pour
$\mathcal{F}$ se déduisent des résultats
correspondants du lemme \ref{lemma-coherent-projective} pour le faisceau cohérent
$i_*\mathcal{F}$ (lemme \ref{lemma-i-star-equivalence}) sur
$\mathbf{P}^n_{A_0}$. Par exemple, ce lemme assure
l'existence d'une surjection
$$
\bigoplus\nolimits_{j = 1, \ldots, r} \mathcal{O}_{\mathbf{P}^n_{A_0}}(d_j)
\longrightarrow i_*\mathcal{F}.
$$
Par adjonction, celle-ci correspond à une application
$\bigoplus_{j = 1, \ldots, r} \mathcal{O}_X(d_j) \longrightarrow \mathcal{F}$,
qui est elle aussi surjective. Les assertions sur la cohomologie résultent
de l'égalité
$H^p(X, \mathcal{F}(d)) = H^p(\mathbf{P}^n_{A_0}, i_*\mathcal{F}(d))$,
d'après le lemme \ref{lemma-relative-affine-cohomology}.
\end{proof}

\begin{lemma}
\label{lemma-recover-tail-graded-module}
Soit $A$ un anneau gradué tel que $A_0$ soit noethérien et que $A$ soit engendré
par un nombre fini d'éléments de $A_1$ sur $A_0$. Soit $M$
un $A$-module gradué de type fini. Posons $X = \text{Proj}(A)$ et soit $\widetilde{M}$
le $\mathcal{O}_X$-module quasi-cohérent sur $X$ associé à $M$.
Les applications
$$
M_n \longrightarrow \Gamma(X, \widetilde{M}(n))
$$
de Constructions, lemme \ref{constructions-lemma-apply-modules},
sont des isomorphismes pour tout $n$ assez grand.
\end{lemma}

\begin{proof}
Puisque $M$ est un $A$-module de type fini,
$\widetilde{M}$ est un $\mathcal{O}_X$-module de
type fini, c'est-à-dire un $\mathcal{O}_X$-module cohérent.
Posons $N = \bigoplus_{n \in \mathbf{Z}} \Gamma(X, \widetilde{M}(n))$.
Il faut montrer que l'application $M \to N$ de $A$-modules
gradués est un isomorphisme en tout degré assez grand.
D'après Propriétés, lemme \ref{properties-lemma-proj-quasi-coherent},
on a un isomorphisme canonique $\widetilde{N} \to \widetilde{M}$
tel que les applications induites $N_n \to N_n = \Gamma(X, \widetilde{M}(n))$
soient les identités. On a donc des applications
$\widetilde{M} \to \widetilde{N} \to \widetilde{M}$
telles que, pour tout $n$, le diagramme
$$
\xymatrix{
M_n \ar[d] \ar[r] & N_n \ar[d] \ar@{=}[rd] \\
\Gamma(X, \widetilde{M}(n)) \ar[r] &
\Gamma(X, \widetilde{N}(n)) \ar[r]^{\cong} &
\Gamma(X, \widetilde{M}(n))
}
$$
soit commutatif. Cela signifie que le composé
$$
M_n \to \Gamma(X, \widetilde{M}(n))
\to \Gamma(X, \widetilde{N}(n)) \to \Gamma(X, \widetilde{M}(n))
$$
est l'application canonique $M_n \to \Gamma(X, \widetilde{M}(n))$.
Cela entraîne clairement que le composé
$\widetilde{M} \to \widetilde{N} \to \widetilde{M}$
est l'identité. Ainsi $\widetilde{M} \to \widetilde{N}$ est un isomorphisme.
Posons $K = \Ker(M \to N)$ et $Q = \Coker(M \to N)$.
Rappelons que le foncteur $M \mapsto \widetilde{M}$ est exact, voir
Constructions, lemme \ref{constructions-lemma-proj-sheaves}.
On a donc $\widetilde{K} = 0$ et $\widetilde{Q} = 0$.
Rappelons que $A$ est un anneau noethérien, que $M$ est un
$A$-module de type fini et que $N$ est un $A$-module gradué tel que
$N' = \bigoplus_{n \geq 0} N_n$ soit de type fini,
d'après la dernière assertion du lemme \ref{lemma-coherent-on-proj}.
Ainsi $K' = \bigoplus_{n \geq 0} K_n$ et
$Q' = \bigoplus_{n \geq 0} Q_n$ sont des $A$-modules de
type fini. Notons que $\widetilde{Q} = \widetilde{Q'}$ et
$\widetilde{K} = \widetilde{K'}$.
Pour achever la démonstration, il suffit donc
de montrer qu'un $A$-module de type fini $K$ tel que $\widetilde{K} = 0$
n'a qu'un nombre fini de composantes homogènes non nulles $K_d$ avec $d \geq 0$.
Pour cela, soient $x_1, \ldots, x_r \in K$ des générateurs homogènes,
de degrés respectifs $d_1, \ldots, d_r$.
Soient $f_1, \ldots, f_n \in A_1$ des éléments engendrant $A$ sur $A_0$.
Pour tout $i$ et tout $j$, il existe $n_{ij} \geq 0$ tel que
$f_i^{n_{ij}} x_j = 0$ dans $K_{d_j + n_{ij}}$ : sinon
$x_i/f_i^{d_i} \in K_{(f_i)}$ serait non nul, donc $\widetilde{K}$
serait non
nul. Alors $K_d$ est nul pour $d > \max_j(d_j + \sum_i n_{ij})$,
car tout élément de $K_d$ est une somme de termes, chacun étant un
monôme en les $f_i$ multiplié par l'un des $x_j$, de degré total $d$.
\end{proof}

\noindent
Soit $A$ un anneau gradué tel que $A_0$ soit noethérien et que $A$ soit engendré
par un nombre fini d'éléments de $A_1$ sur $A_0$. Rappelons que
$A_+ = \bigoplus_{n > 0} A_n$ est l'idéal des éléments de degré strictement
positif. Soit $M$ un $A$-module gradué. Rappelons que
$M$ est un module de torsion pour les puissances de $A_+$ si, pour tout $x \in M$,
il existe $n \geq 1$ tel que $(A_+)^n x = 0$, voir
Compléments d'algèbre, définition \ref{more-algebra-definition-f-power-torsion}.
Si $M$ est de type fini,
cela équivaut à $M_n = 0$ pour $n \gg 0$. Les modules de
torsion pour les puissances de $A_+$ sont parfois appelés modules de torsion. On
dit parfois qu'un $A$-module gradué $M$ est sans torsion si
$x \in M$ avec $(A_+)^n x = 0$, $n > 0$, entraîne $x = 0$.
Notons $\text{Mod}_A$ la catégorie des $A$-modules gradués,
$\text{Mod}^{fg}_A$ la sous-catégorie pleine des modules de type fini,
et $\text{Mod}^{fg}_{A, torsion}$ la sous-catégorie pleine des
modules $M$ tels que $M_n = 0$ pour $n \gg 0$.

\begin{proposition}
\label{proposition-coherent-modules-on-proj}
Soit $A$ un anneau gradué tel que $A_0$ soit noethérien et que $A$ soit engendré
par un nombre fini d'éléments de $A_1$ sur $A_0$.
Posons $X = \text{Proj}(A)$. Le foncteur $M \mapsto \widetilde M$
induit une équivalence
$$
\text{Mod}^{fg}_A/\text{Mod}^{fg}_{A, torsion}
\longrightarrow
\textit{Coh}(\mathcal{O}_X)
$$
dont un quasi-inverse est donné par
$\mathcal{F} \longmapsto \bigoplus_{n \geq 0} \Gamma(X, \mathcal{F}(n))$.
\end{proposition}

\begin{proof}
La sous-catégorie $\text{Mod}^{fg}_{A, torsion}$ est une sous-catégorie de Serre
de $\text{Mod}^{fg}_A$, voir
Homologie, définition \ref{homology-definition-serre-subcategory}.
Cela est clair d'après la description des objets donnée ci-dessus, mais résulte aussi
de Compléments d'algèbre, lemme \ref{more-algebra-lemma-I-power-torsion}.
La catégorie quotient à gauche de la flèche est donc définie
dans Homologie, lemme \ref{homology-lemma-serre-subcategory-is-kernel}.
Pour définir le foncteur de la proposition, il suffit de montrer que
le foncteur $M \mapsto \widetilde M$ envoie les modules de torsion sur $0$.
Cela est clair, car pour tout élément homogène $f \in A_+$, le
module $M_f$ est nul; la valeur $M_{(f)}$ de $\widetilde M$
sur $D_+(f)$ est donc elle aussi nulle.

\medskip\noindent
D'après le lemme \ref{lemma-coherent-on-proj}, le quasi-inverse proposé
est bien défini. En effet, le lemme montre que
$\mathcal{F} \longmapsto \bigoplus_{n \geq 0} \Gamma(X, \mathcal{F}(n))$
est un foncteur $\textit{Coh}(\mathcal{O}_X) \to \text{Mod}^{fg}_A$,
que l'on peut composer avec le foncteur quotient
$\text{Mod}^{fg}_A \to \text{Mod}^{fg}_A/\text{Mod}^{fg}_{A, torsion}$.

\medskip\noindent
D'après le lemme \ref{lemma-recover-tail-graded-module},
le composé de gauche à droite puis de droite à gauche est isomorphe au foncteur identité.

\medskip\noindent
Enfin, soit $\mathcal{F}$ un $\mathcal{O}_X$-module cohérent.
Posons $M = \bigoplus_{n \in \mathbf{Z}} \Gamma(X, \mathcal{F}(n))$,
considéré comme un $A$-module gradué, de sorte que notre foncteur envoie $\mathcal{F}$ sur
$M_{\geq 0} = \bigoplus_{n \geq 0} M_n$.
D'après Propriétés, lemme \ref{properties-lemma-proj-quasi-coherent},
l'application canonique $\widetilde M \to \mathcal{F}$
est un isomorphisme. Puisque l'inclusion
$M_{\geq 0} \to M$ définit un isomorphisme
$\widetilde{M_{\geq 0}} \to \widetilde M$, on en déduit
que le composé de droite à gauche puis de gauche à droite est lui aussi
isomorphe au foncteur identité.
\end{proof}





\section{Faisceaux cohérents sur Proj, II}
\label{section-coherent-proj-general}

\noindent
Dans cette section, nous étudions les faisceaux cohérents sur $\text{Proj}(A)$,
où $A$ est un anneau gradué noethérien. La plupart des résultats
se déduisent astucieusement du résultat correspondant de la section précédente,
où nous avons traité le cas où $A$ est engendré par des
éléments de degré $1$.

\begin{lemma}
\label{lemma-coherent-on-proj-general}
Soit $A$ un anneau gradué noethérien. Posons $X = \text{Proj}(A)$. Alors $X$
est un schéma noethérien. Soit $\mathcal{F}$ un $\mathcal{O}_X$-module cohérent.
\begin{enumerate}
\item Il existe $r \geq 0$, des entiers
$d_1, \ldots, d_r \in \mathbf{Z}$ et une surjection
$$
\bigoplus\nolimits_{j = 1, \ldots, r} \mathcal{O}_X(d_j)
\longrightarrow \mathcal{F}.
$$
\item Pour tout $p$, le groupe de cohomologie $H^p(X, \mathcal{F})$ est un
$A_0$-module de type fini.
\item Si $p > 0$, alors $H^p(X, \mathcal{F}(d)) = 0$ pour tout $d$ suffisamment grand.
\item Pour tout $k \in \mathbf{Z}$, le $A$-module gradué
$$
\bigoplus\nolimits_{d \geq k} H^0(X, \mathcal{F}(d))
$$
est un $A$-module de type fini.
\end{enumerate}
\end{lemma}

\begin{proof}
Nous ramenons l'énoncé au
lemme \ref{lemma-coherent-on-proj}.
D'après Algèbre, lemmes \ref{algebra-lemma-graded-Noetherian} et
\ref{algebra-lemma-S-plus-generated}, l'anneau $A_0$ est noethérien
et $A$ est engendré sur $A_0$ par un nombre fini d'éléments
$f_1, \ldots, f_r$ homogènes de degré strictement positif.
Soit $d$ un entier suffisamment divisible. Posons $A' = A^{(d)}$, avec
les notations d'Algèbre, section \ref{algebra-section-graded}.
Alors $A'$ est engendré sur $A'_0 = A_0$ par des éléments de
degré $1$ ; voir Algèbre, lemme \ref{algebra-lemma-uple-generated-degree-1}.
Le lemme \ref{lemma-coherent-on-proj} s'applique donc à $X' = \text{Proj}(A')$.

\medskip\noindent
Par Constructions, lemme \ref{constructions-lemma-d-uple}, il
existe un isomorphisme de schémas $i : X \to X'$ et des
isomorphismes $\mathcal{O}_X(nd) \to i^*\mathcal{O}_{X'}(n)$
compatibles avec le morphisme $A' \to A$ et les morphismes
$A_n \to H^0(X, \mathcal{O}_X(n))$ et $A'_n \to H^0(X', \mathcal{O}_{X'}(n))$.
Le lemme \ref{lemma-coherent-on-proj} implique donc que $X$ est noethérien et que
(1) et (2) sont satisfaits. Pour établir (3)
et (4), on utilise que, pour $k$, $p$ et $q$ fixés, on a
$$
\bigoplus\nolimits_{dn + q \geq k} H^p(X, \mathcal{F}(dn + q)) =
\bigoplus\nolimits_{dn + q \geq k} H^p(X', (i_*\mathcal{F}(q))(n)
$$
par les compatibilités précédentes. Si $p > 0$, le membre de droite
s'annule pour $k$ suffisamment grand en fonction de $q$, d'après le
lemme \ref{lemma-coherent-on-proj}. Comme il n'existe qu'un nombre fini
de classes de congruence d'entiers modulo $d$, on voit que (3) est satisfait pour
$\mathcal{F}$ sur $X$. Si $p = 0$, le membre de droite
est un $A'$-module de type fini d'après le lemme \ref{lemma-coherent-on-proj}. En
utilisant à nouveau la finitude des classes de congruence, on obtient que
$\bigoplus_{n \geq k} H^0(X, \mathcal{F}(n))$ est lui aussi un $A'$-module de type fini.
Comme la structure de $A'$-module provient de la structure de $A$-module (par
les compatibilités mentionnées plus haut), on conclut qu'il est également de type
fini comme $A$-module.
\end{proof}

\begin{lemma}
\label{lemma-recover-tail-graded-module-general}
Soit $A$ un anneau gradué noethérien et soit $d$ le ppcm des degrés d'un système de générateurs
de $A$ sur $A_0$. Soit $M$ un $A$-module gradué de
type fini. Posons $X = \text{Proj}(A)$ et soit $\widetilde{M}$
le $\mathcal{O}_X$-module quasi-cohérent sur $X$ associé à $M$.
Soit $k \in \mathbf{Z}$.
\begin{enumerate}
\item $N' = \bigoplus_{n \geq k} H^0(X, \widetilde{M(n)})$
est un $A$-module de type fini,
\item $N = \bigoplus_{n \geq k} H^0(X, \widetilde{M}(n))$
est un $A$-module de type fini,
\item il existe un morphisme canonique $N \to N'$,
\item si $k$ est suffisamment petit, il existe un morphisme canonique $M \to N'$,
\item le morphisme $M_n \to N'_n$ est un isomorphisme pour $n \gg 0$,
\item $N_n \to N'_n$ est un isomorphisme si $d | n$.
\end{enumerate}
\end{lemma}

\begin{proof}
Le morphisme $N \to N'$ de (3) provient de
Constructions, Équation (\ref{constructions-equation-multiply-more-generally}),
en prenant les sections globales.

\medskip\noindent
Par
Constructions, Équation
(\ref{constructions-equation-global-sections-more-generally}),
il existe un morphisme de $A$-modules gradués
$M \to \bigoplus_{n \in \mathbf{Z}} H^0(X, \widetilde{M(n)})$.
Si les générateurs de $M$ sont de degrés $\geq k$, alors l'image est
contenue dans le sous-module
$N' \subset \bigoplus_{n \in \mathbf{Z}} H^0(X, \widetilde{M(n)})$,
ce qui donne le morphisme de (4).

\medskip\noindent
Par Algèbre, lemmes \ref{algebra-lemma-graded-Noetherian} et
\ref{algebra-lemma-S-plus-generated}, l'anneau $A_0$ est noethérien
et $A$ est engendré sur $A_0$ par un nombre fini d'éléments
$f_1, \ldots, f_r$ homogènes de degré strictement positif.
Posons $d = \text{lcm}(\deg(f_i))$. Alors (6) résulte, par
exemple, de Constructions,
lemme \ref{constructions-lemma-where-invertible}.

\medskip\noindent
Comme $M$ est un $A$-module de type fini,
$\widetilde{M}$ est un $\mathcal{O}_X$-module de
type fini, c'est-à-dire un $\mathcal{O}_X$-module
cohérent. L'assertion (2) résulte donc du lemme \ref{lemma-coherent-on-proj-general}.

\medskip\noindent
Nous déduisons (1) de (2) par l'astuce suivante. Pour $q \in \{0, \ldots, d - 1\}$,
posons
$$
{}^qN = \bigoplus\nolimits_{n + q \geq k} H^0(X, \widetilde{M(q)}(n))
$$
Ce sont des $A$-modules de type fini par (2). L'anneau noethérien $A$
est fini sur $A^{(d)} = \bigoplus_{n \geq 0} A_{dn}$, car
il est engendré par les $f_i$ sur $A^{(d)}$ et $f_i^d \in A^{(d)}$. Par
conséquent, ${}^qN$ est un $A^{(d)}$-module de type
fini. De plus, $A^{(d)}$ est noethérien (cela résulte
d'Algèbre, lemme \ref{algebra-lemma-dehomogenize-finite-type}).
Il s'ensuit que le sous-$A^{(d)}$-module
${}^qN^{(d)} = \bigoplus_{n \in \mathbf{Z}} {}^qN_{dn}$
est un module de type fini sur $A^{(d)}$. À l'aide des isomorphismes
$\widetilde{M(dn + q)} = \widetilde{M(q)}(dn)$, on peut écrire
$$
N' =
\bigoplus\nolimits_{q \in \{0, \ldots, d - 1\}}
\bigoplus\nolimits_{dn + q \geq k}
H^0(X, \widetilde{M(q)}(dn)) =
\bigoplus\nolimits_{q \in \{0, \ldots, d - 1\}} {}^qN^{(d)}
$$
Ainsi, $N'$ est de type fini sur $A^{(d)}$ et, a fortiori, de type fini sur $A$.
Cela démontre (1).

\medskip\noindent
Soit $K$ un $A$-module de type fini tel que $\widetilde{K} = 0$. Nous affirmons
que $K_n = 0$ si $d|n$ et $n \gg 0$. En raisonnant comme ci-dessus, on voit que
$K^{(d)}$ est un $A^{(d)}$-module de type fini. Soient
$x_1, \ldots, x_m \in K$ des générateurs homogènes de $K^{(d)}$
sur $A^{(d)}$, de degrés respectifs $d_1, \ldots, d_m$, avec $d | d_j$.
Pour tout $i$ et tout $j$, il existe $n_{ij} \geq 0$ tel que
$f_i^{n_{ij}} x_j = 0$ dans $K_{d_j + n_{ij}}$ : sinon,
$x_j/f_i^{d_i/\deg(f_i)} \in K_{(f_i)}$ serait non nul, c'est-à-dire que
$\widetilde{K}$ serait non nul. Nous utilisons ici $\deg(f_i) | d | d_j$
pour tous $i, j$. On conclut que $K_n$ est nul pour tout
$n$ tel que $d | n$ et $n > \max_j (d_j + \sum_i n_{ij} \deg(f_i))$,
car tout élément de $K_n$ est une somme de termes, chacun étant le produit d'un
monôme en les $f_i$ par l'un des $x_j$, de degré total $n$.

\medskip\noindent
Pour achever la démonstration, il reste à montrer que $M \to N'$ est un isomorphisme
en tout degré suffisamment grand. Le morphisme $N \to N'$ induit un
isomorphisme $\widetilde{N} \to \widetilde{N'}$, car les modules correspondants sont isomorphes sur les
ouverts affines $D_+(f_i) = D_+(f_i^d)$ :
$N_{(f_i)} \cong N_{(f_i^d)} \cong N'_{(f_i^d)} \cong N'_{(f_i)}$
par la propriété (6).
D'après Propriétés, lemme \ref{properties-lemma-proj-quasi-coherent},
on a un isomorphisme canonique $\widetilde{N} \to \widetilde{M}$.
Le composé $\widetilde{N} \to \widetilde{M} \to \widetilde{N'}$
est l'isomorphisme précédent (démonstration omise ; indication : vérifier
sur les ouverts affines standards). Ainsi, le morphisme $M \to N'$ induit un isomorphisme
$\widetilde{M} \to \widetilde{N'}$.
Posons $K = \Ker(M \to N')$ et $Q = \Coker(M \to N')$.
Rappelons que le foncteur
$M \mapsto \widetilde{M}$ est exact ; voir
Constructions, lemme \ref{constructions-lemma-proj-sheaves}.
On a donc $\widetilde{K} = 0$ et $\widetilde{Q} = 0$.
D'après le paragraphe précédent, $K_n = 0$ et
$Q_n = 0$ si $d | n$ et $n \gg 0$. On voit ici l'avantage
de l'emploi de $N'$ plutôt que de $N$ : le foncteur $M \leadsto N'$
est compatible aux décalages (cela résulte immédiatement de la construction).
En répétant tout l'argument après avoir remplacé $M$ par $M(q)$,
on obtient $K_n = 0$ et $Q_n = 0$ si $n \equiv q \bmod d$
et $n \gg 0$. Comme il n'existe qu'un
nombre fini de classes de congruence modulo $n$, la démonstration est achevée.
\end{proof}

\noindent
Soit $A$ un anneau gradué noethérien. Rappelons que
$A_+ = \bigoplus_{n > 0} A_n$ est l'idéal non pertinent.
D'après Algèbre, lemmes \ref{algebra-lemma-graded-Noetherian} et
\ref{algebra-lemma-S-plus-generated}, l'anneau $A_0$ est noethérien
et $A$ est engendré sur $A_0$ par un nombre fini d'éléments
$f_1, \ldots, f_r$ homogènes de degré strictement positif.
Posons $d = \text{lcm}(\deg(f_i))$. Soit $M$ un $A$-module
gradué. Dans cette situation, nous dirons qu'un élément homogène $x \in M$
est {\it non pertinent}\footnote{Cette terminologie n'est pas standard.} si
$$
(A_+ x)_{nd} = 0\text{ pour tout }n \gg 0
$$
Si $x \in M$ est homogène et non pertinent et si $f \in A$
est homogène, alors $fx$ est encore non pertinent.
Par conséquent, les éléments non pertinents engendrent un sous-module gradué
$M_{irrelevant} \subset M$. Nous dirons que $M$ est {\it non pertinent}
si tout élément homogène de $M$ est non pertinent, autrement
dit si $M_{irrelevant} = M$.
Lorsque $M$ est de type fini,
cela équivaut à $M_{nd} = 0$ pour $n \gg 0$.
Notons $\text{Mod}_A$ la catégorie des $A$-modules gradués,
$\text{Mod}^{fg}_A$ sa sous-catégorie pleine formée de ceux qui sont de type
fini et $\text{Mod}^{fg}_{A, irrelevant}$ la sous-catégorie pleine des modules
non pertinents.

\begin{proposition}
\label{proposition-coherent-modules-on-proj-general}
Soit $A$ un anneau gradué noethérien. Posons $X = \text{Proj}(A)$. Le foncteur
$M \mapsto \widetilde M$ induit une équivalence
$$
\text{Mod}^{fg}_A/\text{Mod}^{fg}_{A, irrelevant}
\longrightarrow
\textit{Coh}(\mathcal{O}_X)
$$
dont un quasi-inverse est donné par
$\mathcal{F} \longmapsto \bigoplus_{n \geq 0} \Gamma(X, \mathcal{F}(n))$.
\end{proposition}

\begin{proof}
Nous recommandons au lecteur de commencer par lire la démonstration dans le cas où $A$
est engendré en degré $1$ ; voir la
proposition \ref{proposition-coherent-modules-on-proj}.
Soient $f_1, \ldots, f_r \in A$ des éléments homogènes de degré strictement positif
qui engendrent $A$ sur $A_0$. Soit $d$ le ppcm des degrés
$d_i$ des $f_i$. Soit $M$ un $A$-module
de type fini. Montrons que $\widetilde{M}$ est nul si et
seulement si $M$ est un $A$-module gradué non pertinent (au
sens défini avant l'énoncé de la proposition). Soit donc $x \in M$ un élément
homogène.
Choisissons $k \in \mathbf{Z}$ assez petit, et soient
$N \to N'$ et $M \to N'$ comme dans le
lemme \ref{lemma-recover-tail-graded-module-general}.
On peut aussi choisir $l$ assez grand pour que
$M_n \to N_n$ soit un isomorphisme pour $n \geq l$.
Si $\widetilde{M}$ est nul, alors $N = 0$.
Ainsi, pour tout $f \in A_+$ homogène tel que
$\deg(f) + \deg(x) = nd$ et $nd > l$, l'élément $fx$ est nul,
car $N_{nd} \to N'_{nd}$ et $M_{nd} \to N'_{nd}$ sont des isomorphismes.
Donc $x$ est non
pertinent. Réciproquement, supposons $M$ non pertinent. Alors $M_{nd}$ est nul
pour $n \gg 0$ (voir la discussion précédant la
proposition). Cela entraîne manifestement que $M_{(f_i)} = M_{(f_i^{d/\deg(f_i)})} = 0$,
d'où $\widetilde{M} = 0$ par construction.

\medskip\noindent
Il en résulte que la sous-catégorie $\text{Mod}^{fg}_{A, irrelevant}$
est une sous-catégorie de Serre de $\text{Mod}^{fg}_A$, puisqu'elle est
le noyau du foncteur exact $M \mapsto \widetilde M$ ; voir
Homologie, lemme \ref{homology-lemma-kernel-exact-functor},
et Constructions, lemme \ref{constructions-lemma-proj-sheaves}.
La catégorie quotient figurant à gauche de la flèche est donc celle définie
dans Homologie, lemme \ref{homology-lemma-serre-subcategory-is-kernel}.
Pour définir le foncteur de la proposition, il suffit de montrer que
le foncteur $M \mapsto \widetilde M$ envoie les modules non pertinents sur $0$,
ce que nous avons établi ci-dessus.

\medskip\noindent
Le quasi-inverse proposé est bien défini d'après le lemme \ref{lemma-coherent-on-proj-general}. En
effet, ce lemme montre que
$\mathcal{F} \longmapsto \bigoplus_{n \geq 0} \Gamma(X, \mathcal{F}(n))$
est un foncteur $\textit{Coh}(\mathcal{O}_X) \to \text{Mod}^{fg}_A$
que l'on peut composer avec le foncteur quotient
$\text{Mod}^{fg}_A \to \text{Mod}^{fg}_A/\text{Mod}^{fg}_{A, irrelevant}$.

\medskip\noindent
D'après le lemme \ref{lemma-recover-tail-graded-module-general},
le composé de gauche à droite puis de droite à gauche est isomorphe au foncteur identité.
En effet, soit $M$ un $A$-module gradué de type fini. Choisissons
$k \in \mathbf{Z}$ assez petit, et soient
$N \to N'$ et $M \to N'$ comme dans le
lemme \ref{lemma-recover-tail-graded-module-general}.
Le noyau et le conoyau de $M \to N'$ ne sont alors non
nuls que dans un nombre fini de degrés, et sont donc non pertinents.
De plus, le noyau et le conoyau de $N \to N'$ sont nuls en tout degré
assez grand divisible par $d$ ; ce sont donc eux aussi des modules non pertinents.
Ainsi $M \to N'$ et $N \to N'$ sont tous deux des isomorphismes dans la
catégorie quotient, comme voulu.

\medskip\noindent
Enfin, soit $\mathcal{F}$ un $\mathcal{O}_X$-module cohérent.
Posons $M = \bigoplus_{n \in \mathbf{Z}} \Gamma(X, \mathcal{F}(n))$
et considérons-le comme un $A$-module gradué, de sorte que notre foncteur envoie $\mathcal{F}$ sur
$M_{\geq 0} = \bigoplus_{n \geq 0} M_n$.
D'après Propriétés, lemme \ref{properties-lemma-proj-quasi-coherent},
le morphisme canonique $\widetilde M \to \mathcal{F}$
est un isomorphisme. Comme l'inclusion
$M_{\geq 0} \to M$ définit un isomorphisme
$\widetilde{M_{\geq 0}} \to \widetilde M$, on en conclut
que le composé de droite à gauche puis de gauche à droite est lui aussi
isomorphe au foncteur identité.
\end{proof}










\section{Images directes supérieures par les morphismes projectifs}
\label{section-projective-pushforward}

\noindent
Nous énonçons et démontrons d'abord un résultat lorsque la base est affine, puis
nous en déduisons quelques résultats sur les morphismes projectifs.

\begin{lemma}
\label{lemma-coherent-proper-ample}
Soit $R$ un anneau noethérien. Soit $X \to \Spec(R)$ un morphisme propre.
Soit $\mathcal{L}$ un faisceau inversible ample sur $X$. Soit $\mathcal{F}$
un $\mathcal{O}_X$-module cohérent.
\begin{enumerate}
\item L'anneau gradué
$A = \bigoplus_{d \geq 0} H^0(X, \mathcal{L}^{\otimes d})$
est une $R$-algèbre de type fini.
\item Il existe $r \geq 0$, des entiers
$d_1, \ldots, d_r \in \mathbf{Z}$ et une surjection
$$
\bigoplus\nolimits_{j = 1, \ldots, r} \mathcal{L}^{\otimes d_j}
\longrightarrow \mathcal{F}.
$$
\item Pour tout $p$, le groupe de cohomologie $H^p(X, \mathcal{F})$ est un
$R$-module de type fini.
\item Si $p > 0$, alors
$H^p(X, \mathcal{F} \otimes_{\mathcal{O}_X} \mathcal{L}^{\otimes d}) = 0$
pour tout $d$ assez grand.
\item Pour tout $k \in \mathbf{Z}$, le $A$-module gradué
$$
\bigoplus\nolimits_{d \geq k}
H^0(X, \mathcal{F} \otimes_{\mathcal{O}_X} \mathcal{L}^{\otimes d})
$$
est un $A$-module de type fini.
\end{enumerate}
\end{lemma}

\begin{proof}
D'après
Morphismes, lemme \ref{morphisms-lemma-finite-type-over-affine-ample-very-ample},
il existe $d > 0$ et une immersion $i : X \to \mathbf{P}^n_R$
telles que $\mathcal{L}^{\otimes d} \cong i^*\mathcal{O}_{\mathbf{P}^n_R}(1)$.
Puisque $X$ est propre sur $R$, le morphisme $i$ est une immersion fermée
(Morphismes, lemme \ref{morphisms-lemma-image-proper-scheme-closed}).
On a donc $H^i(X, \mathcal{G}) = H^i(\mathbf{P}^n_R, i_*\mathcal{G})$
pour tout faisceau quasi-cohérent $\mathcal{G}$ sur $X$
(d'après le lemme \ref{lemma-relative-affine-cohomology} et le fait que les
immersions fermées sont affines ; voir
Morphismes, lemme \ref{morphisms-lemma-closed-immersion-affine}).
De plus, si $\mathcal{G}$ est cohérent, alors $i_*\mathcal{G}$
l'est aussi (lemme \ref{lemma-i-star-equivalence}).
Nous emploierons désormais ces faits sans autre mention.

\medskip\noindent
Démonstration de (1). Posons $S = R[T_0, \ldots, T_n]$, de sorte que
$\mathbf{P}^n_R = \text{Proj}(S)$.
Remarquons que $A$ est une $S$-algèbre (mais le morphisme d'anneaux $S \to A$ n'est pas
un morphisme d'anneaux gradués, puisque $S_n$ s'envoie dans $A_{dn}$).
La formule de projection
(Cohomologie, lemme \ref{cohomology-lemma-projection-formula})
donne
$$
i_*(\mathcal{L}^{\otimes nd + q}) =
i_*(\mathcal{L}^{\otimes q})
\otimes_{\mathcal{O}_{\mathbf{P}^n_R}}
\mathcal{O}_{\mathbf{P}^n_R}(n)
$$
pour tout $n \in \mathbf{Z}$. On en conclut que $\bigoplus_{n \geq 0} A_{nd + q}$
est un $S$-module gradué de type fini d'après le lemme \ref{lemma-coherent-projective}.
Comme
$A = \bigoplus_{q \in \{0, \ldots, d - 1} \bigoplus_{n \geq 0} A_{nd + q}$
on voit que $A$ est une $S$-algèbre finie, ce qui prouve (1).

\medskip\noindent
Démonstration de (2). Cela résulte de
Propriétés, proposition \ref{properties-proposition-characterize-ample}.

\medskip\noindent
Démonstration de (3). On applique le lemme \ref{lemma-coherent-projective}
et l'égalité $H^p(X, \mathcal{F}) = H^p(\mathbf{P}^n_R, i_*\mathcal{F})$.

\medskip\noindent
Démonstration de (4). Fixons $p > 0$. La formule de projection donne
$$
i_*(\mathcal{F} \otimes_{\mathcal{O}_X} \mathcal{L}^{\otimes nd + q}) =
i_*(\mathcal{F} \otimes_{\mathcal{O}_X} \mathcal{L}^{\otimes q})
\otimes_{\mathcal{O}_{\mathbf{P}^n_R}}
\mathcal{O}_{\mathbf{P}^n_R}(n)
$$
pour tout $n \in \mathbf{Z}$. D'après le lemme \ref{lemma-coherent-projective},
on a donc $H^p(X, \mathcal{F} \otimes \mathcal{L}^{nd + q}) = 0$
pour $n \gg 0$. Comme il n'existe qu'un nombre fini de classes de
congruence d'entiers modulo $d$, ceci prouve (4).

\medskip\noindent
Démonstration de (5). Fixons un entier $k$. Posons
$M = \bigoplus_{n \geq k} H^0(X, \mathcal{F} \otimes \mathcal{L}^{\otimes n})$.
En raisonnant comme ci-dessus, on voit que $\bigoplus_{nd + q \geq k} M_{nd + q}$
est un $S$-module gradué de type fini. Comme
$M = \bigoplus_{q \in \{0, \ldots, d - 1\}}
\bigoplus_{nd + q \geq k} M_{nd + q}$
on voit que $M$ est un $S$-module de type fini. Puisque sa structure de $S$-module se
factorise par le morphisme d'anneaux $S \to A$, on conclut que $M$ est
un $A$-module de type fini.
\end{proof}

\begin{lemma}
\label{lemma-kill-by-twisting}
Soit $f : X \to S$ un morphisme de schémas.
Soit $\mathcal{F}$ un $\mathcal{O}_X$-module quasi-cohérent.
Soit $\mathcal{L}$ un faisceau inversible sur $X$.
Supposons que
\begin{enumerate}
\item $S$ soit noethérien,
\item $f$ soit propre,
\item $\mathcal{F}$ soit cohérent, et
\item $\mathcal{L}$ soit relativement ample sur $X/S$.
\end{enumerate}
Il existe alors $n_0$ tel que, pour tout $n \geq n_0$,
on ait
$$
R^pf_*\left(\mathcal{F} \otimes_{\mathcal{O}_X} \mathcal{L}^{\otimes n}\right)
=
0
$$
pour tout $p > 0$.
\end{lemma}

\begin{proof}
Choisissons un recouvrement ouvert affine fini $S = \bigcup V_j$ et
posons $X_j = f^{-1}(V_j)$.
Il suffit manifestement de résoudre la question pour chacun des systèmes, en nombre
fini, $(X_j \to V_j, \mathcal{L}|_{X_j}, \mathcal{F}|_{V_j})$.
On peut donc supposer $S$ affine. Dans
ce cas, l'annulation de
$R^pf_*(\mathcal{F} \otimes \mathcal{L}^{\otimes n})$
équivaut à celle de
$H^p(X, \mathcal{F} \otimes \mathcal{L}^{\otimes n})$ ; voir le
lemme \ref{lemma-quasi-coherence-higher-direct-images-application}.
L'annulation voulue résulte donc
du lemme \ref{lemma-coherent-proper-ample} (qui s'applique parce que
$\mathcal{L}$ est ample sur $X$ d'après Morphismes, Lemme
\ref{morphisms-lemma-finite-type-over-affine-ample-very-ample}).
\end{proof}

\begin{lemma}
\label{lemma-locally-projective-pushforward}
Soit $S$ un schéma localement noethérien.
Soit $f : X \to S$ un morphisme localement projectif.
Soit $\mathcal{F}$ un $\mathcal{O}_X$-module cohérent.
Alors $R^if_*\mathcal{F}$ est un $\mathcal{O}_S$-module
cohérent pour tout $i \geq 0$.
\end{lemma}

\begin{proof}
Remarquons d'abord qu'un morphisme localement projectif est propre
(Morphismes, lemme \ref{morphisms-lemma-locally-projective-proper})
et donc de type fini.
En particulier, $X$ est localement noethérien
(Morphismes, lemme \ref{morphisms-lemma-finite-type-noetherian}),
de sorte que l'énoncé a un sens. De
plus, d'après le lemme \ref{lemma-quasi-coherence-higher-direct-images},
les faisceaux $R^pf_*\mathcal{F}$ sont quasi-cohérents.

\medskip\noindent
Cela étant, l'énoncé est local sur $S$ (par exemple d'après
Cohomologie, lemme \ref{cohomology-lemma-localize-higher-direct-images}).
On peut donc supposer que $S = \Spec(R)$ est le spectre d'un
anneau noethérien et que $X$ est un sous-schéma fermé de
$\mathbf{P}^n_R$ pour un certain $n$ ; voir
Morphismes, lemme \ref{morphisms-lemma-characterize-locally-projective}.
Dans ce cas, les faisceaux $R^pf_*\mathcal{F}$ sont les faisceaux
quasi-cohérents associés aux $R$-modules $H^p(X, \mathcal{F})$ ; voir le
lemme \ref{lemma-quasi-coherence-higher-direct-images-application}.
Il suffit donc de montrer que les $R$-modules $H^p(X, \mathcal{F})$
sont des $R$-modules de type fini (lemme \ref{lemma-coherent-Noetherian}).
Cela résulte du lemme \ref{lemma-coherent-proper-ample},
car la restriction de $\mathcal{O}_{\mathbf{P}^n_R}(1)$
à $X$ est ample sur $X$.
\end{proof}








\section{Faisceaux inversibles amples et cohomologie}
\label{section-ample-cohomology}

\noindent
Voici un critère d'amplitude sur les schémas propres sur une base affine,
exprimé par l'annulation de la cohomologie des faisceaux tordus.

\begin{lemma}
\label{lemma-vanshing-gives-ample}
\begin{reference}
\cite[III Proposition 2.6.1]{EGA}
\end{reference}
Soit $R$ un anneau noethérien. Soit $f : X \to \Spec(R)$ un morphisme propre.
Soit $\mathcal{L}$ un $\mathcal{O}_X$-module inversible.
Les conditions suivantes sont équivalentes :
\begin{enumerate}
\item $\mathcal{L}$ est ample sur $X$ (cette condition admet bien d'autres formes
équivalentes, voir
Propriétés, proposition \ref{properties-proposition-characterize-ample} et
Morphismes, Lemme
\ref{morphisms-lemma-finite-type-over-affine-ample-very-ample}),
\item pour tout $\mathcal{O}_X$-module cohérent $\mathcal{F}$, il
existe $n_0 \geq 0$ tel que
$H^p(X, \mathcal{F} \otimes \mathcal{L}^{\otimes n}) = 0$ pour tout $n \geq n_0$
et tout $p > 0$, et
\item pour tout faisceau quasi-cohérent d'idéaux
$\mathcal{I} \subset \mathcal{O}_X$, il existe $n \geq 1$
tel que $H^1(X, \mathcal{I} \otimes \mathcal{L}^{\otimes n}) = 0$.
\end{enumerate}
\end{lemma}

\begin{proof}
L'implication (1) $\Rightarrow$ (2) résulte du
lemme \ref{lemma-coherent-proper-ample}.
L'implication (2) $\Rightarrow$ (3) est
triviale. L'implication (3) $\Rightarrow$ (1) est le
lemme \ref{lemma-quasi-compact-h1-zero-invertible}.
\end{proof}

\begin{lemma}
\label{lemma-surjective-finite-morphism-ample}
Soit $R$ un anneau noethérien. Soit $f : Y \to X$ un morphisme de
schémas propres sur $R$. Soit $\mathcal{L}$
un $\mathcal{O}_X$-module inversible. Supposons $f$ fini et surjectif.
Alors $\mathcal{L}$ est ample si et seulement si $f^*\mathcal{L}$ est ample.
\end{lemma}

\begin{proof}
L'image réciproque d'un faisceau inversible ample par un morphisme quasi-affine
est ample, voir Morphismes, Lemme
\ref{morphisms-lemma-pullback-ample-tensor-relatively-ample}.
Cela démontre l'une des implications, car un morphisme fini est affine
par définition.

\medskip\noindent
Supposons $f^*\mathcal{L}$ ample. Soit $P$ la propriété suivante
des $\mathcal{O}_X$-modules cohérents $\mathcal{F}$ : il existe $n_0$
tel que $H^p(X, \mathcal{F} \otimes \mathcal{L}^{\otimes n}) = 0$
pour tout $n \geq n_0$ et tout $p > 0$. Nous allons montrer que $P$ est
vérifiée par tout $\mathcal{O}_X$-module cohérent $\mathcal{F}$, ce qui entraînera que
$\mathcal{L}$ est ample d'après le lemme \ref{lemma-vanshing-gives-ample}.
Nous allons appliquer le lemme \ref{lemma-property-higher-rank-cohomological}.
Il faut donc vérifier les conditions (1), (2) et (3) de ce lemme pour $P$.
La propriété (1) résulte de la suite exacte longue de cohomologie associée
à une suite exacte courte de faisceaux et du fait que le produit tensoriel avec un
faisceau inversible est un foncteur exact. La propriété (2) résulte de ce que
$H^p(X, -)$ est un foncteur additif. Pour vérifier (3), soit $Z \subset X$ un
sous-schéma fermé intègre de point générique $\xi$.
Soit $\mathcal{F}$ un faisceau cohérent sur $Y$ tel que
le support de $f_*\mathcal{F}$ soit égal à $Z$
et que $(f_*\mathcal{F})_\xi$ soit annulé par $\mathfrak m_\xi$,
voir lemme \ref{lemma-finite-morphism-Noetherian}. Nous affirmons
que $\mathcal{G} = f_*\mathcal{F}$ convient. Il suffit de vérifier la
condition (3)(c) du lemme \ref{lemma-property-higher-rank-cohomological}.
Supposons donc que $\mathcal{J} \subset \mathcal{O}_X$ soit un
faisceau quasi-cohérent d'idéaux tel que
$\mathcal{J}_\xi = \mathcal{O}_{X, \xi}$.
Un morphisme fini est affine; d'après le
lemme \ref{lemma-affine-morphism-projection-ideal}, on a donc
$\mathcal{J}\mathcal{G} = f_*(f^{-1}\mathcal{J}\mathcal{F})$.
De plus, comme on l'a remarqué dans la démonstration du
lemme \ref{lemma-affine-morphism-projection-ideal}, le faisceau
$f^{-1}\mathcal{J}\mathcal{F}$ est un $\mathcal{O}_Y$-module cohérent.
Puisque $\mathcal{L}$ est ample, le lemme \ref{lemma-vanshing-gives-ample}
assure l'existence de $n_0$ tel que
$$
H^p(Y, f^{-1}\mathcal{J}\mathcal{F}
\otimes_{\mathcal{O}_Y} f^*\mathcal{L}^{\otimes n}) = 0,
$$
pour $n \geq n_0$ et $p > 0$. Puisque $f$ est fini, donc affine, on a
\begin{align*}
H^p(X, \mathcal{J}\mathcal{G} \otimes_{\mathcal{O}_X}
\mathcal{L}^{\otimes n})
& =
H^p(X, f_*(f^{-1}\mathcal{J}\mathcal{F}) \otimes_{\mathcal{O}_X}
\mathcal{L}^{\otimes n}) \\
& =
H^p(X, f_*(f^{-1}\mathcal{J}\mathcal{F} \otimes_{\mathcal{O}_Y}
f^*\mathcal{L}^{\otimes n})) \\
& =
H^p(Y, f^{-1}\mathcal{J}\mathcal{F} \otimes_{\mathcal{O}_Y}
f^*\mathcal{L}^{\otimes n}) = 0
\end{align*}
On a utilisé ici la formule de projection
(Cohomologie, lemme \ref{cohomology-lemma-projection-formula}) et le
lemme \ref{lemma-relative-affine-cohomology}.
Le sous-faisceau quasi-cohérent $\mathcal{G}' = \mathcal{J}\mathcal{G}$
vérifie donc $P$. Cela établit la propriété (3)(c) du
lemme \ref{lemma-property-higher-rank-cohomological}, comme voulu.
\end{proof}

\noindent
La cohomologie est fonctorielle. En particulier, étant donnés un espace annelé $X$,
un $\mathcal{O}_X$-module inversible $\mathcal{L}$ et une section
$s \in \Gamma(X, \mathcal{L})$, on obtient des applications
$$
H^p(X, \mathcal{F})
\longrightarrow
H^p(X, \mathcal{F} \otimes_{\mathcal{O}_X} \mathcal{L}), \quad
\xi \longmapsto s\xi
$$
induites par l'application
$\mathcal{F} \to \mathcal{F} \otimes_{\mathcal{O}_X} \mathcal{L}$
de multiplication par $s$. On pose
$\Gamma_*(X, \mathcal{L}) =
\bigoplus_{n \geq 0} \Gamma(X, \mathcal{L}^{\otimes n})$,
considéré comme un anneau gradué,
voir Modules, définition \ref{modules-definition-gamma-star}.
Étant donnés un faisceau de $\mathcal{O}_X$-modules $\mathcal{F}$ et un entier
$p \geq 0$, on pose
$$
H^p_*(X, \mathcal{L}, \mathcal{F}) =
\bigoplus\nolimits_{n \in \mathbf{Z}} H^p(X,
\mathcal{F} \otimes_{\mathcal{O}_X} \mathcal{L}^{\otimes n})
$$
La multiplication définie ci-dessus en fait un $\Gamma_*(X, \mathcal{L})$-module gradué. Attention:
la notation $H^p_*(X, \mathcal{L}, \mathcal{F})$
n'est pas usuelle.

\begin{lemma}
\label{lemma-invert-s-cohomology}
Soit $X$ un schéma. Soit $\mathcal{L}$ un faisceau inversible sur $X$.
Soit $s \in \Gamma(X, \mathcal{L})$.
Soit $\mathcal{F}$ un $\mathcal{O}_X$-module quasi-cohérent.
Si $X$ est quasi-compact et quasi-séparé, l'application canonique
$$
H^p_*(X, \mathcal{L}, \mathcal{F})_{(s)}
\longrightarrow
H^p(X_s, \mathcal{F})
$$
qui envoie $\xi/s^n$ sur $s^{-n}\xi$ est un isomorphisme.
\end{lemma}

\begin{proof}
Pour $p = 0$, il s'agit
de Propriétés, lemme \ref{properties-lemma-invert-s-sections}.
Nous allons démontrer l'assertion au moyen du principe de
récurrence (lemme \ref{lemma-induction-principle}), en notant, pour tout ouvert quasi-compact
$U \subset X$, $P(U)$ la propriété :
pour tout $p \geq 0$, l'application
$$
H^p_*(U, \mathcal{L}, \mathcal{F})_{(s)}
\longrightarrow
H^p(U_s, \mathcal{F})
$$
est un isomorphisme.

\medskip\noindent
Si $U$ est affine, les deux membres de la flèche ci-dessus
sont nuls pour $p > 0$, d'après
le lemme \ref{lemma-quasi-coherent-affine-cohomology-zero}
et
Propriétés, lemme \ref{properties-lemma-affine-cap-s-open},
et l'assertion est vraie. Si $P$ est vraie pour $U$, $V$ et $U \cap V$,
on peut utiliser les suites de Mayer-Vietoris
(Cohomologie, lemme \ref{cohomology-lemma-mayer-vietoris}) pour obtenir
un morphisme de suites exactes longues
$$
\xymatrix{
H^{p - 1}_*(U \cap V, \mathcal{L}, \mathcal{F})_{(s)} \ar[r] \ar[d] &
H^p_*(U \cup V, \mathcal{L}, \mathcal{F})_{(s)} \ar[r] \ar[d] &
H^p_*(U, \mathcal{L}, \mathcal{F})_{(s)}
\oplus
H^p_*(V, \mathcal{L}, \mathcal{F})_{(s)} \ar[d] \\
H^{p - 1}(U_s \cap V_s, \mathcal{F}) \ar[r]&
H^p(U_s \cup V_s, \mathcal{F}) \ar[r] &
H^p(U_s, \mathcal{F})
\oplus
H^p(V_s, \mathcal{F})
}
$$
(dont on ne représente qu'un fragment). Notons que $U_s \cap V_s = (U \cap V)_s$ et
que $U_s \cup V_s = (U \cup V)_s$. Les applications verticales de gauche et
de droite sont donc des isomorphismes (ainsi qu'une application supplémentaire à droite et une
à gauche, qui ne figurent pas dans le diagramme). On
en déduit que $P(U \cup V)$ est vérifiée, d'après
le lemme des cinq (Homologie, lemme \ref{homology-lemma-five-lemma}).
Cela achève la démonstration.
\end{proof}

\begin{lemma}
\label{lemma-section-affine-open-kills-classes}
Soit $X$ un schéma.
Soit $\mathcal{L}$ un $\mathcal{O}_X$-module inversible.
Soit $s \in \Gamma(X, \mathcal{L})$ une section.
Supposons que
\begin{enumerate}
\item $X$ soit quasi-compact et quasi-séparé, et
\item $X_s$ soit affine.
\end{enumerate}
Alors, pour tout $\mathcal{O}_X$-module quasi-cohérent $\mathcal{F}$,
tout $p > 0$ et tout $\xi \in H^p(X, \mathcal{F})$, il
existe $n \geq 0$ tel que $s^n\xi = 0$ dans
$H^p(X, \mathcal{F} \otimes_{\mathcal{O}_X} \mathcal{L}^{\otimes n})$.
\end{lemma}

\begin{proof}
Rappelons que $H^p(X_s, \mathcal{G})$ est nul pour tout module
quasi-cohérent $\mathcal{G}$, d'après
le lemme \ref{lemma-quasi-coherent-affine-cohomology-zero}.
Le lemme résulte donc du
lemme \ref{lemma-invert-s-cohomology}.
\end{proof}

\noindent
Pour une version plus générale du lemme suivant, voir
Limites, lemme \ref{limits-lemma-ample-on-reduction}.

\begin{lemma}
\label{lemma-ample-on-reduction}
Soit $i : Z \to X$ une immersion fermée de schémas noethériens induisant
un homéomorphisme des espaces topologiques sous-jacents.
Soit $\mathcal{L}$ un faisceau inversible sur $X$.
Alors $i^*\mathcal{L}$ est ample sur $Z$ si et seulement si
$\mathcal{L}$ est ample sur $X$.
\end{lemma}

\begin{proof}
Si $\mathcal{L}$ est ample, alors $i^*\mathcal{L}$ est ample, par
exemple d'après Morphismes, Lemme
\ref{morphisms-lemma-pullback-ample-tensor-relatively-ample}.
Supposons $i^*\mathcal{L}$ ample. Il faut montrer que $\mathcal{L}$
est ample sur $X$.
Soit $\mathcal{I} \subset \mathcal{O}_X$ le faisceau cohérent d'idéaux
définissant le sous-schéma fermé $Z$. Puisque $i(Z) = X$ ensemblistement,
on a $\mathcal{I}^n = 0$ pour un certain $n$, d'après
le lemme \ref{lemma-power-ideal-kills-sheaf}.
Considérons la suite
$$
X = Z_n \supset Z_{n - 1} \supset Z_{n - 2} \supset \ldots \supset Z_1 = Z
$$
de sous-schémas fermés définis par
$0 = \mathcal{I}^n \subset \mathcal{I}^{n - 1} \subset \ldots \subset
\mathcal{I}$. Chaque immersion fermée $Z_{i - 1} \to Z_{i}$
est alors définie par un faisceau cohérent d'idéaux de carré nul. On se
ramène ainsi au cas où $\mathcal{I}^2 = 0$.

\medskip\noindent
Considérons la suite exacte courte
$$
0 \to \mathcal{I} \to \mathcal{O}_X \to i_*\mathcal{O}_Z \to 0
$$
de $\mathcal{O}_X$-modules quasi-cohérents. En tensorisant par
$\mathcal{L}^{\otimes n}$, on obtient des suites exactes courtes
\begin{equation}
\label{equation-ses}
0 \to \mathcal{I} \otimes_{\mathcal{O}_X} \mathcal{L}^{\otimes n}
\to \mathcal{L}^{\otimes n} \to i_*i^*\mathcal{L}^{\otimes n} \to 0
\end{equation}
Puisque $\mathcal{I}^2 = 0$, on peut utiliser
Morphismes, lemme \ref{morphisms-lemma-i-star-equivalence},
pour considérer $\mathcal{I}$ comme un $\mathcal{O}_Z$-module quasi-cohérent;
on a alors $\mathcal{I} \otimes_{\mathcal{O}_X} \mathcal{L}^{\otimes n} =
\mathcal{I} \otimes_{\mathcal{O}_Z} i^*\mathcal{L}^{\otimes n}$,
avec un abus de notation
évident. De plus, la cohomologie de ce faisceau sur $Z$ s'identifie
canoniquement à sa cohomologie sur $X$ (puisque $i$ est un
homéomorphisme).

\medskip\noindent
Soit $x \in X$ un point, et notons $z \in Z$ le point correspondant.
Puisque $i^*\mathcal{L}$ est ample, il existe $n$ et une section
$s \in \Gamma(Z, i^*\mathcal{L}^{\otimes n})$ tels que $z \in Z_s$
et que $Z_s$ soit affine. L'obstruction au relèvement de $s$ en une section
de $\mathcal{L}^{\otimes n}$ sur $X$ est le cobord
$$
\xi = \partial s \in
H^1(X, \mathcal{I} \otimes_{\mathcal{O}_X} \mathcal{L}^{\otimes n}) =
H^1(Z, \mathcal{I} \otimes_{\mathcal{O}_Z} i^*\mathcal{L}^{\otimes n})
$$
issu de la suite exacte courte de faisceaux (\ref{equation-ses}).
Si l'on remplace $s$ par $s^{e + 1}$, alors $\xi$ est remplacé par
$\partial(s^{e + 1}) = (e + 1) s^e \xi$ dans
$H^1(Z, \mathcal{I} \otimes_{\mathcal{O}_Z} i^*\mathcal{L}^{\otimes (e + 1)n})$,
car l'application de cobord de
$$
0 \to
\bigoplus\nolimits_{m \geq 0}
\mathcal{I} \otimes_{\mathcal{O}_X} \mathcal{L}^{\otimes m} \to
\bigoplus\nolimits_{m \geq 0}
\mathcal{L}^{\otimes m} \to
\bigoplus\nolimits_{m \geq 0}
i_*i^*\mathcal{L}^{\otimes m} \to 0
$$
est une dérivation d'après Cohomologie, Lemme
\ref{cohomology-lemma-boundary-derivation-over-cup-product}. D'après le
lemme \ref{lemma-section-affine-open-kills-classes},
$s^e \xi$ est nul pour $e$ assez grand.
Ainsi, quitte à remplacer $s$ par une puissance, on peut supposer que $s$ est l'image
d'une section $s' \in \Gamma(X, \mathcal{L}^{\otimes n})$.
Alors $X_{s'}$ est un sous-schéma ouvert et $Z_s \to X_{s'}$ est une immersion
fermée surjective de schémas noethériens, avec $Z_s$ affine. Ainsi
$X_s$ est affine d'après le
lemme \ref{lemma-image-affine-finite-morphism-affine-Noetherian}, et
on en déduit que $\mathcal{L}$ est ample.
\end{proof}

\noindent
Pour une version plus générale du lemme suivant, voir
Limites, lemme \ref{limits-lemma-thickening-quasi-affine}.

\begin{lemma}
\label{lemma-thickening-quasi-affine}
Soit $i : Z \to X$ une immersion fermée de schémas noethériens induisant
un homéomorphisme des espaces topologiques sous-jacents.
Alors $X$ est quasi-affine si et seulement si $Z$ est quasi-affine.
\end{lemma}

\begin{proof}
Rappelons qu'un schéma est quasi-affine si
et seulement si son faisceau structural est ample, voir
Propriétés, lemme \ref{properties-lemma-quasi-affine-O-ample}.
Ainsi, si $Z$ est quasi-affine, alors $\mathcal{O}_Z$ est ample,
donc $\mathcal{O}_X$ est ample d'après le
lemme \ref{lemma-ample-on-reduction}, et par suite
$X$ est quasi-affine. La réciproque, que l'on peut
également établir de manière élémentaire, se démontre en lisant le
raisonnement précédent en sens inverse.
\end{proof}

\begin{lemma}
\label{lemma-affine-in-presence-ample}
Soit $X$ un schéma. Soit $\mathcal{L}$ un
$\mathcal{O}_X$-module inversible ample. Soit $n_0$ un entier.
Si $H^p(X, \mathcal{L}^{\otimes -n}) = 0$ pour $n \geq n_0$ et $p > 0$,
alors $X$ est affine.
\end{lemma}

\begin{proof}
Nous affirmons que $H^p(X, \mathcal{F}) = 0$ pour tout
$\mathcal{O}_X$-module quasi-cohérent et tout $p > 0$. Puisque $X$ est quasi-compact
d'après Propriétés, définition \ref{properties-definition-ample},
cette assertion achève la démonstration, en
vertu du lemme \ref{lemma-quasi-compact-h1-zero-covering}.
Le schéma $X$ est séparé d'après
Propriétés, lemme \ref{properties-lemma-ample-separated}. Supposons
que $X$ soit recouvert par $e + 1$ ouverts affines. Alors
$H^p(X, \mathcal{F}) = 0$ pour $p > e$, voir
lemme \ref{lemma-vanishing-nr-affines}. On peut donc procéder par récurrence
descendante sur $p$. En écrivant $\mathcal{F}$
comme limite inductive filtrante de modules quasi-cohérents de type
fini (Propriétés, Lemme
\ref{properties-lemma-quasi-coherent-colimit-finite-type})
et en utilisant Cohomologie, Lemme
\ref{cohomology-lemma-quasi-separated-cohomology-colimit},
on peut supposer $\mathcal{F}$ de type fini. On
peut alors choisir $n > n_0$ tel que
$\mathcal{F} \otimes_{\mathcal{O}_X} \mathcal{L}^{\otimes n}$
soit engendré par ses sections globales, voir Propriétés, Proposition
\ref{properties-proposition-characterize-ample}.
Cela signifie qu'il existe une suite exacte courte
$$
0 \to \mathcal{F}' \to
\bigoplus\nolimits_{i \in I} \mathcal{L}^{\otimes -n}
\to \mathcal{F} \to 0
$$
pour un certain ensemble $I$ (on peut en fait choisir $I$ fini). Par
hypothèse de récurrence, on a $H^{p + 1}(X, \mathcal{F}') = 0$
et, par hypothèse (en utilisant de nouveau la commutation
de la cohomologie aux limites inductives),
on a $H^p(X, \bigoplus_{i \in I} \mathcal{L}^{\otimes -n}) = 0$.
La suite exacte longue de cohomologie donne alors
$H^p(X, \mathcal{F}) = 0$, comme voulu.
\end{proof}

\begin{lemma}
\label{lemma-affine-if-quasi-affine}
Soit $X$ un schéma quasi-affine.
Si $H^p(X, \mathcal{O}_X) = 0$ pour $p > 0$,
alors $X$ est affine.
\end{lemma}

\begin{proof}
Puisque $\mathcal{O}_X$ est ample d'après
Propriétés, lemme \ref{properties-lemma-quasi-affine-O-ample},
cela résulte du lemme \ref{lemma-affine-in-presence-ample}.
\end{proof}







\section{Lemme de Chow}
\label{section-chows-lemma}

\noindent
Dans cette section, nous démontrons le lemme de Chow dans le cas
noethérien (lemme \ref{lemma-chow-Noetherian}).
Dans
Limites, section \ref{limits-section-chows-lemma},
nous en démontrons des variantes dans le cas non noethérien.

\begin{lemma}
\label{lemma-chow-Noetherian}
\begin{reference}
\cite[II théorème 5.6.1(a)]{EGA}
\end{reference}
Soit $S$ un schéma noethérien.
Soit $f : X \to S$ un morphisme séparé de type fini.
Il existe alors $n \geq 0$ et un diagramme
$$
\xymatrix{
X \ar[rd] & X' \ar[d] \ar[l]^\pi \ar[r] & \mathbf{P}^n_S \ar[dl] \\
& S &
}
$$
où $X' \to \mathbf{P}^n_S$ est une immersion et
$\pi : X' \to X$ est propre et surjectif. De plus, on
peut faire en sorte qu'il existe un sous-schéma ouvert dense
$U \subset X$ tel que $\pi^{-1}(U) \to U$ soit un isomorphisme.
\end{lemma}

\begin{proof}
Tous les schémas qui interviendront dans la suite de la
démonstration seront de type fini sur le schéma noethérien $S$,
donc noethériens
(voir Morphismes, lemme \ref{morphisms-lemma-finite-type-noetherian}).
Tous les morphismes entre eux seront automatiquement quasi-compacts, localement de
type fini et quasi-séparés, voir
Morphismes, lemme \ref{morphisms-lemma-permanence-finite-type} et
Propriétés,
lemmes \ref{properties-lemma-locally-Noetherian-quasi-separated} et
\ref{properties-lemma-morphism-Noetherian-schemes-quasi-compact}.

\medskip\noindent
Le schéma $X$ n'a qu'un nombre fini de composantes irréductibles
(Propriétés, lemme \ref{properties-lemma-Noetherian-irreducible-components}).
Soit $X = X_1 \cup \ldots \cup X_r$ la décomposition
de $X$ en composantes irréductibles.
Soit $\eta_i \in X_i$ le point générique.
Pour tout point $x \in X$, il existe un ouvert affine
$U_x \subset X$ contenant $x$ et chacun des points
génériques $\eta_i$. Voir
Propriétés, lemme \ref{properties-lemma-point-and-maximal-points-affine}.
Puisque $X$ est quasi-compact, il admet un recouvrement fini par des ouverts
affines $X = U_1 \cup \ldots \cup U_m$ tel que
chaque $U_i$ contienne $\eta_1, \ldots, \eta_r$.
On en déduit en particulier que l'ouvert
$U = U_1 \cap \ldots \cap U_m \subset X$ est
dense. C'est tout ce que nous utiliserons ci-dessous, avec le fait que les $U_i$ sont des ouverts affines
recouvrant $X$.

\medskip\noindent
Soit $X^* \subset X$ l'adhérence schématique de $U \to X$, voir
Morphismes, définition \ref{morphisms-definition-scheme-theoretic-image}.
Posons $U_i^* = X^* \cap U_i$. Notons que $U_i^*$ est un sous-schéma fermé
de $U_i$. Ainsi $U_i^*$ est affine. Puisque $U$ est dense dans $X$, le
morphisme $X^* \to X$ est une immersion fermée surjective. C'est un isomorphisme
au-dessus de $U$. On peut donc remplacer $X$ par $X^*$ et
$U_i$ par $U_i^*$, et supposer que $U$ est schématiquement dense
dans $X$, voir
Morphismes, définition \ref{morphisms-definition-scheme-theoretically-dense}.

\medskip\noindent
D'après Morphismes, lemme \ref{morphisms-lemma-quasi-projective-finite-type-over-S},
on peut trouver une immersion $j_i : U_i \to \mathbf{P}_S^{n_i}$
pour chaque $i$. D'après
Morphismes, lemme \ref{morphisms-lemma-quasi-compact-immersion}, on peut trouver
des sous-schémas fermés $Z_i \subset \mathbf{P}_S^{n_i}$
tels que $j_i : U_i \to Z_i$ soit une immersion ouverte
schématiquement dense. Notons que $Z_i \to S$ est propre, voir
Morphismes, lemme \ref{morphisms-lemma-locally-projective-proper}.
Considérons le morphisme
$$
j = (j_1|_U, \ldots, j_m|_U) : U \longrightarrow
\mathbf{P}_S^{n_1} \times_S \ldots \times_S \mathbf{P}_S^{n_m}.
$$
D'après le lemme cité ci-dessus, il existe un sous-schéma fermé
$Z$ de $\mathbf{P}_S^{n_1} \times_S \ldots \times_S \mathbf{P}_S^{n_m}$
tel que $j : U \to Z$ soit une immersion ouverte et que $U$
soit schématiquement dense dans $Z$. Le morphisme $Z \to S$
est propre. Considérons la $i$-ième projection
$$
\text{pr}_i|_Z : Z \longrightarrow \mathbf{P}^{n_i}_S.
$$
Ce morphisme se factorise par $Z_i$ (voir Morphismes,
lemme \ref{morphisms-lemma-factor-factor}). Notons $p_i : Z \to Z_i$
le morphisme induit. Ce morphisme est propre, voir par exemple
Morphismes, lemme \ref{morphisms-lemma-image-proper-scheme-closed}.
Nous avons donc des immersions ouvertes
$U \subset U_i \subset Z_i$ schématiquement denses.
De plus, on peut considérer $Z$ comme l'image
schématique du morphisme « diagonal »
$U \to Z_1 \times_S \ldots \times_S Z_m$.

\medskip\noindent
Posons $V_i = p_i^{-1}(U_i)$. Notons que $p_i|_{V_i} : V_i \to U_i$ est propre.
Posons $X' = V_1 \cup \ldots \cup V_m$. Par construction, $X'$ admet une immersion
dans le schéma
$\mathbf{P}^{n_1}_S \times_S \ldots \times_S \mathbf{P}^{n_m}_S$.
Ainsi, grâce au plongement de Segre (voir
Constructions, lemme \ref{constructions-lemma-segre-embedding}
), $X'$ admet
une immersion dans un espace projectif sur $S$.

\medskip\noindent
Nous affirmons que les morphismes $p_i|_{V_i}: V_i \to U_i$ se recollent en un morphisme
$X' \to X$. En effet, $p_i|_U$ est clairement l'application identité
de $U$ dans $U$. Puisque $U \subset X'$ est schématiquement dense
par construction, il est aussi schématiquement dense
dans le sous-schéma ouvert $V_i \cap V_j$. On a donc
$p_i|_{V_i \cap V_j} = p_j|_{V_i \cap V_j}$ comme morphismes dans
le $S$-schéma séparé $X$, voir
Morphismes, lemme \ref{morphisms-lemma-equality-of-morphisms}.
Notons $\pi : X' \to X$ le morphisme ainsi obtenu.

\medskip\noindent
Nous affirmons que $\pi^{-1}(U_i) = V_i$.
Puisque $\pi|_{V_i} = p_i|_{V_i}$, on a
$V_i \subset \pi^{-1}(U_i)$. Considérons le diagramme
$$
\xymatrix{
V_i \ar[r] \ar[rd]_{p_i|_{V_i}} & \pi^{-1}(U_i) \ar[d] \\
& U_i
}
$$
Puisque $V_i \to U_i$ est propre, l'image de la flèche
horizontale est fermée, voir
Morphismes, lemme \ref{morphisms-lemma-image-proper-scheme-closed}.
Puisque $V_i \subset \pi^{-1}(U_i)$ est schématiquement
dense (car il contient $U$),
on en déduit que $V_i = \pi^{-1}(U_i)$, comme annoncé.

\medskip\noindent
Cela montre que $\pi^{-1}(U_i) \to U_i$ s'identifie au morphisme
propre $p_i|_{V_i} : V_i \to U_i$. Ainsi $X$ admet un
recouvrement affine fini $X = \bigcup U_i$ tel que la restriction
de $\pi$ soit propre au-dessus de chaque ouvert du
recouvrement. D'après Morphismes, lemme \ref{morphisms-lemma-proper-local-on-the-base},
$\pi$ est donc propre.

\medskip\noindent
Enfin, il faut montrer que $\pi^{-1}(U) = U$. On raisonne pour cela comme
ci-dessus, à l'aide du diagramme
$$
\xymatrix{
U \ar[r] \ar[rd] & \pi^{-1}(U) \ar[d] \\
& U
}
$$
en utilisant le fait que $\text{id}_U : U \to U$ est propre et que
$U$ est schématiquement dense dans $\pi^{-1}(U)$.
\end{proof}

\begin{remark}
\label{remark-chow-Noetherian}
Dans la situation du Lemme de
Chow \ref{lemma-chow-Noetherian} :
\begin{enumerate}
\item Le morphisme $\pi$ est en fait H-projectif (donc projectif, voir
Morphismes, lemme \ref{morphisms-lemma-H-projective}),
car le morphisme $X' \to \mathbf{P}^n_S \times_S X = \mathbf{P}^n_X$
est une immersion fermée (on utilise le fait que $\pi$ est propre, voir
Morphismes, lemme \ref{morphisms-lemma-image-proper-scheme-closed}).
\item On peut supposer que $\pi^{-1}(U)$ est schématiquement dense
dans $X'$. En effet, il suffit de remplacer $X'$ par l'adhérence schématique
de $\pi^{-1}(U)$. On peut alors considérer $U$ comme un
sous-schéma ouvert schématiquement dense de $X'$.
Voir Morphismes, section \ref{morphisms-section-scheme-theoretic-image}.
\item Si $X$ est réduit, on peut choisir $X'$ réduit. Cela résulte clairement
de (2).
\end{enumerate}
\end{remark}






\section{Images directes supérieures des faisceaux cohérents}
\label{section-proper-pushforward}

\noindent
Dans cette section, nous démontrons le fait fondamental que les images
directes supérieures d'un faisceau cohérent par un morphisme propre
sont cohérentes.

\begin{proposition}
\label{proposition-proper-pushforward-coherent}
\begin{reference}
\cite[III théorème 3.2.1]{EGA}
\end{reference}
Soit $S$ un schéma localement noethérien.
Soit $f : X \to S$ un morphisme propre.
Soit $\mathcal{F}$ un $\mathcal{O}_X$-module cohérent.
Alors $R^if_*\mathcal{F}$ est un $\mathcal{O}_S$-module
cohérent pour tout $i \geq 0$.
\end{proposition}

\begin{proof}
La question étant locale sur $S$, on peut supposer que $S$ est
un schéma noethérien. Un morphisme propre
étant de type fini, $X$ est alors lui aussi un schéma
noethérien. Considérons la propriété $\mathcal{P}$ des faisceaux cohérents
sur $X$ définie par
$$
\mathcal{P}(\mathcal{F}) \Leftrightarrow
R^pf_*\mathcal{F}\text{ est cohérent pour tout }p \geq 0
$$
Nous allons utiliser le
lemme \ref{lemma-property} pour montrer que
$\mathcal{P}$ est vérifiée par tout faisceau cohérent sur $X$.

\medskip\noindent
Soit
$$
0 \to \mathcal{F}_1 \to \mathcal{F}_2 \to \mathcal{F}_3 \to 0
$$
une suite exacte courte de faisceaux cohérents sur $X$.
Considérons la suite exacte longue des images directes supérieures
$$
R^{p - 1}f_*\mathcal{F}_3 \to
R^pf_*\mathcal{F}_1 \to
R^pf_*\mathcal{F}_2 \to
R^pf_*\mathcal{F}_3 \to
R^{p + 1}f_*\mathcal{F}_1
$$
Il est alors clair que, si deux des trois faisceaux $\mathcal{F}_i$
vérifient la propriété $\mathcal{P}$, les images directes supérieures du troisième figurent,
dans ce complexe exact, entre deux faisceaux cohérents. Ces
images directes supérieures sont donc elles aussi cohérentes, d'après les
lemmes \ref{lemma-coherent-abelian-Noetherian} et
\ref{lemma-coherent-Noetherian-quasi-coherent-sub-quotient}.
La propriété $\mathcal{P}$ est donc aussi vérifiée par le troisième.

\medskip\noindent
Soit $Z \subset X$ un sous-schéma fermé intègre.
Il faut trouver un faisceau cohérent $\mathcal{F}$ sur $X$ dont le support soit
contenu dans $Z$ et dont la fibre au point générique $\xi$ de $Z$ soit un espace vectoriel de dimension
$1$ sur $\kappa(\xi)$, tel que $\mathcal{P}$
soit vérifiée par $\mathcal{F}$. Notons $g = f|_Z : Z \to S$ la restriction de $f$.
Supposons que l'on puisse trouver un faisceau cohérent $\mathcal{G}$ sur $Z$ tel
que
(a) $\mathcal{G}_\xi$ soit un espace vectoriel de dimension $1$ sur $\kappa(\xi)$,
(b) $R^pg_*\mathcal{G} = 0$ pour $p > 0$, et
(c) $g_*\mathcal{G}$ soit cohérent. On peut alors considérer
$\mathcal{F} = (Z \to X)_*\mathcal{G}$. Puisque $Z \to X$ est une
immersion fermée, $(Z \to X)_*\mathcal{G}$ est cohérent sur $X$
et $R^p(Z \to X)_*\mathcal{G} = 0$ pour $p > 0$
(lemme \ref{lemma-finite-pushforward-coherent}).
D'après la suite spectrale de Leray relative
(Cohomologie, lemme \ref{cohomology-lemma-relative-Leray}),
on aura donc $R^pf_*\mathcal{F} = R^pg_*\mathcal{G} = 0$ pour $p > 0$
et $f_*\mathcal{F} = g_*\mathcal{G}$ sera cohérent.
Enfin, $\mathcal{F}_\xi = ((Z \to X)_*\mathcal{G})_\xi = \mathcal{G}_\xi$,
ce qui vérifie la condition sur la fibre en $\xi$.
Tout se ramène donc à trouver un faisceau cohérent $\mathcal{G}$
sur $Z$ vérifiant les propriétés (a), (b) et (c).

\medskip\noindent
On peut appliquer le Lemme de Chow \ref{lemma-chow-Noetherian}
au morphisme $Z \to S$. On obtient ainsi un diagramme
$$
\xymatrix{
Z \ar[rd]_g & Z' \ar[d]^-{g'} \ar[l]^\pi \ar[r]_i & \mathbf{P}^m_S \ar[dl] \\
& S &
}
$$
comme dans l'énoncé du lemme de Chow. Soit aussi $U \subset Z$ le
sous-schéma ouvert dense tel que $\pi^{-1}(U) \to U$ soit un
isomorphisme. D'après la remarque \ref{remark-chow-Noetherian},
$i' = (i, \pi) : Z' \to \mathbf{P}^m_Z$ est
une immersion fermée. Ainsi
$$
\mathcal{L} = i^*\mathcal{O}_{\mathbf{P}^m_S}(1) \cong
(i')^*\mathcal{O}_{\mathbf{P}^m_Z}(1)
$$
est relativement ample pour $g'$ et pour $\pi$ (par exemple d'après Morphismes,
lemme \ref{morphisms-lemma-characterize-ample-on-finite-type}).
D'après le lemme \ref{lemma-kill-by-twisting},
il existe donc $n \geq 0$ tel que l'on ait à
la fois $R^p\pi_*\mathcal{L}^{\otimes n} = 0$ pour tout $p > 0$ et
$R^p(g')_*\mathcal{L}^{\otimes n} = 0$ pour tout $p > 0$.
Posons $\mathcal{G} = \pi_*\mathcal{L}^{\otimes n}$.
La propriété (a) est vérifiée, car $\pi_*\mathcal{L}^{\otimes n}|_U$ est
un faisceau inversible (puisque $\pi^{-1}(U) \to U$ est un isomorphisme).
Les propriétés (b) et (c) sont vérifiées, car la suite spectrale de
Leray relative
(Cohomologie, lemme \ref{cohomology-lemma-relative-Leray})
donne
$$
E_2^{p, q} = R^pg_* R^q\pi_*\mathcal{L}^{\otimes n}
\Rightarrow
R^{p + q}(g')_*\mathcal{L}^{\otimes n}
$$
et, par le choix de $n$, les seuls termes non nuls de $E_2^{p, q}$ sont
ceux pour lesquels $q = 0$, tandis que les seuls termes non nuls de
$R^{p + q}(g')_*\mathcal{L}^{\otimes n}$ sont ceux pour lesquels $p = q = 0$.
Cela entraîne que $R^pg_*\mathcal{G} = 0$ pour $p > 0$ et
que $g_*\mathcal{G} = (g')_*\mathcal{L}^{\otimes n}$.
Enfin, en appliquant le
lemme \ref{lemma-locally-projective-pushforward},
on voit que $g_*\mathcal{G} = (g')_*\mathcal{L}^{\otimes n}$ est
cohérent, comme voulu.
\end{proof}

\begin{lemma}
\label{lemma-proper-over-affine-cohomology-finite}
Soit $S = \Spec(A)$, où $A$ est un anneau noethérien.
Soit $f : X \to S$ un morphisme propre.
Soit $\mathcal{F}$ un $\mathcal{O}_X$-module cohérent.
Alors $H^i(X, \mathcal{F})$ est un $A$-module de type fini pour tout $i \geq 0$.
\end{lemma}

\begin{proof}
C'est simplement le cas affine de la
proposition \ref{proposition-proper-pushforward-coherent}. En
effet, d'après les lemmes \ref{lemma-quasi-coherence-higher-direct-images} et
\ref{lemma-quasi-coherence-higher-direct-images-application},
$R^if_*\mathcal{F}$ est le faisceau quasi-cohérent
associé au $A$-module $H^i(X, \mathcal{F})$;
d'après le lemme \ref{lemma-coherent-Noetherian}, ce
faisceau est cohérent si et seulement si $H^i(X, \mathcal{F})$
est un $A$-module de type fini.
\end{proof}

\begin{lemma}
\label{lemma-graded-finiteness}
Soit $A$ un anneau noethérien.
Soit $B$ une $A$-algèbre graduée de
type fini. Soit $f : X \to \Spec(A)$ un morphisme propre.
Posons $\mathcal{B} = f^*\widetilde B$.
Soit $\mathcal{F}$
un $\mathcal{B}$-module gradué quasi-cohérent de type fini.
\begin{enumerate}
\item Pour tout $p \geq 0$, le $B$-module gradué $H^p(X, \mathcal{F})$
est un $B$-module de type fini.
\item Si $\mathcal{L}$ est un $\mathcal{O}_X$-module
inversible ample, il existe un entier $d_0$ tel que
$H^p(X, \mathcal{F} \otimes \mathcal{L}^{\otimes d}) = 0$
pour tout $p > 0$ et tout $d \geq d_0$.
\end{enumerate}
\end{lemma}

\begin{proof}
Pour le démontrer, considérons le carré cartésien
$$
\xymatrix{
X' = \Spec(B) \times_{\Spec(A)} X
\ar[r]_-\pi \ar[d]_{f'} &
X \ar[d]^f \\
\Spec(B) \ar[r] &
\Spec(A)
}
$$
Notons que $f'$ est un morphisme propre, voir
Morphismes, lemme \ref{morphisms-lemma-base-change-proper}. De
plus, $B$ est de type fini comme $A$-algèbre, et est donc
noethérien (Algèbre, lemme \ref{algebra-lemma-Noetherian-permanence}).
Par suite, $X'$ est un schéma noethérien
(Morphismes, lemme \ref{morphisms-lemma-finite-type-noetherian}).
Notons que $X'$ est le spectre relatif de la
$\mathcal{O}_X$-algèbre quasi-cohérente $\mathcal{B}$,
d'après Constructions, lemme \ref{constructions-lemma-spec-properties}.
Puisque $\mathcal{F}$ est un $\mathcal{B}$-module
quasi-cohérent, il existe un unique
$\mathcal{O}_{X'}$-module quasi-cohérent $\mathcal{F}'$ tel que
$\pi_*\mathcal{F}' = \mathcal{F}$, voir
Morphismes, lemme \ref{morphisms-lemma-affine-equivalence-modules}.
Puisque $\mathcal{F}$ est de type fini comme $\mathcal{B}$
-module, $\mathcal{F}'$ est un
$\mathcal{O}_{X'}$-module de type fini (détails omis). Autrement dit,
$\mathcal{F}'$ est un $\mathcal{O}_{X'}$-module cohérent
(lemme \ref{lemma-coherent-Noetherian}).
Puisque le morphisme $\pi : X' \to X$ est affine, on a
$$
H^p(X, \mathcal{F}) = H^p(X', \mathcal{F}')
$$
d'après le lemme \ref{lemma-relative-affine-cohomology}.
Ainsi (1) résulte du
lemme \ref{lemma-proper-over-affine-cohomology-finite}. Étant
donné $\mathcal{L}$ comme dans (2), posons
$\mathcal{L}' = \pi^*\mathcal{L}$. Notons que $\mathcal{L}'$ est
ample sur $X'$,
d'après Morphismes, lemme \ref{morphisms-lemma-pullback-ample-tensor-relatively-ample}.
La formule de projection
(Cohomologie, lemme \ref{cohomology-lemma-projection-formula}) donne
$\pi_*(\mathcal{F}' \otimes \mathcal{L}') = \mathcal{F} \otimes \mathcal{L}$.
L'assertion (2) se déduit donc, par le même raisonnement que ci-dessus, du
lemme \ref{lemma-kill-by-twisting}.
\end{proof}












\section{Le théorème des fonctions formelles}
\label{section-theorem-formal-functions}

\noindent
Dans cette section, nous étudions le comportement de la cohomologie
des suites de faisceaux de la forme $\{I^n\mathcal{F}\}_{n \geq 0}$
ou de la forme $\{\mathcal{F}/I^n\mathcal{F}\}_{n \geq 0}$ lorsque
$n$ varie.

\medskip\noindent
Nous utiliserons ici et dans la suite la notation suivante.
Étant donnés un morphisme de schémas $f : X \to Y$, un faisceau quasi-cohérent
$\mathcal{F}$ sur $X$ et un faisceau quasi-cohérent d'idéaux
$\mathcal{I} \subset \mathcal{O}_Y$, on note
$\mathcal{I}^n\mathcal{F}$ le sous-faisceau quasi-cohérent engendré
par les produits de sections locales de $f^{-1}(\mathcal{I}^n)$ et de
$\mathcal{F}$. Autrement dit,
$$
\mathcal{I}^n\mathcal{F}
=
\Im\left(
f^*(\mathcal{I}^n) \otimes_{\mathcal{O}_X} \mathcal{F}
\longrightarrow
\mathcal{F}
\right).
$$
Notons qu'il existe des applications naturelles
$$
f^{-1}(\mathcal{I}^n) \otimes_{f^{-1}\mathcal{O}_Y} \mathcal{I}^m\mathcal{F}
\longrightarrow
f^*(\mathcal{I}^n) \otimes_{\mathcal{O}_X} \mathcal{I}^m\mathcal{F}
\longrightarrow
\mathcal{I}^{n + m}\mathcal{F}
$$
Ainsi, une section de $\mathcal{I}^n$ définit une
application $R^pf_*(\mathcal{I}^m\mathcal{F}) \to
R^pf_*(\mathcal{I}^{n + m}\mathcal{F})$ par fonctorialité
des images directes supérieures. Par localisation puis faisceautisation, on
obtient des applications de $\mathcal{O}_Y$-modules
$$
\mathcal{I}^n \otimes_{\mathcal{O}_Y} R^pf_*(\mathcal{I}^m\mathcal{F})
\longrightarrow
R^pf_*(\mathcal{I}^{n + m}\mathcal{F}).
$$
Autrement dit,
$\bigoplus_{n \geq 0} R^pf_*(\mathcal{I}^n\mathcal{F})$
est un $\bigoplus_{n \geq 0} \mathcal{I}^n$-module gradué.

\medskip\noindent
Si $Y = \Spec(A)$ et $\mathcal{I} = \widetilde{I}$, on note
$\mathcal{I}^n\mathcal{F}$ simplement $I^n\mathcal{F}$. Les applications
introduites ci-dessus munissent $M = \bigoplus H^p(X, I^n\mathcal{F})$ d'une
structure de module gradué sur $S = \bigoplus I^n$. Si $f$ est propre,
$A$ noethérien et $\mathcal{F}$ cohérent, ce module
est de type fini.

\begin{lemma}
\label{lemma-cohomology-powers-ideal-times-F}
\begin{reference}
\cite[III Cor 3.3.2]{EGA}
\end{reference}
Soit $A$ un anneau noethérien.
Soit $I \subset A$ un idéal.
Posons $B = \bigoplus_{n \geq 0} I^n$.
Soit $f : X \to \Spec(A)$ un morphisme propre.
Soit $\mathcal{F}$ un faisceau cohérent sur $X$.
Alors, pour tout $p \geq 0$, le $B$-module gradué
$\bigoplus_{n \geq 0} H^p(X, I^n\mathcal{F})$ est
un $B$-module de type fini.
\end{lemma}

\begin{proof}
Posons $\mathcal{B} = f^*\widetilde{B}$. Puisque $\mathcal{B}$ admet une
surjection sur $\bigoplus I^n\mathcal{O}_X$, il est clair que
$\bigoplus I^n\mathcal{F}$ est un $\mathcal{B}$-module gradué
de type fini. Le résultat découle donc du lemme \ref{lemma-graded-finiteness} (1).
\end{proof}

\begin{lemma}
\label{lemma-cohomology-powers-ideal-times-sheaf}
Soient un morphisme de schémas $f : X \to Y$, un faisceau quasi-cohérent
$\mathcal{F}$ sur $X$ et un faisceau quasi-cohérent d'idéaux
$\mathcal{I} \subset \mathcal{O}_Y$. Supposons $Y$ localement
noethérien, $f$ propre et $\mathcal{F}$ cohérent.
Alors
$$
\mathcal{M} =
\bigoplus\nolimits_{n \geq 0} R^pf_*(\mathcal{I}^n\mathcal{F})
$$
est un $\mathcal{A} = \bigoplus_{n \geq 0} \mathcal{I}^n$-module
gradué quasi-cohérent de type fini.
\end{lemma}

\begin{proof}
L'assertion est locale sur $Y$; on se ramène donc au
cas où $Y$ est affine. Dans ce cas, le résultat découle
du lemme \ref{lemma-cohomology-powers-ideal-times-F}.
Détails omis.
\end{proof}

\begin{lemma}
\label{lemma-cohomology-powers-ideal-application}
Soit $A$ un anneau noethérien.
Soit $I \subset A$ un idéal.
Soit $f : X \to \Spec(A)$ un morphisme propre.
Soit $\mathcal{F}$ un faisceau cohérent sur $X$.
Alors, pour tout $p \geq 0$, il existe un entier $c \geq 0$
tel que
\begin{enumerate}
\item l'application de multiplication
$I^{n - c} \otimes H^p(X, I^c\mathcal{F}) \to H^p(X, I^n\mathcal{F})$
soit surjective pour tout $n \geq c$;
\item l'image de $H^p(X, I^{n + m}\mathcal{F}) \to H^p(X, I^n\mathcal{F})$
soit contenue dans le sous-module $I^{m - e} H^p(X, I^n\mathcal{F})$,
où $e = \max(0, c - n)$, pour $n + m \geq c$, $n, m \geq 0$;
\item on ait
$$
\Ker(H^p(X, I^n\mathcal{F}) \to H^p(X, \mathcal{F})) =
\Ker(H^p(X, I^n\mathcal{F}) \to H^p(X, I^{n - c}\mathcal{F}))
$$
pour $n \geq c$;
\item il existe des applications $I^nH^p(X, \mathcal{F}) \to H^p(X, I^{n - c}\mathcal{F})$
pour $n \geq c$, telles que les composés
$$
H^p(X, I^n\mathcal{F}) \to
I^{n - c}H^p(X, \mathcal{F}) \to
H^p(X, I^{n - 2c}\mathcal{F})
$$
et
$$
I^nH^p(X, \mathcal{F}) \to
H^p(X, I^{n - c}\mathcal{F}) \to
I^{n - 2c}H^p(X, \mathcal{F})
$$
pour $n \geq 2c$ soient les applications canoniques; et
\item les systèmes projectifs $(H^p(X, I^n\mathcal{F}))$ et
$(I^nH^p(X, \mathcal{F}))$ soient pro-isomorphes.
\end{enumerate}
\end{lemma}

\begin{proof}
Posons $M_n = H^p(X, I^n\mathcal{F})$ pour $n \geq 1$ et
$M_0 = H^p(X, \mathcal{F})$, de sorte que l'on ait des applications
$\ldots \to M_3 \to M_2 \to M_1 \to M_0$. En posant
$B = \bigoplus_{n \geq 0} I^n$, on voit que $M = \bigoplus_{n \geq 0} M_n$
est un $B$-module gradué de type fini,
voir lemme \ref{lemma-cohomology-powers-ideal-times-F}.
Notons que les produits
$B_n \otimes M_m \to M_{m + n}$, $a \otimes m \mapsto a \cdot m$,
sont compatibles avec les applications de notre système projectif, au sens
où les diagrammes
$$
\xymatrix{
B_n \otimes_A M_m \ar[r] \ar[d] & M_{n + m} \ar[d] \\
B_n \otimes_A M_{m'} \ar[r] & M_{n + m'}
}
$$
commutent pour $n, m' \geq 0$ et $m \geq m'$.

\medskip\noindent
Démontrons (1). Choisissons $d_1, \ldots, d_t \geq 0$ et $x_i \in M_{d_i}$
tels que $M$ soit engendré par $x_1, \ldots, x_t$ sur $B$.
Pour tout $c \geq \max\{d_i\}$, on a alors
$B_{n - c} \cdot M_c = M_n$ pour $n \geq c$, ce qui établit (1).

\medskip\noindent
Démontrons (2). Soit $c$ comme dans la démonstration de (1). Soit $n + m \geq c$.
On a $M_{n + m} = B_{n + m - c} \cdot M_c$.
Si $c > n$, on utilise $M_c \to M_n$ et la compatibilité des produits
avec les applications de transition signalée ci-dessus pour en déduire
que l'image de $M_{n + m} \to M_n$ est contenue dans $I^{n + m - c}M_n$.
Si $c \leq n$, on écrit $M_{n + m} = B_m \cdot B_{n - c} \cdot M_c =
B_m \cdot M_n$, ce qui montre que l'image est contenue dans $I^m M_n$.
Cela démontre (2).

\medskip\noindent
Soit $K_n \subset M_n$ le noyau de l'application $M_n \to M_0$. La
compatibilité des produits avec les applications de transition signalée ci-dessus
montre que $K = \bigoplus K_n \subset M$ est un sous-$B$-module gradué.
Puisque $B$ est noethérien et que $M$ est un $B$-module
gradué de type fini, $K$ est un $B$-module gradué de
type fini. Choisissons $d'_1, \ldots, d'_{t'} \geq 0$ et $y_i \in K_{d'_i}$
tels que $K$ soit engendré par $y_1, \ldots, y_{t'}$ sur $B$.
Posons $c = \max(d'_i, d'_j)$. Puisque $y_i \in \Ker(M_{d'_i} \to M_0)$,
on a $B_n \cdot y_i \subset \Ker(M_{n + d'_i} \to M_n)$.
On en déduit que $K_n = \Ker(M_n \to M_{n - c})$ pour $n \geq c$.
Cela démontre (3).

\medskip\noindent
Considérons le diagramme commutatif suivant, dont les flèches pleines sont données :
$$
\xymatrix{
I^n \otimes_A M_0 \ar[r] \ar[d] &
I^nM_0 \ar[r] \ar@{..>}[d] &
M_0 \ar[d] \\
M_n \ar[r] &
M_{n - c} \ar[r] &
M_0
}
$$
Puisque le noyau de la flèche surjective $I^n \otimes_A M_0 \to I^nM_0$
a son image dans $K_n$, ce qui précède fournit la flèche pointillée, et le
composé $I^nM_0 \to M_{n - c} \to M_0$ est l'application canonique. Il est
alors clair que le composé $I^nM_0 \to M_{n - c} \to I^{n - 2c}M_0$
est l'application canonique pour $n \geq 2c$. Considérons le composé
$M_n \to I^{n - c}M_0 \to M_{n - 2c}$. La première application
envoie un élément de la forme $a \cdot m$,
avec $a \in I^{n - c}$ et $m \in M_c$,
sur $a m'$, où $m'$ est l'image de $m$ dans $M_0$.
La seconde application envoie cet élément sur $a \cdot m'$ dans $M_{n - 2c}$,
ce qui établit (4).

\medskip\noindent
L'assertion (5) résulte immédiatement de (4) et de la définition des
morphismes de pro-objets.
\end{proof}

\noindent
Dans la situation des lemmes \ref{lemma-cohomology-powers-ideal-times-F} et
\ref{lemma-cohomology-powers-ideal-application}, considérons le système
projectif
$$
\mathcal{F}/I\mathcal{F} \leftarrow
\mathcal{F}/I^2\mathcal{F} \leftarrow
\mathcal{F}/I^3\mathcal{F} \leftarrow \ldots
$$
Nous voudrions connaître le comportement des groupes de cohomologie.
Voici un premier résultat.

\begin{lemma}
\label{lemma-ML-cohomology-powers-ideal}
Soit $A$ un anneau noethérien.
Soit $I \subset A$ un idéal.
Soit $f : X \to \Spec(A)$ un morphisme propre.
Soit $\mathcal{F}$ un faisceau cohérent sur $X$.
Fixons $p \geq 0$. Il existe $c \geq 0$ tel que
\begin{enumerate}
\item pour tout $n \geq c$, on ait
$$
\Ker(H^p(X, \mathcal{F}) \to H^p(X, \mathcal{F}/I^n\mathcal{F})) \subset
I^{n - c}H^p(X, \mathcal{F}).
$$
\item le système projectif
$$
\left(H^p(X, \mathcal{F}/I^n\mathcal{F})\right)_{n \in \mathbf{N}}
$$
vérifie la condition de Mittag-Leffler (voir
Homologie, définition \ref{homology-definition-Mittag-Leffler}); et
\item on ait
$$
\Im(H^p(X, \mathcal{F}/I^k\mathcal{F})
\to H^p(X, \mathcal{F}/I^n\mathcal{F}))
=
\Im(H^p(X, \mathcal{F})
\to H^p(X, \mathcal{F}/I^n\mathcal{F}))
$$
pour tout $k \geq n + c$.
\end{enumerate}
\end{lemma}

\begin{proof}
Posons $c = \max\{c_p, c_{p + 1}\}$, où $c_p, c_{p + 1}$ sont les entiers obtenus
dans le lemme \ref{lemma-cohomology-powers-ideal-application} pour
$H^p$ et $H^{p + 1}$.

\medskip\noindent
Démontrons (1). Considérons la suite exacte courte
$$
0 \to I^n\mathcal{F} \to \mathcal{F} \to \mathcal{F}/I^n\mathcal{F} \to 0
$$
La suite exacte longue de cohomologie donne
$$
\Ker(
H^p(X, \mathcal{F}) \to H^p(X, \mathcal{F}/I^n\mathcal{F})
)
=
\Im(
H^p(X, I^n\mathcal{F}) \to H^p(X, \mathcal{F})
)
$$
D'après le lemme \ref{lemma-cohomology-powers-ideal-application} (2), ce
sous-module est donc contenu dans $I^{n - c}H^p(X, \mathcal{F})$ pour $n \geq c$.

\medskip\noindent
Notons que (3) entraîne (2), par définition des systèmes de
Mittag-Leffler.

\medskip\noindent
Démontrons (3). Fixons $n$. Considérons le diagramme commutatif
$$
\xymatrix{
0 \ar[r] &
I^n\mathcal{F} \ar[r] &
\mathcal{F} \ar[r] &
\mathcal{F}/I^n\mathcal{F} \ar[r] &
0 \\
0 \ar[r] &
I^{n + m}\mathcal{F} \ar[r] \ar[u] &
\mathcal{F} \ar[r] \ar[u] &
\mathcal{F}/I^{n + m}\mathcal{F} \ar[r] \ar[u] &
0
}
$$
Il induit le diagramme commutatif suivant :
$$
\xymatrix{
H^p(X, \mathcal{F}) \ar[r] &
H^p(X, \mathcal{F}/I^n\mathcal{F}) \ar[r]_\delta &
H^{p + 1}(X, I^n\mathcal{F}) \ar[r] &
H^{p + 1}(X, \mathcal{F}) \\
H^p(X, \mathcal{F}) \ar[r] \ar[u]^1 &
H^p(X, \mathcal{F}/I^{n + m}\mathcal{F}) \ar[r] \ar[u]^\gamma &
H^{p + 1}(X, I^{n + m}\mathcal{F}) \ar[u]^\alpha \ar[r]^-\beta &
H^{p + 1}(X, \mathcal{F}) \ar[u]_1
}
$$
dont les lignes sont exactes. D'après
le lemme \ref{lemma-cohomology-powers-ideal-application} (4), le noyau
de $\beta$ est égal à celui de $\alpha$ pour $m \geq c$.
Une chasse au diagramme montre alors que l'image de $\gamma$ est contenue
dans le noyau de $\delta$, ce qui établit (3) (en
posant $k = n + m$).
\end{proof}

\begin{theorem}[Théorème des fonctions formelles]
\label{theorem-formal-functions}
Soit $A$ un anneau noethérien.
Soit $I \subset A$ un idéal.
Soit $f : X \to \Spec(A)$ un morphisme propre.
Soit $\mathcal{F}$ un faisceau cohérent sur $X$.
Fixons $p \geq 0$.
Le système d'applications
$$
H^p(X, \mathcal{F})/I^nH^p(X, \mathcal{F})
\longrightarrow
H^p(X, \mathcal{F}/I^n\mathcal{F})
$$
définit un isomorphisme de limites
$$
H^p(X, \mathcal{F})^\wedge
\longrightarrow
\lim_n H^p(X, \mathcal{F}/I^n\mathcal{F})
$$
où le membre de gauche est le complété du $A$-module
$H^p(X, \mathcal{F})$ relativement à l'idéal $I$, voir
Algèbre, section \ref{algebra-section-completion}.
De plus, cet isomorphisme est un homéomorphisme pour les topologies de limite projective.
\end{theorem}

\begin{proof}
Cela se déduit du lemme \ref{lemma-ML-cohomology-powers-ideal} comme suit.
Posons $M = H^p(X, \mathcal{F})$, $M_n = H^p(X, \mathcal{F}/I^n\mathcal{F})$,
et notons $N_n = \Im(M \to M_n)$. D'après le
lemme \ref{lemma-ML-cohomology-powers-ideal} (2) et (3),
$(M_n)$ est un système de Mittag-Leffler dans lequel
$N_n \subset M_n$ est l'image de $M_k$ pour tout $k \gg n$.
Il en résulte que $\lim M_n = \lim N_n$ comme modules topologiques (munis des
topologies de limite projective). D'autre part, les $N_n$ forment un système projectif de
quotients du module $M$; par suite, $\lim N_n$ est le complété de $M$
pour la topologie définie par les noyaux $K_n = \Ker(M \to N_n)$.
D'après le lemme \ref{lemma-ML-cohomology-powers-ideal} (1), on
a $K_n \subset I^{n - c}M$ et, puisque $N_n \subset M_n$
est annulé par $I^n$, on a $I^n M \subset K_n$.
La topologie définie par les sous-modules $K_n$
comme système fondamental de voisinages ouverts de $0$
est donc la topologie $I$-adique, et l'application
induite $M^\wedge = \lim M/I^nM \to \lim N_n = \lim M_n$
est un isomorphisme de modules topologiques\footnote{
Précisons que la topologie de limite projective sur $M^\wedge$ coïncide avec sa topologie
$I$-adique, d'après Algèbre,
lemme \ref{algebra-lemma-hathat-finitely-generated}.
Voir Compléments d'algèbre, section \ref{more-algebra-section-topological-ring}.}.
\end{proof}

\begin{lemma}
\label{lemma-spell-out-theorem-formal-functions}
Soit $A$ un anneau. Soit $I \subset A$ un idéal. Supposons $A$ noethérien
et complet relativement à $I$.
Soit $f : X \to \Spec(A)$ un morphisme propre.
Soit $\mathcal{F}$ un faisceau cohérent sur $X$.
Alors
$$
H^p(X, \mathcal{F}) = \lim_n H^p(X, \mathcal{F}/I^n\mathcal{F})
$$
pour tout $p \geq 0$.
\end{lemma}

\begin{proof}
C'est une reformulation du théorème des fonctions formelles
(Théorème \ref{theorem-formal-functions}) dans le
cas d'un anneau de base noethérien complet. En effet, dans ce cas, le
$A$-module $H^p(X, \mathcal{F})$ est de type
fini (lemme \ref{lemma-proper-over-affine-cohomology-finite}), donc complet pour la topologie
$I$-adique (Algèbre, lemme \ref{algebra-lemma-completion-tensor});
il est donc inutile de compléter le membre de gauche.
\end{proof}

\begin{lemma}
\label{lemma-formal-functions-stalk}
Soient un morphisme de schémas $f : X \to Y$ et un faisceau quasi-cohérent
$\mathcal{F}$ sur $X$. Supposons que
\begin{enumerate}
\item $Y$ soit localement noethérien,
\item $f$ soit propre, et
\item $\mathcal{F}$ soit cohérent.
\end{enumerate}
Soit $y \in Y$ un point. Considérons les voisinages infinitésimaux
$$
\xymatrix{
X_n =
\Spec(\mathcal{O}_{Y, y}/\mathfrak m_y^n) \times_Y X
\ar[r]_-{i_n} \ar[d]_{f_n} &
X \ar[d]^f \\
\Spec(\mathcal{O}_{Y, y}/\mathfrak m_y^n) \ar[r]^-{c_n} & Y
}
$$
de la fibre $X_1 = X_y$ et posons $\mathcal{F}_n = i_n^*\mathcal{F}$.
On a alors
$$
\left(R^pf_*\mathcal{F}\right)_y^\wedge
\cong
\lim_n H^p(X_n, \mathcal{F}_n)
$$
comme $\mathcal{O}_{Y, y}^\wedge$-modules.
\end{lemma}

\begin{proof}
C'est simplement une reformulation d'un cas particulier du théorème des
fonctions formelles, Théorème \ref{theorem-formal-functions}.
Détaillons-la. Notons que $\mathcal{O}_{Y, y}$ est un anneau
local noethérien. Considérons le morphisme canonique
$c : \Spec(\mathcal{O}_{Y, y}) \to Y$, voir
Schémas, Équation (\ref{schemes-equation-canonical-morphism}).
Ce morphisme est plat, car il identifie les anneaux locaux.
Notons provisoirement $f' : X' \to \Spec(\mathcal{O}_{Y, y})$
le changement de base de $f$ à cet anneau local. On a
$c^*R^pf_*\mathcal{F} = R^pf'_*\mathcal{F}'$, d'après
le lemme \ref{lemma-flat-base-change-cohomology}.
De plus, les voisinages infinitésimaux
des fibres $X_y$ et $X'_y$ s'identifient (vérification omise; indication :
les morphismes $c_n$ se factorisent par $c$).

\medskip\noindent
On peut donc supposer que $Y = \Spec(A)$ est le spectre d'un
anneau local noethérien $A$ d'idéal maximal $\mathfrak m$,
et que $y \in Y$ est le point fermé (c'est-à-dire celui correspondant à $\mathfrak m$).
Il en résulte en particulier que
$$
\left(R^pf_*\mathcal{F}\right)_y =
\Gamma(Y, R^pf_*\mathcal{F}) =
H^p(X, \mathcal{F}).
$$

\medskip\noindent
Dans ce cas, les morphismes $c_n$ sont aussi des immersions
fermées. Leurs changements de base $i_n$ sont donc eux aussi des immersions
fermées. Notons que $i_{n, *}\mathcal{F}_n = i_{n, *}i_n^*\mathcal{F}
= \mathcal{F}/\mathfrak m^n\mathcal{F}$. La suite spectrale de Leray
pour $i_n$ et le lemme \ref{lemma-finite-pushforward-coherent} donnent
$$
H^p(X_n, \mathcal{F}_n) =
H^p(X, i_{n, *}\mathcal{F}_n) =
H^p(X, \mathcal{F}/\mathfrak m^n\mathcal{F})
$$
On peut donc appliquer le théorème des fonctions formelles
pour calculer la limite de l'énoncé, ce qui conclut.
\end{proof}

\noindent
Voici un lemme que nous généraliserons aux fibres de
dimension $ > 0$ dans le lemme suivant.

\begin{lemma}
\label{lemma-higher-direct-images-zero-finite-fibre}
Soit $f : X \to Y$ un morphisme de schémas.
Soit $y \in Y$. Supposons
que
\begin{enumerate}
\item $Y$ soit localement noethérien,
\item $f$ soit propre, et
\item $f^{-1}(\{y\})$ soit fini.
\end{enumerate}
Alors, pour tout faisceau cohérent $\mathcal{F}$ sur $X$, on a
$(R^pf_*\mathcal{F})_y = 0$ pour tout $p > 0$.
\end{lemma}

\begin{proof}
La fibre $X_y$ est finie et, d'après
Morphismes, lemme \ref{morphisms-lemma-finite-fibre},
c'est un espace discret fini. De
plus, chaque voisinage infinitésimal $X_n$ a le même
espace topologique sous-jacent. Chacun des schémas $X_n$ est donc affine, d'après
Schémas, lemme \ref{schemes-lemma-scheme-finite-discrete-affine}.
Il en résulte que $H^p(X_n, \mathcal{F}_n) = 0$ pour tout
$p > 0$. Ainsi $(R^pf_*\mathcal{F})_y^\wedge = 0$,
d'après le lemme \ref{lemma-formal-functions-stalk}.
Notons que $R^pf_*\mathcal{F}$ est cohérent, d'après la
proposition \ref{proposition-proper-pushforward-coherent}; par
suite, $R^pf_*\mathcal{F}_y$ est un
$\mathcal{O}_{Y, y}$-module de type fini. D'après le lemme
de Nakayama (Algèbre, lemme \ref{algebra-lemma-NAK}),
si le complété d'un module de type fini sur un anneau
local est nul, le module est nul. D'où
$(R^pf_*\mathcal{F})_y = 0$.
\end{proof}

\begin{lemma}
\label{lemma-higher-direct-images-zero-above-dimension-fibre}
Soit $f : X \to Y$ un morphisme de schémas.
Soit $y \in Y$. Supposons
que
\begin{enumerate}
\item $Y$ soit localement noethérien,
\item $f$ soit propre, et
\item $\dim(X_y) = d$.
\end{enumerate}
Alors, pour tout faisceau cohérent $\mathcal{F}$ sur $X$, on a
$(R^pf_*\mathcal{F})_y = 0$ pour tout $p > d$.
\end{lemma}

\begin{proof}
La fibre $X_y$ est de type fini sur $\Spec(\kappa(y))$.
Ainsi $X_y$ est un schéma noethérien, d'après
Morphismes, lemme \ref{morphisms-lemma-finite-type-noetherian}.
L'espace topologique sous-jacent de $X_y$ est donc noethérien, voir
Propriétés, lemme \ref{properties-lemma-Noetherian-topology}.
De plus, chaque voisinage
infinitésimal $X_n$ a le même espace topologique sous-jacent que $X_y$.
Ainsi $H^p(X_n, \mathcal{F}_n) = 0$ pour tout $p > d$,
d'après Cohomologie, proposition \ref{cohomology-proposition-vanishing-Noetherian}.
On en déduit que $(R^pf_*\mathcal{F})_y^\wedge = 0$,
d'après le lemme \ref{lemma-formal-functions-stalk}, pour $p > d$.
Notons que $R^pf_*\mathcal{F}$ est cohérent, d'après la
proposition \ref{proposition-proper-pushforward-coherent}; par
suite, $R^pf_*\mathcal{F}_y$ est un
$\mathcal{O}_{Y, y}$-module de type fini. D'après le lemme
de Nakayama (Algèbre, lemme \ref{algebra-lemma-NAK}),
si le complété d'un module de type fini sur un anneau
local est nul, le module est nul. D'où
$(R^pf_*\mathcal{F})_y = 0$.
\end{proof}








\section{Applications du théorème sur les fonctions formelles}
\label{section-applications-formal-functions}


\noindent
Nous compléterons cette section au besoin. Pour le moment, nous avons besoin de la
caractérisation suivante des morphismes finis dans le cas noethérien.

\begin{lemma}
\label{lemma-characterize-finite}
(Pour une version plus générale, voir Compléments
sur les morphismes, lemme \ref{more-morphisms-lemma-characterize-finite}.)
Soit $f : X \to S$ un morphisme de schémas. Supposons
que $S$ soit localement noethérien. Les
assertions suivantes sont équivalentes :
\begin{enumerate}
\item $f$ est fini, et
\item $f$ est propre à fibres finies.
\end{enumerate}
\end{lemma}

\begin{proof}
Un morphisme fini est propre d'après
Morphismes, lemme \ref{morphisms-lemma-finite-proper}.
Un morphisme fini est quasi-fini d'après
Morphismes, lemme \ref{morphisms-lemma-finite-quasi-finite}.
Un morphisme quasi-fini a des fibres finies, voir
Morphismes, lemme \ref{morphisms-lemma-quasi-finite}.
Ainsi, un morphisme fini est propre et a des fibres finies.

\medskip\noindent
Supposons que $f$ soit propre à fibres finies.
Nous voulons montrer que $f$ est fini.
En fait, il suffit de montrer que $f$ est affine.
En effet, si $f$ est affine, alors
$f$ est intégral d'après
Morphismes, lemme \ref{morphisms-lemma-integral-universally-closed},
et il résulte alors
de Morphismes, lemme \ref{morphisms-lemma-finite-integral},
que $f$ est fini.

\medskip\noindent
Pour montrer que $f$ est affine, nous pouvons supposer que $S$ est affine, et notre
but est alors de montrer que $X$ est également affine.
Puisque $f$ est propre, on voit que $X$ est séparé et quasi-compact.
Nous pouvons donc utiliser le critère du
lemme \ref{lemma-quasi-separated-h1-zero-covering} pour montrer que $X$
est affine. Pour cela, soit $\mathcal{I} \subset \mathcal{O}_X$
un faisceau d'idéaux de type fini. En particulier, $\mathcal{I}$ est
un faisceau cohérent sur $X$. D'après le
lemme \ref{lemma-higher-direct-images-zero-finite-fibre}, nous concluons que
$R^1f_*\mathcal{I}_s = 0$ pour tout $s \in S$.
Autrement dit, $R^1f_*\mathcal{I} = 0$. La suite spectrale
de Leray pour $f$ donne donc
$H^1(X , \mathcal{I}) = H^1(S, f_*\mathcal{I})$.
Comme $S$ est affine et que $f_*\mathcal{I}$ est quasi-cohérent
(Schémas, lemme \ref{schemes-lemma-push-forward-quasi-coherent}),
nous déduisons $H^1(S, f_*\mathcal{I}) = 0$
du lemme \ref{lemma-quasi-coherent-affine-cohomology-zero},
comme voulu. Ainsi, $H^1(X, \mathcal{I}) = 0$, comme voulu.
\end{proof}

\noindent
On en déduit le résultat utile suivant.

\begin{lemma}
\label{lemma-proper-finite-fibre-finite-in-neighbourhood}
\begin{slogan}
Un morphisme propre est fini au voisinage d'une fibre finie.
\end{slogan}
(Pour une version plus générale, voir Compléments
sur les morphismes,
lemme \ref{more-morphisms-lemma-proper-finite-fibre-finite-in-neighbourhood}.)
Soit $f : X \to S$ un morphisme de schémas.
Soit $s \in S$. Supposons
que
\begin{enumerate}
\item $S$ est localement noethérien,
\item $f$ est propre, et
\item $f^{-1}(\{s\})$ est un ensemble fini.
\end{enumerate}
Il existe alors un voisinage ouvert $V \subset S$ de $s$
tel que $f|_{f^{-1}(V)} : f^{-1}(V) \to V$ soit fini.
\end{lemma}

\begin{proof}
Le morphisme $f$ est quasi-fini en tout point de $f^{-1}(\{s\})$
d'après Morphismes, lemme \ref{morphisms-lemma-finite-fibre}.
D'après Morphismes, lemme \ref{morphisms-lemma-quasi-finite-points-open},
l'ensemble des points où $f$ est quasi-fini est un ouvert $U \subset X$.
Posons $Z = X \setminus U$. Alors $s \not \in f(Z)$. Puisque $f$ est propre,
l'ensemble $f(Z) \subset S$ est fermé. Choisissons un voisinage ouvert
$V \subset S$ de $s$ tel que $Z \cap V = \emptyset$. Alors
$f^{-1}(V) \to V$ est localement quasi-fini et propre.
Il est donc quasi-fini
(Morphismes, lemme \ref{morphisms-lemma-quasi-finite-locally-quasi-compact}),
donc à fibres finies
(Morphismes, lemme \ref{morphisms-lemma-quasi-finite}), et par
conséquent fini d'après le lemme \ref{lemma-characterize-finite}.
\end{proof}

\begin{lemma}
\label{lemma-ample-on-fibre}
Soit $f : X \to Y$ un morphisme propre de schémas, avec $Y$ noethérien.
Soit $\mathcal{L}$ un $\mathcal{O}_X$-module inversible.
Soit $\mathcal{F}$ un $\mathcal{O}_X$-module cohérent.
Soit $y \in Y$ un point tel que $\mathcal{L}_y$ soit ample sur $X_y$.
Il existe alors un entier $d_0$ tel que, pour tout $d \geq d_0$, on ait
$$
R^pf_*(\mathcal{F} \otimes_{\mathcal{O}_X} \mathcal{L}^{\otimes d})_y = 0
\text{ pour }p > 0
$$
et l'application
$$
f_*(\mathcal{F} \otimes_{\mathcal{O}_X} \mathcal{L}^{\otimes d})_y
\longrightarrow
H^0(X_y, \mathcal{F}_y \otimes_{\mathcal{O}_{X_y}} \mathcal{L}_y^{\otimes d})
$$
est surjective.
\end{lemma}

\begin{proof}
Remarquons que $\mathcal{O}_{Y, y}$ est un anneau local noethérien.
Considérons le morphisme canonique
$c : \Spec(\mathcal{O}_{Y, y}) \to Y$, voir
Schémas, Équation (\ref{schemes-equation-canonical-morphism}).
Ce morphisme est plat, puisqu'il identifie les anneaux locaux.
Notons provisoirement $f' : X' \to \Spec(\mathcal{O}_{Y, y})$
le changement de base de $f$ à cet anneau local. On a
$c^*R^pf_*\mathcal{F} = R^pf'_*\mathcal{F}'$ d'après le
lemme \ref{lemma-flat-base-change-cohomology}.
De plus, les fibres $X_y$ et $X'_y$ sont identifiées.
Nous pouvons donc supposer que $Y = \Spec(A)$ est le spectre d'un
anneau local noethérien $(A, \mathfrak m, \kappa)$ et que $y \in Y$
correspond à $\mathfrak m$. Dans ce cas,
$R^pf_*(\mathcal{F} \otimes_{\mathcal{O}_X} \mathcal{L}^{\otimes d})_y =
H^p(X, \mathcal{F} \otimes_{\mathcal{O}_X} \mathcal{L}^{\otimes d})$
pour tout $p \geq 0$. Notons $f_y : X_y \to \Spec(\kappa)$ la projection.

\medskip\noindent
Posons $B = \text{Gr}_\mathfrak m(A) =
\bigoplus_{n \geq 0} \mathfrak m^n/\mathfrak m^{n + 1}$.
Considérons le faisceau $\mathcal{B} = f_y^*\widetilde{B}$
de $\mathcal{O}_{X_y}$-algèbres
graduées quasi-cohérentes. Nous utiliserons les notations de la section \ref{section-theorem-formal-functions},
où $I$ est remplacé par $\mathfrak m$.
Comme $X_y$ est le sous-schéma fermé de $X$ défini par
$\mathfrak m\mathcal{O}_X$, nous pouvons considérer
$\mathfrak m^n\mathcal{F}/\mathfrak m^{n + 1}\mathcal{F}$
comme un $\mathcal{O}_{X_y}$-module cohérent, voir
le lemme \ref{lemma-i-star-equivalence}. Alors
$\bigoplus_{n \geq 0} \mathfrak m^n\mathcal{F}/\mathfrak m^{n + 1}\mathcal{F}$
est un $\mathcal{B}$-module gradué quasi-cohérent de type
fini, car il est engendré en degré zéro sur $\mathcal{B}$
et car sa partie de degré zéro est
$\mathcal{F}_y = \mathcal{F}/\mathfrak m \mathcal{F}$
qui est un $\mathcal{O}_{X_y}$-module cohérent.
Ainsi, d'après le lemme \ref{lemma-graded-finiteness}, partie
(2), on voit que
$$
H^p(X_y, \mathfrak m^n \mathcal{F}/ \mathfrak m^{n + 1}\mathcal{F}
\otimes_{\mathcal{O}_{X_y}} \mathcal{L}_y^{\otimes d}) = 0
$$
pour tous $p > 0$, $d \geq d_0$ et $n \geq 0$. D'après le
lemme \ref{lemma-relative-affine-cohomology},
cela revient à dire que
$
H^p(X, \mathfrak m^n \mathcal{F}/ \mathfrak m^{n + 1}\mathcal{F}
\otimes_{\mathcal{O}_X} \mathcal{L}^{\otimes d}) = 0
$
pour tous $p > 0$, $d \geq d_0$ et $n \geq 0$.

\medskip\noindent
Considérons les suites exactes courtes
$$
0 \to \mathfrak m^n\mathcal{F}/\mathfrak m^{n + 1} \mathcal{F}
\to \mathcal{F}/\mathfrak m^{n + 1} \mathcal{F}
\to \mathcal{F}/\mathfrak m^n \mathcal{F} \to 0
$$
de $\mathcal{O}_X$-modules cohérents. La tensorisation par $\mathcal{L}^{\otimes d}$
est un foncteur exact, et nous obtenons les suites exactes courtes
$$
0 \to
\mathfrak m^n\mathcal{F}/\mathfrak m^{n + 1} \mathcal{F}
\otimes_{\mathcal{O}_X} \mathcal{L}^{\otimes d}
\to \mathcal{F}/\mathfrak m^{n + 1} \mathcal{F}
\otimes_{\mathcal{O}_X} \mathcal{L}^{\otimes d}
\to \mathcal{F}/\mathfrak m^n \mathcal{F}
\otimes_{\mathcal{O}_X} \mathcal{L}^{\otimes d} \to 0
$$
À l'aide de la suite exacte longue de cohomologie et de l'annulation
ci-dessus, nous concluons (par récurrence) que
\begin{enumerate}
\item $H^p(X, \mathcal{F}/\mathfrak m^n \mathcal{F}
\otimes_{\mathcal{O}_X} \mathcal{L}^{\otimes d}) = 0$
pour tous $p > 0$, $d \geq d_0$ et $n \geq 0$, et
\item $H^0(X, \mathcal{F}/\mathfrak m^n \mathcal{F}
\otimes_{\mathcal{O}_X} \mathcal{L}^{\otimes d}) \to
H^0(X_y, \mathcal{F}_y \otimes_{\mathcal{O}_{X_y}} \mathcal{L}_y^{\otimes d})$
est surjective pour tous $d \geq d_0$ et $n \geq 1$.
\end{enumerate}
D'après le théorème sur les fonctions formelles (Théorème \ref{theorem-formal-functions}),
on voit que la complétion $\mathfrak m$-adique de
$H^p(X, \mathcal{F} \otimes_{\mathcal{O}_X} \mathcal{L}^{\otimes d})$
est nulle pour tous $d \geq d_0$ et $p > 0$.
Comme $H^p(X, \mathcal{F} \otimes_{\mathcal{O}_X} \mathcal{L}^{\otimes d})$
est un $A$-module de type fini d'après
le lemme \ref{lemma-proper-over-affine-cohomology-finite},
le lemme de Nakayama (Algèbre, lemme \ref{algebra-lemma-NAK}) entraîne
que $H^p(X, \mathcal{F} \otimes_{\mathcal{O}_X} \mathcal{L}^{\otimes d})$
est nul pour tous $d \geq d_0$ et $p > 0$.
Pour $p = 0$, nous déduisons du
lemme \ref{lemma-ML-cohomology-powers-ideal}, partie (3),
que $H^0(X, \mathcal{F} \otimes_{\mathcal{O}_X} \mathcal{L}^{\otimes d}) \to
H^0(X_y, \mathcal{F}_y \otimes_{\mathcal{O}_{X_y}} \mathcal{L}_y^{\otimes d})$
est surjective, ce qui donne la dernière assertion du lemme.
\end{proof}

\begin{lemma}
\label{lemma-ample-in-neighbourhood}
(Pour une version plus générale, voir Compléments
sur les morphismes,
lemme \ref{more-morphisms-lemma-ample-in-neighbourhood}.)
Soit $f : X \to Y$ un morphisme propre de schémas, avec $Y$ noethérien.
Soit $\mathcal{L}$ un $\mathcal{O}_X$-module inversible.
Soit $y \in Y$ un point tel que $\mathcal{L}_y$ soit ample
sur $X_y$. Il existe alors un voisinage ouvert $V \subset Y$
de $y$ tel que $\mathcal{L}|_{f^{-1}(V)}$ soit ample sur $f^{-1}(V)/V$.
\end{lemma}

\begin{proof}
Choisissons $d_0$ comme dans le lemme \ref{lemma-ample-on-fibre} pour
$\mathcal{F} = \mathcal{O}_X$. Choisissons $d \geq d_0$
de sorte que l'on puisse trouver $r \geq 0$ et des sections
$s_{y, 0}, \ldots, s_{y, r} \in H^0(X_y, \mathcal{L}_y^{\otimes d})$
qui définissent une immersion fermée
$$
\varphi_y =
\varphi_{\mathcal{L}_y^{\otimes d}, (s_{y, 0}, \ldots, s_{y, r})} :
X_y \to \mathbf{P}^r_{\kappa(y)}.
$$
Cela est possible d'après Morphismes, Lemme
\ref{morphisms-lemma-finite-type-over-affine-ample-very-ample},
mais nous utilisons également
Morphismes, lemme \ref{morphisms-lemma-image-proper-scheme-closed},
pour voir que $\varphi_y$ est une immersion fermée, et
Constructions, section \ref{constructions-section-projective-space},
pour la description des morphismes vers l'espace
projectif en termes de faisceaux inversibles et de sections.
D'après notre choix de $d_0$, après avoir remplacé $Y$ par un voisinage ouvert
de $y$, nous pouvons choisir
$s_0, \ldots, s_r \in H^0(X, \mathcal{L}^{\otimes d})$
dont les images sont $s_{y, 0}, \ldots, s_{y, r}$.
Soit $X_{s_i} \subset X$ l'ouvert où $s_i$
est un générateur de $\mathcal{L}^{\otimes d}$. Puisque
les $s_{y, i}$ engendrent $\mathcal{L}_y^{\otimes d}$, on voit que
$X_y \subset U = \bigcup X_{s_i}$. Puisque
le morphisme $X \to Y$ est fermé, il existe
un voisinage ouvert $y \in V \subset Y$
tel que $f^{-1}(V) \subset U$.
Après avoir remplacé $Y$ par $V$, nous pouvons supposer que
les $s_i$ engendrent $\mathcal{L}^{\otimes d}$. Nous obtenons
ainsi un morphisme
$$
\varphi = \varphi_{\mathcal{L}^{\otimes d}, (s_0, \ldots, s_r)} :
X \longrightarrow \mathbf{P}^r_Y
$$
tel que $\mathcal{L}^{\otimes d} \cong \varphi^*\mathcal{O}_{\mathbf{P}^r_Y}(1)$
et dont le changement de base à $y$ donne $\varphi_y$.

\medskip\noindent
Nous terminerons la démonstration par une astuce ; la démonstration « correcte »
consiste à montrer directement que $\varphi$ est une immersion
fermée après changement de base à un voisinage ouvert de $y$. En
effet, d'après le lemme \ref{lemma-proper-finite-fibre-finite-in-neighbourhood},
on voit que $\varphi$ est fini au-dessus d'un voisinage ouvert
de la fibre $\mathbf{P}^r_{\kappa(y)}$ de $\mathbf{P}^r_Y \to Y$
au-dessus de $y$. Puisque $\mathbf{P}^r_Y \to Y$ est fermé, quitte à
rétrécir $Y$, nous pouvons supposer que $\varphi$ est fini.
Alors $\mathcal{L}^{\otimes d} \cong \varphi^*\mathcal{O}_{\mathbf{P}^r_Y}(1)$
est ample d'après le résultat très général suivant:
Morphismes, lemme \ref{morphisms-lemma-pullback-ample-tensor-relatively-ample}.
\end{proof}





\section{Cohomologie et changement de base, III}
\label{section-cohomology-and-base-change-perfect}

\noindent
Dans cette section, nous démontrons le cas le plus simple d'un phénomène très
général qui sera étudié dans
Catégories dérivées des schémas, section
\ref{perfect-section-cohomology-and-base-change-perfect}.
Voir la remarque \ref{remark-explain-perfect-direct-image} pour une
traduction algébrique du lemme suivant.

\begin{lemma}
\label{lemma-perfect-direct-image}
Soit $A$ un anneau noethérien et posons $S = \Spec(A)$. Soit $f : X \to S$ un morphisme
propre de schémas. Soit $\mathcal{F}$ un
$\mathcal{O}_X$-module cohérent et plat sur $S$. Alors
\begin{enumerate}
\item $R\Gamma(X, \mathcal{F})$ est un objet parfait de $D(A)$, et
\item pour tout homomorphisme d'anneaux $A \to A'$, le morphisme de changement de base
$$
R\Gamma(X, \mathcal{F}) \otimes_A^{\mathbf{L}} A'
\longrightarrow
R\Gamma(X_{A'}, \mathcal{F}_{A'})
$$
est un isomorphisme.
\end{enumerate}
\end{lemma}

\begin{proof}
Choisissons un recouvrement ouvert affine fini $X = \bigcup_{i = 1, \ldots, n} U_i$.
D'après les lemmes \ref{lemma-separated-case-relative-cech} et
\ref{lemma-base-change-complex}, le complexe de {\v C}ech
$K^\bullet = {\check C}^\bullet(\mathcal{U}, \mathcal{F})$ vérifie
$$
K^\bullet \otimes_A A' = R\Gamma(X_{A'}, \mathcal{F}_{A'})
$$
pour tout homomorphisme d'anneaux $A \to A'$. Soit
$K_{alt}^\bullet = {\check C}_{alt}^\bullet(\mathcal{U}, \mathcal{F})$
le complexe de {\v C}ech alterné. D'après
Cohomologie, lemme \ref{cohomology-lemma-alternating-usual},
il existe une équivalence d'homotopie $K_{alt}^\bullet \to K^\bullet$
de $A$-modules. En particulier, on a
$$
K_{alt}^\bullet \otimes_A A' = R\Gamma(X_{A'}, \mathcal{F}_{A'})
$$
également. Comme $\mathcal{F}$ est plat sur $A$, on voit que chaque $K_{alt}^n$
est plat sur $A$ (voir
Morphismes, lemme \ref{morphisms-lemma-flat-module-characterize}).
Comme, de plus, $K_{alt}^\bullet$ est borné supérieurement (c'est pourquoi nous sommes passés
au complexe de {\v C}ech alterné),
$K_{alt}^\bullet \otimes_A A' = K_{alt}^\bullet \otimes_A^{\mathbf{L}} A'$
d'après la définition des produits tensoriels dérivés (voir
Compléments sur l'algèbre, section \ref{more-algebra-section-derived-tensor-product}).
D'après le
lemme \ref{lemma-proper-over-affine-cohomology-finite},
les groupes de cohomologie $H^i(K_{alt}^\bullet)$ sont des $A$-modules de
type fini. Comme $K_{alt}^\bullet$ est borné, nous en déduisons que $K_{alt}^\bullet$
est pseudo-cohérent, voir
Compléments sur l'algèbre, lemme \ref{more-algebra-lemma-Noetherian-pseudo-coherent}.
Pour tout $A$-module $M$, posons $A' = A \oplus M$, où $M$ est un idéal de carré
nul, c'est-à-dire $(a, m) \cdot (a', m') = (aa', am' + a'm)$. D'après ce qui précède,
on voit que la cohomologie de $K_{alt}^\bullet \otimes_A^\mathbf{L} A'$ est concentrée
dans les degrés $0, \ldots, n$. De même, celle de $K_{alt}^\bullet \otimes_A^\mathbf{L} M$
est concentrée dans les degrés $0, \ldots, n$. Par conséquent, $K_{alt}^\bullet$ a
une dimension de Tor finie,
voir Compléments sur l'algèbre, définition \ref{more-algebra-definition-tor-amplitude}.
La conclusion résulte de Compléments sur l'algèbre, lemme \ref{more-algebra-lemma-perfect}.
\end{proof}

\begin{remark}
\label{remark-explain-perfect-direct-image}
Une conséquence du lemme \ref{lemma-perfect-direct-image} est qu'il
existe un complexe fini de $A$-modules projectifs de type fini $M^\bullet$ tel
que l'on ait
$$
H^i(X_{A'}, \mathcal{F}_{A'}) = H^i(M^\bullet \otimes_A A')
$$
fonctoriellement en $A'$. La condition que $\mathcal{F}$ soit
plat sur $A$ est essentielle, voir \cite{Hartshorne}.
\end{remark}







\section{Modules formels cohérents}
\label{section-coherent-formal}

\noindent
Comme nous ne disposons pas encore de la théorie des schémas formels, nous développons
un peu de terminologie qui remplace la notion de ``module cohérent sur un
schéma formel adique noethérien''.

\medskip\noindent
Soit $X$ un schéma noethérien et soit $\mathcal{I} \subset \mathcal{O}_X$
un faisceau quasi-cohérent d'idéaux. Nous considérerons des systèmes
projectifs $(\mathcal{F}_n)$ de $\mathcal{O}_X$-modules cohérents tels que
\begin{enumerate}
\item $\mathcal{F}_n$ soit annulé par $\mathcal{I}^n$, et
\item les morphismes de transition induisent des isomorphismes
$\mathcal{F}_{n + 1}/\mathcal{I}^n\mathcal{F}_{n + 1} \to \mathcal{F}_n$.
\end{enumerate}
Un morphisme de tels systèmes projectifs est défini de la
manière habituelle. Notons la catégorie de ces systèmes projectifs
$\textit{Coh}(X, \mathcal{I})$. Nous allons établir une série de lemmes
sur les objets de cette catégorie. En fait, la plupart des
lemmes qui suivent sont des conséquences immédiates de la description
suivante de la catégorie dans le cas affine.

\begin{lemma}
\label{lemma-inverse-systems-affine}
Si $X = \Spec(A)$ est le spectre d'un anneau noethérien et si
$\mathcal{I}$ est le faisceau quasi-cohérent d'idéaux associé à l'idéal
$I \subset A$, alors $\textit{Coh}(X, \mathcal{I})$ est équivalente à
la catégorie des $A^\wedge$-modules de type fini, où $A^\wedge$ est le complété
de $A$ par rapport à $I$.
\end{lemma}

\begin{proof}
Soit $\text{Mod}^{fg}_{A, I}$ la catégorie des systèmes projectifs $(M_n)$
de $A$-modules de type fini satisfaisant : (1) $M_n$ est annulé par $I^n$ et (2)
$M_{n + 1}/I^nM_{n + 1} = M_n$. D'après la correspondance entre les faisceaux
cohérents sur $X$ et les $A$-modules de type fini (lemme \ref{lemma-coherent-Noetherian}),
il suffit de montrer que $\text{Mod}^{fg}_{A, I}$ est équivalente à la catégorie
des $A^\wedge$-modules de type fini. Pour le voir, il suffit de prouver qu'étant
donné un objet $(M_n)$ de $\text{Mod}^{fg}_{A, I}$, le module
$$
M = \lim M_n
$$
est un $A^\wedge$-module de type fini et que $M/I^nM = M_n$. Les morphismes de
transition étant surjectifs, nous voyons que $M \to M_1$ est surjectif.
Choisissons $x_1, \ldots, x_t \in M$ dont les images engendrent $M_1$.
Cela induit un morphisme de systèmes $(A/I^n)^{\oplus t} \to M_n$.
D'après le lemme de Nakayama (Algèbre, lemme \ref{algebra-lemma-NAK}), ces morphismes sont
surjectifs. Soit $K_n \subset (A/I^n)^{\oplus t}$ le noyau. La
propriété (2) implique que $K_{n + 1} \to K_n$ est surjectif ; en particulier,
le système $(K_n)$ satisfait la condition de Mittag-Leffler.
D'après Homologie, lemme \ref{homology-lemma-Mittag-Leffler},
nous obtenons une suite exacte
$0 \to K \to (A^\wedge)^{\oplus t} \to M \to 0$
où $K = \lim K_n$.
Ainsi, $M$ est un $A^\wedge$-module de
type fini. Comme $K \to K_n$ est surjectif, il s'ensuit que
$$
M/I^nM = \Coker(K \to (A/I^n)^{\oplus t}) =
(A/I^n)^{\oplus t}/K_n = M_n
$$
comme souhaité.
\end{proof}

\begin{lemma}
\label{lemma-inverse-systems-abelian}
Soit $X$ un schéma noethérien et soit $\mathcal{I} \subset \mathcal{O}_X$
un faisceau quasi-cohérent d'idéaux.
\begin{enumerate}
\item La catégorie $\textit{Coh}(X, \mathcal{I})$ est abélienne.
\item Pour tout ouvert $U \subset X$, le foncteur de restriction
$\textit{Coh}(X, \mathcal{I}) \to \textit{Coh}(U, \mathcal{I}|_U)$
est exact.
\item L'exactitude dans $\textit{Coh}(X, \mathcal{I})$ peut se vérifier
après restriction aux ouverts d'un recouvrement ouvert de $X$.
\end{enumerate}
\end{lemma}

\begin{proof}
Soit $\alpha =(\alpha_n) : (\mathcal{F}_n) \to (\mathcal{G}_n)$ un morphisme dans
$\textit{Coh}(X, \mathcal{I})$. Le conoyau de $\alpha$ est le système projectif
$(\Coker(\alpha_n))$ (détails omis). Pour décrire le noyau, posons
$$
\mathcal{K}'_{l, m} = \Im(\Ker(\alpha_l) \to \mathcal{F}_m)
$$
pour $l \geq m$. Nous
affirmons que :
\begin{enumerate}
\item[(a)] le système projectif $(\mathcal{K}'_{l, m})_{l \geq m}$ est constant
à partir d'un certain rang, disons de valeur $\mathcal{K}'_m$,
\item[(b)] le système $(\mathcal{K}'_m/\mathcal{I}^n\mathcal{K}'_m)_{m \geq n}$
est constant à partir d'un certain rang, disons de valeur $\mathcal{K}_n$,
\item[(c)] le système $(\mathcal{K}_n)$ constitue un objet de
$\textit{Coh}(X, \mathcal{I})$, et
\item[(d)] cet objet est le noyau de $\alpha$.
\end{enumerate}
Pour établir (a), (b) et (c), nous pouvons raisonner localement sur un ouvert affine, disons $X = \Spec(A)$,
et $\mathcal{I}$ correspond à l'idéal $I \subset A$. D'après le
lemme \ref{lemma-inverse-systems-affine},
$\alpha$ correspond à un morphisme $f : M \to N$ de $A^\wedge$-modules de
type fini. Notons $K = \Ker(f)$. Remarquons que $A^\wedge$ est un anneau
noethérien (Algèbre, lemme \ref{algebra-lemma-completion-Noetherian-Noetherian}).
Choisissons un entier $c \geq 0$ tel que
$K \cap I^n M \subset I^{n - c}K$ pour $n \geq c$
(Algèbre, lemme \ref{algebra-lemma-Artin-Rees})
et choisissons-le de sorte que le lemme \ref{algebra-lemma-map-AR}
d'Algèbre s'applique pour le morphisme $f$ et l'idéal $I^\wedge = IA^\wedge$. Alors
$\mathcal{K}'_{l, m}$ correspond au $A$-module
$$
K'_{l, m} = \frac{a^{-1}(I^lN) + I^mM}{I^mM} =
\frac{K + I^{l - c}f^{-1}(I^cN) + I^mM}{I^mM} =
\frac{K + I^mM}{I^mM}
$$
où la dernière égalité a lieu si $l \geq m + c$. Ainsi, $\mathcal{K}'_m$
correspond au $A$-module $K/K \cap I^mM$ et
$\mathcal{K}'_m/\mathcal{I}^n\mathcal{K}'_m$ correspond à
$$
\frac{K}{K \cap I^mM + I^nK} = \frac{K}{I^nK}
$$
pour $m \geq n + c$, d'après le choix de $c$ ci-dessus. Par conséquent, $\mathcal{K}_n$
correspond à $K/I^nK$.

\medskip\noindent
Démontrons (d). Il résulte clairement de la description affine ci-dessus que
la composée $(\mathcal{K}_n) \to (\mathcal{F}_n) \to (\mathcal{G}_n)$
est nulle. Soit $\beta : (\mathcal{H}_n) \to (\mathcal{F}_n)$
un morphisme tel que $\alpha \circ \beta = 0$. Alors l'image de
$\mathcal{H}_l \to \mathcal{F}_l$ est contenue dans $\Ker(\alpha_l)$.
Puisque $\mathcal{H}_m = \mathcal{H}_l/\mathcal{I}^m\mathcal{H}_l$
pour $l \geq m$, nous obtenons un système de morphismes
$\mathcal{H}_m \to \mathcal{K}'_{l, m}$. D'où un morphisme
$\mathcal{H}_m \to \mathcal{K}_m'$. Puisque
$\mathcal{H}_n = \mathcal{H}_m/\mathcal{I}^n\mathcal{H}_m$, nous
obtenons un système de morphismes $\mathcal{H}_n \to \mathcal{K}'_m/\mathcal{I}^n\mathcal{K}'_m$
et donc un morphisme $\mathcal{H}_n \to \mathcal{K}_n$, comme souhaité.

\medskip\noindent
Pour achever la preuve de (1), il reste à montrer que $\Coim = \Im$
dans $\textit{Coh}(X, \mathcal{I})$. Nous avons vu ci-dessus que la formation des
noyaux et des conoyaux est compatible, sur les ouverts affines, avec la description
de $\textit{Coh}(X, \mathcal{I})$ comme catégorie de modules. Puisque l'égalité
$\Im = \Coim$ est vérifiée dans la catégorie des modules,
nous obtenons $\Coim = \Im$ dans $\textit{Coh}(X, \mathcal{I})$.
Les assertions (2) et (3) du lemme découlent immédiatement de notre construction des
noyaux et des conoyaux.
\end{proof}

\begin{lemma}
\label{lemma-inverse-systems-surjective}
Soit $X$ un schéma noethérien et soit $\mathcal{I} \subset \mathcal{O}_X$
un faisceau quasi-cohérent d'idéaux. Un morphisme
$(\mathcal{F}_n) \to (\mathcal{G}_n)$ est surjectif dans
$\textit{Coh}(X, \mathcal{I})$
si et seulement si $\mathcal{F}_1 \to \mathcal{G}_1$ est surjectif.
\end{lemma}

\begin{proof}
Omis. Indication : se placer sur les ouverts affines, utiliser
le lemme \ref{lemma-inverse-systems-affine}, ainsi que
Algèbre, lemme \ref{algebra-lemma-NAK}.
\end{proof}

\noindent
Soit $X$ un schéma noethérien et soit $\mathcal{I} \subset \mathcal{O}_X$
un faisceau quasi-cohérent d'idéaux. Il existe un foncteur
\begin{equation}
\label{equation-completion-functor}
\textit{Coh}(\mathcal{O}_X) \longrightarrow \textit{Coh}(X, \mathcal{I}), \quad
\mathcal{F} \longmapsto \mathcal{F}^\wedge
\end{equation}
qui associe au $\mathcal{O}_X$-module cohérent $\mathcal{F}$
l'objet $\mathcal{F}^\wedge = (\mathcal{F}/\mathcal{I}^n\mathcal{F})$
de $\textit{Coh}(X, \mathcal{I})$.

\begin{lemma}
\label{lemma-exact}
Le foncteur (\ref{equation-completion-functor}) est exact.
\end{lemma}

\begin{proof}
Il suffit de le vérifier localement sur $X$. Nous pouvons donc supposer que $X$ est
affine, c'est-à-dire que nous sommes dans la situation du
lemme \ref{lemma-inverse-systems-affine}.
Le foncteur est le foncteur $\text{Mod}^{fg}_A \to \text{Mod}^{fg}_{A^\wedge}$
qui associe à un $A$-module de type fini $M$ sa complétion $M^\wedge$.
Le résultat découle donc de
Algèbre, lemme \ref{algebra-lemma-completion-flat}.
\end{proof}

\begin{lemma}
\label{lemma-completion-internal-hom}
Soit $X$ un schéma noethérien et soit $\mathcal{I} \subset \mathcal{O}_X$
un faisceau quasi-cohérent d'idéaux. Soient $\mathcal{F}$, $\mathcal{G}$
des $\mathcal{O}_X$-modules cohérents. Posons
$\mathcal{H} = \SheafHom_{\mathcal{O}_X}(\mathcal{G}, \mathcal{F})$.
Alors
$$
\lim H^0(X, \mathcal{H}/\mathcal{I}^n\mathcal{H}) =
\Mor_{\textit{Coh}(X, \mathcal{I})}
(\mathcal{G}^\wedge, \mathcal{F}^\wedge).
$$
\end{lemma}

\begin{proof}
Pour le démontrer, nous pouvons travailler localement sur des ouverts affines de $X$.
Nous pouvons donc supposer $X = \Spec(A)$ et $\mathcal{F}$, $\mathcal{G}$
donnés par les $A$-modules de type fini $M$ et $N$. Alors $\mathcal{H}$
correspond au $A$-module de type fini $H = \Hom_A(M, N)$.
L'énoncé du lemme devient l'égalité
$$
H^\wedge = \Hom_{A^\wedge}(M^\wedge, N^\wedge)
$$
par l'équivalence du lemme \ref{lemma-inverse-systems-affine}.
D'après Algèbre, lemme \ref{algebra-lemma-completion-flat}
(utilisé 3 fois), nous avons
$$
H^\wedge = \Hom_A(M, N) \otimes_A A^\wedge =
\Hom_{A^\wedge}(M \otimes_A A^\wedge, N \otimes_A A^\wedge) =
\Hom_{A^\wedge}(M^\wedge, N^\wedge)
$$
où la deuxième égalité utilise le fait que $A^\wedge$ est plat sur $A$
(voir Compléments d'algèbre, Lemme
\ref{more-algebra-lemma-pseudo-coherence-and-base-change-ext}).
Le lemme en découle.
\end{proof}

\noindent
Soit $X$ un schéma noethérien. Soit $\mathcal{I} \subset \mathcal{O}_X$
un faisceau quasi-cohérent d'idéaux. Nous disons qu'un objet $(\mathcal{F}_n)$ de
$\textit{Coh}(X, \mathcal{I})$ est {\it de torsion annulée par une puissance de $\mathcal{I}$}
ou est {\it annulé par une puissance de $\mathcal{I}$} s'il existe
un $c \geq 1$ tel que $\mathcal{F}_n = \mathcal{F}_c$ pour tout $n \geq c$.
Dans ce cas, nous dirons que $(\mathcal{F}_n)$ est {\it annulé par
$\mathcal{I}^c$}. Si $X = \Spec(A)$ est affine, alors, par l'équivalence du
lemme \ref{lemma-inverse-systems-affine},
ces objets correspondent exactement aux $A$-modules de
type fini annulés par une puissance de $I$ ou par $I^c$.

\begin{lemma}
\label{lemma-existence-easy}
Soit $X$ un schéma noethérien et soit $\mathcal{I} \subset \mathcal{O}_X$
un faisceau quasi-cohérent d'idéaux. Soit $\mathcal{G}$ un
$\mathcal{O}_X$-module cohérent. Soit $(\mathcal{F}_n)$ un objet de
$\textit{Coh}(X, \mathcal{I})$.
\begin{enumerate}
\item Si $\alpha : (\mathcal{F}_n) \to \mathcal{G}^\wedge$ est
un morphisme dont le noyau et le conoyau sont annulés par une puissance de $\mathcal{I}$,
alors il existe un triplet unique (à isomorphisme unique près)
$(\mathcal{F}, a, \beta)$ où
\begin{enumerate}
\item $\mathcal{F}$ est un $\mathcal{O}_X$-module cohérent,
\item $a : \mathcal{F} \to \mathcal{G}$ est un morphisme de $\mathcal{O}_X$-modules dont
le noyau et le conoyau sont annulés par une puissance de $\mathcal{I}$,
\item $\beta : (\mathcal{F}_n) \to \mathcal{F}^\wedge$ est un isomorphisme, et
\item $\alpha = a^\wedge \circ \beta$.
\end{enumerate}
\item Si $\alpha : \mathcal{G}^\wedge \to (\mathcal{F}_n)$ est
un morphisme dont le noyau et le conoyau sont annulés par une puissance de $\mathcal{I}$,
alors il existe un triplet unique (à isomorphisme unique près)
$(\mathcal{F}, a, \beta)$ où
\begin{enumerate}
\item $\mathcal{F}$ est un $\mathcal{O}_X$-module cohérent,
\item $a : \mathcal{G} \to \mathcal{F}$ est un morphisme de $\mathcal{O}_X$-modules dont
le noyau et le conoyau sont annulés par une puissance de $\mathcal{I}$,
\item $\beta : \mathcal{F}^\wedge \to (\mathcal{F}_n)$ est un isomorphisme, et
\item $\alpha = \beta \circ a^\wedge$.
\end{enumerate}
\end{enumerate}
\end{lemma}

\begin{proof}
Démonstration de (1). L'unicité montre qu'il suffit de construire
$(\mathcal{F}, a, \beta)$
localement pour la topologie de Zariski sur $X$. Nous pouvons donc supposer $X = \Spec(A)$ et que $\mathcal{I}$
correspond à l'idéal $I \subset A$. Dans cette situation, le
lemme \ref{lemma-inverse-systems-affine} s'applique.
Soit $M'$ le $A^\wedge$-module de type fini
correspondant à $(\mathcal{F}_n)$. Soit $N$ le $A$-module de type fini correspondant à
$\mathcal{G}$. Alors $\alpha$ correspond à un morphisme
$$
\varphi : M' \longrightarrow N^\wedge
$$
dont le noyau et le conoyau sont annulés par $I^t$ pour un certain $t$. Rappelons que
$N^\wedge = N \otimes_A A^\wedge$
(Algèbre, lemme \ref{algebra-lemma-completion-tensor}).
D'après Compléments d'algèbre, lemme \ref{more-algebra-lemma-application-formal-glueing},
il existe un morphisme de $A$-modules $\psi : M \to N$ dont le noyau et le conoyau sont de torsion annulée par une puissance de
$I$, et un isomorphisme
$M \otimes_A A^\wedge = M'$ compatible avec $\varphi$.
Comme $N$ et $M'$ sont des modules de type fini, nous concluons que $M$
est un $A$-module de type
fini, voir Compléments d'algèbre, remarque \ref{more-algebra-remark-formal-glueing-algebras}.
Ainsi $M \otimes_A A^\wedge = M^\wedge$. Nous omettons la vérification
que le triplet $(M, N \to M, M^\wedge \to M')$ ainsi obtenu
est unique à isomorphisme unique près.

\medskip\noindent
La démonstration de (2) est exactement la même et nous l'omettons.
\end{proof}

\begin{lemma}
\label{lemma-torsion-hom-ext}
Soit $X$ un schéma noethérien et soit $\mathcal{I} \subset \mathcal{O}_X$
un faisceau quasi-cohérent d'idéaux. Tout objet de
$\textit{Coh}(X, \mathcal{I})$ qui est annulé
par une puissance de $\mathcal{I}$ appartient à l'image essentielle de
(\ref{equation-completion-functor}).
De plus, si $\mathcal{F}$, $\mathcal{G}$ appartiennent à $\textit{Coh}(\mathcal{O}_X)$
et si $\mathcal{F}$ ou $\mathcal{G}$ est annulé par une puissance de
$\mathcal{I}$, alors les applications
$$
\xymatrix{
\Hom_X(\mathcal{F}, \mathcal{G}) \ar[d] &
\Ext_X(\mathcal{F}, \mathcal{G}) \ar[d] \\
\Hom_{\textit{Coh}(X, \mathcal{I})}(\mathcal{F}^\wedge, \mathcal{G}^\wedge) &
\Ext_{\textit{Coh}(X, \mathcal{I})}(\mathcal{F}^\wedge, \mathcal{G}^\wedge)
}
$$
sont des isomorphismes.
\end{lemma}

\begin{proof}
Supposons que $(\mathcal{F}_n)$ soit un objet de $\textit{Coh}(X, \mathcal{I})$
annulé par $\mathcal{I}^c$ pour un certain $c \geq 1$. Alors
$\mathcal{F}_n \to \mathcal{F}_c$ est un isomorphisme pour $n \geq c$.
Ainsi, si nous posons $\mathcal{F} = \mathcal{F}_c$, nous voyons que
$\mathcal{F}^\wedge \cong (\mathcal{F}_n)$. Cela prouve la première assertion.

\medskip\noindent
Soient $\mathcal{F}$, $\mathcal{G}$ des objets de $\textit{Coh}(\mathcal{O}_X)$
tels que $\mathcal{F}$ ou $\mathcal{G}$ soit annulé par
$\mathcal{I}^c$ pour un certain $c \geq 1$. Alors
$\mathcal{H} = \SheafHom_{\mathcal{O}_X}(\mathcal{G}, \mathcal{F})$
est un $\mathcal{O}_X$-module cohérent annulé par $\mathcal{I}^c$.
Nous voyons donc que
$$
\Hom_X(\mathcal{G}, \mathcal{F}) =
H^0(X, \mathcal{H}) =
\lim H^0(X, \mathcal{H}/\mathcal{I}^n\mathcal{H}) =
\Mor_{\textit{Coh}(X, \mathcal{I})}
(\mathcal{G}^\wedge, \mathcal{F}^\wedge).
$$
Cela résulte du lemme \ref{lemma-completion-internal-hom}.
Cela prouve l'énoncé concernant les homomorphismes.

\medskip\noindent
La notation $\Ext$ désigne les extensions définies dans
Homologie, section \ref{homology-section-extensions}.
L'injectivité de l'application sur les $\Ext$ découle immédiatement de
la bijectivité de l'application sur les $\Hom$. Pour
la surjectivité, supposons que $\mathcal{F}$ soit annulé
par une puissance de $I$. Alors la partie (1) du lemme \ref{lemma-existence-easy}
montre que, étant donnée une extension
$$
0 \to \mathcal{G}^\wedge \to (\mathcal{E}_n) \to \mathcal{F}^\wedge \to 0
$$
dans $\textit{Coh}(U, I\mathcal{O}_U)$,
le morphisme $\mathcal{G}^\wedge \to (\mathcal{E}_n)$ est
isomorphe à $\mathcal{G} \to \mathcal{E}^\wedge$
pour un certain $\mathcal{G} \to \mathcal{E}$ dans $\textit{Coh}(\mathcal{O}_U)$.
On procède de même dans l'autre cas.
\end{proof}

\begin{lemma}
\label{lemma-finite-over-rees-algebra}
Soit $X$ un schéma noethérien et soit $\mathcal{I} \subset \mathcal{O}_X$
un faisceau quasi-cohérent d'idéaux. Si $(\mathcal{F}_n)$ est un objet de
$\textit{Coh}(X, \mathcal{I})$, alors
$\bigoplus \Ker(\mathcal{F}_{n + 1} \to \mathcal{F}_n)$ est
un
$\bigoplus \mathcal{I}^n/\mathcal{I}^{n + 1}$-module gradué quasi-cohérent de type fini.
\end{lemma}

\begin{proof}
La question est locale sur $X$ ; nous pouvons donc supposer $X$ affine, c'est-à-dire que
nous sommes dans la situation du lemme \ref{lemma-inverse-systems-affine}.
Dans ce cas, si $(\mathcal{F}_n)$ correspond au $A^\wedge$
-module de type fini $M$, alors $\bigoplus \Ker(\mathcal{F}_{n + 1} \to \mathcal{F}_n)$
correspond à $\bigoplus I^nM/I^{n + 1}M$, qui est manifestement un module de
type fini sur $\bigoplus I^n/I^{n + 1}$.
\end{proof}

\begin{lemma}
\label{lemma-inverse-systems-pullback}
Soit $f : X \to Y$ un morphisme de schémas noethériens.
Soit $\mathcal{J} \subset \mathcal{O}_Y$ un faisceau quasi-cohérent
d'idéaux et posons $\mathcal{I} = f^{-1}\mathcal{J} \mathcal{O}_X$.
Il existe alors un foncteur exact à droite
$$
f^* : \textit{Coh}(Y, \mathcal{J}) \longrightarrow \textit{Coh}(X, \mathcal{I})
$$
qui envoie $(\mathcal{G}_n)$ sur $(f^*\mathcal{G}_n)$. Si $f$ est plat,
alors $f^*$ est un foncteur exact.
\end{lemma}

\begin{proof}
Puisque $f^* : \textit{Coh}(\mathcal{O}_Y) \to \textit{Coh}(\mathcal{O}_X)$
est exact à droite, nous avons
$$
f^*\mathcal{G}_n =
f^*(\mathcal{G}_{n + 1}/\mathcal{I}^n\mathcal{G}_{n + 1}) =
f^*\mathcal{G}_{n + 1}/f^{-1}\mathcal{I}^nf^*\mathcal{G}_{n + 1} =
f^*\mathcal{G}_{n + 1}/\mathcal{J}^nf^*\mathcal{G}_{n + 1}
$$
donc l'image inverse d'un système est un système. La construction des
conoyaux dans la démonstration du lemme \ref{lemma-inverse-systems-abelian}
montre que
$f^* : \textit{Coh}(Y, \mathcal{J}) \to \textit{Coh}(X, \mathcal{I})$
est toujours exact à droite. Si $f$ est plat, alors
$f^* : \textit{Coh}(\mathcal{O}_Y) \to \textit{Coh}(\mathcal{O}_X)$
est un foncteur exact. Il résulte de la construction des noyaux
dans la démonstration du lemme \ref{lemma-inverse-systems-abelian}
que, dans ce cas,
$f^* : \textit{Coh}(Y, \mathcal{J}) \to \textit{Coh}(X, \mathcal{I})$
envoie également les noyaux sur les noyaux.
\end{proof}

\begin{lemma}
\label{lemma-inverse-systems-pullback-equivalence}
Soit $f : X' \to X$ un morphisme de schémas noethériens. Soit $Z \subset X$
un sous-schéma fermé et notons $Z' = f^{-1}Z$ l'image inverse
schématique. Soient $\mathcal{I} \subset \mathcal{O}_X$,
$\mathcal{I}' \subset \mathcal{O}_{X'}$ les faisceaux
quasi-cohérents d'idéaux correspondants.
Si $f$ est plat et que le morphisme induit $Z' \to Z$
est un isomorphisme, alors le foncteur d'image inverse
$f^* : \textit{Coh}(X, \mathcal{I}) \to \textit{Coh}(X', \mathcal{I}')$
(lemme \ref{lemma-inverse-systems-pullback})
est une équivalence.
\end{lemma}

\begin{proof}
Si $X$ et $X'$ sont affines, cela découle immédiatement de
Compléments d'algèbre, lemme \ref{more-algebra-lemma-neighbourhood-equivalence}.
Pour le démontrer en général, soient $Z_n \subset X$, $Z'_n \subset X'$
les $n$-ièmes voisinages infinitésimaux de $Z$, $Z'$.
Le morphisme $Z_n \to Z'_n$ induit un homéomorphisme entre
les espaces topologiques sous-jacents. D'autre part, si $z' \in Z'$
s'envoie sur $z \in Z$, alors le morphisme d'anneaux
$\mathcal{O}_{X, z} \to \mathcal{O}_{X', z'}$ est plat
et induit un isomorphisme
$\mathcal{O}_{X, z}/\mathcal{I}_z \to \mathcal{O}_{X', z'}/\mathcal{I}'_{z'}$.
Il induit donc un isomorphisme
$\mathcal{O}_{X, z}/\mathcal{I}_z^n \to
\mathcal{O}_{X', z'}/(\mathcal{I}'_{z'})^n$
pour tout $n \geq 1$, par exemple
d'après Compléments d'algèbre, lemme \ref{more-algebra-lemma-neighbourhood-isomorphism}.
Ainsi $Z'_n \to Z_n$ est un isomorphisme de schémas.
Ainsi $f^*$ induit une équivalence entre
la catégorie des $\mathcal{O}_X$-modules cohérents annulés par $\mathcal{I}^n$
et la
catégorie des $\mathcal{O}_{X'}$-modules cohérents annulés par
$(\mathcal{I}')^n$, voir le
lemme \ref{lemma-i-star-equivalence}.
Cela entraîne clairement le lemme.
\end{proof}

\begin{lemma}
\label{lemma-inverse-systems-ideals-equivalence}
Soit $X$ un schéma noethérien. Soient
$\mathcal{I}, \mathcal{J} \subset \mathcal{O}_X$
des faisceaux quasi-cohérents d'idéaux.
Si $V(\mathcal{I}) = V(\mathcal{J})$ comme sous-ensembles fermés
de $X$, alors $\textit{Coh}(X, \mathcal{I})$ et $\textit{Coh}(X, \mathcal{J})$
sont équivalentes.
\end{lemma}

\begin{proof}
Supposons d'abord que $X = \Spec(A)$ soit affine. Soient $I, J \subset A$ les idéaux
correspondant à $\mathcal{I}, \mathcal{J}$. Alors $V(I) = V(J)$
implique que $I^c \subset J$ et $J^d \subset I$ pour certains $c, d \geq 1$
par les propriétés élémentaires de la topologie de
Zariski (voir Algèbre, section \ref{algebra-section-spectrum-ring} et
lemme \ref{algebra-lemma-Noetherian-power}).
Les complétions $I$-adique et $J$-adique de $A$ coïncident donc,
voir Algèbre, lemme \ref{algebra-lemma-change-ideal-completion}.
L'équivalence découle ainsi du lemme \ref{lemma-inverse-systems-affine}
dans ce cas.

\medskip\noindent
En général, d'après ce qui précède et puisque
$X$ est quasi-compact, on peut choisir $c, d \geq 1$ tels que
$\mathcal{I}^c \subset \mathcal{J}$ et $\mathcal{J}^d \subset \mathcal{I}$.
Alors, étant donné un objet $(\mathcal{F}_n)$ de
$\textit{Coh}(X, \mathcal{I})$, nous affirmons que le
système projectif
$$
(\mathcal{F}_{cn}/\mathcal{J}^n\mathcal{F}_{cn})
$$
appartient à $\textit{Coh}(X, \mathcal{J})$. Cela peut se vérifier sur
les éléments d'un recouvrement affine ; nous omettons les détails.
De la même manière, nous pouvons construire un objet de
$\textit{Coh}(X, \mathcal{I})$ à partir d'un objet de
$\textit{Coh}(X, \mathcal{J})$. Nous omettons la vérification que
ces constructions définissent des foncteurs quasi-inverses l'un de l'autre.
\end{proof}





\section{Théorème d'existence de Grothendieck, I}
\label{section-existence}

\noindent
Dans cette section, nous étudions le théorème d'existence de Grothendieck dans
le cas projectif. Nous utiliserons la notion de modules formels cohérents développée
à la section \ref{section-coherent-formal}. Le lecteur familier des
schémas formels est invité à lire l'énoncé et la démonstration
du théorème dans \cite{EGA}.

\begin{lemma}
\label{lemma-fully-faithful}
Soit $A$ un anneau noethérien complet pour la topologie $I$-adique.
Soit $f : X \to \Spec(A)$ un morphisme propre. Posons
$\mathcal{I} = I\mathcal{O}_X$.
Alors le foncteur (\ref{equation-completion-functor}) est pleinement fidèle.
\end{lemma}

\begin{proof}
Soient $\mathcal{F}$ et $\mathcal{G}$ des $\mathcal{O}_X$-modules cohérents.
Alors $\mathcal{H} = \SheafHom_{\mathcal{O}_X}(\mathcal{G}, \mathcal{F})$
est un $\mathcal{O}_X$-module cohérent,
voir Modules, lemme \ref{modules-lemma-internal-hom-locally-kernel-direct-sum}.
D'après le lemme \ref{lemma-completion-internal-hom}, l'application
$$
\lim_n H^0(X, \mathcal{H}/\mathcal{I}^n\mathcal{H})
\to
\Mor_{\textit{Coh}(X, \mathcal{I})}
(\mathcal{G}^\wedge, \mathcal{F}^\wedge)
$$
est bijective. Par conséquent, la pleine fidélité de
(\ref{equation-completion-functor}) résulte du théorème des fonctions
formelles (lemme \ref{lemma-spell-out-theorem-formal-functions})
appliqué au faisceau cohérent $\mathcal{H}$.
\end{proof}

\begin{lemma}
\label{lemma-vanishing-projective}
Soit $A$ un anneau noethérien et soit $I \subset A$ un idéal.
Soit $f : X \to \Spec(A)$ un morphisme propre et soit
$\mathcal{L}$ un faisceau inversible $f$-ample. Posons
$\mathcal{I} = I\mathcal{O}_X$. Soit $(\mathcal{F}_n)$ un objet
de $\textit{Coh}(X, \mathcal{I})$. Il existe alors un
entier $d_0$ tel que
$$
H^1(X, \Ker(\mathcal{F}_{n + 1} \to \mathcal{F}_n)
\otimes \mathcal{L}^{\otimes d} )
= 0
$$
pour tout $n \geq 0$ et tout $d \geq d_0$.
\end{lemma}

\begin{proof}
Posons $B = \bigoplus I^n/I^{n + 1}$ et
$\mathcal{B} = \bigoplus \mathcal{I}^n/\mathcal{I}^{n + 1} = f^*\widetilde{B}$.
D'après le lemme \ref{lemma-finite-over-rees-algebra}, le
$\mathcal{B}$-module gradué quasi-cohérent
$\mathcal{G} = \bigoplus \Ker(\mathcal{F}_{n + 1} \to \mathcal{F}_n)$
est de type fini. Le lemme découle donc de la partie (2) du
lemme \ref{lemma-graded-finiteness}.
\end{proof}

\begin{lemma}
\label{lemma-existence-projective}
Soit $A$ un anneau noethérien complet pour la topologie $I$-adique.
Soit $f : X \to \Spec(A)$ un morphisme projectif. Posons
$\mathcal{I} = I\mathcal{O}_X$.
Alors le foncteur (\ref{equation-completion-functor}) est une équivalence.
\end{lemma}

\begin{proof}
Nous avons déjà vu que (\ref{equation-completion-functor}) est
pleinement fidèle (lemme \ref{lemma-fully-faithful}). Il suffit donc
de montrer que ce foncteur est essentiellement surjectif.

\medskip\noindent
Montrons d'abord que tout objet $(\mathcal{F}_n)$ de
$\textit{Coh}(X, \mathcal{I})$ est le quotient d'un objet
appartenant à l'image de (\ref{equation-completion-functor}).
Soit $\mathcal{L}$ un faisceau inversible $f$-ample sur $X$.
Choisissons $d_0$ comme dans le lemme \ref{lemma-vanishing-projective}.
Choisissons $d \geq d_0$ tel que
$\mathcal{F}_1 \otimes \mathcal{L}^{\otimes d}$
soit globalement engendré par des sections $s_{1, 1}, \ldots, s_{t, 1}$.
Puisque les applications de transition du système
$$
H^0(X, \mathcal{F}_{n + 1} \otimes \mathcal{L}^{\otimes d})
\longrightarrow
H^0(X, \mathcal{F}_n \otimes \mathcal{L}^{\otimes d})
$$
sont surjectives par l'annulation de $H^1$, nous pouvons relever
$s_{1, 1}, \ldots, s_{t, 1}$ en un système compatible de sections globales
$s_{1, n}, \ldots, s_{t, n}$ de
$\mathcal{F}_n \otimes \mathcal{L}^{\otimes d}$.
Celles-ci déterminent un système compatible d'applications
$$
(s_{1, n}, \ldots, s_{t, n}) :
(\mathcal{L}^{\otimes -d})^{\oplus t} \longrightarrow \mathcal{F}_n
$$
En utilisant le lemme \ref{lemma-inverse-systems-surjective},
nous en déduisons une application surjective
$$
\left((\mathcal{L}^{\otimes -d})^{\oplus t}\right)^\wedge
\longrightarrow
(\mathcal{F}_n)
$$
comme voulu.

\medskip\noindent
Le résultat du paragraphe précédent et le fait que
$\textit{Coh}(X, \mathcal{I})$ est abélienne
(lemme \ref{lemma-inverse-systems-abelian})
impliquent que
tout objet de $\textit{Coh}(X, \mathcal{I})$ est le conoyau
d'une application entre des objets provenant de $\textit{Coh}(\mathcal{O}_X)$.
Comme (\ref{equation-completion-functor}) est pleinement fidèle et exact par les
lemmes \ref{lemma-fully-faithful} et \ref{lemma-exact},
nous concluons.
\end{proof}








\section{Théorème d'existence de Grothendieck, II}
\label{section-existence-proper}

\noindent
Dans cette section, nous étudions le théorème d'existence de Grothendieck dans le
cas propre. Avant d'en donner l'énoncé et la démonstration, nous devons développer quelque
peu la théorie des catégories $\textit{Coh}(X, \mathcal{I})$
de modules formels cohérents introduites
à la section \ref{section-coherent-formal}.

\begin{remark}
\label{remark-inverse-systems-kernel-cokernel-annihilated-by}
Soit $X$ un schéma noethérien et soient
$\mathcal{I}, \mathcal{K} \subset \mathcal{O}_X$
des faisceaux quasi-cohérents d'idéaux. Soit
$\alpha : (\mathcal{F}_n) \to (\mathcal{G}_n)$ un morphisme de
$\textit{Coh}(X, \mathcal{I})$.
Étant donné un ouvert affine $\Spec(A) = U \subset X$ tel que
$\mathcal{I}|_U, \mathcal{K}|_U$ correspondent aux idéaux $I, K \subset A$,
notons $\alpha_U : M \to N$ le morphisme de $A^\wedge$-modules finis
qui correspond à $\alpha|_U$ via le lemme \ref{lemma-inverse-systems-affine}.
Nous affirmons que les conditions suivantes sont équivalentes :
\begin{enumerate}
\item il existe un entier $t \geq 1$ tel que
$\Ker(\alpha_n)$ et $\Coker(\alpha_n)$
soient annulés par $\mathcal{K}^t$ pour tout $n \geq 1$,
\item pour tout ouvert affine $\Spec(A) = U \subset X$ comme ci-dessus,
les modules $\Ker(\alpha_U)$ et $\Coker(\alpha_U)$
sont annulés par $K^t$ pour un certain entier $t \geq 1$, et
\item il existe un recouvrement affine fini $X = \bigcup U_i$
tel que la conclusion de (2) soit vérifiée pour $\alpha_{U_i}$.
\end{enumerate}
Si ces conditions équivalentes sont vérifiées, nous dirons que
$\alpha$ est une
{\it application dont le noyau et le conoyau sont annulés par une puissance de
$\mathcal{K}$}.
Pour voir l'équivalence, utilisons le fait d'algèbre commutative suivant :
supposons donnée une suite exacte
$$
0 \to T \to M \to N \to Q \to 0
$$
de $A$-modules, où $T$ et $Q$ sont annulés par $K^t$ pour un certain
idéal $K \subset A$. Alors, pour tous $f, g \in K^t$, il existe une
application canonique $"fg": N \to M$ telle que $M \to N \to M$ soit la
multiplication par $fg$. En effet, pour $y \in N$, nous pouvons choisir $x \in M$
ayant pour image $fy$ dans $N$, puis poser $"fg"(y) = gx$. Il est alors
clair que $\Ker(M/JM \to N/JN)$ et $\Coker(M/JM \to N/JN)$
sont annulés par $K^{2t}$ pour tout idéal $J \subset A$.

\medskip\noindent
En appliquant ce fait d'algèbre commutative à $\alpha_{U_i}$ et $J = I^n$,
nous voyons que (3) implique (1). Réciproquement,
supposons (1) vérifiée et $M \to N$ égal à $\alpha_U$. Il existe
alors $t \geq 1$ tel que
$\Ker(M/I^nM \to N/I^nN)$ et $\Coker(M/I^nM \to N/I^nN)$
soient annulés par $K^t$ pour tout $n$. Nous obtenons des applications
$"fg" : N/I^nN \to M/I^nM$ qui induisent, par passage à la limite, une application $N \to M$,
car $N$ et $M$ sont complets pour la topologie $I$-adique. Comme le composé
$N \to M \to N$ est la multiplication par $fg$, nous concluons que $fg$
annule $T$ et $Q$. Autrement dit, $T$ et $Q$ sont annulés par
$K^{2t}$, comme voulu.
\end{remark}

\begin{lemma}
\label{lemma-existence-tricky}
Soit $X$ un schéma noethérien. Soient
$\mathcal{I}, \mathcal{K} \subset \mathcal{O}_X$
des faisceaux quasi-cohérents d'idéaux.
Soit $X_e \subset X$ le sous-schéma fermé défini par $\mathcal{K}^e$.
Posons $\mathcal{I}_e = \mathcal{I}\mathcal{O}_{X_e}$.
Soit $(\mathcal{F}_n)$ un objet de $\textit{Coh}(X, \mathcal{I})$. Supposons
que
\begin{enumerate}
\item le foncteur
$\textit{Coh}(\mathcal{O}_{X_e}) \to \textit{Coh}(X_e, \mathcal{I}_e)$
est une équivalence pour tout $e \geq 1$, et
\item il existe un faisceau cohérent $\mathcal{H}$ sur $X$ et une application
$\alpha : (\mathcal{F}_n) \to \mathcal{H}^\wedge$ dont
le noyau et le conoyau sont annulés par une puissance de $\mathcal{K}$.
\end{enumerate}
Alors $(\mathcal{F}_n)$ appartient à l'image essentielle de
(\ref{equation-completion-functor}).
\end{lemma}

\begin{proof}
Dans toute cette démonstration, nous utiliserons sans autre mention le fait que, pour une immersion
fermée $i : Z \to X$, le foncteur $i_*$ donne une équivalence entre la
catégorie des modules cohérents sur $Z$ et celle des modules cohérents sur $X$ annulés
par le faisceau d'idéaux de $Z$, voir le lemme \ref{lemma-i-star-equivalence}.
En particulier, nous pouvons identifier $\textit{Coh}(\mathcal{O}_{X_e})$
à la catégorie des $\mathcal{O}_X$-modules cohérents annulés par
$\mathcal{K}^e$, et $\textit{Coh}(X_e, \mathcal{I}_e)$ à la sous-catégorie pleine
de $\textit{Coh}(X, \mathcal{I})$ formée des objets annulés par $\mathcal{K}^e$.
De plus, (1) nous dit que le foncteur de
complétion (\ref{equation-completion-functor}) induit une équivalence entre ces catégories.

\medskip\noindent
En appliquant cette équivalence, nous obtenons un $\mathcal{O}_X$-module cohérent
$\mathcal{G}_e$, annulé par $\mathcal{K}^e$, qui correspond au système
$(\mathcal{F}_n/\mathcal{K}^e\mathcal{F}_n)$ de
$\textit{Coh}(X, \mathcal{I})$. Les applications
$\mathcal{F}_n/\mathcal{K}^{e + 1}\mathcal{F}_n \to
\mathcal{F}_n/\mathcal{K}^e\mathcal{F}_n$ correspondent à des applications canoniques
$\mathcal{G}_{e + 1} \to \mathcal{G}_e$ qui induisent des isomorphismes
$\mathcal{G}_{e + 1}/\mathcal{K}^e\mathcal{G}_{e + 1} \to \mathcal{G}_e$.
Ainsi, $(\mathcal{G}_e)$ est un objet de $\textit{Coh}(X, \mathcal{K})$.
L'application $\alpha$ induit un système d'applications
$$
\mathcal{F}_n/\mathcal{K}^e\mathcal{F}_n
\longrightarrow
\mathcal{H}/(\mathcal{I}^n + \mathcal{K}^e)\mathcal{H}
$$
et donc des applications $\mathcal{G}_e \to \mathcal{H}/\mathcal{K}^e\mathcal{H}$
(en utilisant à nouveau l'équivalence de catégories).
Soit $t \geq 1$ un entier, qui existe par l'hypothèse (2),
tel que $\mathcal{K}^t$ annule le noyau et le conoyau de toutes les applications
$\mathcal{F}_n \to \mathcal{H}/\mathcal{I}^n\mathcal{H}$.
Alors $\mathcal{K}^{2t}$ annule le noyau et le conoyau des applications
$\mathcal{F}_n/\mathcal{K}^e\mathcal{F}_n \to
\mathcal{H}/(\mathcal{I}^n + \mathcal{K}^e)\mathcal{H}$, voir la
remarque \ref{remark-inverse-systems-kernel-cokernel-annihilated-by}.
Nous en concluons que $\mathcal{K}^{4t}$ annule le noyau et
le conoyau des applications
$$
\mathcal{G}_e
\longrightarrow
\mathcal{H}/\mathcal{K}^e\mathcal{H},
$$
voir la remarque \ref{remark-inverse-systems-kernel-cokernel-annihilated-by}.
Nous appliquons le lemme \ref{lemma-existence-easy} pour obtenir un
$\mathcal{O}_X$-module cohérent $\mathcal{F}$, une application
$a : \mathcal{F} \to \mathcal{H}$ et un isomorphisme
$\beta : (\mathcal{G}_e) \to (\mathcal{F}/\mathcal{K}^e\mathcal{F})$
dans $\textit{Coh}(X, \mathcal{K})$. En remontant les constructions, pour $n$
donné, le triplet
$(\mathcal{F}/\mathcal{I}^n\mathcal{F}, a \bmod \mathcal{I}^n, \beta
\bmod \mathcal{I}^n)$ est un triplet comme dans le lemme pour le morphisme
$\alpha_n \bmod \mathcal{K}^e :
(\mathcal{F}_n/\mathcal{K}^e\mathcal{F}_n) \to
(\mathcal{H}/(\mathcal{I}^n + \mathcal{K}^e)\mathcal{H})$
de $\textit{Coh}(X, \mathcal{K})$. L'unicité dans le
lemme \ref{lemma-existence-easy}
donne donc un isomorphisme canonique
$\mathcal{F}/\mathcal{I}^n\mathcal{F} \to \mathcal{F}_n$
compatible avec tous les morphismes considérés. Ceci achève la
démonstration du lemme.
\end{proof}

\begin{lemma}
\label{lemma-inverse-systems-push-pull}
Soit $Y$ un schéma noethérien. Soient
$\mathcal{J}, \mathcal{K} \subset \mathcal{O}_Y$
des faisceaux quasi-cohérents d'idéaux.
Soit $f : X \to Y$ un morphisme propre qui est un isomorphisme au-dessus
de $V = Y \setminus V(\mathcal{K})$.
Posons $\mathcal{I} = f^{-1}\mathcal{J} \mathcal{O}_X$.
Soit $(\mathcal{G}_n)$ un objet de $\textit{Coh}(Y, \mathcal{J})$,
soit $\mathcal{F}$ un $\mathcal{O}_X$-module cohérent, et soit
$\beta : (f^*\mathcal{G}_n)  \to \mathcal{F}^\wedge$ un isomorphisme dans
$\textit{Coh}(X, \mathcal{I})$. Il existe alors une application
$$
\alpha :
(\mathcal{G}_n)
\longrightarrow
(f_*\mathcal{F})^\wedge
$$
dans $\textit{Coh}(Y, \mathcal{J})$ dont le noyau et le
conoyau sont annulés par une puissance de $\mathcal{K}$.
\end{lemma}

\begin{proof}
Puisque $f$ est un morphisme propre, $f_*\mathcal{F}$
est un $\mathcal{O}_Y$-module cohérent
(proposition \ref{proposition-proper-pushforward-coherent}).
L'énoncé du lemme a donc bien un sens.
Considérons les composés
$$
\gamma_n : \mathcal{G}_n \to
f_*f^*\mathcal{G}_n \to
f_*(\mathcal{F}/\mathcal{I}^n\mathcal{F}).
$$
Ici, la première application est l'application d'adjonction, et la seconde est $f_*\beta_n$.
Nous affirmons qu'il existe une unique $\alpha$ comme dans le lemme
telle que les composés
$$
\mathcal{G}_n \xrightarrow{\alpha_n}
f_*\mathcal{F}/\mathcal{J}^nf_*\mathcal{F} \to
f_*(\mathcal{F}/\mathcal{I}^n\mathcal{F})
$$
soient égaux à $\gamma_n$ pour tout $n$. En raison de l'unicité, nous pouvons supposer
que $Y = \Spec(B)$ est affine. Soit $J \subset B$ l'idéal correspondant
à $\mathcal{J}$. Posons
$$
M_n = H^0(X, \mathcal{F}/\mathcal{I}^n\mathcal{F})
\quad\text{et}\quad
M = H^0(X, \mathcal{F})
$$
D'après le lemme \ref{lemma-ML-cohomology-powers-ideal} et le
Théorème \ref{theorem-formal-functions},
la limite inverse des modules
$M_n$ est égale à la complétion $M^\wedge = \lim M/J^nM$.
Posons $N_n = H^0(Y, \mathcal{G}_n)$ et $N = \lim N_n$.
Par l'équivalence de catégories du
lemme \ref{lemma-inverse-systems-affine},
les $B^\wedge$-modules finis $N$ et $M^\wedge$ correspondent
à $(\mathcal{G}_n)$ et $f_*\mathcal{F}^\wedge$.
Il s'ensuit que $\alpha$ doit être le morphisme de
$\textit{Coh}(Y, \mathcal{J})$ correspondant à l'homomorphisme
$$
\lim \gamma_n : N = \lim_n N_n \longrightarrow \lim M_n = M^\wedge
$$
de $B^\wedge$-modules finis.

\medskip\noindent
Il reste à montrer que le noyau et le conoyau de $\alpha$ sont
annulés par une puissance de $\mathcal{K}$. Posons $Y' = \Spec(B^\wedge)$
et $X' = Y' \times_Y X$. Soient $\mathcal{K}'$, $\mathcal{J}'$, $\mathcal{G}'_n$
et $\mathcal{I}'$, $\mathcal{F}'$ les images inverses de
$\mathcal{K}$, $\mathcal{J}$, $\mathcal{G}_n$ et
$\mathcal{I}$, $\mathcal{F}$, sur $Y'$ et $X'$.
Le morphisme de projection $f' : X' \to Y'$ est le changement de base de
$f$ par $Y' \to Y$. Notons que $Y' \to Y$ est un morphisme plat de schémas,
car $B \to B^\wedge$ est plat d'après
Algèbre, lemme \ref{algebra-lemma-completion-flat}.
Ainsi, $f'_*\mathcal{F}'$, resp.\ $f'_*(f')^*\mathcal{G}_n'$,
est l'image inverse de $f_*\mathcal{F}$, resp.\ $f_*f^*\mathcal{G}_n$,
sur $Y'$ par le lemme \ref{lemma-flat-base-change-cohomology}.
L'unicité de notre construction montre que l'image inverse de $\alpha$ sur $Y'$
est l'application correspondante $\alpha'$ construite dans la situation sur $Y'$.
De plus, pour vérifier que le noyau et le conoyau de $\alpha$ sont
annulés par $\mathcal{K}^t$, il suffit de vérifier que le
noyau et le conoyau de $\alpha'$ sont annulés par
$(\mathcal{K}')^t$. En effet, il suffit pour cela de le vérifier pour
les noyaux et conoyaux des applications $\alpha_n$ et $\alpha'_n$
(voir la remarque \ref{remark-inverse-systems-kernel-cokernel-annihilated-by}),
et l'homomorphisme d'anneaux $B \to B^\wedge$ induit
une équivalence de catégories entre les modules annulés par
$J^n$ et ceux annulés par $(J')^n$, voir
Compléments sur l'algèbre, lemme \ref{more-algebra-lemma-neighbourhood-equivalence}.
Nous pouvons donc supposer que $B$ est complet pour la topologie $J$-adique.

\medskip\noindent
Supposons que $Y = \Spec(B)$ soit affine, que $\mathcal{J}$ corresponde à l'idéal
$J \subset B$, et que $B$ soit complet pour la topologie $J$-adique.
Dans ce cas, $(\mathcal{G}_n)$ appartient à l'image essentielle du foncteur
$\textit{Coh}(\mathcal{O}_Y) \to \textit{Coh}(Y, \mathcal{J})$.
Soit $\mathcal{G}$ un $\mathcal{O}_Y$-module cohérent tel que
$(\mathcal{G}_n) = \mathcal{G}^\wedge$. Notons que
$f^*(\mathcal{G}^\wedge) = (f^*\mathcal{G})^\wedge$. Le
lemme \ref{lemma-fully-faithful}
nous dit donc que $\beta$ provient d'un isomorphisme
$b : f^*\mathcal{G} \to \mathcal{F}$,
et que $\alpha$ s'obtient en appliquant le foncteur de complétion à
$$
\mathcal{G} \to f_*f^*\mathcal{G} \cong f_*\mathcal{F}
$$
Il s'agit donc de vérifier que le noyau et le conoyau de
l'application d'adjonction $c : \mathcal{G} \to f_*f^*\mathcal{G}$ sont annulés par
une puissance de $\mathcal{K}$. Or, puisque la restriction
$f|_{f^{-1}(V)} : f^{-1}(V) \to V$ est un isomorphisme,
nous voyons que $c|_V$ est un isomorphisme. Ainsi, les faisceaux cohérents
$\Ker(c)$ et $\Coker(c)$ sont supportés sur $V(\mathcal{K})$,
et sont donc annulés par une puissance de $\mathcal{K}$
(lemme \ref{lemma-power-ideal-kills-sheaf}), comme voulu.
\end{proof}

\noindent
La proposition suivante est la forme du théorème d'existence de
Grothendieck qui est la plus souvent utilisée en pratique.

\begin{proposition}
\label{proposition-existence-proper}
Soit $A$ un anneau noethérien complet pour la topologie $I$-adique.
Soit $f : X \to \Spec(A)$ un morphisme propre de schémas.
Posons $\mathcal{I} = I\mathcal{O}_X$.
Alors le foncteur (\ref{equation-completion-functor}) est une équivalence.
\end{proposition}

\begin{proof}
Nous avons déjà vu que (\ref{equation-completion-functor}) est
pleinement fidèle (lemme \ref{lemma-fully-faithful}). Il suffit donc
de montrer que ce foncteur est essentiellement surjectif.

\medskip\noindent
Considérons l'ensemble $\Xi$ des faisceaux quasi-cohérents d'idéaux
$\mathcal{K} \subset \mathcal{O}_X$ tels que tout objet
$(\mathcal{F}_n)$ annulé par $\mathcal{K}$ appartienne à l'image essentielle.
Nous voulons montrer que $(0)$ appartient à $\Xi$. Sinon, puisque $X$ est noethérien,
il existe un faisceau quasi-cohérent d'idéaux maximal $\mathcal{K}$
n'appartenant pas à $\Xi$, voir le
lemme \ref{lemma-acc-coherent}.
Après avoir remplacé $X$ par le sous-schéma fermé de $X$
correspondant à $\mathcal{K}$, nous pouvons supposer que tout faisceau quasi-cohérent d'idéaux
$\mathcal{K}$ non nul appartient à $\Xi$. (On utilise ici la correspondance entre
les modules cohérents annulés par $\mathcal{K}$ et les modules
cohérents sur le sous-schéma fermé correspondant à $\mathcal{K}$, voir le
lemme \ref{lemma-i-star-equivalence}.)
Soit $(\mathcal{F}_n)$ un objet de
$\textit{Coh}(X, \mathcal{I})$.
Nous allons montrer que cet objet appartient à l'image essentielle du
foncteur (\ref{equation-completion-functor}), achevant ainsi la
démonstration de la proposition.

\medskip\noindent
Appliquons le lemme de Chow (lemme \ref{lemma-chow-Noetherian}) pour obtenir un
morphisme propre surjectif $f : X' \to X$ qui soit un isomorphisme
au-dessus d'un ouvert dense $U \subset X$ et tel que $X'$ soit projectif sur $A$.
Soit $\mathcal{K}$ le faisceau quasi-cohérent d'idéaux définissant
le complémentaire réduit $X \setminus U$. Par le cas
projectif du théorème d'existence de Grothendieck
(lemme \ref{lemma-existence-projective}),
il existe un module cohérent $\mathcal{F}'$ sur $X'$ tel
que $(\mathcal{F}')^\wedge \cong (f^*\mathcal{F}_n)$. D'après la
proposition \ref{proposition-proper-pushforward-coherent},
le $\mathcal{O}_X$-module $\mathcal{H} = f_*\mathcal{F}'$ est cohérent
et, d'après le lemme \ref{lemma-inverse-systems-push-pull},
il existe un morphisme $(\mathcal{F}_n) \to \mathcal{H}^\wedge$
de $\textit{Coh}(X, \mathcal{I})$ dont le noyau et le conoyau
sont annulés par une puissance de $\mathcal{K}$. Toutes les puissances $\mathcal{K}^e$
appartiennent à $\Xi$, si bien que (\ref{equation-completion-functor})
est une équivalence pour les sous-schémas fermés $X_e = V(\mathcal{K}^e)$.
Nous concluons par le lemme \ref{lemma-existence-tricky}.
\end{proof}





\section{Propreté sur une base}
\label{section-proper-over-base}

\noindent
Il s'agit simplement d'une courte section destinée à signaler quelques propriétés utiles des sous-ensembles
fermés propres sur une base et des modules quasi-cohérents de type fini
dont le support est propre sur une base.

\begin{lemma}
\label{lemma-closed-proper-over-base}
Soit $f : X \to S$ un morphisme de schémas localement de type fini.
Soit $Z \subset X$ un sous-ensemble fermé. Les conditions suivantes sont équivalentes :
\begin{enumerate}
\item le morphisme $Z \to S$ est propre lorsque $Z$ est muni de la structure
de sous-schéma fermé réduit induit
(Schémas, définition \ref{schemes-definition-reduced-induced-scheme}),
\item pour une certaine structure de sous-schéma fermé sur $Z$, le morphisme $Z \to S$
est propre,
\item pour toute structure de sous-schéma fermé sur $Z$, le morphisme
$Z \to S$ est propre.
\end{enumerate}
\end{lemma}

\begin{proof}
Les implications (3) $\Rightarrow$ (1) et (1) $\Rightarrow$ (2)
sont immédiates. Il suffit donc de prouver que (2) implique (3).
Nous invitons le lecteur à trouver sa propre démonstration de ce fait.
Soient $Z'$ et $Z''$ des structures de sous-schémas fermés sur $Z$
telles que $Z' \to S$ soit propre. Il faut montrer que $Z'' \to S$ est propre.
Soit $Z''' = Z' \cup Z''$ la réunion schématique, voir
Morphismes, Définition
\ref{morphisms-definition-scheme-theoretic-intersection-union}.
Alors $Z'''$ est une autre structure de sous-schéma fermé sur $Z$.
Cela résulte par exemple de la description des réunions schématiques dans
Morphismes, lemme \ref{morphisms-lemma-scheme-theoretic-union}.
Puisque $Z'' \to Z'''$ est une immersion fermée, il suffit de prouver
que $Z''' \to S$ est propre (voir
Morphismes, lemmes \ref{morphisms-lemma-closed-immersion-proper} et
\ref{morphisms-lemma-composition-proper}).
Le morphisme $Z' \to Z'''$ est une immersion fermée bijective,
donc en particulier surjective et universellement
fermée. Le fait que $Z' \to S$ soit séparé implique alors que
$Z''' \to S$ est séparé, voir
Morphismes, lemme \ref{morphisms-lemma-image-universally-closed-separated}. De
plus, $Z''' \to S$ est localement de type fini
puisque $X \to S$ est localement de type fini
(Morphismes, lemmes \ref{morphisms-lemma-immersion-locally-finite-type} et
\ref{morphisms-lemma-composition-finite-type}).
Puisque $Z' \to S$ est quasi-compact et $Z' \to Z'''$ est un homéomorphisme,
nous voyons que $Z''' \to S$ est quasi-compact.
Enfin, puisque $Z' \to S$ est universellement fermé, nous voyons qu'il
en va de même pour $Z''' \to S$ d'après
Morphismes, lemme \ref{morphisms-lemma-image-proper-is-proper}.
Cela achève la démonstration.
\end{proof}

\begin{definition}
\label{definition-proper-over-base}
Soit $f : X \to S$ un morphisme de schémas localement de type fini.
Soit $Z \subset X$ un sous-ensemble fermé. Nous
disons que {\it $Z$ est propre sur $S$}
si les conditions équivalentes du lemme \ref{lemma-closed-proper-over-base}
sont satisfaites.
\end{definition}

\noindent
Le lemme utilisé dans la définition ci-dessus est faux si le morphisme
$f : X \to S$ n'est pas localement de type fini. Nous invitons donc le
lecteur à ne pas employer cette terminologie si $f$ n'est pas localement de
type fini.

\begin{lemma}
\label{lemma-closed-closed-proper-over-base}
Soit $f : X \to S$ un morphisme de schémas localement de type fini.
Soient $Y \subset Z \subset X$ des sous-ensembles fermés.
Si $Z$ est propre sur $S$, alors il en va de même pour $Y$.
\end{lemma}

\begin{proof}
Démonstration omise.
\end{proof}

\begin{lemma}
\label{lemma-base-change-closed-proper-over-base}
Considérons un diagramme cartésien de schémas
$$
\xymatrix{
X' \ar[d]_{f'} \ar[r]_{g'} & X \ar[d]^f \\
S' \ar[r]^g & S
}
$$
où $f$ est localement de type fini.
Si $Z$ est un sous-ensemble fermé de $X$ propre sur $S$, alors
$(g')^{-1}(Z)$ est un sous-ensemble fermé de $X'$ propre sur $S'$.
\end{lemma}

\begin{proof}
Remarquons que l'énoncé a un sens puisque $f'$ est localement de
type fini d'après Morphismes, lemme \ref{morphisms-lemma-base-change-finite-type}.
Munissons $Z$ de la structure de sous-schéma fermé réduit induit.
Notons $Z' = (g')^{-1}(Z)$ l'image inverse schématique (Schémas,
définition \ref{schemes-definition-inverse-image-closed-subscheme}).
Alors $Z' = X' \times_X Z = (S' \times_S X) \times_X Z = S' \times_S Z$
est propre sur $S'$ en tant que changement de base de $Z$ sur $S$
(Morphismes, lemme \ref{morphisms-lemma-base-change-proper}).
\end{proof}

\begin{lemma}
\label{lemma-functoriality-closed-proper-over-base}
Soit $S$ un schéma. Soit $f : X \to Y$ un morphisme entre schémas
localement de type fini sur $S$.
\begin{enumerate}
\item Si $Y$ est séparé sur $S$ et si $Z \subset X$ est un sous-ensemble fermé
propre sur $S$, alors $f(Z)$ est un sous-ensemble fermé de $Y$ propre sur $S$.
\item Si $f$ est universellement fermé et si $Z \subset X$ est un
sous-ensemble fermé propre sur $S$, alors $f(Z)$ est un sous-ensemble fermé
de $Y$ propre sur $S$.
\item Si $f$ est propre et si $Z \subset Y$ est un sous-ensemble fermé
propre sur $S$, alors $f^{-1}(Z)$ est un sous-ensemble fermé de $X$ propre sur $S$.
\end{enumerate}
\end{lemma}

\begin{proof}
Démonstration de (1). Supposons $Y$ séparé sur $S$ et $Z \subset X$
fermé et propre sur $S$. Munissons $Z$ de sa structure réduite
induite de sous-schéma fermé et appliquons
Morphismes, lemme \ref{morphisms-lemma-scheme-theoretic-image-is-proper},
à $Z \to Y$ sur $S$ pour conclure.

\medskip\noindent
Démonstration de (2). Supposons $f$ universellement fermé et $Z \subset X$ fermé
et propre sur $S$. Munissons $Z$ et $Z' = f(Z)$ de leurs structures
réduites induites de sous-schémas fermés. On obtient un morphisme
induit $Z \to Z'$.
Notons $Z'' = f^{-1}(Z')$ l'image inverse schématique (Schémas,
définition \ref{schemes-definition-inverse-image-closed-subscheme}).
Alors $Z'' \to Z'$ est universellement fermé, comme changement de base de $f$
(Morphismes, lemme \ref{morphisms-lemma-base-change-proper}).
Ainsi, $Z \to Z'$ est universellement fermé, comme composé de
l'immersion fermée $Z \to Z''$ et de $Z'' \to Z'$
(Morphismes, Lemmes
\ref{morphisms-lemma-closed-immersion-proper} et
\ref{morphisms-lemma-composition-proper}).
On en déduit que $Z' \to S$ est séparé d'après
Morphismes, lemme \ref{morphisms-lemma-image-universally-closed-separated}.
Puisque $Z \to S$ est quasi-compact et que $Z \to Z'$ est surjectif,
le morphisme $Z' \to S$ est quasi-compact.
Puisque $Z' \to S$ est le composé de $Z' \to Y$ et de $Y \to S$,
le morphisme $Z' \to S$ est localement de type fini
(Morphismes, lemmes \ref{morphisms-lemma-immersion-locally-finite-type} et
\ref{morphisms-lemma-composition-finite-type}).
Enfin, puisque $Z \to S$ est universellement fermé, il
en va de même pour $Z' \to S$ d'après
Morphismes, lemme \ref{morphisms-lemma-image-proper-is-proper}.
Cela achève la démonstration.

\medskip\noindent
Démonstration de (3). Supposons $f$ propre et $Z \subset Y$ fermé et propre
sur $S$. Munissons $Z$ de sa structure réduite induite de sous-schéma
fermé. Notons $Z' = f^{-1}(Z)$ l'image inverse schématique (Schémas,
définition \ref{schemes-definition-inverse-image-closed-subscheme}).
Alors $Z' \to Z$ est propre, comme changement de base de $f$
(Morphismes, lemme \ref{morphisms-lemma-base-change-proper}).
Ainsi $Z' \to S$ est propre, comme composé de $Z' \to Z$
et de $Z \to S$
(Morphismes, lemme \ref{morphisms-lemma-composition-proper}).
Cela achève la démonstration.
\end{proof}

\begin{lemma}
\label{lemma-union-closed-proper-over-base}
Soit $f : X \to S$ un morphisme de schémas localement de type fini.
Soient $Z_i \subset X$, $i = 1, \ldots, n$, des sous-ensembles fermés.
Si les $Z_i$, $i = 1, \ldots, n$, sont propres sur $S$, il en va de
même pour $Z_1 \cup \ldots \cup Z_n$.
\end{lemma}

\begin{proof}
Munissons les $Z_i$ de leurs structures réduites induites de sous-schémas fermés.
Le morphisme
$$
Z_1 \amalg \ldots \amalg Z_n \longrightarrow X
$$
est fini d'après Morphismes, Lemmes
\ref{morphisms-lemma-closed-immersion-finite} et
\ref{morphisms-lemma-finite-union-finite}.
Les morphismes finis étant universellement fermés
(Morphismes, lemme \ref{morphisms-lemma-finite-proper}),
et $Z_1 \amalg \ldots \amalg Z_n$ étant propre sur $S$,
on déduit du
lemme \ref{lemma-functoriality-closed-proper-over-base} (2)
que l'image $Z_1 \cup \ldots \cup Z_n$ est propre sur $S$.
\end{proof}

\noindent
Soit $f : X \to S$ un morphisme de schémas localement de
type fini. Soit $\mathcal{F}$ un
$\mathcal{O}_X$-module quasi-cohérent de type fini. Alors le support $\text{Supp}(\mathcal{F})$
de $\mathcal{F}$ est un sous-ensemble fermé de $X$, voir
Morphismes, lemme \ref{morphisms-lemma-support-finite-type}.
On peut donc dire que «
le support de $\mathcal{F}$ est propre sur $S$ ».

\begin{lemma}
\label{lemma-module-support-proper-over-base}
Soit $f : X \to S$ un morphisme de schémas localement de
type fini. Soit $\mathcal{F}$ un
$\mathcal{O}_X$-module quasi-cohérent de type fini. Les conditions suivantes sont équivalentes :
\begin{enumerate}
\item le support de $\mathcal{F}$ est propre sur $S$,
\item le support schématique de $\mathcal{F}$
(Morphismes, définition \ref{morphisms-definition-scheme-theoretic-support})
est propre sur $S$, et
\item il existe un sous-schéma fermé $Z \subset X$ et
un $\mathcal{O}_Z$-module quasi-cohérent de type fini
$\mathcal{G}$ tels que (a) $Z \to S$ soit propre et (b)
$(Z \to X)_*\mathcal{G} = \mathcal{F}$.
\end{enumerate}
\end{lemma}

\begin{proof}
Le support $\text{Supp}(\mathcal{F})$ de $\mathcal{F}$ est un sous-ensemble fermé
de $X$, voir Morphismes, lemme \ref{morphisms-lemma-support-finite-type}.
On peut donc appliquer la définition \ref{definition-proper-over-base}.
Puisque le support schématique de $\mathcal{F}$ est un sous-schéma
fermé dont le sous-ensemble fermé sous-jacent est $\text{Supp}(\mathcal{F})$,
les conditions (1) et (2) sont équivalentes d'après la
définition \ref{definition-proper-over-base}.
Il est clair que (2) implique (3).
Réciproquement, si (3) est vérifiée, alors
$\text{Supp}(\mathcal{F}) \subset Z$
(inclusion de sous-ensembles fermés de $X$),
et donc $\text{Supp}(\mathcal{F})$
est propre sur $S$, par exemple d'après
le lemme \ref{lemma-closed-closed-proper-over-base}.
\end{proof}

\begin{lemma}
\label{lemma-base-change-module-support-proper-over-base}
Considérons un diagramme cartésien de schémas
$$
\xymatrix{
X' \ar[d]_{f'} \ar[r]_{g'} & X \ar[d]^f \\
S' \ar[r]^g & S
}
$$
où $f$ est localement de type fini. Soit $\mathcal{F}$
un $\mathcal{O}_X$-module quasi-cohérent
de type fini. Si le support de $\mathcal{F}$ est propre sur $S$, alors
le support de $(g')^*\mathcal{F}$ est propre sur $S'$.
\end{lemma}

\begin{proof}
Remarquons que l'énoncé a un sens, car
$(g')*\mathcal{F}$ est de type fini d'après
Modules, lemme \ref{modules-lemma-pullback-finite-type}.
On a $\text{Supp}((g')^*\mathcal{F}) = (g')^{-1}(\text{Supp}(\mathcal{F}))$
d'après Morphismes, lemme \ref{morphisms-lemma-support-finite-type}.
Le lemme résulte donc du
lemme \ref{lemma-base-change-closed-proper-over-base}.
\end{proof}

\begin{lemma}
\label{lemma-cat-module-support-proper-over-base}
Soit $f : X \to S$ un morphisme de schémas localement de
type fini. Soient $\mathcal{F}$, $\mathcal{G}$
des $\mathcal{O}_X$-modules quasi-cohérents de type fini.
\begin{enumerate}
\item Si les supports de $\mathcal{F}$, $\mathcal{G}$
sont propres sur $S$, il en va de même
pour $\mathcal{F} \oplus \mathcal{G}$, pour toute extension
de $\mathcal{G}$ par $\mathcal{F}$, pour $\Im(u)$ et $\Coker(u)$,
quel que soit le morphisme de $\mathcal{O}_X$-modules $u : \mathcal{F} \to \mathcal{G}$,
et pour tout quotient quasi-cohérent de $\mathcal{F}$ ou de $\mathcal{G}$.
\item Si $S$ est localement noethérien, la catégorie
des $\mathcal{O}_X$-modules cohérents à support propre sur
$S$ est une sous-catégorie de Serre (Homologie, Définition
\ref{homology-definition-serre-subcategory})
de la catégorie abélienne
des $\mathcal{O}_X$-modules cohérents.
\end{enumerate}
\end{lemma}

\begin{proof}
Démonstration de (1). Soient $Z$, $Z'$ les supports de $\mathcal{F}$
et de $\mathcal{G}$. Tous les faisceaux mentionnés dans (1) ont
leur support contenu dans $Z \cup Z'$. L'assertion résulte
donc des lemmes \ref{lemma-closed-closed-proper-over-base} et
\ref{lemma-union-closed-proper-over-base}
dès que l'on a vérifié que ces faisceaux sont de type
fini et quasi-cohérents. Pour la quasi-cohérence, nous renvoyons le lecteur à
Schémas, section \ref{schemes-section-quasi-coherent}.
Pour la propriété « de type fini », nous invitons le lecteur à consulter
Modules, section \ref{modules-section-finite-type}.

\medskip\noindent
Démonstration de (2). On raisonne comme pour (1). Notons que les assertions
ont un sens, puisque $X$ est localement noethérien d'après
Morphismes, lemme \ref{morphisms-lemma-finite-type-noetherian},
et compte tenu de la description de la catégorie des modules
cohérents dans la section \ref{section-coherent-sheaves}.
\end{proof}

\begin{lemma}
\label{lemma-support-proper-over-base-pushforward}
Soit $S$ un schéma localement noethérien.
Soit $f : X \to S$ un morphisme de schémas localement de type fini.
Soit $\mathcal{F}$ un $\mathcal{O}_X$-module
cohérent à support propre sur $S$. Alors $R^pf_*\mathcal{F}$
est un $\mathcal{O}_S$-module cohérent pour tout $p \geq 0$.
\end{lemma}

\begin{proof}
D'après le lemme \ref{lemma-module-support-proper-over-base},
il existe une immersion fermée $i : Z \to X$ et
un $\mathcal{O}_Z$-module quasi-cohérent de type fini
$\mathcal{G}$ tels que (a) $g = f \circ i : Z \to S$ soit propre et (b)
$i_*\mathcal{G} = \mathcal{F}$.
Le faisceau $R^pg_*\mathcal{G}$ est cohérent sur $S$ d'après la
proposition \ref{proposition-proper-pushforward-coherent}.
D'autre part, $R^qi_*\mathcal{G} = 0$ pour $q > 0$
(lemme \ref{lemma-finite-pushforward-coherent}).
D'après Cohomologie, lemme \ref{cohomology-lemma-relative-Leray},
on obtient $R^pf_*\mathcal{F} = R^pg_*\mathcal{G}$, ce qui achève la démonstration.
\end{proof}

\begin{lemma}
\label{lemma-systems-with-proper-support}
Soit $S$ un schéma noethérien. Soit $f : X \to S$ un morphisme de type fini.
Soit $\mathcal{I} \subset \mathcal{O}_X$ un
faisceau quasi-cohérent d'idéaux. Les catégories suivantes sont des sous-catégories de Serre
de $\textit{Coh}(X, \mathcal{I})$ :
\begin{enumerate}
\item la sous-catégorie pleine de $\textit{Coh}(X, \mathcal{I})$
formée des objets $(\mathcal{F}_n)$ tels que
le support de $\mathcal{F}_1$ soit propre sur $S$,
\item la sous-catégorie pleine de $\textit{Coh}(X, \mathcal{I})$
formée des objets $(\mathcal{F}_n)$ tels qu'il
existe un sous-schéma fermé $Z \subset X$ propre sur $S$
avec $\mathcal{I}_Z \mathcal{F}_n = 0$ pour tout $n \geq 1$.
\end{enumerate}
\end{lemma}

\begin{proof}
Utilisons le critère de
Homologie, lemme \ref{homology-lemma-characterize-serre-subcategory}.
Nous utiliserons aussi le fait que, si
$0 \to (\mathcal{G}_n) \to (\mathcal{F}_n) \to (\mathcal{H}_n) \to 0$
est une suite exacte courte de $\textit{Coh}(X, \mathcal{I})$, alors
(a) $\mathcal{G}_n \to \mathcal{F}_n \to \mathcal{H}_n \to 0$
est exacte pour tout $n \geq 1$ et
(b) $\mathcal{G}_n$ est un quotient de $\Ker(\mathcal{F}_m \to \mathcal{H}_m)$
pour un certain $m \geq n$. Voir la démonstration du lemme \ref{lemma-inverse-systems-abelian}.

\medskip\noindent
Démonstration de (1). Soit $(\mathcal{F}_n)$ un objet de
$\textit{Coh}(X, \mathcal{I})$. Alors
$\text{Supp}(\mathcal{F}_n) = \text{Supp}(\mathcal{F}_1)$ pour tout $n \geq 1$.
Les remarques (a) et (b) ci-dessus montrent donc que,
pour toute suite exacte courte
$0 \to (\mathcal{G}_n) \to (\mathcal{F}_n) \to (\mathcal{H}_n) \to 0$
de $\textit{Coh}(X, \mathcal{I})$, on a
$\text{Supp}(\mathcal{G}_1) \cup \text{Supp}(\mathcal{H}_1) =
\text{Supp}(\mathcal{F}_1)$.
Cela prouve que la catégorie définie dans (1) est
une sous-catégorie de Serre de $\textit{Coh}(X, \mathcal{I})$.

\medskip\noindent
Démonstration de (2). On raisonne de même. Soit
$0 \to (\mathcal{G}_n) \to (\mathcal{F}_n) \to (\mathcal{H}_n) \to 0$
une suite exacte courte de $\textit{Coh}(X, \mathcal{I})$.
Si $Z \subset X$ est un sous-schéma fermé et si $\mathcal{I}_Z$
annule $\mathcal{F}_n$ pour tout $n$, alors
$\mathcal{I}_Z$ annule $\mathcal{G}_n$ et $\mathcal{H}_n$
pour tout $n$ d'après (a) et (b) ci-dessus.
Si donc $Z \to S$ est propre, on en déduit que la catégorie
définie dans (2) est stable par passage aux sous-objets et aux objets
quotients dans $\textit{Coh}(X, \mathcal{I})$.
Enfin, supposons que $Z \subset X$ et $Y \subset X$ soient
des sous-schémas fermés propres sur $S$ tels que
$\mathcal{I}_Z \mathcal{G}_n = 0$ et
$\mathcal{I}_Y \mathcal{H}_n = 0$ pour tout $n \geq 1$.
Il résulte alors de (a) ci-dessus que
$\mathcal{I}_{Z \cup Y} = \mathcal{I}_Z \cdot \mathcal{I}_Y$
annule $\mathcal{F}_n$ pour tout $n$.
D'après le lemme \ref{lemma-union-closed-proper-over-base}
(et la définition \ref{definition-proper-over-base},
qui permet de choisir une structure de schéma quelconque sur la
réunion), le morphisme $Z \cup Y \to S$ est propre, ce qui achève la démonstration.
\end{proof}








\section{Théorème d'existence de Grothendieck, III}
\label{section-existence-proper-support}

\noindent
Pour énoncer la version générale du théorème d'existence
de Grothendieck, introduisons quelques notations supplémentaires. Soit $A$ un anneau noethérien complet
pour la topologie définie par un idéal $I$. Soit $f : X \to \Spec(A)$
un morphisme de schémas séparé et de type fini. Posons
$\mathcal{I} = I\mathcal{O}_X$. Dans cette situation, notons
$$
\textit{Coh}_{\text{support propre sur } A}(\mathcal{O}_X)
$$
la sous-catégorie pleine de $\textit{Coh}(\mathcal{O}_X)$
formée des $\mathcal{O}_X$-modules cohérents dont
le support est propre sur $\Spec(A)$. C'est une sous-catégorie de Serre de
$\textit{Coh}(\mathcal{O}_X)$, voir
lemme \ref{lemma-cat-module-support-proper-over-base}.
De même, notons
$$
\textit{Coh}_{\text{support propre sur } A}(X, \mathcal{I})
$$
la sous-catégorie pleine de $\textit{Coh}(X, \mathcal{I})$
formée des objets $(\mathcal{F}_n)$ tels que
le support de $\mathcal{F}_1$ soit propre sur $\Spec(A)$.
C'est une sous-catégorie de Serre de $\textit{Coh}(X, \mathcal{I})$
d'après le lemme \ref{lemma-systems-with-proper-support} (1). Puisque
le support d'un module quotient est contenu dans celui
du module, le foncteur (\ref{equation-completion-functor})
induit un foncteur
\begin{equation}
\label{equation-completion-functor-proper-over-A}
\textit{Coh}_{\text{support propre sur }A}(\mathcal{O}_X)
\longrightarrow
\textit{Coh}_{\text{support propre sur }A}(X, \mathcal{I})
\end{equation}
Nous pouvons maintenant énoncer le théorème principal de cette section.

\begin{theorem}[Théorème d'existence de Grothendieck]
\label{theorem-grothendieck-existence}
\begin{reference}
\cite[III théorème 5.1.5]{EGA}
\end{reference}
Soit $A$ un anneau noethérien complet pour la topologie définie par un idéal $I$.
Soit $X$ un schéma séparé et de type fini sur $A$. Alors
le foncteur
(\ref{equation-completion-functor-proper-over-A})
$$
\textit{Coh}_{\text{support propre sur }A}(\mathcal{O}_X)
\longrightarrow
\textit{Coh}_{\text{support propre sur }A}(X, \mathcal{I})
$$
est une équivalence.
\end{theorem}

\begin{proof}
Nous utiliserons l'équivalence de catégories du
lemme \ref{lemma-i-star-equivalence}
sans autre mention.
Pour un sous-schéma fermé $Z \subset X$ propre sur $A$,
nous dirons, dans cette démonstration, qu'un module cohérent sur $X$ est
« porté par $Z$ » s'il est annulé par le faisceau
d'idéaux de $Z$, ou, de manière équivalente, s'il est l'image
directe d'un module cohérent sur $Z$.
D'après la proposition \ref{proposition-existence-proper}, le
résultat est vrai pour le
foncteur entre modules cohérents et systèmes de modules
cohérents portés par $Z$. Il suffit donc de montrer que
tout objet de
$\textit{Coh}_{\text{support propre sur }A}(\mathcal{O}_X)$
et tout objet de
$\textit{Coh}_{\text{support propre sur }A}(X, \mathcal{I})$ est
porté par un sous-schéma fermé $Z \subset X$ propre sur $A$.
C'est vrai par définition pour les objets de
$\textit{Coh}_{\text{support propre sur }A}(\mathcal{O}_X)$.
Nous allons le démontrer pour les objets de
$\textit{Coh}_{\text{support propre sur }A}(X, \mathcal{I})$
en utilisant la méthode de démonstration de la proposition \ref{proposition-existence-proper}.
Nous invitons le lecteur à lire d'abord cette démonstration.

\medskip\noindent
Considérons l'ensemble $\Xi$ des faisceaux quasi-cohérents d'idéaux
$\mathcal{K} \subset \mathcal{O}_X$ tels que l'assertion soit vraie
pour tout objet $(\mathcal{F}_n)$ de
$\textit{Coh}_{\text{support propre sur }A}(X, \mathcal{I})$
annulé par $\mathcal{K}$. Montrons que $(0)$ appartient à $\Xi$.
Sinon, puisque $X$ est noethérien, il existe un
faisceau quasi-cohérent d'idéaux $\mathcal{K}$ maximal parmi ceux qui n'appartiennent pas à $\Xi$, voir
lemme \ref{lemma-acc-coherent}.
En remplaçant $X$ par le sous-schéma fermé de $X$
correspondant à $\mathcal{K}$, on peut supposer que tout faisceau non nul
$\mathcal{K}$ appartient à $\Xi$. Soit $(\mathcal{F}_n)$ un objet de
$\textit{Coh}_{\text{support propre sur }A}(X, \mathcal{I})$.
Nous allons montrer que cet objet est porté par un sous-schéma fermé
$Z \subset X$ propre sur $A$, ce qui achèvera
la démonstration du théorème.

\medskip\noindent
Appliquons le lemme de Chow (lemme \ref{lemma-chow-Noetherian}) pour obtenir un
morphisme propre surjectif $f : Y \to X$ qui est un isomorphisme
au-dessus d'un ouvert dense $U \subset X$ et tel que $Y$ soit H-quasi-projectif
sur $A$. Choisissons une immersion ouverte $j : Y \to Y'$, où
$Y'$ est projectif sur $A$, voir
Morphismes, lemme \ref{morphisms-lemma-H-quasi-projective-open-H-projective}.
Remarquons que
$$
\text{Supp}(f^*\mathcal{F}_n) = f^{-1}\text{Supp}(\mathcal{F}_n) =
f^{-1}\text{Supp}(\mathcal{F}_1)
$$
La première égalité résulte de
Morphismes, lemme \ref{morphisms-lemma-support-finite-type}.
Par hypothèse et d'après le
lemme \ref{lemma-functoriality-closed-proper-over-base} (3),
le sous-ensemble $f^{-1}\text{Supp}(\mathcal{F}_1)$ est propre sur $A$.
L'image de $f^{-1}\text{Supp}(\mathcal{F}_1)$
par $j$ est donc fermée dans $Y'$ d'après le
lemme \ref{lemma-functoriality-closed-proper-over-base} (1).
Ainsi $\mathcal{F}'_n = j_*f^*\mathcal{F}_n$ est cohérent sur
$Y'$ d'après le lemme \ref{lemma-pushforward-coherent-on-open}.
Il en résulte que $(\mathcal{F}_n')$
est un objet de $\textit{Coh}(Y', I\mathcal{O}_{Y'})$.
D'après le cas projectif du théorème d'existence de Grothendieck
(lemme \ref{lemma-existence-projective}),
il existe un $\mathcal{O}_{Y'}$-module cohérent
$\mathcal{F}'$ et un isomorphisme
$(\mathcal{F}')^\wedge \cong (\mathcal{F}'_n)$ dans
$\textit{Coh}(Y', I\mathcal{O}_{Y'})$.
Puisque $\mathcal{F}'/I\mathcal{F}' = \mathcal{F}'_1$, on a
$$
\text{Supp}(\mathcal{F}') \cap V(I\mathcal{O}_{Y'}) =
\text{Supp}(\mathcal{F}'_1) = j(f^{-1}\text{Supp}(\mathcal{F}_1))
$$
Le morphisme structural $p' : Y' \to \Spec(A)$ est propre, donc
$p'(\text{Supp}(\mathcal{F}') \setminus j(Y))$
est fermé dans $\Spec(A)$. Un sous-ensemble fermé non vide de $\Spec(A)$
contient un point de $V(I)$, car $I$ est contenu dans le radical de Jacobson
de $A$ d'après Algèbre, lemme \ref{algebra-lemma-radical-completion}.
L'égalité ci-dessus montre que
$\text{Supp}(\mathcal{F}') \cap (p')^{-1}V(I) \subset j(Y)$
et l'on en déduit que $\text{Supp}(\mathcal{F}') \subset j(Y)$.
Ainsi $\mathcal{F}'|_Y = j^*\mathcal{F}'$
est porté par un sous-schéma fermé $Z'$ de $Y$ propre sur $A$,
et $(\mathcal{F}'|_Y)^\wedge = (f^*\mathcal{F}_n)$.

\medskip\noindent
Soit $\mathcal{K}$ le faisceau quasi-cohérent d'idéaux définissant
le complémentaire réduit $X \setminus U$. D'après la
proposition \ref{proposition-proper-pushforward-coherent},
le $\mathcal{O}_X$-module $\mathcal{H} = f_*(\mathcal{F}'|_Y)$ est cohérent,
et d'après le lemme \ref{lemma-inverse-systems-push-pull},
il existe un morphisme $\alpha : (\mathcal{F}_n) \to \mathcal{H}^\wedge$
de $\textit{Coh}(X, \mathcal{I})$ dont le noyau et le conoyau
sont annulés par une puissance $\mathcal{K}^t$ de $\mathcal{K}$.
On obtient une suite exacte
$$
0 \to \Ker(\alpha) \to (\mathcal{F}_n) \to
\mathcal{H}^\wedge \to \Coker(\alpha) \to 0
$$
dans $\textit{Coh}(X, \mathcal{I})$. Si $Z_0 \subset X$ est le support
schématique de $\mathcal{H}$, il est clair que $Z_0 \subset f(Z')$
ensemblistement. Ainsi $Z_0$ est propre sur $A$ d'après le
lemme \ref{lemma-closed-closed-proper-over-base} et le
lemme \ref{lemma-functoriality-closed-proper-over-base} (2). Par
conséquent, $\mathcal{H}^\wedge$ appartient à la sous-catégorie définie dans le
lemme \ref{lemma-systems-with-proper-support} (2), et
a fortiori à
$\textit{Coh}_{\text{support propre sur }A}(X, \mathcal{I})$.
On en déduit que $\Ker(\alpha)$ et $\Coker(\alpha)$
appartiennent à $\textit{Coh}_{\text{support propre sur }A}(X, \mathcal{I})$
d'après le lemme \ref{lemma-systems-with-proper-support} (1). Par
l'hypothèse de récurrence, plus précisément parce que $\mathcal{K}^t$ appartient à $\Xi$,
les objets $\Ker(\alpha)$ et $\Coker(\alpha)$ appartiennent à
la sous-catégorie définie dans le
lemme \ref{lemma-systems-with-proper-support} (2). Celle-ci
étant une sous-catégorie de Serre d'après le lemme, il en va
de même pour $(\mathcal{F}_n)$, comme voulu.
\end{proof}

\begin{remark}[Explicitation du théorème d'existence de Grothendieck]
\label{remark-reformulate-existence-theorem}
Soit $A$ un anneau noethérien complet pour la topologie définie par un idéal $I$.
Écrivons $S = \Spec(A)$ et $S_n = \Spec(A/I^n)$.
Soit $X \to S$ un morphisme séparé de type fini.
Pour $n \geq 1$, posons $X_n = X \times_S S_n$. On a le
diagramme suivant :
$$
\xymatrix{
X_1 \ar[r]_{i_1} \ar[d] & X_2 \ar[r]_{i_2} \ar[d] & X_3 \ar[r] \ar[d] &
\ldots & X \ar[d] \\
S_1 \ar[r] & S_2 \ar[r] & S_3 \ar[r] & \ldots & S
}
$$
Dans cette situation, considérons les systèmes $(\mathcal{F}_n, \varphi_n)$
tels que
\begin{enumerate}
\item $\mathcal{F}_n$ soit un $\mathcal{O}_{X_n}$-module cohérent,
\item $\varphi_n : i_n^*\mathcal{F}_{n + 1} \to \mathcal{F}_n$
soit un isomorphisme, et
\item $\text{Supp}(\mathcal{F}_1)$ soit propre sur $S_1$.
\end{enumerate}
Le Théorème \ref{theorem-grothendieck-existence} affirme que le foncteur
de complétion
$$
\begin{matrix}
\text{les }\mathcal{O}_X\text{-modules cohérents }\mathcal{F} \\
\text{à support propre sur }A
\end{matrix}
\quad
\longrightarrow
\quad
\begin{matrix}
\text{systèmes }(\mathcal{F}_n) \\
\text{comme ci-dessus}
\end{matrix}
$$
est une équivalence de catégories. Dans le cas particulier où $X$ est
propre sur $A$, on peut omettre les conditions portant sur les supports.
\end{remark}



\section{Théorème d'algébrisation de Grothendieck}
\label{section-algebraization}

\noindent
Notre premier résultat est une reformulation du théorème d'existence de
Grothendieck en termes de sous-schémas fermés et de morphismes finis.

\begin{lemma}
\label{lemma-algebraize-formal-closed-subscheme}
Soit $A$ un anneau noethérien complet pour la topologie $I$-adique.
Posons $S = \Spec(A)$ et $S_n = \Spec(A/I^n)$.
Soit $X \to S$ un morphisme séparé de type fini.
Pour $n \geq 1$, posons $X_n = X \times_S S_n$.
Supposons donné un diagramme commutatif
$$
\xymatrix{
Z_1 \ar[r] \ar[d] & Z_2 \ar[r] \ar[d] & Z_3 \ar[r] \ar[d] & \ldots \\
X_1 \ar[r]^{i_1} & X_2 \ar[r]^{i_2} & X_3 \ar[r] & \ldots
}
$$
de schémas dont les carrés sont cartésiens. Supposons que
\begin{enumerate}
\item $Z_1 \to X_1$ est une immersion fermée, et
\item $Z_1 \to S_1$ est propre.
\end{enumerate}
Il existe alors une immersion fermée de schémas $Z \to X$ telle que
$Z_n = Z \times_S S_n$. De plus, $Z$ est propre sur $S$.
\end{lemma}

\begin{proof}
Notons $j_n : Z_n \to X_n$ les morphismes verticaux.
Les carrés de l'énoncé étant cartésiens,
le changement de base de $j_n$ à $X_1$ est $j_1$.
Ainsi, Morphismes, lemme \ref{morphisms-lemma-check-closed-infinitesimally},
montre que $j_n$ est une immersion fermée.
Posons $\mathcal{F}_n = j_{n, *}\mathcal{O}_{Z_n}$, de sorte que
$j_n^\sharp$ soit une surjection $\mathcal{O}_{X_n} \to \mathcal{F}_n$.
En utilisant à nouveau le caractère cartésien des carrés, on voit que
l'image inverse de $\mathcal{F}_{n + 1}$ sur $X_n$ est $\mathcal{F}_n$.
Le théorème d'existence de Grothendieck, sous la forme de la
remarque \ref{remark-reformulate-existence-theorem},
assure donc l'existence d'un morphisme
$\mathcal{O}_X \to \mathcal{F}$
de $\mathcal{O}_X$-modules cohérents dont la restriction à
$X_n$ redonne $\mathcal{O}_{X_n} \to \mathcal{F}_n$.
De plus, le support de $\mathcal{F}$ est propre sur $S$.
Comme le foncteur de complétion est exact (lemme \ref{lemma-exact}),
le conoyau $\mathcal{Q}$ de $\mathcal{O}_X \to \mathcal{F}$
a une complétion nulle. Comme le support de $\mathcal{F}$ est propre
sur $S$, et qu'il en est de même de $\mathcal{Q}$, il en résulte que
$\mathcal{Q} = 0$ ; cela découle par exemple du fait que le
foncteur (\ref{equation-completion-functor-proper-over-A}) est une équivalence par
le théorème d'existence de Grothendieck.
Ainsi $\mathcal{F} = \mathcal{O}_X/\mathcal{J}$
pour un certain faisceau quasi-cohérent d'idéaux $\mathcal{J}$. Il
suffit de poser $Z = V(\mathcal{J})$ pour achever la démonstration.
\end{proof}

\noindent
Dans le lemme suivant, il suffit en fait de supposer que $Y_1 \to X_1$
est fini, car cela entraîne que $Y_n \to X_n$ est également fini
(voir Compléments sur les morphismes, Lemme
\ref{more-morphisms-lemma-thicken-property-morphisms-cartesian}).

\begin{lemma}
\label{lemma-algebraize-formal-scheme-finite-over-proper}
Soit $A$ un anneau noethérien complet pour la topologie $I$-adique.
Posons $S = \Spec(A)$ et $S_n = \Spec(A/I^n)$.
Soit $X \to S$ un morphisme séparé de type fini.
Pour $n \geq 1$, posons $X_n = X \times_S S_n$.
Supposons donné un diagramme commutatif
$$
\xymatrix{
Y_1 \ar[r] \ar[d] & Y_2 \ar[r] \ar[d] & Y_3 \ar[r] \ar[d] & \ldots \\
X_1 \ar[r]^{i_1} & X_2 \ar[r]^{i_2} & X_3 \ar[r] & \ldots
}
$$
de schémas dont les carrés sont cartésiens. Supposons que
\begin{enumerate}
\item $Y_n \to X_n$ est un morphisme fini, et
\item $Y_1 \to S_1$ est propre.
\end{enumerate}
Il existe alors un morphisme fini de schémas $Y \to X$ tel que
$Y_n = Y \times_S S_n$. De plus, $Y$ est propre sur $S$.
\end{lemma}

\begin{proof}
Notons $f_n : Y_n \to X_n$ les morphismes verticaux.
Posons $\mathcal{F}_n = f_{n, *}\mathcal{O}_{Y_n}$. C'est
un $\mathcal{O}_{X_n}$-module cohérent puisque $f_n$ est fini
(lemme \ref{lemma-finite-pushforward-coherent}).
Le caractère cartésien des carrés montre que
l'image inverse de $\mathcal{F}_{n + 1}$ sur $X_n$ est $\mathcal{F}_n$.
Le théorème d'existence de Grothendieck, sous la forme de la
remarque \ref{remark-reformulate-existence-theorem},
assure donc l'existence d'un $\mathcal{O}_X$-module cohérent
$\mathcal{F}$ dont la restriction à $X_n$ redonne $\mathcal{F}_n$.
De plus, le support de $\mathcal{F}$ est propre sur $S$.
Comme le foncteur de complétion est pleinement fidèle
(Théorème \ref{theorem-grothendieck-existence}),
les applications de multiplication
$\mathcal{F}_n \otimes_{\mathcal{O}_{X_n}} \mathcal{F}_n \to
\mathcal{F}_n$ se recollent en une structure d'algèbre sur $\mathcal{F}$. Il
suffit de poser $Y = \underline{\Spec}_X(\mathcal{F})$ pour achever la démonstration.
\end{proof}

\begin{lemma}
\label{lemma-algebraize-morphism}
Soit $A$ un anneau noethérien complet pour la topologie $I$-adique.
Posons $S = \Spec(A)$ et $S_n = \Spec(A/I^n)$. Soient $X$ et $Y$ des schémas
sur $S$. Pour $n \geq 1$, posons $X_n = X \times_S S_n$ et
$Y_n = Y \times_S S_n$. Supposons donné un système compatible de
diagrammes commutatifs
$$
\xymatrix{
& & X_{n + 1} \ar[rd] \ar[rr]_{g_{n + 1}} & & Y_{n + 1} \ar[ld] \\
X_n \ar[rru] \ar[rd] \ar[rr]_{g_n} & & Y_n \ar[rru] \ar[ld] & S_{n + 1} \\
& S_n \ar[rru]
}
$$
Supposons que
\begin{enumerate}
\item $X \to S$ est propre, et
\item $Y \to S$ est séparé de type fini.
\end{enumerate}
Il existe alors un unique morphisme de schémas $g : X \to Y$
au-dessus de $S$ tel que $g_n$ soit le changement de base de $g$ à $S_n$.
\end{lemma}

\begin{proof}
Les morphismes $(1, g_n) : X_n \to X_n \times_S Y_n$ sont des immersions fermées,
car $Y_n \to S_n$ est séparé
(Schémas, lemme \ref{schemes-lemma-section-immersion}). Le
lemme \ref{lemma-algebraize-formal-closed-subscheme}
fournit donc un sous-schéma fermé $Z \subset X \times_S Y$,
propre sur $S$, dont le changement de base à $S_n$ redonne
$X_n \subset X_n \times_S Y_n$. La première projection $p : Z \to X$
est un morphisme propre (car $Z$ est propre sur $S$, voir
Morphismes, lemme \ref{morphisms-lemma-image-proper-scheme-closed})
dont le changement de base à $S_n$ est un isomorphisme
pour tout $n$. En particulier, $p : Z \to X$ est fini au-dessus d'un voisinage
ouvert de $X_0$, d'après
le lemme \ref{lemma-proper-finite-fibre-finite-in-neighbourhood}.
Comme $X$ est propre sur $S$, ce voisinage ouvert est $X$
tout entier, et nous concluons que $p : Z \to X$ est fini.
En appliquant l'équivalence de la proposition \ref{proposition-existence-proper},
on obtient $p_*\mathcal{O}_Z = \mathcal{O}_X$, puisque cette égalité est vraie
modulo $I^n$ pour tout $n$. Par conséquent, $p$ est un isomorphisme, et
le morphisme $g$ est le composé $X \cong Z \to Y$.
Nous omettons la démonstration de l'unicité.
\end{proof}

\noindent
Pour démontrer un théorème d'algébrisation « abstrait », nous
devons supposer l'existence d'un faisceau inversible ample : sans cette hypothèse,
le résultat est faux.

\begin{theorem}[Théorème d'algébrisation de Grothendieck]
\label{theorem-algebraization}
Soit $A$ un anneau noethérien complet pour la topologie $I$-adique.
Posons $S = \Spec(A)$ et $S_n = \Spec(A/I^n)$. Considérons un diagramme
commutatif
$$
\xymatrix{
X_1 \ar[r]_{i_1} \ar[d] & X_2 \ar[r]_{i_2} \ar[d] & X_3 \ar[r] \ar[d] &
\ldots \\
S_1 \ar[r] & S_2 \ar[r] & S_3 \ar[r] & \ldots
}
$$
de schémas dont les carrés sont cartésiens. Supposons donnés des couples $(\mathcal{L}_n, \varphi_n)$,
où chaque $\mathcal{L}_n$ est un faisceau inversible sur $X_n$ et
$\varphi_n : i_n^*\mathcal{L}_{n + 1} \to \mathcal{L}_n$
est un isomorphisme. Si
\begin{enumerate}
\item $X_1 \to S_1$ est propre, et
\item $\mathcal{L}_1$ est ample sur $X_1$,
\end{enumerate}
alors il existe un morphisme propre de schémas $X \to S$,
un $\mathcal{O}_X$-module inversible ample $\mathcal{L}$,
ainsi que des isomorphismes $X_n \cong X \times_S S_n$ et
$\mathcal{L}_n \cong \mathcal{L}|_{X_n}$ compatibles avec
les morphismes $i_n$ et $\varphi_n$.
\end{theorem}

\begin{proof}
Les carrés du diagramme étant cartésiens et les morphismes
$S_n \to S_{n + 1}$ étant des immersions fermées, les morphismes
$i_n$ sont eux aussi des immersions fermées. Nous pouvons donc considérer
$X_m$ comme un sous-schéma fermé de $X_n$ pour $m < n$. Plus précisément, $X_m$ est
le sous-schéma fermé défini par le faisceau quasi-cohérent d'idéaux
$I^m\mathcal{O}_{X_n}$. De plus, les espaces topologiques
sous-jacents aux schémas $X_1, X_2, X_3, \ldots$ s'identifient tous ; nous pouvons
donc considérer, et considérerons, les faisceaux $\mathcal{O}_{X_n}$ comme définis sur un même
espace topologique sous-jacent, et de même pour les
$\mathcal{O}_{X_n}$-modules cohérents. Posons
$$
\mathcal{F}_n =
\Ker(\mathcal{O}_{X_{n + 1}} \to \mathcal{O}_{X_n})
$$
de sorte que l'on obtienne des suites exactes courtes
$$
0 \to \mathcal{F}_n \to \mathcal{O}_{X_{n + 1}} \to \mathcal{O}_{X_n} \to 0
$$
On a, d'après ce qui précède, $\mathcal{F}_n = I^n\mathcal{O}_{X_{n + 1}}$.
Il s'ensuit que $\mathcal{F}_n$ est un faisceau cohérent sur $X_{n + 1}$
annulé par $I$ ; nous pouvons donc le considérer, et le considérerons, comme
un $\mathcal{O}_{X_1}$-module cohérent. Remarquons que, pour $m > n$, le faisceau
$$
I^n\mathcal{O}_{X_m}/I^{n + 1}\mathcal{O}_{X_m}
$$
s'envoie isomorphiquement sur $\mathcal{F}_n$ par le morphisme
$\mathcal{O}_{X_m} \to \mathcal{O}_{X_{n + 1}}$. Étant donnés
$n_1, n_2 \geq 0$, on peut donc choisir $m > n_1 + n_2$ et considérer le morphisme
de multiplication
$$
I^{n_1}\mathcal{O}_{X_m} \times I^{n_2}\mathcal{O}_{X_m}
\longrightarrow
I^{n_1 + n_2}\mathcal{O}_{X_m} \to \mathcal{F}_{n_1 + n_2}
$$
Celui-ci induit une application $\mathcal{O}_{X_1}$-bilinéaire
$$
\mathcal{F}_{n_1} \times \mathcal{F}_{n_2} \longrightarrow
\mathcal{F}_{n_1 + n_2}
$$
qui définit à son tour une structure de $\mathcal{O}_{X_1}$-algèbre graduée
sur $\mathcal{F} = \bigoplus_{n \geq 0} \mathcal{F}_n$.

\medskip\noindent
Posons $B = \bigoplus I^n/I^{n + 1}$ ; c'est
une $A/I$-algèbre graduée de type fini. Posons $\mathcal{B} = (X_1 \to S_1)^*\widetilde{B}$.
La discussion précédente fournit une surjection canonique
$$
\mathcal{B} \longrightarrow \mathcal{F}
$$
de $\mathcal{O}_{X_1}$-algèbres graduées. En particulier,
$\mathcal{F}$ est un $\mathcal{B}$-module gradué quasi-cohérent
de type fini. Le lemme \ref{lemma-graded-finiteness} permet de choisir un entier $d_0$
tel que $H^1(X_1, \mathcal{F} \otimes \mathcal{L}^{\otimes d}) = 0$
pour tout $d \geq d_0$. Choisissons $d \geq d_0$ de sorte qu'il existe des sections
$s_{0, 1}, \ldots, s_{N, 1} \in \Gamma(X_1, \mathcal{L}_1^{\otimes d})$
induisant une immersion
$$
\psi_1 : X_1 \to \mathbf{P}^N_{S_1}
$$
au-dessus de $S_1$, voir
Morphismes, lemme \ref{morphisms-lemma-finite-type-over-affine-ample-very-ample}.
Comme $X_1$ est propre sur $S_1$, $\psi_1$
est une immersion fermée, voir
Morphismes, lemme \ref{morphisms-lemma-image-proper-scheme-closed},
ainsi que
Schémas, lemme \ref{schemes-lemma-immersion-when-closed}.
Nous allons « relever » $\psi_1$ en un système compatible
d'immersions fermées des $X_n$ dans $\mathbf{P}^N$.

\medskip\noindent
En tensorisant les suites exactes courtes du premier paragraphe
de la démonstration
par $\mathcal{L}_{n + 1}^{\otimes d}$, on obtient les suites exactes courtes
$$
0 \to \mathcal{F}_n \otimes \mathcal{L}_{n + 1}^{\otimes d} \to
\mathcal{L}_{n + 1}^{\otimes d} \to \mathcal{L}_{n + 1}^{\otimes d} \to 0
$$
À l'aide des isomorphismes $\varphi_n$, on obtient des isomorphismes
$\mathcal{L}_{n + 1} \otimes \mathcal{O}_{X_l} = \mathcal{L}_l$
pour $l \leq n$. La suite précédente devient donc
$$
0 \to \mathcal{F}_n \otimes \mathcal{L}_1^{\otimes d} \to
\mathcal{L}_{n + 1}^{\otimes d} \to \mathcal{L}_n^{\otimes d} \to 0
$$
L'annulation de $H^1(X, \mathcal{F}_n \otimes \mathcal{L}_1^{\otimes d})$
permet de relever par récurrence
$s_{0, 1}, \ldots, s_{N, 1} \in \Gamma(X_1, \mathcal{L}_1^{\otimes d})$
en des sections
$s_{0, n}, \ldots, s_{N, n} \in \Gamma(X_n, \mathcal{L}_n^{\otimes d})$.
On obtient ainsi un diagramme commutatif
$$
\xymatrix{
X_1 \ar[r]_{i_1} \ar[d]_{\psi_1} &
X_2 \ar[r]_{i_2} \ar[d]_{\psi_2} &
X_3 \ar[r] \ar[d]_{\psi_3} &
\ldots \\
\mathbf{P}^N_{S_1} \ar[r] &
\mathbf{P}^N_{S_2} \ar[r] &
\mathbf{P}^N_{S_3} \ar[r] & \ldots
}
$$
où
$\psi_n = \varphi_{(\mathcal{L}_n, (s_{0, n}, \ldots, s_{N, n}))}$
avec les notations de Constructions, section
\ref{constructions-section-projective-space}.
Les carrés de l'énoncé du théorème étant
cartésiens, ceux du diagramme ci-dessus le sont
aussi. Le lemme \ref{lemma-algebraize-formal-closed-subscheme} permet alors de conclure.
\end{proof}

















\begin{multicols}{2}[\section{Autres chapitres}]
\noindent
Préliminaires
\begin{enumerate}
\item \hyperref[introduction-section-phantom]{Introduction}
\item \hyperref[conventions-section-phantom]{Conventions}
\item \hyperref[sets-section-phantom]{Théorie des ensembles}
\item \hyperref[categories-section-phantom]{Catégories}
\item \hyperref[topology-section-phantom]{Topologie}
\item \hyperref[sheaves-section-phantom]{Faisceaux sur les espaces}
\item \hyperref[sites-section-phantom]{Sites et faisceaux}
\item \hyperref[stacks-section-phantom]{Champs}
\item \hyperref[fields-section-phantom]{Corps}
\item \hyperref[algebra-section-phantom]{Algèbre commutative}
\item \hyperref[brauer-section-phantom]{Groupes de Brauer}
\item \hyperref[homology-section-phantom]{Algèbre homologique}
\item \hyperref[derived-section-phantom]{Catégories dérivées}
\item \hyperref[simplicial-section-phantom]{Méthodes simpliciales}
\item \hyperref[more-algebra-section-phantom]{Compléments d'algèbre}
\item \hyperref[smoothing-section-phantom]{Lissage des morphismes d'anneaux}
\item \hyperref[modules-section-phantom]{Faisceaux de modules}
\item \hyperref[sites-modules-section-phantom]{Modules sur les sites}
\item \hyperref[injectives-section-phantom]{Objets injectifs}
\item \hyperref[cohomology-section-phantom]{Cohomologie des faisceaux}
\item \hyperref[sites-cohomology-section-phantom]{Cohomologie sur les sites}
\item \hyperref[dga-section-phantom]{Algèbre différentielle graduée}
\item \hyperref[dpa-section-phantom]{Algèbre à puissances divisées}
\item \hyperref[sdga-section-phantom]{Faisceaux différentiels gradués}
\item \hyperref[hypercovering-section-phantom]{Hyperrecouvrements}
\end{enumerate}
Schémas
\begin{enumerate}
\setcounter{enumi}{25}
\item \hyperref[schemes-section-phantom]{Schémas}
\item \hyperref[constructions-section-phantom]{Constructions de schémas}
\item \hyperref[properties-section-phantom]{Propriétés des schémas}
\item \hyperref[morphisms-section-phantom]{Morphismes de schémas}
\item \hyperref[coherent-section-phantom]{Cohomologie des schémas}
\item \hyperref[divisors-section-phantom]{Diviseurs}
\item \hyperref[limits-section-phantom]{Limites de schémas}
\item \hyperref[varieties-section-phantom]{Variétés}
\item \hyperref[topologies-section-phantom]{Topologies sur les schémas}
\item \hyperref[descent-section-phantom]{Descente}
\item \hyperref[perfect-section-phantom]{Catégories dérivées de schémas}
\item \hyperref[more-morphisms-section-phantom]{Compléments sur les morphismes}
\item \hyperref[flat-section-phantom]{Compléments sur la platitude}
\item \hyperref[groupoids-section-phantom]{Schémas en groupoïdes}
\item \hyperref[more-groupoids-section-phantom]{Compléments sur les schémas en groupoïdes}
\item \hyperref[etale-section-phantom]{Morphismes étales de schémas}
\end{enumerate}
Thèmes de théorie des schémas
\begin{enumerate}
\setcounter{enumi}{41}
\item \hyperref[chow-section-phantom]{Homologie de Chow}
\item \hyperref[intersection-section-phantom]{Théorie de l'intersection}
\item \hyperref[pic-section-phantom]{Schémas de Picard des courbes}
\item \hyperref[weil-section-phantom]{Théories cohomologiques de Weil}
\item \hyperref[adequate-section-phantom]{Modules adéquats}
\item \hyperref[dualizing-section-phantom]{Complexes dualisants}
\item \hyperref[duality-section-phantom]{Dualité pour les schémas}
\item \hyperref[discriminant-section-phantom]{Discriminants et différentes}
\item \hyperref[derham-section-phantom]{Cohomologie de de Rham}
\item \hyperref[local-cohomology-section-phantom]{Cohomologie locale}
\item \hyperref[algebraization-section-phantom]{Géométrie algébrique et formelle}
\item \hyperref[curves-section-phantom]{Courbes algébriques}
\item \hyperref[resolve-section-phantom]{Résolution des surfaces}
\item \hyperref[models-section-phantom]{Réduction semi-stable}
\item \hyperref[functors-section-phantom]{Foncteurs et morphismes}
\item \hyperref[equiv-section-phantom]{Catégories dérivées de variétés}
\item \hyperref[pione-section-phantom]{Groupes fondamentaux de schémas}
\item \hyperref[etale-cohomology-section-phantom]{Cohomologie étale}
\item \hyperref[crystalline-section-phantom]{Cohomologie cristalline}
\item \hyperref[proetale-section-phantom]{Cohomologie pro-étale}
\item \hyperref[relative-cycles-section-phantom]{Cycles relatifs}
\item \hyperref[more-etale-section-phantom]{Compléments de cohomologie étale}
\item \hyperref[trace-section-phantom]{La formule des traces}
\end{enumerate}
Espaces algébriques
\begin{enumerate}
\setcounter{enumi}{64}
\item \hyperref[spaces-section-phantom]{Espaces algébriques}
\item \hyperref[spaces-properties-section-phantom]{Propriétés des espaces algébriques}
\item \hyperref[spaces-morphisms-section-phantom]{Morphismes d'espaces algébriques}
\item \hyperref[decent-spaces-section-phantom]{Espaces algébriques décents}
\item \hyperref[spaces-cohomology-section-phantom]{Cohomologie des espaces algébriques}
\item \hyperref[spaces-limits-section-phantom]{Limites d'espaces algébriques}
\item \hyperref[spaces-divisors-section-phantom]{Diviseurs sur les espaces algébriques}
\item \hyperref[spaces-over-fields-section-phantom]{Espaces algébriques sur des corps}
\item \hyperref[spaces-topologies-section-phantom]{Topologies sur les espaces algébriques}
\item \hyperref[spaces-descent-section-phantom]{Descente et espaces algébriques}
\item \hyperref[spaces-perfect-section-phantom]{Catégories dérivées d'espaces}
\item \hyperref[spaces-more-morphisms-section-phantom]{Compléments sur les morphismes d'espaces}
\item \hyperref[spaces-flat-section-phantom]{Platitude sur les espaces algébriques}
\item \hyperref[spaces-groupoids-section-phantom]{Groupoïdes dans les espaces algébriques}
\item \hyperref[spaces-more-groupoids-section-phantom]{Compléments sur les groupoïdes dans les espaces}
\item \hyperref[bootstrap-section-phantom]{Amorçage}
\item \hyperref[spaces-pushouts-section-phantom]{Sommes amalgamées d'espaces algébriques}
\end{enumerate}
Thèmes de géométrie
\begin{enumerate}
\setcounter{enumi}{81}
\item \hyperref[spaces-chow-section-phantom]{Groupes de Chow des espaces}
\item \hyperref[groupoids-quotients-section-phantom]{Quotients de groupoïdes}
\item \hyperref[spaces-more-cohomology-section-phantom]{Compléments sur la cohomologie des espaces}
\item \hyperref[spaces-simplicial-section-phantom]{Espaces simpliciaux}
\item \hyperref[spaces-duality-section-phantom]{Dualité pour les espaces}
\item \hyperref[formal-spaces-section-phantom]{Espaces algébriques formels}
\item \hyperref[restricted-section-phantom]{Algébrisation des espaces formels}
\item \hyperref[spaces-resolve-section-phantom]{Résolution des surfaces revisitée}
\end{enumerate}
Théorie des déformations
\begin{enumerate}
\setcounter{enumi}{89}
\item \hyperref[formal-defos-section-phantom]{Théorie formelle des déformations}
\item \hyperref[defos-section-phantom]{Théorie des déformations}
\item \hyperref[cotangent-section-phantom]{Le complexe cotangent}
\item \hyperref[examples-defos-section-phantom]{Problèmes de déformation}
\end{enumerate}
Champs algébriques
\begin{enumerate}
\setcounter{enumi}{93}
\item \hyperref[algebraic-section-phantom]{Champs algébriques}
\item \hyperref[examples-stacks-section-phantom]{Exemples de champs}
\item \hyperref[stacks-sheaves-section-phantom]{Faisceaux sur les champs algébriques}
\item \hyperref[criteria-section-phantom]{Critères de représentabilité}
\item \hyperref[artin-section-phantom]{Axiomes d'Artin}
\item \hyperref[quot-section-phantom]{Espaces de Quot et de Hilbert}
\item \hyperref[stacks-properties-section-phantom]{Propriétés des champs algébriques}
\item \hyperref[stacks-morphisms-section-phantom]{Morphismes de champs algébriques}
\item \hyperref[stacks-limits-section-phantom]{Limites de champs algébriques}
\item \hyperref[stacks-cohomology-section-phantom]{Cohomologie des champs algébriques}
\item \hyperref[stacks-perfect-section-phantom]{Catégories dérivées de champs}
\item \hyperref[stacks-introduction-section-phantom]{Introduction aux champs algébriques}
\item \hyperref[stacks-more-morphisms-section-phantom]{Compléments sur les morphismes de champs}
\item \hyperref[stacks-geometry-section-phantom]{La géométrie des champs}
\end{enumerate}
Thèmes de théorie des espaces de modules
\begin{enumerate}
\setcounter{enumi}{107}
\item \hyperref[moduli-section-phantom]{Champs de modules}
\item \hyperref[moduli-curves-section-phantom]{Espaces de modules de courbes}
\end{enumerate}
Divers
\begin{enumerate}
\setcounter{enumi}{109}
\item \hyperref[examples-section-phantom]{Exemples}
\item \hyperref[exercises-section-phantom]{Exercices}
\item \hyperref[guide-section-phantom]{Guide bibliographique}
\item \hyperref[desirables-section-phantom]{Desiderata}
\item \hyperref[coding-section-phantom]{Style de programmation}
\item \hyperref[obsolete-section-phantom]{Obsolète}
\item \hyperref[fdl-section-phantom]{Licence de documentation libre GNU}
\item \hyperref[index-section-phantom]{Index généré automatiquement}
\end{enumerate}
\end{multicols}


\bibliography{my}
\bibliographystyle{amsalpha}

\end{document}
